% PHP Gallery Vollständige Produktdokumentation
% Project: PHP Gallery
% Repository: https://github.com/klusik/PHP_gallery
% File: docs/PHP_Gallery_Manual_DE.tex
% Module Type: LaTeX-Produkthandbuch
% Purpose: Administrative Arbeitsabläufe, Betrieb und Wartung von PHP Gallery erläutern.
% Responsibilities:
%   - Nutzungsleitfaden und technische Referenz unter Erhalt reproduzierbarer PDF-Typografie pflegen.
% Author: Rudolf Klusal
% In diesem Verzeichnis erstellen mit:
%   pdflatex PHP_Gallery_Manual_DE.tex
%   makeindex PHP_Gallery_Manual_DE.idx
%   pdflatex PHP_Gallery_Manual_DE.tex
%   pdflatex PHP_Gallery_Manual_DE.tex
%
% Die in den Quellen aufgeführten Markdown- und Implementierungsdateien sind die
% redaktionellen und technischen Quellen dieser Ausgabe. Prüfen Sie das lokale Projekt,
% wenn sich die Anwendungsversion ändert.

\documentclass[11pt,a4paper]{article}

\usepackage[utf8]{inputenc}
\usepackage[T1]{fontenc}
\usepackage{lmodern}
\usepackage[ngerman]{babel}
\usepackage{geometry}
\usepackage{microtype}
\usepackage{xcolor}
\usepackage{graphicx}
\usepackage{booktabs}
\usepackage{array}
\usepackage{ragged2e}
\usepackage{longtable}
\usepackage{tabularx}
\usepackage{enumitem}
\usepackage{listings}
\usepackage{fancyhdr}
\usepackage{titlesec}
\usepackage[bottom,hang,flushmargin]{footmisc}
\usepackage{makeidx}
\usepackage{hyperref}
\usepackage{bookmark}

\geometry{top=23mm,bottom=25mm,left=24mm,right=24mm,headheight=15pt}
\setlength{\parindent}{0pt}
\setlength{\emergencystretch}{1.5em}
\setlength{\parskip}{0.65em}
\setcounter{tocdepth}{3}
\setcounter{secnumdepth}{3}
\makeatletter
% Zweistellige Unterabschnittsnummern von ihren Titeln im Inhaltsverzeichnis getrennt halten.
\renewcommand*\l@subsection{\@dottedtocline{2}{1.5em}{4.1em}}
\makeatother
\setlist{itemsep=0.3em,topsep=0.45em}
\renewcommand{\arraystretch}{1.2}

\definecolor{GalleryBlue}{HTML}{245A86}
\definecolor{GalleryLightBlue}{HTML}{EAF2F8}
\definecolor{GalleryDark}{HTML}{263238}
\definecolor{GalleryGray}{HTML}{F3F5F7}
\definecolor{GalleryCode}{HTML}{1F2933}
\definecolor{GalleryGreen}{HTML}{2E6B45}
\definecolor{GalleryRed}{HTML}{8A2F2F}

% Ausgabenmetadaten stehen hier, damit eine Versionsaktualisierung nur eine Versionsänderung benötigt.
\newcommand{\version}{0.124.4}
\newcommand{\manualdate}{8.~Oktober~2026}

\hypersetup{
  colorlinks=true,
  linkcolor=GalleryBlue,
  citecolor=GalleryGreen,
  urlcolor=GalleryBlue,
  pdftitle={PHP Gallery Vollständige Produktdokumentation},
  pdfauthor={Rudolf Klusal},
  pdfsubject={Referenz zu Installation, Administration, Architektur, Sicherheit, Wartung und Entwicklung},
  pdfkeywords={PHP Gallery, Galerie-CMS, Administration, Architektur, Dokumentation},
  pdfdisplaydoctitle=true,
  bookmarksopen=true,
  bookmarksnumbered=true
}

\pagestyle{fancy}
\fancyhf{}
\fancyhead[L]{\small PHP Gallery}
\fancyhead[R]{\small Vollständige Produktdokumentation}
\fancyfoot[L]{\small Version~\version}
\fancyfoot[C]{\thepage}
\fancyfoot[R]{\small \manualdate}
\renewcommand{\headrulewidth}{0.4pt}
\renewcommand{\footrulewidth}{0.2pt}

\titleformat{\section}{\Large\bfseries\color{GalleryBlue}}{\thesection}{0.75em}{}
\titleformat{\subsection}{\large\bfseries\color{GalleryDark}}{\thesubsection}{0.75em}{}
\titleformat{\subsubsection}{\normalsize\bfseries\color{GalleryDark}}{\thesubsubsection}{0.75em}{}
\pretocmd{\section}{\clearpage}{}{}

\lstdefinestyle{gallery}{
  basicstyle=\ttfamily\small\color{GalleryCode},
  backgroundcolor=\color{GalleryGray},
  frame=single,
  rulecolor=\color{GalleryLightBlue},
  breaklines=true,
  breakatwhitespace=false,
  columns=fullflexible,
  keepspaces=true,
  showstringspaces=false,
  tabsize=4,
  xleftmargin=0.8em,
  xrightmargin=0.8em,
  framexleftmargin=0.4em,
  aboveskip=0.8em,
  belowskip=0.8em
}
\lstset{style=gallery}

\newcommand{\product}{\textbf{PHP Gallery}}
\newcommand{\setting}[1]{\nolinkurl{#1}}
\newcommand{\filepath}[1]{\path{#1}}
\newcommand{\ui}[1]{\textbf{#1}}
\newcommand{\codeword}[1]{\nolinkurl{#1}}
\newcommand{\important}[1]{%
  \begin{center}\fcolorbox{GalleryBlue}{GalleryLightBlue}{\parbox{0.91\linewidth}{\textbf{Wichtig.} #1}}\end{center}}
\newcommand{\securitynote}[1]{%
  \begin{center}\fcolorbox{GalleryRed}{GalleryGray}{\parbox{0.91\linewidth}{\textbf{Sicherheitshinweis.} #1}}\end{center}}
\newcommand{\performancenote}[1]{%
  \begin{center}\fcolorbox{GalleryGreen}{GalleryGray}{\parbox{0.91\linewidth}{\textbf{Leistungshinweis.} #1}}\end{center}}

\makeindex

\begin{document}

\pagenumbering{gobble}
\begin{titlepage}
  \centering
  \vspace*{24mm}
  {\Huge\bfseries\color{GalleryBlue} PHP Gallery\par}
  \vspace{6mm}
  {\LARGE Vollständige Produktdokumentation\par}
  \vspace{14mm}
  {\large Referenz zu Installation, Administration, Architektur, Sicherheit, Wartung und Entwicklung\par}
  \vfill
  \begin{tabular}{rl}
    \textbf{Anwendungsversion:} & \version\\
    \textbf{Handbuchausgabe:} & \manualdate\\
    \textbf{Dokumentformat:} & A4, PDF mit Index und Verknüpfungen\\
    \textbf{Projektautor:} & Rudolf Klusal\\
    \textbf{Lizenz:} & MIT-Lizenz\\
  \end{tabular}
  \vfill
  {\small Leitfaden zum Erstellen, Organisieren, Präsentieren, Schützen, Teilen und Bewahren einer Fotosammlung mit PHP Gallery Version~\version.\par}
\end{titlepage}

\pagenumbering{roman}
\begingroup
\small
\setlength{\parskip}{0pt}
\tableofcontents
\endgroup
\clearpage
\pagenumbering{arabic}

\section{Zweck und Lesehinweise}
\index{d00079@Dokumentation}
\index{h00146@Handbuch!u00246@Umfang}

\product{} ist eine selbst gehostete Fotogalerie. Die Fotos bleiben auf einem Speicher, den Sie kontrollieren; die Anwendung ergänzt Katalog, Navigation, Zugriffsregeln, Vorschaubilder, Karten, Suche, Downloads und Verwaltungswerkzeuge. Eine Installation eignet sich für ein öffentliches Portfolio, ein Familienarchiv, die Übergabe an Kunden, Veranstaltungsfotos, Reisesammlungen oder eine große persönliche Bibliothek.

Das Handbuch orientiert sich an den Aufgaben, die ein Betreiber tatsächlich ausführt. Die ersten Abschnitte behandeln die Installation und die erste Galerie; der mittlere Teil behandelt Veröffentlichung, Medien, Datenschutz, Präsentation und Wartung. Architektur und Entwicklerthemen stehen weiter hinten, damit sie die alltägliche Administration nicht unterbrechen.

Beginnen Sie bei einer neuen Installation mit Abschnitt~\ref{sec:first-hour}. Die Planung Ihrer Sammlung behandelt Abschnitt~\ref{sec:planning}; die regelmäßige Bearbeitung von Galerien und Fotos finden Sie in den Abschnitten~\ref{sec:create-edit-gallery} und~\ref{sec:photo-workflow}. Abschnitt~\ref{sec:audiences} ist vor der Veröffentlichung hilfreich, weil er verschiedene typische Zielgruppen vergleicht. Inhaltsverzeichnis und Index ermöglichen den direkten Sprung zu einer Einstellung oder einem Arbeitsablauf, wenn die Website bereits genutzt wird.





\subsection{So verwenden Sie dieses Handbuch}
\index{h00146@Handbuch!v00257@Verwendung}

Bezeichnungen aus der Administration wie \ui{Design} und \ui{Systemzustand} erscheinen fett. Nichtproportionale Schrift kennzeichnet Dateinamen, Einstellungsnamen, Werte oder Befehle. Bei Abläufen mit verbindlicher Reihenfolge werden nummerierte Listen verwendet; Checklisten und Referenzinformationen stehen in Aufzählungen oder Tabellen. Technische Hinweise sind nur enthalten, wenn sie ein sichtbares Verhalten, eine Wiederherstellungsentscheidung oder eine Wartungsanforderung erklären.

\subsection{Weitere Projektreferenzen}

Referenzen für Entwicklung und Hosting bleiben von diesem Betreiberhandbuch getrennt. Der Abschnitt mit den Quellen führt die Repository-Dokumente zu Architektur, Datenbank, Quellcodeverantwortung, Tests und Beitragsregeln auf \cite{readme,architecture,database,codemap,testing,agents}.\footnote{Diese Dateien sind die bessere Quelle, wenn Sie Code ändern, Schemadetails prüfen oder eine Veröffentlichung vorbereiten.}

\subsection{Begriffe}

\begin{description}
  \item[Galerie] Eine benannte Fotosammlung. Eine Galerie kann Fotos und kleinere Galerien enthalten.
  \item[Foto oder Bild] Ein einzelner Eintrag einer Galerie mit Titel, Beschreibung, Schlagwörtern, Sichtbarkeit und verfügbaren Kamerainformationen.
  \item[Vorschaubild] Eine kleinere Kopie für die schnelle Anzeige in Galerierastern, Suchergebnissen, Titelbildern und Vorschauen.
  \item[Besucher] Eine Person, die den öffentlichen Bereich der Website betrachtet.
  \item[Administrator] Eine vertrauenswürdige Person, die sich anmeldet, um Inhalte und Einstellungen zu erstellen und zu verwalten.
  \item[Ausgewählte Galerie] Die derzeit auf dem Bildschirm geöffnete Galerie.
  \item[Galeriezweig] Eine Galerie einschließlich aller darin untergeordneten Galerien.
\end{description}

\section{Produktübersicht}
\label{sec:overview}
\index{p00193@Produktübersicht}
\index{s00220@Shared Hosting}

\product{} macht aus einer Ordnersammlung eine durchsuchbare Fotowebsite. Besucher sehen Titelbilder, Titel, Beschreibungen, Fotoraster, Vollbildansichten, Karten, Schlagwörter, Suche und Downloads entsprechend den Entscheidungen des Betreibers. Administratoren erhalten private Werkzeuge zum Hinzufügen von Fotos, Ordnen von Sammlungen, Steuern der Privatsphäre, Verbessern von Beschreibungen und Erhalten einer funktionsfähigen Installation.

Gewöhnliches PHP/MySQL-Webhosting ist die primäre Zielumgebung. Der Betreiber hält die Originalfotos in nachvollziehbaren Ordnern; die Datenbank liefert den Katalog und den für das Durchsehen benötigten Anwendungszustand.\footnote{Die Projektübersicht enthält eine kurze Zusammenfassung von Installation und Funktionen \cite{readme}.}

\subsection{Wesentliche Funktionen}

\begin{itemize}
  \item Erstellen Sie Galerien innerhalb anderer Galerien, beispielsweise \emph{2026 / Italien / Rom} oder \emph{Kunden / Northwind / Endauswahl}.
  \item Geben Sie jeder Galerie einen Titel, eine Beschreibung, einen Datumsbereich, ein Titelbild, eine öffentliche Adresse, eine Datenschutzeinstellung und ein individuelles Erscheinungsbild.
  \item Laden Sie Fotos über den Browser hoch oder lassen Sie Dateien erkennen, die in Galerieordner kopiert wurden.
  \item Ergänzen Sie Fotos um Titel, Geschichten, Schlagwörter, Sichtbarkeit, Reihenfolge, Datums- und Ortsinformationen.
  \item Ermöglichen Sie Besuchern die Suche, das Durchsehen von Schlagwörtern, die Vollbildanzeige von Fotos, die Kartennavigation, Abstimmungen und Downloads erlaubter Sammlungen.
  \item Schützen Sie Familien-, Kunden- oder sensible Sammlungen mit Zugriffskontrollen und Freigabemethoden.
  \item Finden Sie wahrscheinliche Fotoduplikate und prüfen Sie sie über direkte Links zu beiden Speicherorten.
  \item Halten Sie die Installation mit geführten Aktualisierungen, Zustandsprüfungen, Sicherungen und Wartungswerkzeugen aktuell.
  \item Automatisieren Sie auf bestimmte Bereiche beschränkte Uploads, tauschen Sie Galerien mit kompatiblen Installationen aus und bereiten Sie Downloads ausgewählter Fotos oder ganzer Galerien vor.
  \item Erzeugen Sie optional KI-Metadaten und Übersetzungsentwürfe, wobei jeder erzeugte Wert vor der Veröffentlichung überprüfbar bleibt.
  \item Veröffentlichen Sie crawlerfreundliche Robots- und Bild-Sitemap-Ausgaben für öffentliche Inhalte, auf die Besucher tatsächlich zugreifen können.
\end{itemize}

\subsection{Zusammenspiel von Dateisystem und Datenbank}
\index{d00061@dateisystemorientiertes Konzept}
\index{d00062@Datenbank}

PHP Gallery speichert dauerhafte Daten sowohl im Dateisystem als auch in der Datenbank. Galerieverzeichnisse enthalten die Originalfotos in einem nachvollziehbaren Verzeichnisbaum. Die Datenbank speichert Informationen, die nicht zur Datei selbst gehören: Titel, Beschreibungen, Schlagwörter, Reihenfolge, Zugriffseinstellungen, Stimmen, Datumsangaben, Datensätze erzeugter Medien und weitere Anwendungszustände.

Keine der beiden Seiten stellt allein eine vollständige Sicherung dar. Eine Wiederherstellung nur der Datenbank hinterlässt Katalogeinträge ohne Originaldateien; eine Wiederherstellung nur des Galeriebaums verliert den Verwaltungs- und Präsentationszustand dieser Dateien. Sicherung und Wiederherstellung sollten daher Datenbank, Galerieverzeichnisse und lokale Konfiguration als zusammengehörigen Satz behandeln. Abschnitt~\ref{sec:user-data} geht darauf ausführlicher ein.

\section{Anforderungen und Umgebung}
\label{sec:requirements}
\index{a00015@Anforderungen}
\index{p00186@PHP}
\index{m00175@MySQL}
\index{m00167@MariaDB}

Verwenden Sie diesen Abschnitt bei der Auswahl eines Hostingpakets oder der Prüfung eines vorhandenen Kontos. Die Anforderungstabelle eignet sich auch zur direkten Weitergabe an einen Hostinganbieter.

\subsection{Mindestanforderungen an die Laufzeitumgebung}

\begin{longtable}{>{\bfseries}>{\hspace{0pt}}p{0.27\linewidth}>{\hspace{0pt}}p{0.65\linewidth}}
\toprule
Komponente & Anforderung\\
\midrule
\endhead
PHP & Version~8.1 oder neuer.\\
Datenbank & MySQL~8.4 LTS, MariaDB~10.11 LTS oder MariaDB~11.4 LTS.\\
PDO & PDO-MySQL-Erweiterung für Anwendungsabfragen und Migrationen.\\
Bildverarbeitung & GD für die übliche Vorschaubilderzeugung. Zusätzliche lokale Bildwerkzeuge können spezielle Formate unterstützen.\\
Archive & ZipArchive für Paket-, Download- und Übertragungsabläufe, die ZIP-Daten erzeugen oder verarbeiten.\\
Dateisystem & Beschreibbare Speicherorte für Anwendungsdaten, Galerien, Cache und erzeugte Medien.\\
Webserver & Umschreibregeln nach Apache-Art werden für sprechende URLs empfohlen; das Routing über Abfrageparameter unterstützt kompatible Umgebungen ohne Umschreibregeln.\\
\bottomrule
\end{longtable}

Die gepflegte Supportmatrix qualifiziert MySQL 8.4 LTS, MariaDB 10.11 LTS und MariaDB 11.4 LTS als unmittelbar in der vorgeschriebenen CI getestete Vertreter. Das kanonische Dokument führt MySQL 8.0 und MariaDB 10.4--10.10 als ältere Versionen mit Unterstützung nach bestem Bemühen auf; für sie gibt es weder einen vorgeschriebenen Datenbank-CI-Lauf noch eine projektseitige Regressions- oder Sicherheitszusage. Sie liegen außerhalb der gepflegten Qualifizierung. MySQL 5.7 und älter sowie MariaDB 10.3 und älter werden nicht unterstützt. Andere Datenbankserien bleiben unqualifiziert, bis sie direkt durch die vorgeschriebene CI abgedeckt sind; eine höhere Versionsnummer allein belegt keine Kompatibilität. Einzelheiten enthält \filepath{docs/DATABASE_SUPPORT.md}.

Ausgehender HTTPS-Zugriff unterstützt Anwendungsaktualisierungen und entfernte Veröffentlichungsmetadaten. Lokaler Betrieb und manuell bereitgestellte Versionen können auch bei stärker eingeschränktem ausgehendem Zugriff funktionieren.\footnote{Die Netzwerkrichtlinie sollte den Aktualisierungs- und Datenschutzanforderungen des Installationsbetreibers folgen. Der Aktualisierungscode prüft Pakete und bewahrt Wiederherstellungsdaten entsprechend der lokalen Aktualisierungsimplementierung auf.}

\subsection{Berechtigungen}
\index{b00031@Berechtigungen!d00060@Dateisystem}

Der Webprozess benötigt Lesezugriff auf den Laufzeitcode und Schreibzugriff auf die vorgesehenen veränderlichen Speicherorte. Halten Sie den Laufzeitcode schreibgeschützt, soweit das Hosting dies erlaubt, außer während eines kontrollierten Aktualisierungsvorgangs. Originaldateien der Galerien benötigen denselben Sicherungsumfang wie die Datenbank, weil sich der vollständige Anwendungszustand über beide Bereiche erstreckt.

\important{Sichern Sie Datenbank, Galeriedateien und Installationskonfiguration gemeinsam. Eine reine Datenbanksicherung bewahrt Metadaten, lässt aber den maßgeblichen Medienbaum aus.}

\section{Installation und Ersteinrichtung}
\label{sec:installation}
\index{i00149@Installation}
\index{e00082@Einrichtung}

Interne Verzeichnisse für Anwendung, Werkzeuge, Dokumentation, Cache, Daten und Protokolle besitzen eigene Apache-Zugriffssperren unabhängig von URL-Umschreibungen. Der Hoster muss diese Regeln beachten oder gleichwertige Serverregeln bereitstellen; der PHP-Entwicklungsserver liest sie nicht. Direkte HTTP-Aufrufe von CLI-Befehlen erhalten vor dem Anwendungsstart eine leere Antwort 404. Verwenden Sie für Browserrouten den öffentlichen Einstiegspunkt und für Wartungsbefehle die Kommandozeile. Siehe \filepath{docs/HTTP_ENTRYPOINTS.md}.

Erstellen Sie Bereitstellungsordner oder ZIP-Dateien bei Bedarf aus der geprüften Liste \filepath{app/production-files.json}. Die Werkzeuge prüfen die exakte Dateimenge und aktuelle Integritätsdaten; lokale Agentendateien, aktive Medien, Konfiguration, eigenes CSS und Laufzeitdaten gehören nicht zum normalen Paket. Freigegebene WinApp-Quellen werden mitgeliefert, der Installer bleibt unabhängig. Enthält ein interaktiv gewähltes Windows-Ziel bereits ein Paket, wird ein neuer benachbarter Pfad mit Zeitstempel gewählt; ein ausdrücklich angegebenes Automatisierungsziel lehnt Kollisionen ab. Lesen Sie vor manueller Bereitstellung \filepath{docs/PRODUCTION_FILES.md}.

\subsection{Installation im Browser}

\subsubsection{Installationsablauf}

Das browserbasierte Installationsprogramm erfasst Datenbankeinstellungen, initialisiert das Schema durch Migrationen, erstellt den ersten Administrator, bereitet beschreibbare Speicherorte vor und schreibt die lokale Konfiguration. Fehlt eine erforderliche Konfiguration, leitet eine neue Anfrage zur Installation weiter.\footnote{Installationszustand und erzeugte Konfiguration sollten ebenso geschützt werden wie Zugangsdaten. Quellcodearchive sollten die lokale \filepath{config.php} weiterhin ausschließen.}

Migrationen werden vor~dem ersten Administratorkonto, \filepath{config.php} und~der Installationssperre abgeschlossen. Reparaturen verwenden die~geprüfte Installer-Verbindung und~die~gemeinsame Schreibsperre für~Galerien; vor~der Konfiguration liegt der~Galeriespeicher im~Installationsordner \filepath{galleries/}. Beheben Sie~bei einer fehlgeschlagenen Migration die~Ursache und~senden Sie~den Administratorschritt erneut. Abgeschlossene Migrationen bleiben protokolliert, der~fehlgeschlagene Callback wird erneut ausgeführt. Bewahren Sie~das Migrationsprotokoll sowie Konfiguration und~Sperren abgeschlossener Installationen auf.

Typische Schritte sind:

\begin{enumerate}
  \item Erstellen Sie im Hosting-Kontrollbereich eine leere MySQL- oder MariaDB-Datenbank.
  \item Laden Sie das Anwendungspaket in das vorgesehene Webverzeichnis hoch.
  \item Öffnen Sie die Website und folgen Sie dem Installationsprogramm.
  \item Geben Sie die Datenbankverbindungswerte und die Zugangsdaten des ersten Administrators ein.
  \item Bestätigen Sie die beschreibbaren Verzeichnisse und die Erweiterungsprüfungen.
  \item Schließen Sie die Installation ab und öffnen Sie den Administrationsbereich.
\end{enumerate}

\subsection{Einrichtung über die Befehlszeile}

\subsubsection{Lokale Befehle}

Das Repository bietet einfache PHP-Skripte für Migrationen und die Administratorerstellung:

\begin{lstlisting}[language=bash]
php scripts/migrate.php
php scripts/create_admin.php <username> <password>
\end{lstlisting}

Die Einrichtung über die Befehlszeile folgt demselben dauerhaften Schemamodell wie die Browsereinrichtung. Prüfen Sie \filepath{config.example.php}, erstellen Sie die lokale Konfiguration, richten Sie die Datenbank ein, führen Sie Migrationen aus und erstellen Sie den ersten Administrator.

\subsection{Erste Prüfung der Administration}

Prüfen Sie nach der Installation:

\begin{itemize}
  \item Websitename, Sprache, Design, Seitenbreite und öffentliche Darstellungspräferenzen.
  \item Berechtigungen des Galeriestammverzeichnisses und Erkennungsergebnisse.
  \item Vorschaubildformate, Abmessungen, Qualität und Wartungszustand.
  \item Einstellungen für Zugriff, Passwortzurücksetzung, E-Mail, Telemetrie, Aktualisierungskanal und Sicherung.
  \item Kompatibilität der Umschreibregeln und kanonische öffentliche URLs.
  \item Diagnose des Systemzustands und ausstehende Migrationen.
\end{itemize}

\subsection{Sprechende URLs und Kompatibilität der Umschreibregeln}
\index{u00250@URL-Umschreibung}
\index{r00208@Routing über Abfrageparameter}
\index{s00232@sprechende URLs}

PHP Gallery unterstützt sprechende öffentliche Pfade, wenn die Hostingumgebung die Umschreibregeln des Projekts anwendet. Sprechende URLs sind eine Präsentations- und Routingpräferenz, keine Anforderung des Galeriedatenmodells. Routen mit Abfrageparametern bleiben der Kompatibilitätsweg und funktionieren weiter, wenn Umschreibungen deaktiviert sind oder nicht zuverlässig überprüft werden können.

Die Einstellung zur URL-Umschreibung steuert, welche Form PHP Gallery in erzeugten Links ausgibt. Ist die Umschreibung ausgeschaltet, verwenden Hilfsfunktionen Adressen mit Abfrageparametern für Galerien, Fotos, Vorschaubilder und andere Anwendungsrouten, statt die Übersetzung eines sprechenden Pfads durch den Webserver vorauszusetzen. Slugs, Galerieordner, Datenbankzeilen und Originalfotos müssen bei einer Änderung dieser Einstellung nicht umbenannt werden.

Scheitern sprechende Pfade bei einem neuen Hoster, deaktivieren Sie zunächst die URL-Umschreibung in der Administration und prüfen Sie, ob die Form mit Abfrageparametern, etwa \codeword{?page=gallery}, funktioniert. Prüfen Sie auf Apache oder LiteSpeed anschließend, ob die \filepath{.htaccess} des Projekts berücksichtigt wird und Umschreibunterstützung verfügbar ist. Behalten Sie auf anderen Webservern den Kompatibilitätsmodus bei, sofern keine entsprechenden Routingregeln eingerichtet wurden.

\section{Ihre erste Stunde mit PHP Gallery}
\label{sec:first-hour}
\index{e00092@erste Stunde}
\index{e00085@Einstieg}
\index{g00122@Galeriebetreiber}
\index{a00007@Administrationsübersicht}

Halten Sie den Umfang der ersten Sitzung klein: Erstellen Sie eine Galerie, fügen Sie einige Fotos hinzu und öffnen Sie das Ergebnis in einem Browser, in dem Sie nicht als Administrator angemeldet sind. Designarbeit und eine tiefere Strukturierung können warten, bis dieser grundlegende Ablauf nachweislich funktioniert.

\subsection{Administrationsbereich öffnen}
\index{a00008@Administratoranmeldung}

Wählen Sie \ui{Administratoranmeldung} in der öffentlichen Kopfzeile und geben Sie das während der Installation erstellte Konto ein. Die Übersicht ist der Ausgangspunkt für die alltägliche Verwaltung. Sie fasst Galerien, Fotos, Systemzustand und verfügbare Wartungsaktionen zusammen.\footnote{Die genauen Inhalte der Übersicht hängen von aktivierten Funktionen, abgeschlossenen Migrationen und Entwicklungseinstellungen ab. Auch eine fehlerfreie Installation kann informative Hinweise anzeigen, die einen Administrator zur Erledigung optionaler Aufgaben auffordern.}

Nach der Anmeldung kann ein Rücksprunglink eine lokale Seite dieser Installation öffnen. Externe URLs und Rücksprungpfade mit umgekehrten Schrägstrichen werden abgelehnt; stattdessen verwendet die Anwendung das sichere lokale Ersatzziel des Anmeldeablaufs.

Nehmen Sie sich einige Minuten, um diese Bereiche kennenzulernen:

\begin{itemize}
  \item \ui{Galerien}, in denen Sammlungen erstellt, erkannt, bearbeitet und geordnet werden;
  \item \ui{Design}, in dem Erscheinungsbild und Präsentationsvorgaben der öffentlichen Website ausgewählt werden;
  \item \ui{Schlagwörter}, in denen wiederverwendbare Themen und Kategorien gepflegt werden;
  \item \ui{Systemzustand}, in dem Migrationen, Vorschaubilder, Speicher und Konsistenz geprüft werden;
  \item \ui{Konto}, in dem Zugangsdaten, Wieder\-her\-stel\-lungs-E-Mail und E-Mail-Zustellung verwaltet werden.
\end{itemize}

\subsection{Erste Galerie erstellen}
\index{g00117@Galerie!e00091@erste Galerie erstellen}

Öffnen Sie die Galerieverwaltung und erstellen Sie eine Galerie mit einem kurzen Titel wie \emph{Sommerspaziergang}. Prüfen Sie den vorgeschlagenen öffentlichen Pfad und Ordnernamen und ergänzen Sie dann einen Satz zum Inhalt der Sammlung.

Machen Sie die Galerie für diese erste Prüfung öffentlich und belassen Sie erweitertes Branding und geerbte Überschreibungen bei ihren Vorgaben. Fügen Sie drei bis zehn Fotos hinzu. Das genügt, um Reihenfolge, Vorschaubilderzeugung, Titel und Lightbox zu prüfen, ohne aus der Übung einen großen Import zu machen.

\subsection{Vorschau als Besucher}
\index{a00017@anonyme Vorschau}
\index{b00039@Besuchervorschau}

Verwenden Sie den Link zur öffentlichen Galerie oder die Besuchervorschau. Öffnen Sie für eine realistische Prüfung die öffentliche Adresse in einem privaten Browserfenster ohne Administratorsitzung. Prüfen Sie, ob Titel, Beschreibung, Fotos und Navigation wie vorgesehen erscheinen. Öffnen Sie ein Foto, wechseln Sie im Betrachter vor und zurück und kehren Sie zur Galerie zurück.

Administratoransichten können Bearbeitungssteuerungen und zusätzliche Ressourcen enthalten. Die anonyme Vorschau vermittelt das klarste Bild dessen, was Familie, Kunden oder Öffentlichkeit erleben.\footnote{Geschützte Galerien und Inhalte mit Altersbeschränkung verdienen einen gesonderten anonymen Test, weil eine aktive Administratorsitzung bereits weitergehenden Zugriff besitzt.}

\subsection{Checkliste für die erste Stunde abschließen}

\begin{enumerate}
  \item Legen Sie den öffentlichen Websitenamen und die Sprache fest.
  \item Erstellen Sie eine Galerie mit Titel und Beschreibung.
  \item Laden Sie eine kleine Gruppe von Fotos hoch.
  \item Geben Sie mindestens einem Foto einen aussagekräftigen Titel.
  \item Ändern Sie die Reihenfolge zweier Fotos und prüfen Sie die neue Reihenfolge öffentlich.
  \item Öffnen Sie die Galerie als anonymer Besucher.
  \item Prüfen Sie die Galerie auf einem Bildschirm in Telefongröße.
  \item Bestätigen Sie, dass der Systemzustand keine dringenden Migrations- oder Berechtigungsprobleme meldet.
  \item Erstellen Sie nach Abschluss der Ersteinrichtung eine abgestimmte Sicherung.
\end{enumerate}

\section{Eine Galeriesammlung planen}
\label{sec:planning}
\index{g00131@Galerieplanung}
\index{s00210@Sammlungsgestaltung}
\index{g00127@Galeriehierarchie!p00189@Planung}

Eine durchdachte Struktur hilft Besuchern, eine Sammlung zu verstehen, und dem Betreiber, sie langfristig zu pflegen. Beginnen Sie damit, wie Menschen nach Fotos suchen. Ordnernamen und Datenbankdetails können dieser menschlichen Ordnung folgen.

\subsection{Eine nachvollziehbare Struktur wählen}

Übliche Strukturen sind:

\begin{description}
  \item[Nach Jahr und Ereignis] Ein Jahr auf oberster Ebene enthält Hochzeiten, Reisen, Feiern und Projekte. Das eignet sich für Familien- und historische Archive.
  \item[Nach Ort] Länder enthalten Regionen und Städte. Das eignet sich für Reise-, Luftfahrt-, Landschafts- und Dokumentarfotografie.
  \item[Nach Kunde oder Projekt] Jeder Kunde enthält Aufträge und Übergabegalerien. Das eignet sich für professionelle Arbeit mit getrennten Zugriffsregeln.
  \item[Nach Thema] Wildtiere, Architektur, Porträts, Flugzeuge und ähnliche Themen bilden dauerhafte Kategorien. Schlagwörter bieten einen zweiten, kategorienübergreifenden Weg zum Durchsehen.
  \item[Nach Arbeitsablauf] Eingangs-, Auswahl-, Übergabe- und Archivgalerien stehen für unterschiedliche Phasen. Sichtbarkeitseinstellungen halten Arbeitssammlungen privat, während fertige Auswahlen öffentlich werden.
\end{description}

Verwenden Sie Galerien für deutliche Navigationsgrenzen und Schlagwörter für Beziehungen über diese Grenzen hinweg. Ein Foto aus Prag kann in \emph{Reisen / Tschechische Republik / Prag} liegen und Schlagwörter wie \emph{Architektur}, \emph{Nacht} und \emph{Straßenbahn} tragen.\footnote{Im normalen Inhaltsmodell besitzt ein Foto einen Galerieort, während wiederverwendbare Schlagwörter mehrere semantische Verbindungen schaffen. Kopier- und Verschiebewerkzeuge unterstützen bewusste Änderungen der physischen Organisation.}

\subsection{Tiefe der Hierarchie bestimmen}
\index{g00127@Galeriehierarchie!t00242@Tiefe}

Verschachtelte Galerien unterstützen tiefe Sammlungen; die meisten Besucher können jedoch leichter navigieren, wenn jede Ebene einen klaren Zweck hat. Zwei bis vier Ebenen eignen sich für viele Websites. Auch ein tieferes Archiv bleibt praktikabel, wenn Brotkrümelnavigation, aussagekräftige Titel und einheitliche Datumsangaben den Leser führen.

Erstellen Sie eine Untergalerie, wenn sie einen eigenen Titel, eine Beschreibung, ein Titelbild, eine Zugriffsrichtlinie, einen Datumsbereich oder eine öffentliche Adresse verdient. Lassen Sie Fotos zusammen, wenn eine neue Ebene nur einen zusätzlichen Klick erzeugen würde.

\subsection{Titel und Beschreibungen vorbereiten}
\index{g00134@Galerietitel}
\index{g00121@Galeriebeschreibung}

Ein Galerietitel sollte auch außerhalb der Administrationsoberfläche verständlich sein. Eine Beschreibung kann drei Fragen beantworten:

\begin{itemize}
  \item Was zeigt diese Sammlung?
  \item Wo und wann wurde sie erstellt?
  \item Welcher Zusammenhang verleiht den Fotos Bedeutung?
\end{itemize}

Beschreibungen können bei zwanglosen Alben kurz und bei dokumentarischer Arbeit ausführlich sein. Verwenden Sie den ersten Satz als wesentliche Zusammenfassung, weil Besucher häufig zunächst überfliegen, bevor sie den ganzen Text lesen.

\subsection{Privatsphäre vor dem Hochladen planen}
\index{d00069@Datenschutzplanung}
\index{g00133@Galeriesichtbarkeit}

Entscheiden Sie, ob eine Sammlung öffentlich, zur administrativen Vorbereitung nicht veröffentlicht oder für eine bekannte Zielgruppe geschützt sein soll. Die Richtlinie der übergeordneten Galerie kann sich auf Nachkommen auswirken; legen Sie daher Sammlungen mit ähnlichen Zielgruppen nach Möglichkeit zusammen ab.

Beispiele sind:

\begin{itemize}
  \item ein öffentliches Portfolio mit sorgfältig ausgewählten Fotos;
  \item ein passwortgeschützter Familienzweig;
  \item eine nicht veröffentlichte Kundengalerie zur Bildauswahl während der Vorbereitung;
  \item eine öffentliche Veranstaltungsgalerie, in der einzelne ausgewählte Fotos als privat markiert sind;
  \item ein Zweig mit Altersprüfung für Inhalte, die eine Besucherbestätigung erfordern.
\end{itemize}
\section{Galerien erstellen und bearbeiten}
\label{sec:create-edit-gallery}
\index{g00117@Galerie!e00093@Erstellung}
\index{g00117@Galerie!b00027@Bearbeitung}
\index{g00125@Galerieeditor}

\subsection{Eine Galerie in der Administration erstellen}

Verwenden Sie die Schaltfläche \ui{+} einer öffentlichen Galerie oder \ui{Galerie hier erstellen} in deren Editor, um eine Untergalerie anzulegen. Der rechte Administrations-Seitenbereich fragt nur nach dem neuen Titel; die übergeordnete Galerie, von der die Schaltfläche stammt, ist bereits bekannt. Die neue Galerie beginnt als \ui{Nicht veröffentlicht}. Nach der Erstellung öffnet derselbe Seitenbereich den vollständigen Editor für Beschreibung, Datum, Sprache, Schlagwörter, Zugriff, Anzeige, Medien und Fotos. Die öffentliche Seite bleibt während dieses erweiterten Ablaufs geöffnet.

Wenn Sie auf der eigenständigen Administrationsseite zur Erstellung beginnen, geben Sie einen Titel ein und ergänzen Sie vor dem Absenden optional SimBrief-Angaben, Beschreibung, Ausgangssprache und Schlagwörter. Der eingeklappte Abschnitt \ui{Weitere Optionen} enthält Ordnername, anfängliche Sichtbarkeit, Datumsbereich, übergeordnete Galerie, Abstimmung, Dateinamensanzeige und Kennzeichnung enthaltener Bilder. Eine kurze Zusammenfassung darunter zeigt die aktuelle Sichtbarkeit und die übergeordnete Galerie. Das spezielle Formular zum Erstellen und Hochladen bleibt verfügbar, wenn Fotos bereits während der Erstellung hochgeladen werden sollen. Bei einem Validierungsfehler bleiben übermittelte, nicht geheime Werte und der Kontext der übergeordneten Galerie zur Korrektur erhalten.

Sobald Sie mindestens zwei Zeichen in das Titelfeld eingeben, kann die Administration den restlichen Text eines passenden vorhandenen Galerietitels als blassen Vorschlag direkt im Feld anzeigen. Neuere Titel unter der ausgewählten übergeordneten Galerie werden bevorzugt; Treffer an anderen Stellen werden erst berücksichtigt, nachdem die begrenzte Suche unter Geschwistergalerien ausgeschöpft ist. Ältere Titel können fehlen; kein Vorschlag beweist daher nicht die Einzigartigkeit eines Titels. Das Formular fordert höchstens acht Kandidaten an, statt den vollständigen Titelkatalog zu laden.

Wenn der Cursor am Ende steht und kein Text ausgewählt ist, drücken Sie \codeword{Tab} oder \ui{Pfeil nach rechts} ohne Zusatztaste oder wählen Sie den vorgeschlagenen Text aus, um ihn zu übernehmen. Tastenkombinationen mit Zusatztasten und die Texteingabe über Eingabemethoden behalten ihr übliches Verhalten. Der erste Druck auf \codeword{Escape} verwirft die Vervollständigung; ein späterer Druck kann den umgebenden Seitenbereich erreichen. Übersetzte Hilfe und eine zurückhaltende Ansage machen Vorschläge auffindbar, ohne den Feldnamen zu ändern.

Der Unicode-Abgleich akzeptiert nur sichere Grenzen zwischen vollständigen Zeichen. Einige Vorschläge außerhalb des ASCII-Zeichensatzes entfallen, wenn Servererweiterungen oder die Segmentierungsunterstützung des Browsers fehlen. Eine langsame, fehlgeschlagene oder veraltete Abfrage verhindert niemals die normale Eingabe. Die Übernahme verändert nur das Titelfeld: Sie sendet das Formular nicht ab, benennt keine vorhandene Galerie um und verändert deren Zugriffsregeln nicht. Dasselbe Verhalten gilt nach dem Ersetzen eines Formulars im Seitenbereich. Abgleichgrenzen und technische Einzelheiten finden Sie in \filepath{docs/TITLE_COMPLETION.md}.

Nach der Erstellung mit nur einem Namen ist die Galerie bereits eine echte Sammlung im Galeriebaum, auch bevor weitere Editorfelder gespeichert wurden. Ihr physischer Ordner kann Fotos über Browser, Dateisystemübertragung, Erkennung oder einen anderen unterstützten Uploadablauf aufnehmen. Füllen und speichern Sie die Editorfelder vor der Veröffentlichung. Die Datenbank speichert beschreibende Einstellungen und Präsentationseinstellungen; die normale Zugriffsrichtlinie bestimmt, was Besucher sehen können.

Jedes erzeugte Erstellungsformular und klassische Uploadformular enthält außerdem eine private Vorgangskennung. Bricht eine Verbindung ab, nachdem der Server die Arbeit möglicherweise bereits beendet hat, versuchen Sie es mit demselben Formular erneut, statt ein neues zu öffnen: Für denselben Administrator, Vorgang und Inhalt erhalten Sie das ursprüngliche abgeschlossene Ergebnis. Ein frisch geöffnetes Formular steht für einen neuen beabsichtigten Vorgang und kann weiterhin ein bewusstes Duplikat erstellen. Diese Kennung ist kein Passwort und ersetzt weder Anmeldung noch CSRF-Schutz.

\subsubsection{Erstellungsablauf}

\begin{enumerate}
  \item Öffnen Sie die vorgesehene übergeordnete Galerie und wählen Sie deren Schaltfläche \ui{+}, oder beginnen Sie eine Stammgalerie in der Administration.
  \item Geben Sie einen besucherfreundlichen Titel ein; übernehmen Sie bei Bedarf einen Titelvorschlag.
  \item Erstellen Sie die nicht veröffentlichte Galerie. Bei verfügbarem JavaScript öffnet sich der vollständige Editor im selben Seitenbereich.
  \item Ergänzen Sie unter \ui{Identität} Beschreibung, Datum oder Datumsbereich, Ausgangssprache, Schlagwörter und optionale SimBrief-Daten.
  \item Prüfen Sie \ui{Zugriff}, \ui{Anzeige}, \ui{Medien} sowie gegebenenfalls die Werte für Galerieordner und übergeordnete Galerie unter den erweiterten Optionen.
  \item Speichern Sie die Galerie, fügen Sie Fotos hinzu und wählen Sie ein Titelbild, sobald geeignete Fotos vorhanden sind.
  \item Veröffentlichen Sie die fertige Galerie und prüfen Sie anschließend das Ergebnis anonym.
\end{enumerate}

\subsubsection{Gespeicherte Vorgaben und SimBrief-Entwürfe}
\index{g00126@Galerieerstellung!g00141@gespeicherte Vorgaben}
\index{s00224@SimBrief!e00088@Entwurf vor der Erstellung}

Der Editor kann die aktuelle SimBrief-Kennung und Ausgangssprache für den angemeldeten Administrator speichern. Wählen Sie die benachbarte Option \ui{Für zukünftige Galerien merken} nur bei einem Wert, den Sie wiederverwenden möchten. Bei späteren Galerien ist dieser Wert bereits ausgewählt; ein kleiner Hinweis kennzeichnet ihn als gespeicherte Vorgabe. Der Wert bleibt bearbeitbar; wenn Sie die Merken-Option nicht aktivieren, bleibt die bisher gespeicherte Auswahl erhalten. Vorgaben gehören zu einem einzelnen Administrator und veröffentlichen weder Zugangsdaten noch kopieren sie Zugriffsregeln. Die normale Migrationsaktion muss \codeword{202609250001\_gallery\_creation\_preferences.php} anwenden, bevor die Merken-Funktion eine neue Vorgabe speichern kann; die Galerieerstellung funktioniert auch ohne diese optionale Tabelle.

Das sichtbare SimBrief-Feld akzeptiert eine Pilot ID oder einen Pilotennamen. Ausschließlich aus Ziffern bestehende Eingaben gelten als Pilot ID, anderer Text als Name. \ui{Beschreibungsentwurf erzeugen} kann bei einer neuen oder bestehenden Galerie den neuesten OFP abrufen und einen bearbeitbaren Markdown-Text in der ausgewählten Inhaltssprache einfügen. Wenn die Galerielokalisierung aktiviert und ihr Übersetzungsspeicher bereit ist, füllt der Entwurf auch die Beschreibungsfelder der gepflegten Sprachen. Prüfen und bearbeiten Sie die Beschreibungen vor~dem~Speichern. Der Import hält für 30~Minuten eine private Referenz vor, die auf denselben Administrator und dieselbe Sitzung beschränkt ist; eine Änderung der Kennung macht diese Referenz ungültig. Bei einer bestehenden Galerie werden OFP und~verfügbare Routenkoordinaten erst nach dem Speichern der Galerie angehängt. Scheitert ein optionaler Anhang, bleibt die Galerie gespeichert und der Seitenbereich meldet eine Warnung. Importieren Sie erneut, wenn der Entwurf abläuft oder sich die Sitzung ändert.

\subsubsection{Im kompakten Editor arbeiten}
\index{g00125@Galerieeditor!s00214@Seitenbereich}

Die üblichen Arbeitsbereiche sind \ui{Identität}, \ui{API}, \ui{Zugriff}, \ui{Anzeige} und \ui{Medien}; anschließend folgen je nach Verfügbarkeit Bild- und Spezialwerkzeuge. Der rechte Seitenbereich hält \ui{Galerie speichern} unten sichtbar, während das Einstellungsformular gescrollt wird. Ein Registerwechsel speichert nicht automatisch. Verwenden Sie nach der Bearbeitung die Schaltfläche und öffnen Sie die Galerie anschließend erneut, wenn Sie die gespeicherten Werte prüfen möchten. Längere Erklärungen stehen hinter per Tastatur erreichbaren Fragezeichen-Schaltflächen, seltenere Optionen unter aufklappbaren \ui{Erweitert}-Bereichen. Ein direkter Administrationseditor verwendet dieselben Felder; ein Browser ohne JavaScript nutzt die normale Formularübermittlung.

\ui{Identität} hält Titel, Beschreibung, Datum, Ausgangssprache, SimBrief und Schlagwörter direkt erreichbar. Der erweiterte Abschnitt enthält Slug, Ordner, übergeordnete Galerie, Sortierreihenfolge und, sofern aktiviert, KI-Galerietext. \ui{API} zeigt einen kopierbaren Uploadendpunkt, Schlüsselbezeichnung und -aktion sowie einen neu erzeugten Schlüssel genau einmal; vorhandene Schlüssel erscheinen nur, wenn welche existieren. Die Neuerzeugung von KI-Metadaten, Migration und die globale API-Verwaltung liegen unter \ui{Erweiterte API-Werkzeuge}. \ui{Zugriff} bündelt Sichtbarkeit, Passwort, Freigabelink-Steuerungen und das NSFW-Kontrollkästchen; längere Richtlinienerklärungen stehen hinter Fragezeichen. \ui{Anzeige} beginnt mit dem Raster und den üblichen Schaltern für Abstimmung und Dateinamen; EXIF/GPS, Galeriekartenformat, Anzahlkennzeichnung, Lightbox-Modus, Picture Game, Vorschaubildgrenzen liegen unter \ui{Erweiterte Anzeigeeinstellungen}; Flugroutentext ist neben SimBrief unter \ui{Identität} verfügbar.

\subsection{Identität und Kontext bearbeiten}
\index{g00129@Galeriemetadaten}
\index{d00070@Datumsbereich}

Mit dem Editor können Sie eine Galerie im Laufe der Zeit verbessern. Titel und Beschreibungen können sich mit der Sammlung entwickeln. Datumsfelder können einen einzelnen Tag, eine Reise, eine Jahreszeit oder einen längeren historischen Zeitraum darstellen. EXIF-Datumsvorschläge können Fotos und untergeordnete Galerien untersuchen und anschließend einen bearbeitbaren Bereich zur Bestätigung anbieten.

Der vollständige Editor vermerkt die genaue Galerierevision, aus der seine Felder stammen. Ändert zuerst ein anderes Register, ein Administrator, Upload, Verschiebevorgang, Scan oder Wartungsablauf die Galerie, wird ein älterer Speicherversuch abgelehnt, statt den neueren Zustand still zu überschreiben. Die Konfliktansicht bewahrt eingegebenen, nicht geheimen Text und zeigt einen sicheren aktuellen Vergleich. Prüfen und kopieren Sie die gewünschten Änderungen in den aktuellen Editor; geben Sie Passwörter erneut ein und wählen Sie Dateien erneut aus, weil private Werte niemals in Konfliktdetails kopiert werden.

Verwenden Sie Datumsvorschläge als hilfreichen Ausgangspunkt. Gescannte Fotos, Bildschirmfotos, bearbeitete Exporte und Kameras mit falsch eingestellter Uhr können Datumsangaben liefern, die menschliche Prüfung erfordern.\footnote{Die Übernahme eines EXIF-basierten Vorschlags aktualisiert die vorgeschlagenen Galeriefelder über den bestehenden Editorablauf. Der Administrator bleibt für die endgültigen, Besuchern angezeigten Datumsangaben verantwortlich.}

\subsubsection{Datumsvorschläge galerieübergreifend prüfen}
\index{g00124@Galeriedaten!e00097@EXIF-Prüfung}

Die spezielle Wartungsseite \ui{Galeriedaten} wendet dieselben EXIF-Anhaltspunkte auf viele Galerien an. Sie kann die ganze Installation oder einen ausgewählten Galeriezweig prüfen. Jede Zeile zeigt den aktuellen manuellen Datumsbereich, den vorgeschlagenen Bereich, die Anzahl der EXIF-Fotos und Galerieknoten, die Anhaltspunkte beigetragen haben, sowie bearbeitbare Felder \ui{Von}/\ui{Bis}. Zeilen ohne vorhandenes manuelles Datum sind standardmäßig zur Übernahme ausgewählt; Zeilen mit manuellem Datum bleiben unmarkiert, damit eine Sammelprüfung die vom Betreiber eingetragene Chronologie nicht still überschreibt.

Nur markierte Zeilen werden gespeichert. Unmarkierte Vorschläge werden ignoriert, und der Administrator kann beide Datumswerte vor der Übernahme der Auswahl bearbeiten. Falls ältere importierte Fotos noch keine gescannten EXIF-Aufnahmedaten besitzen, führen Sie zuerst den normalen Bildscan beziehungsweise Import der Galerie aus. Dieselbe gezielte Vorschlagsaktion steht auch im einzelnen Galerieeditor zur Verfügung und kann diesen an Ort und Stelle aktualisieren.

Datumsfelder in der Administration verwenden weiterhin das native Datumsfeld des Browsers als übermittelten Wert, ergänzen aber einen kompakten Kalenderöffner sowie die Hilfen \ui{Heute} und \ui{Löschen}. \ui{Heute} trägt das aktuelle lokale Kalenderdatum des Browsers ein; \ui{Löschen} leert das Feld. Diese Erweiterung wird auch auf Formulare angewendet, die in den Administrations-Seitenbereich geladen werden; das zugrunde liegende native Datumsfeld bleibt ohne die zusätzlichen JavaScript-Steuerungen nutzbar.

\subsection{Referenz der Galerieeditor-Steuerungen}
\index{g00125@Galerieeditor!s00236@Steuerungen}
\index{g00117@Galerie!s00223@Sichtbarkeit}
\index{g00117@Galerie!a00020@Anzeigeüberschreibungen}

Der Galerieeditor ist der maßgebliche Ort für Einstellungen, die zu einer physischen Galerie gehören. Die üblichen Register sind Identität, API, Zugriff, Anzeige und Medien, gefolgt von Bilder und optionalen Organizer-/Umbenennungswerkzeugen, sofern verfügbar. Aufklappbare erweiterte Bereiche behalten selten genutzte Einstellungen bei, ohne sie aus dem gemeinsamen Speicherformular zu entfernen. Die folgende Tabelle fasst die aktuellen galeriebezogenen Steuerungen in Version~\version{} zusammen.

\begin{longtable}{>{\RaggedRight\arraybackslash}>{\hspace{0pt}}p{0.21\textwidth} >{\RaggedRight\arraybackslash}>{\hspace{0pt}}p{0.27\textwidth} >{\RaggedRight\arraybackslash}>{\hspace{0pt}}p{0.43\textwidth}}
\toprule
\textbf{Bereich} & \textbf{Steuerung} & \textbf{Aktuelles Verhalten}\\
\midrule
\endfirsthead
\toprule
\textbf{Bereich} & \textbf{Steuerung} & \textbf{Aktuelles Verhalten}\\
\midrule
\endhead
Identität & Titel und Beschreibung & Öffentlicher redaktioneller Text. Gepflegte Übersetzungen können separat gespeichert werden, ohne den Ausgangstext zu ändern.\\
Identität & Manuelles Datum und Datumsbereich & Optionales Ereignis- oder Aufnahmedatum und unterstützter Anfangs-/Endbereich. EXIF-basierte Vorschläge bleiben vor~dem~Speichern bearbeitbar.\\
Identität & Slug & Bestimmt die öffentliche URL-Kennung der Galerie. Halten Sie bestehende Slugs nach Möglichkeit stabil, weil Lesezeichen und externe Links davon abhängen können.\\
Identität & Ordnername & Benennt den physischen Galerieordner auf dem Datenträger über den normalen Änderungsablauf um. Es handelt sich nicht nur um eine Anzeigebezeichnung.\\
Identität & Übergeordnete Galerie und Sortierreihenfolge & Ändert Hierarchie und Reihenfolge unter Geschwistergalerien. Beim Verschieben einer Galerie wird ihr Ordnerbaum physisch verlagert, soweit das konfigurierte Galeriestammverzeichnis dies erlaubt.\\
Identität & Schlagwörter & Wiederverwendbare, durch Kommas getrennte Schlagwörter. Vorschläge werden anhand benachbarter Galerien, Fotos und des Ordnerkontexts gewichtet und bilden keine ungeordnete Zufallsliste.\\
Identität & SimBrief-Kennung und Vorgaben & Ein einziges sichtbares Feld akzeptiert eine numerische Pilot ID oder einen Pilotennamen. Der Administrator kann diese Kennung und die Ausgangssprache für zukünftige Galerien speichern; vorbelegte Werte werden gekennzeichnet.\\
API & Uploadendpunkt und Schlüssel & Kopieren Sie den galeriebezogenen Endpunkt, erzeugen Sie einen Schlüssel mit Bezeichnung und kopieren Sie den neuen Rohwert sofort. Aktive Schlüssel erscheinen nur, wenn vorhanden; der Rohwert wird nicht erneut angezeigt. KI-Neuerzeugung und Migration liegen unter den erweiterten Optionen.\\
Zugriff & Sichtbarkeit & \ui{Öffentlich} wird normal aufgelistet. \ui{Nicht veröffentlicht} bleibt in normalen Listen verborgen, kann aber über die normale URL geöffnet werden, wenn die Zugriffsregeln es erlauben. \ui{Privat} ist auf Administratoren beschränkt, außer bei unterstütztem direktem Tokenzugriff.\\
Zugriff & Passwortschutz & Unabhängig von der Sichtbarkeit. Eine öffentliche, nicht veröffentlichte oder private Galerie kann eine eigene Passwortrichtlinie besitzen, wenn das Zugriffsschema verfügbar ist. Ein leeres Feld für das neue Passwort erhält das bestehende Passwort.\\
Zugriff & Freigabelink & Kann erzeugt, neu erzeugt, widerrufen und optional mit einem Ablaufdatum versehen werden. Gespeichertes Tokenmaterial ist durch Hashes geschützt; ein älteres oder nicht entschlüsselbares Token muss möglicherweise neu erzeugt werden, bevor wieder ein kopierbarer Link angezeigt werden kann.\\
Zugriff & NSFW / 18+ & Kennzeichnet den Galeriezweig für die Altersbestätigung anonymer Besucher. Der Schutz gilt gemäß der wirksamen Zugriffsrichtlinie auch für Untergalerien und ausgelieferte Medien.\\
Anzeige & Picture Game & Erweiterte Option für das Vergleichsspiel in diesem Galeriezweig. Die Aktivierung des Spiels erfordert außerdem verfügbare Bildabstimmung.\\
Anzeige & Bildabstimmung & Aktiviert öffentliche Likes für Bilder dieser Galerie. Bei Deaktivierung bleiben gespeicherte Bewertungen erhalten; neue Abstimmungssteuerungen werden entfernt und neue Einsendungen blockiert.\\
Anzeige & Dateinamen anzeigen & Standardmäßig aus. Ursprünglich hochgeladene Dateinamen bleiben ohne Aktivierung verborgen; benutzerdefinierte Fototitel und Beschreibungen werden weiterhin angezeigt.\\
Anzeige & EXIF/GPS-Anzeige & \ui{Globale Vorgabe verwenden}, \ui{Einschalten erzwingen} oder \ui{Ausschalten erzwingen}. Der wirksame Zustand steuert öffentliche GPS-Markierungen, EXIF-Galeriekarten und die Anzeige von GPS-Koordinaten.\\
Anzeige & Breadcrumb-Stil & \ui{Standard / vom Design übernehmen} folgt der Designvorgabe; eine Galerie kann einen der neun registrierten Stile wählen. Ein ungültiger Override übernimmt die Designvorgabe, danach den integrierten Fallback \codeword{chevron}. Alle Vorgängerlinks umbrechen auf schmalen Bildschirmen und bleiben auch ohne JavaScript zugänglich.\\
Anzeige & Beschreibungsformat der Galeriekarten & Übernimmt standardmäßig die Designvorgabe und kann, soweit unterstützt, das Kartenlayout dieser Galerie überschreiben.\\
Anzeige & Kennzeichnung enthaltener Bilder & Übernimmt die Designvorgabe oder überschreibt das Symbol gestapelter Bilder und die auf Galeriekarten gezeigte Anzahl enthaltener Bilder.\\
Anzeige & Lightbox-Browsingmodus & Designvorgabe übernehmen, klassisches Einzelbild, Bildstreifen oder 3D-Karussell. Vollbild- und Diashow-Verhalten bleiben in allen Modi verfügbar.\\
Anzeige & Benutzerdefiniertes Raster & Erste sichtbare Anzeigesteuerung: optionale Spalten und Zeilen für diese Galerie. Ohne ausdrückliches Raster gilt das nächste übergeordnete Raster, das Vererbung an Untergalerien zulässt, danach die Designvorgabe.\\
Anzeige & Grenzen responsiver Vorschaubilder & Optionale Mindest-/Höchstgrenzen für Vorschaubildgrößen. Beim Speichern können die Galeriegrenzen rekursiv auf Nachkommen kopiert werden; rekursives Speichern ist standardmäßig aus.\\
Anzeige & Flugroute & Erweiterter manueller Routentext unterstützt die lokale Kennungssuche und ausdrückliche \codeword{NAME@latitude,longitude}-Punkte. SimBrief kann einen OFP, einen Beschreibungsentwurf und aufgelöste Routenkoordinaten speichern.\\
Medien & Titelbild & Automatisch oder ein ausgewähltes Foto, einschließlich geeigneter Fotos aus Nachkommen, wenn die Auswahl sie anbietet.\\
Medien & Hochgeladenes Galerievorschaubild & Separates Medium für die Galeriekarte, unabhängig von gewöhnlichen Galeriefotos gespeichert.\\
Medien & Hintergrundquelle & Designvorgabe übernehmen, einen neuen Hintergrund hochladen, ein vorhandenes Galeriebild wählen oder eine Collage aus öffentlichen Galerien erzeugen.\\
Medien & Branding-Ressourcen & Optionale galeriespezifische Banner-, Logo- und Trennlinienressourcen überschreiben konfigurierte Designersatzwerte nur für diese Galerie.\\
Bilder & Scan/Import & Durchsucht den physischen Galerieordner erneut und synchronisiert unterstützte Medien in den Katalog.\\
Bilder & Alle Vorschaubilder erstellen & Erstellt die Vorschaubilder dieser Galerie neu. Wählen Sie \ui{Untergalerien einschließen}, um ihre Nachkommen einzubeziehen; andere Galeriezweige bleiben außerhalb des Auftrags. Die Browsererzeugung ist standardmäßig ausgewählt, wenn verfügbar; es gibt eine Option zur Servererzeugung.\\
Werkzeuge & Organizer, Umbenenner, API & Datumsorganizer, physischer Medienumbenenner, Verwaltung von Upload-API-Schlüsseln, KI-Neuerzeugung und Galeriemigration erscheinen nur bei verfügbaren Schemata/Funktionen.\\
\bottomrule
\end{longtable}

\important{Galeriespezifische Steuerungen werden bewusst nicht in der globalen Einstellungszentrale zusammengeführt. Vererbung gehört zum Datenmodell; ein globales Suchergebnis kann daher zum Galerieeditor verlinken, ohne den galeriebezogenen Wert zu einer websiteweiten Einstellung zu machen.}

\subsection{Titelbild und visuelle Identität wählen}
\index{g00135@Galerietitelbild}
\index{g00123@Galeriebranding}

Ein Titelbild lädt visuell in die Galerie ein. Wählen Sie ein Bild, das als Vorschaubild erkennbar bleibt und die Sammlung angemessen vertritt. Gesichter, starke Silhouetten, klare Motive und übersichtliche Kompositionen funktionieren bei kleinen Größen häufig gut.

Je nach verfügbaren Steuerungen kann Galeriebranding Logo, Banner, Hintergrund, Trennlinie oder Präsentationsüberschreibungen bereitstellen. Beginnen Sie mit dem geerbten Designstil und ergänzen Sie galeriespezifische Ressourcen dort, wo eine Sammlung eine eigene Identität besitzt, etwa eine Hochzeit, Ausstellung, ein Kunde oder ein langfristiges Projekt.

\subsection{Galerien verschieben und neu ordnen}
\index{g00117@Galerie!v00255@Verschieben}
\index{g00117@Galerie!n00178@Neuordnen}

Das Verschieben einer Galerie verändert ihre Position in der Hierarchie und kann geerbten Zugriff oder Präsentation beeinflussen. Prüfen Sie das Ziel, bevor Sie den Vorgang bestätigen. Neuordnen verändert die Reihenfolge unter Geschwistergalerien und ermöglicht eine kuratierte Betrachtungsreihenfolge.

Beim Verschieben von Fotos zwischen Galerien wird ein dauerhaftes Vorgangsprotokoll verwendet. Der Organizer verschiebt jedes Original physisch und prüft seine aufgezeichnete Identität, bevor er die Katalogzuordnung ändert. Hält ein Arbeitsprozess oder eine Verbindung mitten im Ablauf an, zeigen Systemzustand und Laufzeitdiagnose einen begrenzten Aktionszustand, ohne private Pfade offenzulegen. Ein fehlgeschlagener Verschiebevorgang zeigt einen sicheren Grund auf dem Bildschirm; vertrauenswürdige Administrationsprotokolle bewahren ausführlichere private Diagnosen. Stoppen Sie konkurrierende Speicherarbeiten, bewahren Sie eine Sicherung auf und folgen Sie \filepath{docs/IMAGE_MOVE_RECOVERY.md}; die Wiederherstellung behält geänderte oder nicht überprüfbare Dateien, statt sie zu überschreiben.

Die Reihenfolge kann chronologisch, geografisch, erzählerisch, alphabetisch oder manuell hervorgehoben sein. Halten Sie die Entscheidung innerhalb einer übergeordneten Galerie ausreichend einheitlich, damit Besucher den nächsten Eintrag erwarten können.

Ein angemeldeter Administrator kann auch die direkten Untergaleriekarten auf der öffentlichen Galerieseite neu ordnen. Die kompakte öffentliche Sortierleiste ist auf die Karten der aktuellen Seite der Seitennavigation beschränkt. Die Speicheranfrage enthält den Versatz dieser Seite und ihre sichtbare Anzahl; der Server lehnt die Änderung ab, wenn die übermittelten IDs nicht mehr zum selben aktuellen Ausschnitt passen. Sie ändert nur die \codeword{sort_order}-Werte der Geschwistergalerien; übergeordnete Galerien werden nicht geändert, und Galerien werden nicht verschachtelt.

Betrachtet ein angemeldeter Administrator eine öffentliche Galerie, können direkte Untergalerien mit ausgefülltem \ui{Von}-Datum auch vorläufig vom ältesten zum neuesten oder vom neuesten zum ältesten Datum sortiert werden. Untergalerien ohne brauchbares Datum behalten ihre Positionen. Die Vorschau ist vorübergehend, bis \ui{Speichern} gewählt wird; das Speichern schreibt das Ergebnis als tatsächliche Geschwisterreihenfolge für alle Besucher. Mindestens zwei datierte direkte Untergalerien sind erforderlich, bevor die Datumsreihenfolge gespeichert werden kann.

\subsection{Eine Galerie bedacht löschen}
\index{g00117@Galerie!l00166@Löschen}

Das Löschen einer Galerie kann Originaldateien, Nachkommen, Metadaten, Vorschaubilder, Links und Besucherlesezeichen betreffen. Prüfen Sie die genau ausgewählte Galerie und ihre Untergalerien, erstellen Sie eine Sicherung und verwenden Sie den normalen administrativen Löschablauf. Eine vorübergehende Sichtbarkeitsänderung kann eine sicherere Vorbereitungsphase bieten, wenn die Galerie später noch benötigt werden könnte.

Der Galeriepapierkorb ist standardmäßig aktiviert. Das Löschen über eine normale administrative Galerie-Löschaktion verschiebt den ausgewählten Ordner und seinen vollständigen untergeordneten Baum aus dem aktiven Galeriestammverzeichnis, zeichnet die zur Wiederherstellung benötigten Metadaten auf und führt den Eintrag unter \ui{Wartung > Papierkorb} auf. Öffentliche Seiten erkennen diesen Teilbaum sofort nicht mehr; ein Administrator kann ihn jedoch am ursprünglichen Pfad wiederherstellen, solange dieser frei bleibt. Die Wiederherstellung überschreibt niemals eine neue aktive Galerie am selben Speicherort.

Wählen Sie im Papierkorb-Seitenbereich \ui{Wiederherstellen}, \ui{Endgültig löschen} oder \ui{Papierkorb leeren}. Endgültiges Löschen und das Leeren des Papierkorbs sind nicht rückgängig zu machen. Das Leeren arbeitet in begrenzten Chargen, damit große Papierkorbsammlungen keine unbegrenzte einzelne Anfrage erfordern. Das Löschen einzelner Fotos gehört nicht zu dieser Wiederherstellungsfunktion auf Galerieebene und bleibt endgültig. Eine Sicherung bleibt als Schutz gegen Speicherausfall, administrative endgültige Löschung oder Schäden außerhalb von PHP Gallery notwendig.
\section{Smart Galleries}
\index{s00226@Smart Galleries}
\index{v00260@virtuelle Galerien}

Eine Smart Gallery ist eine gespeicherte Abfrage, kein Ordner. Öffnen Sie \ui{Administration > Galerien > Smart Galleries}, fügen Sie Bedingungen hinzu und kombinieren Sie diese mit verschachtelten Gruppen \ui{AND}, \ui{OR} und \ui{NOT}. Die Felder umfassen Quellgalerien und deren Nachkommen, Schlagwörter, Aufnahmedaten, EXIF/GPS, Text, KI-Metadaten, Duplikatstatus, Dateieigenschaften und private redaktionelle Bewertungen.

Mit \ui{Vorschau} sehen Sie die aktuelle Trefferzahl und die tatsächlich passenden Fotokarten; speichern Sie anschließend und veröffentlichen Sie die Galerie bei Bedarf. Bilder und Dateien bleiben in ihren physischen Galerien, sodass Metadaten- oder Bewertungsänderungen die Zugehörigkeit sofort verändern. Die redaktionelle Bewertung ist privat und unabhängig von Besucherabstimmungen. Administratoren können außerdem in der öffentlichen Suche \ui{Suche als Smart Gallery speichern} verwenden.

Wählen Sie \ui{Nicht gelistet} für Zugriff nur über die URL, \ui{Stammgalerie}, um die Smart Gallery auf der Startseite zu zeigen, oder \ui{Platzierung als Untergalerie}, damit physische Galerien sie einbinden können. Wählen Sie in jedem physischen Galerieeditor, ob eine Einbindung \ui{Über dem Galerieinhalt} oder \ui{Unter dem Galerieinhalt} erscheinen soll, und geben Sie einen Reihenfolgewert an. Einbindungen oberhalb erscheinen zuerst, danach gewöhnliche Untergalerien und Fotos und schließlich Einbindungen unterhalb. Gleiche Reihenfolgewerte werden einheitlich anhand der Smart-Gallery-ID aufgelöst. Vorhandene Einbindungen stehen standardmäßig unterhalb mit Reihenfolge~0.

Dieselbe virtuelle Galerie kann unter mehreren physischen Galerien eingebunden werden, jeweils mit eigener Platzierung und Reihenfolge. Der Abschnitt \ui{In physischen Galerien verwendet} des Smart-Gallery-Editors zeigt alle übergeordneten Galerien und erlaubt, deren Platzierung/Reihenfolge zu speichern oder die Einbindung mit \ui{Trennen} zu lösen, ohne andere übergeordnete Galerien zu verändern. Wird der Editor im Administrations-Seitenbereich geöffnet, bleiben diese Aktionen dort; bei deaktiviertem JavaScript steht die gewöhnliche Formularübermittlung zur Verfügung. Die Karte einer öffentlichen Smart Gallery unter einer privaten oder anderweitig eingeschränkten übergeordneten Galerie erbt weiterhin deren öffentliche Zugriffsgrenze.

Öffentliche Ergebnisse und Zählungen wenden stets dieselbe Quellzugriffsrichtlinie an wie gewöhnliche öffentliche Galerieinhalte. Ein passendes Foto zählt nur dann, wenn seine physische Galerie öffentlich gelistet und für den aktuellen Besucher zugänglich ist und das Foto selbst öffentlich ist. Bestehende Regeln für Passwort, Freigabelink, geerbte Sichtbarkeit, nicht veröffentlichte Galerien und Altersbeschränkung gelten daher weiterhin. Routen mit Abfrageparametern funktionieren ohne Umschreibung; sprechende \texttt{/smart/<slug>}-URLs sind optional.

\subsection{Sichere Beziehungen und Reparatur}
\index{s00226@Smart Galleries!z00282@Zyklenschutz}

Eine Smart Gallery darf nicht zu einer physischen Galerie zurückführen, die sie einbindet. Der Server prüft bei jeder Änderung von Regeln, Einbindungen oder Galeriehierarchie den vollständigen Beziehungsgraphen; Browserfilterung gilt nicht als Sicherheitsmaßnahme. Das umfasst indirekte Wege über physische Nachkommen und andere eingebundene Smart Galleries. Die aus dem Dateisystem abgeleitete Hierarchiesynchronisierung und die Reparatur übergeordneter Galerien anhand öffentlicher Pfade prüfen ihre vollständige vorgeschlagene Zuordnung, bevor eine neue Elternverknüpfung geschrieben wird. Der aktuelle Regeleditor kann physische Galerien referenzieren, bietet aber keine direkte Bedingung zur Referenzierung einer Smart Gallery.

Ein vollständiger administrativer Drag-and-drop-Baum wird einmal in seiner endgültigen Form geprüft. Die anschließenden Ordnerverschiebungen verwenden diese vorab geprüfte Zuordnung und verschieben die Hierarchiesynchronisierung bis zum Vorliegen des endgültigen Baums, sodass ein vorübergehender Zwischenzustand keine ansonsten gültige Stapeländerung ablehnen kann. Jede festgeschriebene Änderung der physischen Hierarchie leert den anfragelokalen Smart-Gallery-Graphcache, bevor spätere graphabhängige Arbeit fortgesetzt wird.

Enthalten ältere oder manuell veränderte Datenbankdaten bereits einen Zyklus, beendet die öffentliche Darstellung diese ungültige Beziehung, statt rekursiv fortzufahren. Die administrative Einbindungsliste kennzeichnet die Beziehung zur Reparatur; \ui{Trennen} bleibt verfügbar. Validierung und Laufzeitdurchlauf verwenden konservative Tiefen-/Knotengrenzen und anfragelokale Deduplizierung, damit fehlerhafte Daten dieselbe Smart Gallery nicht wiederholt erweitern können. Die normale Ergebnisabfrage bleibt seitenweise und kopiert nicht sämtliche passenden Fotos nach PHP, nur um Zyklen zu erkennen.

Migration \texttt{202608170002\_smart\_gallery\_attachment\_ordering.php} ergänzt die Felder für Platzierung und Reihenfolge je übergeordneter Galerie. Solange diese Migration nicht verfügbar ist, behalten vorhandene Einbindungen ihr bisheriges Verhalten unterhalb des Inhalts; administrative Platzierungsänderungen werden abgelehnt, statt teilweise gespeichert zu werden.

\subsection{Präsentationseinstellungen}
\index{s00226@Smart Galleries!p00192@Präsentation}

Standardmäßig folgt eine Smart Gallery dem aktuellen Design und der Websitepräsentation. Aktivieren Sie \ui{Designvorgaben für diese Smart Gallery überschreiben}, wenn diese virtuelle Sammlung eigene Spalten, Zeilen, Seitennavigation, Vorschaubildgrößengrenzen, Vorschaubildrenderer, Galeriekartenlayout, Metadatenüberlagerungen, Lightbox-Verhalten, Diashow-, Download-, Abstimmungs- oder Breadcrumb-Stile benötigt. Der Breadcrumb-Selektor steht bei den Präsentationseinstellungen; \ui{Standard / vom Design übernehmen} folgt der Designvorgabe. Sortierung und Richtung bleiben Teil der Smart-Gallery-Definition.

Eine Smart Gallery erbt keine Präsentation von der physischen Galerie, in der ihre Karte erscheint. Eine Smart Gallery kann unter mehreren physischen Galerien erscheinen; daher gibt es keine einzelne übergeordnete Präsentation mit eindeutigem Vorrang. Fehlt eine Überschreibung oder sind gespeicherte Präsentationsdaten nicht vorhanden oder ungültig, bleiben die aktuellen Design-/Websitevorgaben der Ersatzwert. Vorschaubildgrenzen physischer Galerien und Fotos haben weiterhin Vorrang, wenn eine Smart-Gallery-Grenze ihnen widerspricht.

\subsection{Lightbox über Ergebnisseiten hinweg}
\index{s00226@Smart Galleries!l00164@Lightbox}

Der große Betrachter ist nicht auf die auf der aktuellen HTML-Seite dargestellten Fotos beschränkt. Smart-Gallery-Karten merken sich die Position jedes Fotos in der vollständigen autorisierten Ergebnisreihenfolge. Wenn der Besucher vor- oder zurückwechselt, fordert der Browser benachbarte Metadaten in kleinen autorisierten Gruppen vom Server an. Der Server wiederholt dieselben Smart-Gallery-Regeln, Quellzugriffsprüfungen, Sortierung und stabile Auflösung gleicher Sortierwerte wie die Seite.

So bleiben große virtuelle Sammlungen praktikabel: Das Öffnen einer Smart Gallery lädt nicht alle Originale oder sämtliche Ergebnismetadaten vor. Bestehendes Verhalten für Tastatur, Berührung, Vollbild, Zoom, Schließen/erneutes Öffnen und Diashow wird mit gewöhnlichen Galerien geteilt. Ist die Diashow für diese Smart Gallery deaktiviert, entfallen ihre Diashow-Steuerungen, während normale Galerien ihre üblichen Vorgaben behalten.

\subsection{Downloads und Sicherheitsgrenzen}
\index{s00226@Smart Galleries!d00080@Downloads}

Sind sowohl Website-Downloads als auch die Downloadoption der Smart Gallery aktiviert, kann der Server eine ZIP-Datei aus den aktuell autorisierten passenden Originalen vorbereiten. Er wertet gespeicherte Regeln und Zugriffsrichtlinie erneut aus, statt Dateisystempfade gegenüber dem Browser offenzulegen. Der Archivablauf lehnt Sammlungen mit mehr als 5.000 passenden Bildern ab und erzwingt außerdem eine Gesamtobergrenze für die Bytes der Originaldateien. Bestehende Installationen verwenden eine Grenze von 2~GiB, sofern \texttt{smart\_gallery\_zip\_max\_source\_bytes} nicht in \texttt{config.php} gesetzt ist. Gleichzeitige Anfragen für dieselbe Archivsignatur werden durch eine Dateisperre serialisiert; nach Erhalt der Sperre wird der Cache erneut geprüft, eine eindeutige vorläufige ZIP-Datei geschrieben und das vollständige Archiv durch atomare Umbenennung bereitgestellt. Fehlerprotokolle enthalten nur begrenzte Gründe-/Klassendiagnosen statt unverarbeiteter Ausnahmemeldungen oder privater Serverpfade.

Regeln sind lesbares versioniertes JSON, niemals vom Benutzer geliefertes SQL: Felder/Operatoren stehen auf einer Positivliste, Werte sind vorbereitete Parameter, die Tiefe ist auf fünf begrenzt und höchstens 50 Bedingungen werden akzeptiert. Gelöschte Referenzen treffen nichts; fehlerhafte oder unbekannte Regelversionen werden abgelehnt. Der normale Migrationsablauf ergänzt den Speicher für die Smart-Gallery-Präsentation; vorhandene Smart Galleries übernehmen Design-/Websitevorgaben, bis eine Überschreibung gespeichert wird.

\section{Kooperative Galerien}
\index{kooperative Galerien}
Kooperative Galerien verbinden unabhängige, öffentlich erreichbare HTTPS-Installationen. Jeder Betreiber behält seine Alben und Originale. Wenden Sie ausstehende Migrationen über den normalen Verwaltungsablauf an und aktivieren Sie kooperative Galerien unter Funktionen auf jedem Server. Die Funktion ist standardmäßig ausgeschaltet; Ausschalten bewahrt Identitäten, Freundschaften, Vorschläge und Einstellungen.

Koppeln Sie unter Einstellungen > Erweitert > Befreundete Galerien jede Teilnehmerpaarung direkt. Freundschaft allein teilt keine Fotos. Wählen Sie bei Albumzusammenarbeit ein öffentliches, gelistetes Album ohne Passwort und NSFW und erstellen Sie seinen Referenzcode. Ein Teilnehmer kombiniert die Referenzen zu einem unveränderlichen Vorschlag. Jeder Administrator prüft Mitglieder und Berechtigungen und stimmt auf seiner eigenen Installation ausdrücklich zu. Zustellung und Prüfung laufen in begrenzten Schritten im Seitenpanel. E-Mail-Einladungen verwenden den konfigurierten Mailtransport; Öffnen eines Links ist keine Zustimmung.

Nach Aktivierung öffnen Sie gemeinsame Fotos im Panel oder über den Link des lokalen Albums. Quellen behalten Zuordnung und Seitennavigation. Nur vorhandene autorisierte JPEG/WebP-Vorschaubilder werden geteilt; Originale, geschützte Alben und entfernte Schreibzugriffe bleiben ausgeschlossen. Links funktionieren ohne JavaScript. Fehlende Vorschaubilder benötigen die normale Erzeugung. Neue Fotos erhalten vollständige Pfadidentitäten einschließlich Erweiterung. Alte Namen bleiben ohne Massenerneuerung erhalten, mehrdeutige Eigentümerschaft wird jedoch abgelehnt.

Rücknahme, geänderter Quellschutz und lokale Freundschaftswiderrufe sperren nachfolgende autorisierte Zugriffe. Unerreichbare Teilnehmer bleiben durch höchstens 120 Sekunden Prüfungsdauer begrenzt; bereits heruntergeladene Pixel lassen sich nicht zurückrufen. Neue Teilnehmer benötigen einen separaten Vorschlag mit neuer Zustimmung aller. Technische Erneuerung kann bei Nutzung erfolgen; ein Zeitplan ist optional. Fehlendes oder unbekanntes erforderliches Schema verweigert Freigaben und Änderungen. Betriebsgrenzen beschreibt \filepath{docs/COOPERATIVE_GALLERIES.md}.

\subsection{Simulatorposition unter Windows}
WinApp 0.3.2 verwendet automatisch die Kameraweltposition in MSFS 2024 und fällt bei begrenztem Kamerafehler auf die Flugzeugposition zurück. MSFS 2020 verwendet direkt die Flugzeugposition über dieselbe SimConnect-Laufzeit. Koordinaten sind optional; ein fehlender Simulator verhindert keine Uploads. Die Diagnose unterscheidet die Positionsquelle. Das Installationsprogramm fordert das Beenden der Anwendung vor dem Dateiaustausch an. Schließen Sie laufende Aufgaben vorher ab. Prüfen Sie echte Simulatoren und das Upgrade auf dem Ziel-PC.

Der gebündelte Aktualisierer lädt Windows-APIs ausdrücklich aus System32 und vermeidet damit eine Kollision der VERSION-Textdatei mit der Windows-Versionsbibliothek. Starten Sie beim Upgrade von 0.3.1 den geprüften Installer 0.3.2 manuell. Die laufende ältere Anwendung kopiert ihren eigenen Hilfsprozess vor der Installation; dieser alte Helfer kann den bisherigen DLL-Fehler auch nach dem Ersetzen der Anwendung melden. Künftige Updates aus der korrigierten Anwendung 0.3.2 verwenden den reparierten Helfer; prüfen Sie die installierte Übergabe unter Windows.

Flugzeugantworten werden anhand der Anfrage- und Definitionskennungen zugeordnet; der USER-Selektor der Anfrage verlangt keine Objektkennung null in der Antwort. Begrenzte Diagnosen erfassen Empfangsmetadaten, Verarbeitungsreihenfolge und konkrete Gründe für verworfene Pakete. Fehlergründe bleiben nach der Pfadredaktion lesbar, und Aktivitätsmeldungen zeigen nur eine SimConnect-Markierung. Der Objektkennungsfilter war ein bestätigter Codefehler, doch der gemeldete FS2020-Timeout enthielt die zurückgegebene Objektkennung nicht. Prüfen Sie die Korrektur auf dem betroffenen PC; simulierte Tests belegen keinen tatsächlichen Simulatorerfolg.

Der installierte Uploader prüft seine eigene stabile GitHub-Installerversion höchstens einmal pro Stunde; \ui{Check for updates} startet eine manuelle Prüfung. Der letzte automatische Versuch wird unabhängig für den aktuellen Benutzer gespeichert und bleibt nach einem Anwendungsneustart erhalten. CMS-Tags bestimmen nicht seine Version. Die höchste numerische Version benötigt SHA256 vom GitHub-Asset oder aus einer passenden \filepath{winapp-update.json} desselben Releases. Fehlende oder widersprüchliche Daten verhindern das Update; identische Installer in mehreren CMS-Releases werden akzeptiert. Galerie-Zugangsdaten werden nie an GitHub gesendet.

\ui{Update} lädt den Installer mit Fortschritt und Abbruchmöglichkeit herunter, prüft ihn, sperrt neue Arbeit und wartet bis zu 120 Sekunden auf den sicheren Abschluss der Hintergrundaufgaben. Misslingt dieser Abschluss, wird die Installation verweigert und der Uploader bleibt verfügbar. Ein externer Hilfsprozess prüft die Datei erneut und fordert Windows-UAC-Zustimmung an. Das Selbstupdate beendet keine Arbeiter gewaltsam; normal manuell gestartete Installer behalten ihre bisherige Beendigungsregel. Der tatsächliche Installationsfortschritt kommt über einen lokalen sitzungsspezifischen Kanal. Nach Erfolg startet der Uploader als ursprünglicher Benutzer mit gespeicherten Einstellungen und Aufgaben. Es gibt kein automatisches Rollback. Native Python- oder Quellcode-Ausführungen bieten den geprüften Installer zum manuellen Start an. Wenn Windows einen Neustart verlangt, starten Sie Windows vor dem erneuten Öffnen des Uploaders neu; ein UAC-Abbruch öffnet die vorhandene Installation erneut.


\section{Fotos hinzufügen, beschreiben und organisieren}
\label{sec:photo-workflow}
\index{f00106@Fotoarbeitsablauf}
\index{b00044@Bildmetadaten}
\index{f00108@Fotoorganisation}

\subsection{Uploadmethode wählen}
\index{u00248@Uploadmethoden}

Der Browserupload eignet sich für alltägliche Ergänzungen. Wählen Sie eine Galerie und Dateien aus, prüfen Sie die Charge und verfolgen Sie den Fortschritt, während der Server Medien validiert und Datensätze und Vorschaubilder vorbereitet. Die Dateisystemerkennung eignet sich für große vorhandene Ordnerbäume oder Übertragungen per FTP und Hostingwerkzeugen. Automatisierung und WebDAV eignen sich nach Einrichtung ihrer bereichsgebundenen Zugangsdaten für wiederholte oder mobile Abläufe.

Im rechten Administrationsbereich \ui{Fotos hochladen} lassen sich nacheinander Bildschirmfotos mit Strg+V unter Windows/Linux oder Cmd+V unter macOS einfügen, weitere Dateien auswählen und Bilder in den Uploadbereich ziehen. Jedes Bild erscheint vor dem Absenden in einer lokalen Miniaturvorschau-Liste mit Dateiname und Größe. Über die Pfeiltasten der Liste lässt sich die Reihenfolge ändern; falsche Bilder können entfernt oder die gesamte Auswahl geleert werden. Die vorhandene Schaltfläche \ui{Bilder hochladen} übermittelt die Originaldateien einmal über den normalen Uploadablauf. Miniaturbilder ersetzen nicht die Originale. Die Liste ist kein auf dem Server gespeicherter Entwurf: Beim Verlassen des Bereichs muss das Verwerfen ungesendeter Dateien bestätigt werden; nach einem Fehler bleiben sie für einen erneuten Versuch erhalten, nach bestätigtem Erfolg wird die Liste geleert. Nicht darstellbare Formate behalten ihr Dateisymbol und durchlaufen weiterhin die serverseitige Validierung. Die eigenständige Admin-Uploadseite und das Formular ohne JavaScript bleiben unverändert.

Laden Sie bei einer großen Veranstaltung in sinnvollen Gruppen hoch. Kleinere Gruppen erleichtern das Erkennen von Fehlern und ermöglichen eine schrittweise Prüfung von Titeln, Reihenfolge und Sichtbarkeit.

\subsection{Nützliche Fototitel schreiben}
\index{f00110@Fototitel}
\index{b00046@Bildunterschriften}
\index{a00011@Alternativtext}

Ein Titel identifiziert das Foto rasch. Eine Beschreibung oder Bildunterschrift kann Menschen, Ort, Tätigkeit, technischen Kontext oder eine Geschichte außerhalb des Bildausschnitts erklären. Nützliche Beispiele sind:

\begin{itemize}
  \item \emph{Karlsbrücke vor Sonnenaufgang};
  \item \emph{Ankunft des ersten Zuges am restaurierten Bahnhof};
  \item \emph{Familiengruppe im Garten, Juni 1998};
  \item \emph{Endanflug während der Vorführung am Samstag}.
\end{itemize}

Aussagekräftiger Text verbessert Suche, Barrierefreiheit, spätere Identifizierung und das Erlebnis von Lesern, die mehr als ein visuelles Raster wünschen. Dateinamen können betriebliche Details bleiben, wenn öffentliche Anzeigeeinstellungen kuratierte Titel bevorzugen.

\subsection{Schlagwörter als zweite Organisationsebene verwenden}
\index{s00212@Schlagwörter!p00191@praktische Verwendung}
\index{v00256@Verschlagwortungsstrategie}

Erstellen Sie einen kleinen Wortschatz, bevor Sie viele nahezu identische Schlagwörter hinzufügen. Bevorzugen Sie einen stabilen Begriff wie \emph{Flugzeug} gegenüber mehreren unbeabsichtigten Varianten. Schlagwörter können Motive, Personen, Orte, Techniken, Ausrüstung, Jahreszeiten oder redaktionellen Status darstellen.

Bei einem Reisearchiv können Galerien Reisen und Schlagwörter wiederkehrende Themen repräsentieren. Bei einem Familienarchiv können Galerien Jahre und Ereignisse repräsentieren, während Schlagwörter Personen identifizieren. Bei einem professionellen Portfolio können Galerien Kunden und Schlagwörter die Art des Ergebnisses, den Ort und die Spezialisierung kennzeichnen.

Schlagwortfelder können beim Tippen vorhandene Schlagwörter vorschlagen. Die Wiederverwendung eines Vorschlags vermeidet unbeabsichtigte Schreibvarianten und hält Schlagwort-Zielseiten einheitlich. Der Vorschlag ist nur eine Bearbeitungshilfe; weiterhin entscheidet die normale Speicheraktion, welche Schlagwörter für Galerie oder Foto gespeichert werden.

\subsection{Betrachtungsreihenfolge gestalten}
\index{f00109@Fotoreihenfolge}
\index{e00094@Erzählen}

Die Reihenfolge macht aus einer Sammlung eine Abfolge. Beginnen Sie mit einem einladenden Bild, führen Sie Ort oder Kontext ein, variieren Sie weite und nahe Ansichten, halten Sie eng zusammengehörige Details beieinander und enden Sie mit einem natürlichen Abschluss. Chronologische Reihenfolge eignet sich für Veranstaltungen; ein visueller Rhythmus kann zu Portfolios und Ausstellungen passen.

Die öffentlichen Sortiersteuerungen helfen Administratoren, sichtbare Karten anzupassen. Auch die Fotosortierung auf einer öffentlichen Seite ist auf den aktuell sichtbaren Ausschnitt der Seitennavigation begrenzt und speichert über den authentifizierten Bildreihenfolge-Endpunkt. Wenn Picture Manager aktiv ist, verändert das Ziehen ausgewählter Fotokarten auf eine sichtbare Untergaleriekarte die Geste vom Neuordnen zu einem validierten Verschieben in diese Untergalerie. Die Seitennavigation kann eine große Sammlung unterteilen; prüfen Sie daher neben der ersten Seite auch die Übergänge an Seitengrenzen.

\subsection{EXIF- und Ortsinformationen prüfen}
\index{e00096@EXIF!b00030@Benutzerprüfung}
\index{s00234@Standortdatenschutz}

EXIF kann ein Foto um Aufnahmedatum, Kamera, Objektiv, Belichtung und GPS-Koordinaten ergänzen. Diese Informationen unterstützen Karten, Datumsvorschläge, Duplikathinweise und technischen Kontext. Prüfen Sie die GPS-Sichtbarkeit, bevor Sie Fotos veröffentlichen, die an Wohnhäusern, Schulen, privaten Arbeitsplätzen oder empfindlichen Naturorten aufgenommen wurden.

\section{Veröffentlichen, Teilen und Zielgruppenszenarien}
\label{sec:audiences}
\index{v00252@Veröffentlichungsablauf}
\index{z00277@Zielgruppen}
\index{a00019@Anwendungsfälle}

\subsection{Öffentliches Portfolio}
\index{p00190@Portfolio}

Ein öffentliches Portfolio profitiert von wenigen starken Galerien auf oberster Ebene, einheitlichen Titelbildern, knappen Beschreibungen und sorgfältiger Reihenfolge. Verwenden Sie Schlagwörter für projektübergreifende Spezialgebiete. Halten Sie Quellarchive in privaten oder nicht veröffentlichten Zweigen und kopieren Sie ausgewählte Arbeiten in öffentliche Präsentationsgalerien, wenn der Ablauf eine getrennte Kuratierung verlangt.

Testen Sie das Portfolio auf einem Telefon, einem hochauflösenden Display und einer langsameren Verbindung. Progressive Schärfung ist Standard und sicherer Ersatz für Fotokarten in der geöffneten Galerie; responsive Auswahl bleibt eine unterstützte Alternative.

\subsection{Privates Familienarchiv}
\index{f00098@Familienarchiv}

Ordnen Sie nach Jahr, Familienzweig oder Ereignis. Verwenden Sie passwortgeschützte übergeordnete Galerien, um einen verständlichen privaten Bereich zu schaffen, und legen Sie zusammengehörige Nachkommen darin ab. Ergänzen Sie Namen und Datumsangaben, solange Erinnerungen vorhanden sind. Gescannte Fotos profitieren besonders von menschlichen Beschreibungen, weil eingebettete Kamerametadaten fehlen können.

Teilen Sie geschützte Adresse und Passwort über einen geeigneten privaten Kanal. Erklären Sie, dass Empfänger Untergalerien ohne Administratorkonto durchsehen können. Prüfen Sie regelmäßig, ob alte Links und Passwörter noch der vorgesehenen Zielgruppe dienen.

\subsection{Übergabe an Kunden}
\index{k00159@Kundengalerie}
\index{b00042@Bildauswahlablauf}

Erstellen Sie pro Kunde oder Auftrag eine geschützte Galerie. Laden Sie Auswahlbilder hoch, ordnen Sie sie in nützlichen Gruppen und erläutern Sie in Beschreibungen Erwartungen an Auswahl oder Download. Abstimmungen können eine einfache Sammlung von Präferenzen unterstützen, wenn sie für diese Galerie aktiviert sind. Eine endgültige Übergabegalerie kann freigegebene Dateien in der vorgesehenen Reihenfolge enthalten.

Freigabelinks bieten, sofern eingerichtet, einen weiteren kontrollierten Zugangsweg. Testen Sie den genauen Kundenablauf in einem privaten Browserfenster, bevor Sie die Adresse versenden.

\subsection{Veranstaltungs- und Gemeinschaftsgalerie}
\index{v00251@Veranstaltungsgalerie}

Eine Veranstaltungsgalerie kann zunächst nicht veröffentlicht bleiben, während Fotos eintreffen. Verwenden Sie Untergalerien für Tage, Bühnen, Teams, Orte oder Aktivitäten. Veröffentlichen Sie die übergeordnete Galerie, sobald eine erste nützliche Auswahl fertig ist, und ergänzen Sie später weitere Gruppen, ohne die vertraute öffentliche Adresse zu ändern.

Öffentliche Abstimmungen können Publikumslieblinge hervorheben. Schlagwörter verbinden wiederkehrende Teilnehmer oder Motive. Downloads bieten eine bequeme Übertragung der Sammlung, wenn die Veranstaltungsrichtlinie sie erlaubt.

\subsection{Reise- und Ortsarchiv}
\index{r00201@Reisearchiv}
\index{k00155@Karten!r00202@Reisenutzung}

Länder-, Regionen-, Städte- und Reisegalerien bieten eine natürliche Navigation. Datumsbereiche schaffen Chronologie, während GPS-Karten einen geografischen Betrachtungsweg eröffnen. Prüfen Sie Koordinaten von Unterkünften und privaten Orten. Eine beschreibende Einführung kann Route, Zweck und historischen Kontext einer Reise erläutern.

\subsection{Historisches und institutionelles Archiv}
\index{h00147@historisches Archiv}
\index{i00150@institutionelles Archiv}

Historische Sammlungen profitieren von sorgfältigen Quellenhinweisen, Formulierungen bei unsicheren Datumsangaben, Namen, Orten und Herkunftsnachweisen. Verwenden Sie Beschreibungen, um bestätigte Fakten von familiären oder institutionellen Erinnerungen zu unterscheiden. Schlagwörter können Menschen, Gebäude, Organisationen und Zeiträume über viele Galerien hinweg verbinden.

Erstellen Sie vor großen Metadatenprojekten Sicherungen und exportieren Sie bei geplanter externer Bewahrung Informationen über verfügbare Werkzeuge. Die Galerie kann eine zugängliche Präsentationsebene bilden, während daneben die ursprünglichen Erhaltungsverfahren fortgeführt werden.

\section{Ihre Galeriedaten verstehen}
\label{sec:user-data}
\index{d00067@Datenhoheit}
\index{d00062@Datenbank!f00116@für Galeriebetreiber}
\index{d00060@Dateisystem!f00116@für Galeriebetreiber}

Galeriebetreiber arbeiten mit zwei zusammenwirkenden Speichern: gewöhnlichen Dateien und einer Datenbank. Das Verständnis ihrer Rollen erleichtert Sicherung, Migration und Fehlersuche.

\subsection{Was in Galerieordnern liegt}

Die Galerieordner enthalten Fotos und Ordnerhierarchie, sodass die Kernsammlung per FTP, Hosting-Dateimanager, Sicherungswerkzeug oder lokaler Kopie erkennbar bleibt. Das Verschieben oder Umbenennen von Dateien außerhalb der Anwendung kann Erkennung oder Synchronisierung erfordern, damit die Datenbank den neuen Zustand erfährt.

\subsection{Was in der Datenbank liegt}

Die Datenbank merkt sich Titel, Beschreibungen, Datumsangaben, Sichtbarkeit, Passwörter und Zugriffskonfiguration, öffentliche Pfade, Reihenfolge, Schlagwörter, Stimmen, Bildabmessungen, EXIF-Indizes, Prüfsummen, Vorschaubilddatensätze, Administratorkonten, Einstellungen, Auditinformationen und Ablaufzustände.\footnote{Passwortwerte verwenden zweckgerechte Hashes oder geschützte Konfiguration. Ein Datenbankexport enthält dennoch sensible Betriebsinformationen und verdient sichere Aufbewahrung.}

Die Datenbank macht Suche und Administration schnell. Sie ermöglicht außerdem einen verständlichen Dateititel, eine ausführliche Beschreibung, Schlagwörter, eine festgelegte Reihenfolge und Zugriffsregeln, ohne die ursprünglichen Bildbytes zu verändern.

\subsection{Wie Benutzer von der Datenbank profitieren}
\index{d00062@Datenbank!v00267@Vorteile}

Alltägliche Vorteile sind:

\begin{itemize}
  \item Wörtersuche in Galerie- und Fotobeschreibungen;
  \item Öffnen einer Galerie über einen stabilen öffentlichen Pfad;
  \item Anzeige von Fotos in einer kuratierten Reihenfolge;
  \item Speichern von Schlagwörtern und Besucherstimmen;
  \item Schutz einer Galerie durch geerbte Zugriffsregeln;
  \item Finden exakt gleicher Inhalte anhand gespeicherter Prüfsummen;
  \item Vorbereiten geeigneter Vorschaubilder, ohne auf jeder Seite jede Originaldatei erneut zu lesen;
  \item Aufzeichnen von Migrationen und Wartungsentscheidungen mit einem Auditprotokoll.
\end{itemize}

\subsection{Sicherung als vollständiger Satz}

Eine vollständige Sicherung enthält Galeriebaum, Datenbankexport, lokale Konfiguration und eigene benutzerdefinierte Ressourcen der Installation. Geben Sie diesen Teilen dasselbe Sicherungsdatum, damit sie einen zusammenhängenden Zustand beschreiben. Bewahren Sie eine Kopie außerhalb des Webservers auf und testen Sie die Wiederherstellung an einem getrennten Ort.
\section{Alltägliche Verwaltung und Pflege}
\label{sec:daily-care}
\index{t00239@tägliche Administration}
\index{m00174@monatliche Wartung}
\index{b00040@Betreibercheckliste}

\subsection{Nach jedem wichtigen Upload}

Prüfen Sie Zielgalerie, Fotoanzahl, Reihenfolge, Titel, Sichtbarkeit und Erscheinungsbild der Vorschaubilder. Öffnen Sie mehrere Fotos in der Lightbox und prüfen Sie die Galerie anonym. Wiederholen Sie bei einer geschützten Sammlung den Zugang per Passwort oder Freigabelink in einer neuen privaten Sitzung.

\subsection{Monatliche Prüfung}

Eine monatliche Wartungsprüfung sollte Folgendes umfassen:

\begin{enumerate}
  \item Prüfen Sie Systemzustand und ausstehende Migrationen.
  \item Prüfen Sie die jüngsten administrativen Auditereignisse.
  \item Prüfen Sie verfügbare Anwendungsaktualisierungen und deren Veröffentlichungsinformationen.
  \item Bestätigen Sie den Abschluss geplanter oder manueller Sicherungen.
  \item Öffnen Sie repräsentative öffentliche und geschützte Galerien.
  \item Prüfen Sie die Vorschaubildwartung, wenn sich Originaldateien oder Darstellungseinstellungen wesentlich geändert haben.
  \item Entfernen oder widerrufen Sie Zugangsdaten und Freigabemethoden, die ihren Zweck erfüllt haben.
\end{enumerate}

\subsection{Vor einem größeren öffentlichen Start}
\index{s00235@Startcheckliste}

Prüfen Sie Titel, Beschreibungen, Titelbilder, Rechtschreibung, mobiles Layout, Privatsphäre, Offenlegung von GPS-Daten, Downloads, Karten, Abstimmung, Suchergebnisse und Erwartungen an die Kontaktaufnahme. Testen Sie mit einem langsamen Netzwerkprofil, wenn die Zielgruppe möglicherweise mobile Verbindungen nutzt. Erstellen Sie unmittelbar vor dem Startzustand eine Sicherung.

\subsection{Wenn eine weitere Person bei der Administration hilft}

Geben Sie jedem Administrator ein eigenes Konto, sofern das Kontomodell den vorgesehenen Ablauf unterstützt. Erläutern Sie Galeriebereich, Datenschutzerwartungen, Namenskonventionen, Schlagwortvokabular, Sicherungsverantwortlichkeiten und Löschverfahren. Einzelkonten verbessern die Nachvollziehbarkeit in Auditprotokollen und halten persönliche Duplikatprüfungslisten getrennt.

\section{Öffentliches Besuchererlebnis}
\label{sec:public}
\index{o00181@öffentliche Galerie}
\index{r00207@Routing}

\subsection{Was ein Besucher sieht}
\index{b00035@Besuchererlebnis}

Ein Besucher öffnet die Startseite oder folgt einem direkten Link zu einer Galerie oder einem Foto. Die Seite zeigt Websitename, Navigation, Galerietitel, Beschreibung, Brotkrümelnavigation, Schlagwörter, Untergalerien und Fotos, die der Besucher sehen darf. Geschützte Inhalte verlangen den passenden Zugangsschritt, bevor privates Material offengelegt wird.

Jede Galerie besitzt eine öffentliche Adresse, die als Lesezeichen gespeichert oder geteilt werden kann. Eine direkte Fotoadresse öffnet den Galeriekontext und führt den Besucher zu diesem Foto. Aussagekräftige Titel und stabile öffentliche Pfade halten Links in Nachrichten, Lesezeichen und Suchergebnissen verständlich.

\subsubsection{Typischer Besuch}

Ein üblicher Besucherablauf ist:

\begin{enumerate}
  \item Öffnen Sie die Startseite und wählen Sie ein interessantes Galerietitelbild.
  \item Lesen Sie die Galerieeinführung und folgen Sie Untergalerien oder Schlagwörtern.
  \item Scrollen Sie durch die kleineren Vorschaubilder.
  \item Öffnen Sie ein Foto im großen Betrachter.
  \item Wechseln Sie mit Schaltflächen, Tasten, einem Bildstreifen oder dem ausgewählten Betrachtungsmodus durch benachbarte Fotos.
  \item Öffnen Sie eine Karte, stimmen Sie ab, suchen Sie oder laden Sie Inhalte herunter, wenn die Galerie diese Aktionen anbietet.
  \item Kehren Sie über die Brotkrümelnavigation zu einer übergeordneten Sammlung zurück.
\end{enumerate}

\subsection{Galeriehierarchie und Seitennavigation}
\index{s00216@Seitennavigation}
\index{g00127@Galeriehierarchie}

\subsubsection{Direkte Inhalte und stabile Reihenfolge}

Eine Galerie kann zuerst kleinere Galerien und danach Fotos anzeigen. Beispielsweise kann \emph{Italien 2026} die Galerien \emph{Rom}, \emph{Florenz} und \emph{Venedig} enthalten und zugleich einige Übersichtsfotos direkt auf der Italien-Seite zeigen.

Die Seitennavigation teilt eine lange Sammlung in überschaubare Seiten. Dadurch erscheint die erste Seite rasch, und das Scrollen bleibt angenehm. Eine kleine Ausstellung wirkt möglicherweise auf einer Seite am besten. Eine Hochzeit mit 1.500 Fotos profitiert meist von Seiten oder Untergalerien. Der große Fotobetrachter kann die umfassendere Galerieabfolge weiter durchlaufen und zusätzliche Informationen bei Bedarf laden.\footnote{Die Seitennavigation verändert, wie viele Karten gleichzeitig erscheinen. Die zugrunde liegenden Fotos und deren Reihenfolge bleiben unverändert.}

Lange öffentliche Listen enthalten eine schwebende Schaltfläche \ui{Nach oben}, sobald der Besucher in die Galerieliste gescrollt hat. Sie verwendet sanftes Browserscrolling, verschwindet außerhalb der aktiven Liste und bleibt verborgen, solange die Lightbox oder der Vollbildbetrachter geöffnet ist, damit sie nicht mit den Fotosteuerungen konkurriert.

\subsection{Suche, Schlagwörter und Navigation}
\index{s00237@Suche}
\index{s00212@Schlagwörter}

Die Suche kann die gesamte öffentliche Website durchsuchen oder auf die aktuelle Galerie und ihre Untergalerien beschränkt bleiben. Ein Besucher im \emph{Familienarchiv} kann nur diesen Zweig nach einem Personennamen durchsuchen; ein Besucher auf der Startseite kann über alle öffentlichen Sammlungen hinweg suchen. Der aktuelle Suchdienst beginnt bei zwei sichtbaren Suchzeichen, liefert eine begrenzte Mischung aus Galerie- und Fotoergebnissen und ordnet stärkere Titel-/Namens-/Schlagworttreffer vor schwächeren Texttreffern ein. Die Standardanfrage liefert bis zu 12 zusammengeführte Ergebnisse; der Dienst begrenzt eine Anfrage auf höchstens 30 Ergebnisse.

Durchsuchbare Informationen umfassen öffentliche Galerie-/Fototitel und Beschreibungen, Bilddateinamen, Galerie- und Fotoschlagwörter und deren öffentliche Beschreibungen, aktive gepflegte Sprachinhalte bei verfügbarem Schema und internen durchsuchbaren KI-Text bei bereiten KI-Metadaten. Die Zugriffsfilterung erfolgt als Teil der Abfrage: Private, unzugängliche oder nicht öffentliche Bildinhalte werden nicht allein wegen eines passenden Texts ausgegeben. Eine Zweigsuche beschränkt außerdem den Listenkontext auf diesen Galerieteilbaum, statt den gesamten Katalog zu durchsuchen und unpassende Treffer nachträglich auszublenden.

Schlagwörter bieten einen weiteren Weg durch das Archiv. Ein Schlagwort \emph{Nachtfotografie} kann Prag, London und Tokio verbinden, auch wenn diese Fotos in getrennten Reisegalerien liegen. Die Brotkrümelnavigation zeigt die aktuelle Position in der Hierarchie; Favoritenverknüpfungen platzieren wichtige Galerien direkt in der Websitekopfzeile.
\section{Medien- und Galerieverwaltung}
\label{sec:admin}
\index{a00005@Administration}

\subsection{Galerieübersicht und geprüfte Änderungen}
\index{Galerie!Funktionsplanung}
Der Galeriebereich zeigt direkte und untergeordnete Fotos, Untergalerien sowie eingeschaltete, ausgeschaltete und gemischte Funktionszustände. Klappen Sie einen Zweig auf und planen Sie Karten, Dateinamen, Abstimmung oder Picture Game für dessen Wurzel und Nachfolger. Die Planung speichert noch nichts. Prüfen Sie den genauen Umfang und die Abhängigkeiten und wenden Sie erst danach den geprüften Plan an. Überlappende Absichten behalten ihre Reihenfolge. Picture Game benötigt Abstimmung; deren Abschaltung deaktiviert auch das Spiel. Abbrechen oder Verwerfen ändert keine gespeicherten Werte. Nach zwischenzeitlichen Änderungen muss neu geprüft werden.

Die schnelle Passwortaktion ändert ausschließlich das eigene Passwort der aktuellen Galerie. Geben Sie ein neues Passwort ein oder entfernen Sie das eigene Passwort. Token- und übergeordneter Schutz, Sichtbarkeit und andere Einstellungen bleiben erhalten. Veraltete Revisionen werden abgelehnt. Ohne JavaScript bleiben normale Formulare verfügbar; Seitenpanel-Aktionen werden vor Ort abgeschlossen und lassen das Panel geöffnet.

Angemeldete Administratoren können auf der öffentlichen Startseite eine Wurzelgalerie erstellen. Der Hinweis auf die erste Galerie verlangt einen nachweislich leeren Katalog. Inhalt und Ort/Sichtbarkeit sind getrennt, weitere Optionen aufklappbar. Inhaltssprache, Tags, SimBrief-Entwürfe und gespeicherte Vorgaben bleiben verfügbar.

\subsubsection{Upload-Einstellungen und mobile Verbindungen}
Öffnen Sie \ui{Einstellungen > Uploads} für Formate, Umbenennung, Worker und Grenzen, ZIP-Stapel, Quelldatenblöcke für Miniaturen und mobile WebDAV-Verbindungen. Zusammengehörige Grenzen werden gemeinsam normalisiert; Größen erscheinen in MB. Die alten Menülinks können bei Bedarf eingeschaltet werden. Ausgeblendete Links deaktivieren keine bestehenden Routen oder Dienste.

Mobile Aktionen aktualisieren ihren eigenen Bereich. Kopieren Sie ein neu ausgegebenes Einmalpasswort sofort. Scheitert die anschließende Listenabfrage, bleibt die Verbindung gespeichert und das Passwort in dieser Antwort sichtbar; aktualisieren Sie die Liste getrennt. Deaktivierte mobile Uploads lesen die Verbindungsliste nicht.

\subsubsection{Darstellung, Sprache und eigene Stylesheets}
Darstellung verbindet gruppierte Felder mit einer Live-Vorschau und einem begrenzten Maus-/Tastaturteiler auf breiten Bildschirmen. Seine Breite ist eine lokale Browserpräferenz. Website-Bilder gruppiert Branding/Hintergründe; Layout trennt Verknüpfungen, Raster, Karten und Viewer-Vorgaben. Sprache bietet Einstellungen, Design, Editor und Diagnose. Die Designvorschau verwendet die angebotenen Sprachen und die öffentliche Vorgabe desselben Formulars und speichert selbst nichts.

Lassen Sie die Ersatzauswahl leer, um das installierte CSS zu behalten. Wählen Sie ausdrücklich eine Vorlage oder laden Sie eine gültige CSS-Datei hoch; der Upload hat Vorrang. Ersetzen/Zurücksetzen prüft zuerst den Speicher und stellt bei fehlgeschlagener Markerspeicherung die vorherige Datei wieder her. Bei unvollständiger Wiederherstellung behalten Sie die verfügbare Kopie und folgen dem sicheren Hinweis. Deployment erhält die installationsspezifische \filepath{public/assets/custom.css}.

\subsubsection{Kartenorientierung und Betriebszustand}
Vertikale Karten zeigen das Foto über dem Text, horizontale daneben. Die normale Orientierungsmigration erhält das frühere Aussehen von Einstellungen, physischen und Smart Galleries, Sidecars, importierten Metadaten und Papierkorb. Sichern Sie Datenbank, Galerien und Papierkorb gemeinsam. Beheben Sie nach Unterbrechung den gemeldeten Fehler und wiederholen Sie die Migration; Dokumentmarker verhindern doppelte Umwandlung. Neuinstallationen verwenden vertikale Karten. Alle Tags bleiben ohne JavaScript vorhanden; Aufklappen und Scrollen funktionieren auch auf Karten.

Öffentliche Suche ist nur ohne gespeicherte Wahl und bei verfügbarer Funktion standardmäßig an. Updates aktualisieren Versionen, Notizen und API-Zustand gemeinsam; automatische Metadatenprüfung erfolgt höchstens stündlich, deaktivierte automatische Updates lassen alte Hintergrundaufträge pausiert. OurAirports verwendet wöchentliche Aktualität und eine Stunde Fehlerabstand. Der Papierkorb zeigt Restdauer und sichere Löschbarkeit. Auch Präsentationsschema-Probleme verlangen eine Aktion im Systemstatus.

Der unabhängige Uploader bleibt bei 0.3.2. Der Build liefert Installer und passende generierte \filepath{winapp-update.json} gemeinsam unter \filepath{winapp/dist/0.3.2/}. Beide gehören zu derselben Veröffentlichung. Der Build ersetzt keine Prüfung der installierten Windows-Anwendung oder des Simulators.

\subsection{Administrativer Arbeitsbereich}

Nach der Anmeldung sieht ein Administrator private Steuerungen für Übersicht, Galerien, Fotos, Design, Schlagwörter, Wartung, Aktualisierungen, Protokolle und Kontoeinstellungen. Viele Bearbeitungen öffnen sich im rechten Seitenbereich, sodass die Galerie sichtbar bleibt, während ein Eintrag geändert und gespeichert wird.\footnote{Eine vollständige Administrationsseite bleibt für Abläufe verfügbar, die mehr Platz benötigen, oder wenn die Browsererweiterung nicht verfügbar ist.}

Auch die öffentliche Galerie zeigt kontextbezogene Administrationssteuerungen, solange eine normale Administratorsitzung aktiv ist. Die Stiftschaltflächen von Galerien und Fotos öffnen dieselben vollständigen Editoren im Seitenbereich. Die kompakte Galerieentfernung verwendet bei aktiviertem Papierkorb den wiederherstellbaren Papierkorbablauf; bei deaktiviertem Papierkorb erfolgen zukünftige Galerieentfernungen wieder über den bisherigen Dienst zur endgültigen Teilbaum-Löschung. Die kompakte Fotoentfernung ist bewusst als \ui{Foto ... aus dem CMS entfernen} beschriftet: Diese Schnellaktion auf der öffentlichen Seite entfernt die Bildzeile aus dem Katalog und führt nicht den vollständigen Dienst zur Originaldateilöschung aus. Verwenden Sie den normalen Foto-/Sammellöschablauf, wenn auch die Originaldatei und erzeugte Ableitungen gelöscht werden sollen. Der anonyme Vorschaumodus blendet diese kontextbezogenen Administrationssteuerungen aus.

Die administrative Galerietabelle kann Zeilen nach Sichtbarkeit filtern: alle, nicht veröffentlicht, öffentlich oder privat. Die Filterung ist eine browserseitige Ansichtshilfe, keine Änderung der Zugriffsrichtlinie. Bei einer durch den Filter ausgeblendeten Zeile wird auch die Auswahl aufgehoben und die Gesamtauswahl geleert, damit ein verborgenes veraltetes Kontrollkästchen nicht versehentlich in die nächste Sammelaktion eingeht. Die Zusammenfassung meldet die Anzahl angezeigter/berücksichtigungsfähiger Galerien.

Der Galeriebaum unterstützt das Auf-/Zuklappen je Zweig sowie \ui{Alle zuklappen} und \ui{Alle aufklappen}. Er merkt sich, welche Galeriezweige der Administrator zugeklappt hat. Der Browser sendet die IDs der zugeklappten Galerien über eine authentifizierte, CSRF-geschützte Aktion, und die Anwendung speichert diesen Oberflächenzustand. Bei der Rückkehr zum Administrationsbaum bleiben somit dieselben Zweige geschlossen, statt die gesamte Hierarchie wieder aufzuklappen. Bei Nachkommen, die durch einen zugeklappten Zweig verborgen sind, wird vor Sammelarbeiten ebenfalls die Auswahl aufgehoben. Neuordnen aktualisiert die Baumsteuerungen, wenn eine Galerie Untergalerien erhält oder verliert.

\subsection{Zuverlässige Änderungen direkt im Seitenbereich}
\index{a00005@Administration!a00012@Änderungen im Seitenbereich}
\index{s00214@Seitenbereich!a00013@Änderungsabschluss}

Bei verfügbarem JavaScript wird eine dauerhafte Aktion, die im rechten Administrations-Seitenbereich beginnt, dort abgeschlossen. Der Bereich bleibt eingebunden und offen, die Browseradresse unverändert, und die Seite wird nicht allein zur Ergebnisanzeige neu geladen. Das gilt für Erstellen/Bearbeiten/Verschieben/Löschen von Galerien, Bildaktionen für Hochladen/Bearbeiten/Löschen/Verschieben/Neuordnen/Sichtbarkeit/NSFW/Titelbild, Scan/Import, EXIF-Datumsvorschläge, Vorschaubilderstellung, Medienumbenenner, Metadatenorganizer, Prüf-/Löschaktionen des Fotoduplikatdetektors, Zugangsdaten der Uploadautomatisierung, Zielabrufe der Galeriemigration, Schlagwortbearbeitung sowie Smart-Gallery-Aktionen für Erstellen/Bearbeiten/Platzieren/Trennen/Duplizieren/Löschen. Eine eigenständige Administrationsseite und gewöhnliches POST-/Weiterleitungsverhalten bleiben bei deaktiviertem JavaScript oder direkt geöffneter Route verfügbar.

Nur die neueste Anfrage zum Öffnen des Seitenbereichs darf dessen Inhalt ersetzen; Schließen oder Öffnen eines anderen Eintrags entzieht einer älteren Antwort die Zuständigkeit. Der Tastaturfokus bleibt im geöffneten Bereich und kehrt zu dessen öffnender Steuerung zurück. Ungespeicherter gewöhnlicher Text kann für die aktuelle Seite in begrenztem Browserspeicher vorgehalten und bei der Rückkehr desselben Formulars wiederhergestellt werden. Passwörter, Tokens, API-Schlüssel, Dateifelder, Vorgangskennungen und andere private Werte sind ausgeschlossen; nichts wird dauerhaft im Browserspeicher abgelegt.

Der gespeicherte Datenbank-/Dateisystemzustand bleibt maßgeblich. Nach der Speicherung beschreibt der Server die geänderte Entität, das gegebenenfalls neu darzustellende Seitenbereichsfragment und jeden betroffenen öffentlichen Kontext. Stabile Galerie-, Bild-, Schlagwort- und Smart-Gallery-Kennungen werden verwendet, statt anzunehmen, dass die alte URL weiterhin das gespeicherte Objekt bezeichnet. Das ist nach einer Umbenennung von Galerieordner/Slug, einem Wechsel der übergeordneten Galerie oder einer Schlagwort-Slug-Änderung wichtig: Die Anwendung kann die neue kanonische Darstellungsquelle abrufen und dabei die sichtbare Browseradresse unverändert lassen. Bei Verschiebungen werden Quell- und Zielgalerie beschrieben; bei einem Wechsel der übergeordneten Galerie die alte und neue übergeordnete Galerie sowie die verschobene Galerie; Smart-Gallery-Platzierungsänderungen beschreiben jeden betroffenen Stamm- oder physischen Elternkontext.

Eine erfolgreiche Fragmentanfrage beweist allein noch nicht die Sichtbarkeit der Änderung. Das aktualisierte serverseitig dargestellte HTML muss eine vorgangsspezifische beobachtbare Bedingung erfüllen, sofern eine existiert. Beispiele sind Vorhandensein oder Fehlen von Galerien/Karten und vollständige Zählungen, genaue wirksame Sichtbarkeit, gespeicherte Galerierevision, Bildabwesenheit/-sichtbarkeit/NSFW-Zustand, galerieübergreifende Bildrevision bei Änderungen außerhalb der Seite, Titelbildkennung, gespeicherte Bildreihenfolge, Schlagwortkennung und Smart-Gallery-Platzierung/-Reihenfolge. Seitennavigationsbewusste Prüfungen akzeptieren eine direkt sichtbare passende Karte, bevor sie Gesamtzählungen verwenden; bei Ergebnissen außerhalb der aktuellen Seite wird weiterhin darauf gewartet, dass die maßgebliche Gesamtzahl übereinstimmt. Die Reparatur des Vorschaubildcaches ist die beabsichtigte Ausnahme: Sie aktualisiert betroffene Kontexte, besitzt aber keine stabile datenbankgestützte Bedingung im öffentlichen HTML.

Unabhängige Darstellungsanfragen können bei Shared Hosting kurzzeitig veraltete Daten liefern. PHP Gallery führt daher eine einzige zentralisierte, begrenzte Wiederholungsfolge nur dann aus, wenn die Änderung selbst erfolgreich war und die deklarierte Bedingung noch nicht beobachtbar ist. Jeder Abschluss erhält eine neue Vorgangsgeneration; eine neuere Änderung bricht ältere Abrufe, Wiederholungsaktivierungen und Seitenbereichsantworten für denselben Kontext ab oder ersetzt sie. Strukturell gültiges, aber veraltetes HTML wird vor dem Ersetzen geprüft, sodass eine ältere Antwort keinen neueren sichtbaren Zustand überschreiben kann. Aktualisierungen umgehen Browsercaches und bewahren anwendbare Zustände von Seitennavigation, Sortierung, Sprache, Filter, aktivem Register sowie Scrollpositionen von Bereich und Seite.

Gelingt die Speicherung, lässt sich aber die aktualisierte Ansicht innerhalb des begrenzten Wiederholungsbudgets nicht prüfen, meldet der Seitenbereich ein Synchronisierungsproblem und erhält die gespeicherte Serveränderung. Er meldet die Änderung nicht fälschlich als fehlgeschlagen, lädt die Seite nicht still neu, navigiert nicht weg und schreibt den Browserverlauf nicht um. Fehler bei Authentifizierung, Autorisierung, CSRF, Schemabereitschaft, Pfadsicherheit und Validierung liefern einen erwarteten strukturierten Fehler an den erweiterten Ablauf und werden nicht wiederholt. Dynamisch aktualisierte Steuerungen bleiben aktiv, weil der Seitenbereich lebenszyklussichere delegierte Behandlung verwendet statt Bindungen, die nach dem Fragmentaustausch verschwinden.

Browsergestütztes Erstellen/Hochladen besitzt zusätzliche Schutzmaßnahmen. Die Auswahl der Browserverarbeitung ist eine ausdrückliche Ausführungsentscheidung: Eine fehlende Browserfähigkeit oder fehlgeschlagene Arbeitsprozessvorbereitung stoppt vor der Speicherung, statt still auf serverseitige Vorschaubilderzeugung umzuschalten. PHP prüft vor der Speicherung der Originale, ob jede erforderliche vorbereitete Vorschaubildgröße und jedes erforderliche Format vorliegt. Sobald die browsergestützte Speicherung beginnt, werden Ersatzmetadaten für die direkte Seite niemals als Erlaubnis zur Wiederholung der klassischen Erstellungs-/Uploadanfrage interpretiert; dies verhindert doppelte Galerien und Fotos. Die konfigurierte ZIP-Chargengröße ist ein weiches Packziel: Ein unteilbares Fotopaket oberhalb dieses Ziels kann allein in einer Charge übertragen werden; die wirksame PHP-Uploadgrenze bleibt jedoch eine harte Grenze. Die Windows-Uploadautomatisierung kann die vollständige JPEG+WebP-Kompatibilitätsmatrix senden; ein moderner reiner WebP-Server akzeptiert die erforderlichen WebP-Varianten und überspringt gültige zusätzliche JPEG-Varianten sicher. Fehlerhafte, unpassende, unbekannt große oder nicht unterstützte Vorschaubilddaten bleiben harte Fehler.

\subsection{Zentrale Einstellungen}
\index{e00084@Einstellungen!z00275@Zentrale}

Die Administrationsnavigation enthält eine zentrale Seite \textbf{Einstellungen} für websiteweite Konfiguration. Allgemein bündelt Websiteadresse, Namen, Administrations- und öffentliche Sprache sowie Navigation. Weitere Abschnitte umfassen Öffentliches Erscheinungsbild, Inhalte, Medien und Navigation, Uploads und Automatisierung, Datenschutz und Diagnose sowie Erweitert. Frühere Websiteadresslinks führen zu Allgemein. Stabile Abschnittslinks und normale Formulare funktionieren ohne JavaScript.

Mit JavaScript werden Abschnitte beim Öffnen geladen und direkt gespeichert. Eine feste Speicherleiste zeigt geänderte Einstellungen und bietet Speichern oder Zurücksetzen des Entwurfs. Erweiterte Zusammenfassungen bleiben eingeklappt. Adress- oder Administrationssprachänderungen bieten einen ausdrücklichen Link zur neuen Adresse oder Sprache. Das Dashboard zeigt zuerst seine Oberfläche und lädt Übersichtssummen oder die Galerieliste bei Bedarf; fehlgeschlagene Abfragen können wiederholt werden. Die Übersicht zählt indexierte Originale ohne Dateiscan oder Ermittlung von Titelbildern.

Ein Feld nach Spotlight-Art durchsucht Einstellungen während der Eingabe. Es gleicht Namen, Beschreibungen, Kategorien und technische Namen ab, einschließlich Steuerungen spezieller Administrationsseiten. Die Auswahl eines Ergebnisses öffnet den passenden Abschnitt oder die Spezialseite und fokussiert die entsprechende Steuerung. Pfeiltasten, Eingabetaste, Escape und die Löschschaltfläche sind ohne Maus nutzbar.

Häufig verwendete globale Einstellungen lassen sich direkt in der Zentrale bearbeiten, darunter Websitename, öffentliche Sprache, URL-Umschreibung, öffentliche Suche, soweit verfügbar, Vorschaubilddarstellung, globale EXIF/GPS-Anzeigevorgabe und Entwicklungsdiagnose. Andere Einstellungen bleiben auf ihren Spezialseiten, wenn sie mehr Kontext, Zugangsdaten, Dateiuploads, destruktive Aktionen oder größere Editoren benötigen. Beispiele sind Designlayout, galeriegebundene Upload-API-Schlüssel, Telemetrie, Kontosicherheit, Google-/OpenAI-Konfiguration, benutzerdefiniertes CSS, Branding, Sprachpakete und Datenbankwartung.

Die Zentrale hebt bestehende Vererbung nicht auf. Das Schlagwortseitenlayout kann weiterhin globales Raster und Kartendesign erben; galeriespezifische Karten-, Lightbox- und EXIF/GPS-Entscheidungen bleiben bei ihrer Galerie. Kontextbezogene Galerie- und Fotowerte werden nicht zu globalen Einstellungen.

Geheimnisse erscheinen niemals in Suchergebnissen oder Zusammenfassungen. Passwörter, Tokens, API-Zugangsdaten und ähnliche Werte erscheinen nur als sichere Zustände wie Konfiguriert, Nicht konfiguriert oder als Link zu ihrer Spezialseite.

\subsubsection{Geführter Einrichtungsassistent}
\index{e00084@Einstellungen!e00083@Einrichtungsassistent}
\index{e00083@Einrichtungsassistent}

Wählen Sie oben in \ui{Einstellungen} den \ui{Einrichtungsassistenten}, um einen geführten Konfigurationsentwurf zu beginnen oder fortzusetzen. Der Assistent folgt denselben acht Abschnitten wie die Zentrale. Kurze Unterabschnitte halten zusammengehörige Optionen beieinander; \ui{Erweitert} und \ui{Experten- und Spezialeinstellungen} zeigen seltenere Optionen. Bei deaktiviertem JavaScript bleiben alle Unterabschnittsinhalte lesbar, und die normale Formularnavigation funktioniert weiterhin.

Jede unterstützte Einstellung besitzt Erklärung und Beispiel. Lassen Sie ihr Einbeziehungskontrollkästchen ausgewählt, um einen Wert vorzumerken, oder deaktivieren Sie es, um die ursprüngliche Einstellung zu erhalten. \ui{Überspringen} erhält den gesamten Abschnitt. \ui{Weiter}, \ui{Zurück} und die Unterabschnittsnavigation speichern keine Anwendungseinstellungen. Unterstützte Designoptionen teilen die vorhandene Live-Vorschau; eine Vorschau ist vorübergehend und verändert die öffentliche Website nicht.

Der Assistent unterstützt sichere Werte der zentralen Einstellungen, grundlegende Designfarben und Layouts, Raster-/Karten- und Hero-Schlagwortoptionen, begrenzte Upload- und Telemetriepräferenzen, Vorschaubildvorbereitung, SEO-Schutz, Galeriepapierkorb, automatische Aktualisierungen und Einstellungen der geplanten Wartung. Zugangsdaten, API-Schlüsselaktionen, hochgeladenes Branding, unverarbeitetes CSS, die Verlagerung des Galeriespeichers und destruktive Wartung bleiben rein informativ. Ihre Spezialseiten werden nicht aus dem Entwurf geöffnet; verlassen Sie den Assistenten, bevor Sie diese Abläufe verwenden. Überschreibungen je Galerie bleiben unverändert.

Die Abschlussprüfung zeigt zuerst vorgeschlagene Änderungen. Klappen Sie unveränderte Details auf oder verwenden Sie \ui{Alle Einstellungen anzeigen}, um erhaltene und informative Einträge zu prüfen. Wählen Sie das Zustimmungskontrollkästchen und übernehmen Sie erst nach Prüfung der Zusammenfassung. \ui{Abbrechen} oder \ui{Neu starten} verwirft den Entwurf. Ein Entwurf gehört dem angemeldeten Administrator und der aktuellen Browsersitzung; er ist kein gemeinsam verwendetes Konfigurationsdokument.

Das Speichern prüft die aktuellen Einstellungen erneut. Hat eine andere Sitzung einen betroffenen Wert geändert, verweigert der Assistent dessen Überschreibung. Der erforderliche Datenbankspeicher muss überprüft sein; fehlender oder unbekannter Speicher blockiert die betroffene Speicherung. Anwendungs- und Telemetriezeilen verwenden eine Transaktion; Änderungen der Websiteadresse besitzen einen reversiblen Konfigurationsschreibvorgang zur Fehlerwiederherstellung. Papierkorb- und geplante Wartungspräferenzen werden über ihre bestehenden Verantwortungsbereiche als zusammenhängende Gruppen gespeichert. Geheimnisse gelangen niemals in Entwurf, Zusammenfassung oder Assistentenprotokoll.

\subsubsection{Websiteadresse korrigieren}
\index{e00084@Einstellungen!w00272@Websiteadresse}
\index{k00158@Konfiguration!b00026@Basis-URL}

Öffnen Sie \ui{Einstellungen > Allgemein}, um nach der Einrichtung die öffentliche Installations-URL zu korrigieren. Geben Sie eine absolute HTTP- oder HTTPS-Adresse wie \url{https://example.com} ein. Fügen Sie nur dann einen Pfad hinzu, wenn die Galerie tatsächlich in diesem Unterverzeichnis installiert ist. Hat die Hostingerkennung ein unerwünschtes Verzeichnis ergänzt, entfernen Sie es aus diesem Feld und speichern Sie.

Das Speichern verändert nur die \codeword{base_url}-Zeichenfolge auf oberster Ebene in \filepath{config.php}; es erstellt oder verändert keine Datenbankeinstellung. Andere Konfigurationswerte, Kommentare und Formatierung bleiben erhalten. Sowohl die Datei als auch ihr Verzeichnis müssen beschreibbar sein. Die Anwendung lehnt Adressen mit Zugangsdaten, Abfrageparametern oder Fragmenten ab und verweigert das Umschreiben mehrfach vorhandener oder berechneter Konfigurationswerte. Bearbeiten Sie eine benutzerdefinierte Konfiguration manuell, wenn sie nicht sicher aktualisiert werden kann.

Nach erfolgreichem Speichern öffnet der Browser denselben Einstellungsabschnitt unter der neuen Adresse. Nachfolgende Anfragen verwenden die aktualisierte Konfiguration. Diese Einstellung korrigiert erzeugte Adressen; sie verschiebt weder Galeriedateien noch konfiguriert sie DNS oder Webserverrouting.

\subsection{Schlagwortmetadaten verwalten}
\index{s00212@Schlagwörter!a00005@Administration}
\index{s00212@Schlagwörter!m00170@Metadaten}

\ui{Administration > Schlagwörter bearbeiten} ist der zentrale Editor für das wiederverwendbare Schlagwortvokabular. Die Liste kann nach Gesamtnutzung oder alphabetisch sortiert werden und zeigt für jedes Schlagwort die aktuelle gemeinsame Nutzungszahl in Galerien und Fotos. Die Auswahl eines Schlagworts zeigt seinen kanonischen kleingeschriebenen Namen, den Slug seiner sprechenden öffentlichen URL und eine optionale öffentliche Beschreibung für die Schlagwort-Zielseite. Dieselbe Ansicht führt aktuell zugeordnete Galerien und Fotos auf, mit Links zurück zum passenden öffentlichen oder Fotoeditorkontext.

Das Umbenennen eines Schlagworts oder Ändern seines Slugs aktualisiert den gemeinsamen Schlagwortdatensatz, sodass jede Galerie-/Fotozuordnung weiterhin dasselbe wiederverwendbare Schlagwort referenziert. Namen und Slugs werden beim Speichern gemäß der sicheren Schlagwortkonvention normalisiert; ein doppelter Slug wird abgelehnt. Das Löschen eines Schlagworts ist anders: Die Anwendung löst es transaktional von Galerien und Fotos und entfernt anschließend den Schlagwortdatensatz. Verwenden Sie \ui{Öffentliches Schlagwort anzeigen}, um die Zielseite zu prüfen, und \ui{Schlagwortanzeige konfigurieren}, um die Designsteuerungen für Schlagwortseiten/-karten und Hero-Schlagwortpräsentation zu öffnen.

\subsection{Funktionsschalter}
\index{f00114@Funktionsschalter}
\index{f00112@Funktionen!a00005@Administration}

\ui{Erscheinungsbild > Funktionen} bietet globale Schalter für optionale Anwendungsteile. Registrierte Funktionen sind standardmäßig aktiviert, damit eine Aktualisierung nicht still Verhalten entfernt, das eine vorhandene Installation bereits nutzt. Ein Administrator kann einzelne Besucherfunktionen, Karten- und Flugsimulationswerkzeuge, KI-/Uploadintegrationen oder spezielle Administrationswerkzeuge deaktivieren, wenn eine einfachere Installation gewünscht ist.

Eine deaktivierte Funktion wird aus der normalen Navigation ausgeblendet; ihre geschützten Routen verweigern den Funktionsablauf. Das Ausschalten eines Schalters löscht weder Funktionscode noch Fotos, Datenbankdatensätze, Zugangsdaten oder Konfiguration. Erneutes Aktivieren stellt daher die normalen Einstiegspunkte unter denselben Schema-, Zugriffs- und Konfigurationsprüfungen wie zuvor wieder her. Schaltbare Funktionen umfassen beispielsweise öffentliche Suche, Lightbox-Browsingmodi, Picture Manager, Inline-Administration auf öffentlichen Seiten, Smart Galleries, Abstimmung, Picture Game, Downloads, Galerie- und Flugkarten, Navigationsdaten, SimBrief, OpenAI-Unterstützung, lokale KI-Metadaten, API-Uploads, mobiles WebDAV, Galeriemigration, Medienumbenenner und anonyme Telemetrie.

\subsection{Galerieverwaltung}

Die Galerieverwaltung umfasst den gesamten Lebenszyklus einer Sammlung: Erstellung, Erkennung, Benennung, Beschreibung, Datumsangaben, Schlagwörter, Privatsphäre, Passwörter, Branding, Titelbilder, Layout, Reihenfolge, Verschieben und schließlich Löschen. Eine Untergalerie kann nützliche Entscheidungen ihrer übergeordneten Galerie erben. Beispielsweise kann jede Galerie innerhalb eines privaten Familienzweigs dessen Schutz übernehmen; eine einzelne Untergalerie kann eine eigene Präsentation wählen, soweit der Editor eine Überschreibung anbietet. Die Sammelgalerieansicht bietet entsprechende Aktionen über mehrere Galerien, wenn dieselbe Änderung für mehrere Sammlungen zugleich sinnvoll ist.

Galeriebeschreibungen können sichere externe Links darstellen, geschrieben als gewöhnliche Markdown-Links, \codeword{[link=URL]label[/link]}, \codeword{[url=URL]label[/url]} oder Kurzformen, deren sichtbarer Text zugleich das Ziel ist. Nur HTTP- und HTTPS-Ziele werden verlinkt; andere Schemata bleiben inaktiver Text. Erkannte Dienste verwenden gebündelte lokale Markensymbole. Bei anderen öffentlichen Hosts kann das Speichern der Galerie ein kleines validiertes Favicon in den lokalen Galeriecache abrufen. Abrufe werden nach Hostnamen dedupliziert, durch Größen-, Zeit-, Weiterleitungs- und Wiederholungsgrenzen beschränkt, lehnen private oder reservierte Netzwerkziele ab und erfolgen niemals während der Darstellung einer öffentlichen Beschreibung für einen Besucher. Ist der Cache oder seine Datenbanktabelle nicht verfügbar, bleibt der Link ohne Symbol nutzbar.

Linkziele behalten Unterstriche und Abfrageparameter. Formatierte Beschriftungen und wörtlicher Inline-Code bleiben erhalten, und Links bleiben im umgebenden Absatz.

Ist die Speicherung sprechender öffentlicher Pfade verfügbar, erstellt \ui{Sprechende öffentliche Pfade neu erzeugen} die gespeicherten Galerie- und Bildpfade aus der aktuellen Hierarchie neu und meldet getrennte Galerie-/Bildanzahlen. Neuordnen und ausgewählte Editorspeicherungen aktualisieren Pfade ebenfalls, wenn ihre Änderungen es erfordern. Diese Wartungsaktion ist nach Hierarchiereparatur, Migration oder älteren Importen nützlich; sie aktiviert die URL-Umschreibung nicht selbst, sodass Routing über Abfrageparameter der Ersatzweg bleibt, wenn Umschreibung deaktiviert oder vom Hoster nicht unterstützt wird.
\subsection{Metadatenorganizer}
\index{m00171@Metadatenorganizer}
\index{e00096@EXIF!d00071@Datumsorganizer}
\index{g00117@Galerie!a00025@automatische Organisation}

Der Galerieeditor enthält einen \ui{Organizer}-Ablauf, um eine flache Sammlung direkter Fotos in datumsbasierte Untergalerien umzuwandeln. Die aktuelle Implementierung gruppiert nach dem bereits in der Bilddatenbank gespeicherten EXIF-Aufnahmedatum. Eine sekundäre Gruppierung nach Ort ist für die Zukunft vorgesehen und in Version~\version{} keine aktive Organizerstrategie.

Wählen Sie frühestes und spätestes EXIF-Datum und erstellen Sie anschließend eine Vorschau. Die Vorschau liest Datenbankmetadaten in kleinen Chargen und scannt, benennt oder verschiebt keine Originaldateien. Der Entwurf zeigt die Anzahl direkter Fotos, passender Kandidaten, vorgeschlagener Untergalerien und ignorierter Fotos ohne brauchbares EXIF-Datum oder außerhalb des gewählten Bereichs. Ein frühestes Datum hilft besonders dabei, Fotos einer Kamera mit nicht eingestellter Uhr auszuschließen.

Für jeden akzeptierten Kalendertag ist der vorgeschlagene Unterordnername das stabile ISO-Datum \codeword{YYYY-MM-DD}. Sein Anzeigetitel verwendet die Formatierung für Galeriedaten und fällt auf einen lesbaren Tag-/Monat-/Jahr-Wert zurück, wenn diese keine Bezeichnung liefern kann. Bevor der Organizer eine neue Untergalerie vorschlägt, sucht er nach einer vorhandenen direkten Untergalerie mit demselben Organizertitel und verwendet sie erneut. Dadurch ergänzt ein wiederholter Vorschau-/Übernahmezyklus dasselbe Tagesziel, statt eine weitere gleichwertige Tagesgalerie zu erstellen.

Die Übernahme des Entwurfs erfordert eine gesonderte Bestätigung. PHP Gallery erstellt oder verwendet die vorgeschlagenen Datumsuntergalerien erneut und verschiebt passende Fotos über denselben physischen Verschiebedienst wie der gewöhnliche Sammelbildverschieber. Originale und bekannte erzeugte Ableitungen werden daher gemeinsam mit der Aktualisierung der Datenbankzuordnung verschoben. Große Übernahmen werden in begrenzte browsergesteuerte Anfragen aufgeteilt statt in einer langen PHP-Anfrage ausgeführt. Prüfen Sie den Entwurf und halten Sie eine Sicherung vor, bevor Sie eine große Galerie neu organisieren.

\subsection{Medienumbenenner}
\index{m00168@Medienumbenenner}
\index{d00059@Dateinamen!u00245@Umbenennen}
\index{m00169@Medienwartung}

\ui{Galerien > Medienumbenenner} und der Umbenennungs-Seitenbereich einer einzelnen Galerie können physische Bilddateinamen anhand des Galeriekontexts und der aktuellen Fotoreihenfolge normalisieren. Dies ist ein Dateisystemvorgang, keine kosmetische Titeländerung; daher beginnt das Werkzeug immer mit einem Probelaufplan. Die Vorschau führt jeden alten und vorgeschlagenen Dateinamen auf und kennzeichnet Dateien, die bereits passen, fehlen, kollidieren würden oder aufgrund einer Sicherheitsregel übersprungen werden müssen.

Nur direkte Bilddateien im ausgewählten Galerieordner kommen infrage. Fehlende Originale, Pfade außerhalb der Galerie und unsichere Ziele werden nicht umbenannt. Ist der bevorzugte Zielname bereits belegt, kann aber ein sicherer deterministischer Zusatz die Kollision lösen, zeigt die Vorschau diese Anpassung. Die Übernahme eines Plans erfordert die Bestätigung, dass die Vorschau geprüft wurde.

Das Standardmuster ist \codeword{\{gallery_path\}_\{seq4\}}. Ein benutzerdefiniertes Muster darf \codeword{\{gallery_path\}}, \codeword{\{gallery_title\}}, \codeword{\{photo_title\}}, \codeword{\{old_name\}}, \codeword{\{old_stem\}}, \codeword{\{image_id\}} sowie Sequenzplatzhalter von \codeword{\{seq\}} bis \codeword{\{seq5\}} verwenden. Nummerierte Sequenzvarianten füllen auf zwei, drei, vier oder fünf Stellen auf. Textbestandteile werden zu kleingeschriebenen dateisystemsicheren ASCII-artigen Namensstämmen mit Unterstrichen normalisiert, die ursprüngliche Erweiterung bleibt in bereinigter Form erhalten, und ein leeres benutzerdefiniertes Muster fällt auf die Vorgabe zurück.

Nach erfolgreicher physischer Umbenennung hält PHP Gallery Datenbankzeile und öffentliche Pfadreferenzen synchron, aktualisiert abgeleitete Titel, soweit der Umbenenner dafür zuständig ist, verschiebt oder verwirft erzeugte Ableitungen nach Bedarf und entfernt veraltete ZIP-Archivcachedatensätze mit Bezug auf die alte Originalkennung. Warnungen bei der Ableitungsbereinigung rechtfertigen keine Rücknahme einer ansonsten abgeschlossenen Originalumbenennung, weil reproduzierbare Ableitungen erneut erzeugt werden können. Bei websiteweiter Auswahl führt der Browser die Arbeit in begrenzten Teilen weiter und zeigt das Ergebnis für jede Galerie.

\subsection{Bildverwaltung}
\label{sec:media}
\index{b00043@Bilder}
\index{u00249@Uploads}

Fotos können im Browser hochgeladen, in Ordner kopiert und erkannt, durch konfigurierte Automatisierung empfangen, von einer anderen kompatiblen Installation übertragen oder über einen bereichsgebundenen mobilen Ablauf hinzugefügt werden. Nach dem Eingang vermerkt die Galerie Größe und verfügbare Kamerainformationen und bereitet kleinere Anzeigekopien vor, die auf der gesamten Website verwendet werden.

Der einzelne Fotoeditor umfasst derzeit Titel, Beschreibung, gepflegte Sprachvarianten, Sichtbarkeit, Sortierreihenfolge, Schlagwörter, NSFW-/18+-Zustand, die private redaktionelle Bewertung für Smart-Gallery-Regeln, optionale fotoeigene Grenzen responsiver Vorschaubilder und schreibgeschützte EXIF/GPS-Details. Sind lokale KI-Metadaten vorhanden, zeigt der Editor das erzeugende Modell/die Quelle, durchsuchbaren internen Text und unverarbeitetes Metadaten-JSON getrennt von der öffentlichen Bildunterschrift. OpenAI-Steuerungen setzen prüfbare Entwürfe in bearbeitbare Textfelder ein und ersetzen gespeicherten redaktionellen Text nicht automatisch.

Sammelbildaktionen sind bewusst enger gefasst als der einzelne Editor und wiederholen serverseitige Zuordnungsprüfungen. Aktuelle Sammelaktionen umfassen:

\begin{itemize}
  \item ausgewählte Fotos in eine vorhandene Galerie verschieben;
  \item eine neue Zielgalerie erstellen und die Auswahl hinein verschieben, mit Bereinigung eines leeren neu erstellten Ziels, falls der Vorgang vor der Festschreibung von Dateien scheitert;
  \item ausgewählte Fotodatensätze mitsamt ihren zugehörigen Originalen und bekannten erzeugten Ableitungen über den normalen Löschdienst löschen;
  \item ein ausgewähltes Foto als Titelbild der Galerie festlegen;
  \item die Fotosichtbarkeit zwischen den unterstützten Zuständen Entwurf/öffentlich/privat ändern;
  \item fotoeigenen NSFW-Schutz bei bekanntem Schema ein- oder ausschalten;
  \item Vorschaubilder für ausgewählte Fotos erstellen.
\end{itemize}

Die Fotoreihenfolge lässt sich durch Ziehen des Zeilengriffs ändern; sortiert werden kann nach Dateiname oder EXIF-Aufnahmedatum. Fotos ohne Aufnahmedatum bleiben am Ende; Änderungen werden über die normale Reihenfolgeroute gespeichert.

Der Reiter \ui{Bilder} im Galeriebearbeiter bündelt häufige Fotoarbeiten in einer kompakten Liste. Zugängliche Vorschaubilder lassen sich mit Mauszeiger, Tastaturfokus oder Klick öffnen; die Zeilenauswahl unterstützt mit Umschalt-Klick einen Bereich. Jede Zeile bietet direkte Steuerungen für Sichtbarkeit, Titelbild der Galerie, Bearbeitung und Löschung des Fotos. Diese Aktionen verwenden weiterhin die vorhandenen authentifizierten Serverabläufe. \ui{Dateinamen anzeigen} ändert nur die Darstellung im Editor; Administratoren können eine optionale browserweite Vorgabe speichern, während jede Galerie ihre Auswahl für die aktuelle Editorsitzung merkt.

Ist die optionale Funktion \ui{Picture Manager} aktiviert, kann ein normal angemeldeter Administrator Fotos und direkte physische Untergalerien im öffentlichen Galerieraster auswählen. Häkchen wählen einzelne Karten; Umschalt-Klick wählt einen Bereich; Strg-Klick oder Cmd-Klick schaltet die Auswahl einer Karte um; \ui{Alle auswählen}, \ui{Auswahl aufheben} und Strg/Cmd+A unterstützen größere Auswahlen auf der sichtbaren Seite. Smart-Gallery-Karten gehören nicht zu dieser dateibasierten Auswahl. Das Ziehen eines nicht ausgewählten Fotos beginnt mit diesem Foto als Auswahl; das Ziehen einer vorhandenen Fotoauswahl verschiebt die ausgewählten Fotos gemeinsam. Eine sichtbare Untergaleriekarte wird zum Ablageziel; das Ablegen führt denselben validierten Fotoverschiebevorgang aus wie die Zielauswahl. Physische Galeriekarten selbst sind keine Ziehquellen.

Über die Werkzeugleiste kann der Administrator eine gemischte Auswahl löschen oder daraus eine echte physische Galerie erstellen. Die Erstellung akzeptiert eine durchsuchbare Auswahl der übergeordneten Galerie, kopiert ausgewählte Fotos in die neue Wurzel und ausgewählte physische Untergaleriebäume darunter, während die Quellen unverändert bleiben. Der Server schreibt aktuelle Galeriebegleitdateien, lehnt rekursive Ziele und vorhandene Zielordner ab, reserviert jede Zielwurzel vor dem Kopieren, erstellt kopierte Galeriezeilen aus dem Dateisystem neu und entfernt bei einem Fehler das neue Teilergebnis. Gemischtes Löschen entfernt ausgewählte Fotos über die normale Bildbereinigung und wendet auf ausgewählte Galeriebäume die konfigurierte wiederherstellbare Papierkorbrichtlinie an, sofern verfügbar.

Verschieben und Kopieren in eine vorhandene Galerie bleiben reine Fotoaktionen. Ist eine physische Galerie ausgewählt, heben Sie deren Auswahl auf, bevor Sie diese Steuerungen verwenden. Fotoauswahlen können auch geteilt oder heruntergeladen werden. \ui{Teilen} ruft die autorisierten browserdarstellbaren Bildversionen als \codeword{File}-Objekte im Browser ab und öffnet den nativen Web-Share-Dialog, wenn Browser/Gerät das Teilen mehrerer Dateien unterstützen. Der Browser bestimmt die verfügbaren empfangenden Apps; PHP Gallery kann daher keine Instagram Story, keinen Reel und kein anderes bestimmtes Ziel erzwingen. Ist natives Teilen nicht verfügbar, kann es die Dateien nicht annehmen oder scheitert es nach der Vorbereitung, startet der Manager den authentifizierten ZIP-Download der ausgewählten Fotos als Ersatz. Derselbe ZIP-Endpunkt bedient somit sowohl einen ausdrücklichen Download ausgewählter Fotos als auch den Ersatz beim Teilen. Gewöhnliche Besucher erhalten diese Steuerungen niemals; der Server prüft bei jeder Änderung oder ZIP-Anfrage erneut angemeldetes Konto, CSRF-Token, Quellzuordnung, Zugehörigkeit direkter Untergalerien und Zielgalerie.

\subsection{Uploads und Erkennung}
\index{e00089@Erkennung}

Die Dateisystemerkennung ist nützlich, wenn Fotos bereits per FTP, Hosting-Dateimanager oder wiederhergestellter Sicherung auf den Server kopiert wurden. Die Erkennungsansicht zeigt die Funde und erlaubt dem Administrator, diese Ordner und Dateien in den Katalog zu übernehmen. Die Erkennung läuft als sitzungsgestützter Browserauftrag in kleinen AJAX-Chargen, statt während der ersten Administrationsseitenanfrage den gesamten Baum zu durchsuchen. Die aktuelle Vorgabe untersucht bis zu 80 Verzeichnisse je Charge, akzeptiert eine begrenzte Chargengröße bis 300 und lässt den Zustand aufgegebener Erkennungsaufträge nach zwei Stunden ablaufen. So bleibt ein großer Ordnerbaum auf langsamerem Hosting reaktionsfähig und zeigt trotzdem Fortschritt und eine abschließende bearbeitbare Ergebnisliste.

Der Scanner durchsucht unterhalb des konfigurierten Galeriestammverzeichnisses und folgt keinen symbolischen Verzeichnisverknüpfungen. Verborgene Verzeichnisse und bekannte Namen erzeugter/Cacheverzeichnisse wie \filepath{cache}, \filepath{thumbs}, \filepath{thumbnail}, \filepath{thumbnails}, \filepath{preview} und \filepath{previews} werden ignoriert. Ein neuer Kandidatenzweig muss irgendwo unterhalb mindestens ein unterstütztes Bild enthalten. Leere strukturelle Elternordner können dennoch zu Galerien werden, wenn sie zum Erhalt der Hierarchie über bildhaltigen Nachkommen benötigt werden. Ein Ordner nur mit einer \filepath{gallery.json}-Begleitdatei und ohne unterstütztes Bild irgendwo in seinem Zweig wird vom aktuellen Erkennungsablauf als reiner Metadatenordner gemeldet statt als Galerie importiert. Bereits indizierte Galerieordner werden aus der Kandidatenliste ausgeblendet; ein Importkandidat, dessen Anzeigetitel einen vorhandenen Geschwistertitel duplizieren würde, wird nicht still erstellt.

Bei jedem nicht verwalteten Ergebnis kann der Administrator den Ordnerbaum an Ort und Stelle importieren, unterstützte Fotos aus ausgewählten nicht verwalteten Ordnern in eine vorhandene Galerie verschieben oder ausgewählte nicht verwaltete Ordner vom Datenträger löschen. Der Import erweitert die Auswahl in Eltern-vor-Kind-Reihenfolge und kann fehlende bildhaltige Vorfahren und Nachkommen einbeziehen. Die Verschiebeaktion vermeidet das Überschreiben eines vorhandenen Zieldateinamens durch einen deterministisch gewählten eindeutigen Namen, scannt nach erfolgreichem Verschieben die Zielgalerie und entfernt Quellverzeichnisse nur, wenn sie leer werden. Die erkennungsspezifische Datenträgerlöschung ist bewusst auf sichere nicht verwaltete Pfade unterhalb der Galeriewurzel beschränkt, die selbst keine Wurzel sind, und lehnt Ordner ab, die bereits durch Datenbankgalerien vertreten sind. Verwenden Sie für eine importierte CMS-Galerie den normalen Galerielöschablauf.

\subsubsection{Die Metadatenbegleitdatei gallery.json}
\index{g00136@gallery.json}
\index{b00028@Begleitmetadaten}

PHP Gallery hält aktiven Katalog und Beziehungen in der Datenbank; eine optionale \filepath{gallery.json}-Datei bei einem Galerieordner bewahrt nützliche bearbeitbare Metadaten im Dateisystem. Normale Galeriespeicherungen schreiben die Begleitdatei zurück in den Galerieordner. Die aktuelle Ausgabe enthält Titel, Beschreibung, normalisierte Galerieschlagwörter, Sichtbarkeit, Reihenfolge, Abstimmungs- und Dateinamensanzeigezustand, gepflegte Inhaltssprache/Übersetzungen, soweit verfügbar, manuelle Galeriedaten, ausgewählte Karten-/Anzahl-/Lightbox-/Rasterüberschreibungen, Grenzen responsiver Vorschaubilder, Zugriffs-/Listungsrichtlinie, Titelbildreferenzen und konfigurierte Galeriebranding-Ressourcen bei verfügbaren Schemata. Das Schreiben der Begleitdatei folgt denselben Schemaprüfungen wie die entsprechenden Galeriefunktionen.

Bei Erkennung und Reparatur fehlender übergeordneter Galerien kann PHP Gallery unterstützte Werte aus der Begleitdatei lesen, statt alles aus dem Ordnernamen abzuleiten. Importierte Zeilen werden vorsichtshalber nicht veröffentlicht, wenn keine Sichtbarkeit angegeben ist. Der aktuelle Importweg stellt die ausdrücklich vom Importer unterstützten Begleitdateifelder wieder her, darunter Titel/Beschreibung/Lokalisierung, Sichtbarkeit, manuelle Datumsangaben, Abstimmungs-/Dateinamenspräsentation, ausgewählte Karten-/Anzahl-/Lightbox-/Raster-/Vorschaubildüberschreibungen, Zugriffs-/Listungsrichtlinie, Reihenfolge und Brandingpfade. Die Begleitdatei ist daher als portable Galerieordnermetadaten zu behandeln, nicht als vollständige Datenbanksicherung: Beziehungszustand, Zugangsdaten, Protokolle, Stimmen, Spielzustand, Smart-Gallery-Definitionen und andere datenbankeigene Datensätze benötigen weiterhin das normale Sicherungsverfahren.

\securitynote{Die Aktion \ui{Vom Datenträger löschen} in der Erkennungsansicht ist eine echte rekursive Dateisystemlöschung nicht verwalteter Kandidatenordner. Prüfen Sie die ausgewählten relativen Pfade, bevor Sie bestätigen. Der Schutz gegen bekannte importierte Galeriepfade ist eine zusätzliche Sicherheitsgrenze, kein Ersatz für eine Sicherung.}

Der Browserupload ist für Routinearbeiten die zugänglichste Wahl. Er kann in eine vorhandene Galerie hochladen oder eine Untergalerie erstellen und optionale Fotos im selben Seitenbereichsablauf hochladen. Wählen Sie Ziel und Dateien und verfolgen Sie die Fortschrittsanzeige. Die Anwendung prüft Ziel und Dateiinformationen vor der Annahme jedes Ergebnisses.

Bei aktiviertem JavaScript fokussieren Sie den Uploadbereich und fügen ein Bild oder Bildschirmfoto mit Strg+V unter Windows/Linux oder Cmd+V unter macOS ein. Das fokussierte oder zuletzt verwendete sichtbare Uploadformular erhält die Bilder; ein offener Seitenbereich bestimmt sein eigenes Uploadziel. Alle vom Browser bereitgestellten unterstützten Bilder werden derselben Dateiauswahl wie beim Dateiwähler hinzugefügt. JPEG, PNG, GIF und WebP folgen der Wählerrichtlinie; HEIC, HEIF und DNG benötigen zusätzlich einen ausdrücklichen Hinweis auf Serverunterstützung. Vom Browser gelieferte Dateinamen bleiben erhalten; namenlose Bilddaten erhalten einen eindeutigen Namen mit Zeitstempel, laufender Nummer und passender Erweiterung. Senden Sie das normale Formular über den gewählten Browser-/Serverablauf ab; Validierung, Größenlimits, Metadaten, Vorschaubilder, Fortschritt und Fehler bleiben gleich. Textfelder behalten normales Einfügen; nicht unterstützte Inhalte und Einfügen während eines laufenden Uploads werden ignoriert. Kann der Browser Zwischenablagebilder nicht bereitstellen oder hinzufügen, verwenden Sie Dateiwähler oder Dateiablage. Uploads im Seitenbereich erhalten den offenen Bereich und die aktuelle URL.

Im rechten Uploadbereich ergänzen wiederholtes Einfügen, Dateiauswahl und Ablegen dieselbe lokale Vorschauwarteschlange. Prüfen Sie Anzahl und Gesamtgröße; Entfernen, Auswahl leeren und Pfeiltasten ändern die tatsächliche Dateiauswahl und Uploadreihenfolge. Betätigen Sie anschließend einmal die normale Uploadschaltfläche. Sowohl der Upload in eine vorhandene Galerie als auch das Erstellen mit Upload im Seitenbereich verwenden diese Warteschlange; die eigenständige Uploadseite behält ihr bisheriges Auswahlverhalten. Es entstehen weder Serverentwürfe noch eine Veröffentlichungsaktion Commit.

Vorschauen werden bedarfsgesteuert mit einem Decoder und höchstens zwölf Miniaturressourcen erzeugt; PNG-Miniaturen haben höchstens 256 Pixel an der längsten Seite. Vor dem Dekodieren werden nur die ersten 256 KiB der Header gelesen; die Quelle muss innerhalb von 24 MiB, 8192 Pixeln pro Seite und 16 * 1024 * 1024 Pixeln insgesamt liegen. Nicht unterstützte, animierte, unlesbare oder größere Bilder erhalten ein Symbol oder einen Hinweis, bleiben aber als Original ausgewählt. Diese Vorschaugrenzen ersetzen keine serverseitigen Uploadlimits. Die Auswahl entfernt keine EXIF-Daten und ersetzt kein Original durch eine Miniatur.

Wiederholte Dateinamen erhalten einen Hinweis, bleiben aber getrennte Auswahlpositionen; nach dem Entfernen eines Duplikats wird der Hinweis aktualisiert. Leere Dateien, fehlende Vorschauen und Format-/Größenhinweise entfernen keine intakten Dateien. Eine deaktivierte Dateiauswahl blockiert auch Einfügen und Ablegen. Die angezeigte Summe beschreibt die lokale Auswahl: klassische Uploads senden ein Bild pro Anfrage, die Browservorbereitung verwendet die konfigurierten Stapellimits.

Vor Schließen oder Wechseln des Seitenbereichs gibt es eine ausdrückliche Entscheidung für ungesendete Dateien; Escape schließt die Warnung und stellt den Fokus wieder her. Normales Zurücksetzen leert nur eine noch nicht gesendete Auswahl. Nach einem unklaren oder teilweise fehlgeschlagenen Upload sind Entfernen, Ersetzen und Zurücksetzen gesperrt, damit die exakten Wiederholungsdaten und Operationsschlüssel erhalten bleiben. Wiederholen Sie im selben Formular; erst die kanonische Bestätigung leert die Warteschlange. Navigation gibt Vorschauressourcen frei. Ein Rücksprung über den Browser-Verlaufscache kann den Seitenzustand wiederherstellen, normales Neuladen bewahrt ungesendete Dateien jedoch nicht.

Die komprimierte Dateigröße allein beweist nicht, dass ein Bild sicher dekodiert werden kann. Bevor GD oder optional Imagick ein vollständiges Raster zuweist, prüft PHP Gallery Format, Abmessungen, Pixelarithmetik, konfiguriertes/Laufzeit-Speicherbudget und Pipelinegrenzen. Eine nicht sicher zulässige Datei wird abgelehnt, bevor ihr Original oder erzeugte Ableitungen in der Zielgalerie festgeschrieben werden. Diese Richtlinie gilt für klassische und browservorbereitete Aufnahme; sie behauptet nicht, dass jeder externe RAW-Dekoder dasselbe prozessinterne Speichermodell besitzt.

Die Uploadeinstellungen unterscheiden den gewöhnlichen Serverweg von browsergestützter Vorbereitung. \ui{Alle serverunterstützten Formate zulassen} hält die Dateiauswahl passend zu den erkannten Serverfähigkeiten. \ui{Vom Telefon gerendertes JPG/PNG/WebP bevorzugen, kein RAW/DNG} fordert mobile Browser nach Möglichkeit zu einer gerenderten Bilddarstellung auf; das hilft, wenn ein iPhone-ProRAW-/DNG-Ablauf ungeeignete Farben erzeugt. Der Browserhinweis ist kein garantierter Transcoder; daher prüft der Server weiterhin die tatsächlich hochgeladenen Medien.

Browsergestützte Vorbereitung wird unabhängig aktiviert und ist im Uploadablauf standardmäßig ausgewählt. Der Browser kann vor dem Versand optimierte Vorschaubilder und begrenzte ZIP-Chargen vorbereiten. Scheitert die Browservorbereitung vor dem Beginn eines serverseitigen Schreibvorgangs, können gewöhnliche Bildauswahlen automatisch auf den klassischen Serverupload zurückfallen. Administratoren können Standard- und Höchstzahl der Browserarbeitsprozesse, Bilder je Charge, Verhältnis von ZIP-Größe zu PHP-Uploadgrenze, maximale ZIP-Chargenbytes und die beim browserseitigen Vorschaubildneuaufbau verwendete Quellblockgröße begrenzen. Die Normalisierung erzwingt außerdem eine harte Arbeitsprozessgrenze, damit eine versehentliche Einstellung keine unbegrenzten parallelen Anfragen auf Shared Hosting erzeugt.

Bei der Annahme einer im Browser vorbereiteten ZIP-Charge prüft der Server außerdem die CRC-32-Prüfsumme jeder gespeicherten Datei sowie ein vollständiges Zentralverzeichnis und einen Enddatensatz mit übereinstimmenden Dateinamen, Offsets, Größen und Eintragszahlen. Eine beschädigte oder widersprüchliche Charge wird vor dem Schreiben von Galeriedateien abgelehnt; erstellen und senden Sie die Charge im Browser erneut.

Aktuelle Browserpipeline-Vorgaben sind 8 Arbeitsprozesse, ein Ziel von 80~\% der erkannten PHP-Uploadgrenze für jede unkomprimierte ZIP-Charge, höchstens 8 Bilder je ZIP-Charge und eine absolute ZIP-Chargengrenze von 24~MiB. Validierte Bereiche sind 1--32 Arbeitsprozesse, 10--95~\% der PHP-Uploadgrenze, 1--64 Bilder je Charge und 1--128~MiB für die absolute ZIP-Grenze. Der browserseitige Vorschaubildneuaufbau verwendet standardmäßig einen Quell-ZIP-Block von 512~MiB, begrenzt auf 16~MiB--3~GiB; der Quellerzeuger begrenzt einen Block normalerweise auf 96 Einträge mit einer internen harten Grenze von 512. Dies sind Sicherheitsobergrenzen und Vorgaben, keine Empfehlung, einen langsamen gemeinsam genutzten Host auf die Höchstwerte anzuheben.

Bei aktiviertem browsergestütztem Upload akzeptiert die Auswahl auch ZIP-Archive, etwa einen iCloud-Photos-Export. Der Browser untersucht das Archiv lokal, ergänzt unterstützte JPEG-, PNG-, WebP- und GIF-Bilder in der gewöhnlichen Vorbereitungswarteschlange und überspringt Ordner, verborgene Metadaten und nicht unterstützte Einträge. Die ausgewählte ZIP-Datei selbst wird niemals zu PHP hochgeladen. Archive mit Stored- und Deflate-Verfahren werden unterstützt; verschlüsselte, mehrteilige, ZIP64-, beschädigte, verdächtig komprimierte oder übermäßig große Archive werden vor Beginn eines serverseitigen Uploads abgelehnt. ZIP-Import besitzt keinen klassischen PHP-Ersatzweg; lassen Sie daher browsergestützte Vorbereitung aktiviert, wenn Sie ein Archiv auswählen.

Die globale Option zur automatischen Umbenennung beim Upload kann nach Browser-, API- oder WebDAV-Aufnahme dieselbe Benennungsrichtlinie nach dem Scan ausführen wie der Medienumbenenner. Lassen Sie sie ausgeschaltet, wenn externe Dateinamen zu Ihrer Archivkonvention gehören; aktivieren Sie sie, wenn alle Eingangswege auf das normalisierte Namensschema der Galerie zusammenlaufen sollen.
\subsection{Erzeugte Metadaten und Automatisierung}
\index{k00157@KI-Metadaten}
\index{o00182@OpenAI}
\index{u00247@Uploadautomatisierung}
\index{w00270@WebDAV}

Die optionale lokale oder arbeitsprozessbasierte Bildanalysewarteschlange kann visuelle Metadaten getrennt von menschlich verfasstem Titel, Beschreibung, Schlagwörtern und EXIF-Feldern aufzeichnen. Administratoren können diese Metadaten im Fotoeditor prüfen und als durchsuchbare Anhaltspunkte verwenden, ohne redaktionellen Text still zu ersetzen. Konfigurierte OpenAI-Unterstützung kann eine Foto- oder Galeriebeschreibung erzeugen, vorhandenen Text bereinigen, Kontext zusammenfassen, sichtbare Inhalte anhand sicherer erzeugter Vorschaubilder beschreiben oder eine Übersetzung vorschlagen. Diese Aktionen setzen prüfbare Entwürfe ein; sie veröffentlichen oder speichern Text nicht automatisch, solange der Administrator den normalen Speichervorgang nicht abschließt.

Erstellen Sie für wiederholte Übertragungen ein auf die vorgesehene Galerie beschränktes Uploadautomatisierungstoken und senden Sie unterstützte Dateien an den Uploadendpunkt. Tokens lassen sich unabhängig widerrufen und müssen wie Passwörter behandelt werden. Der API-Schlüssel im Klartext wird nur bei seiner Erstellung angezeigt; danach wird nur sein Hash gespeichert. Kopieren Sie daher den neuen Schlüssel bei der Erstellung in den vorgesehenen Client. Der globale \ui{API-Manager} führt Automatisierungszugangsdaten galerieübergreifend auf und bietet einen zentralen Widerrufsablauf, ohne die ursprünglichen Geheimnisse offenzulegen. Steuerungen auf Galerieebene bleiben für dieselben bereichsgebundenen Zugangsdaten verfügbar.

Die Windows-Begleitanwendung verwendet denselben galeriegebundenen Uploadvertrag. Sie besitzt zwei normale Uploadmodi. \ui{Ordner überwachen} ignoriert bereits beim Überwachungsbeginn vorhandene Dateien, wartet auf die Stabilisierung neuer Dateien, merkt sich bestätigte Uploads lokal und wiederholt vorübergehende Fehler mit zunehmenden Warteabständen. Optional kann sie unmittelbar vor jedem überwachten Upload den aktuellen Breitengrad/Längengrad/die Höhe der Kamera aus Microsoft Flight Simulator über SimConnect anhängen; eine überwachte Quelldatei kann sie erst löschen, nachdem die Galerie einen erfolgreichen Upload bestätigt. \ui{Manueller Upload} akzeptiert eine Dateiauswahl unabhängig vom überwachten Ordner und kann responsive Vorschaubilder entweder auf dem PC mit Arbeitsprozessen erzeugen oder die Erzeugung durch den Galerieserver anfordern. Vorschaubildarbeit und Netzwerkübertragungen laufen als Pipeline, statt zunächst eine gesamte große Auswahl vorzubereiten. Die Unterstützung des Infobereichs lässt Überwachung oder manuellen Auftrag bei verborgenem Hauptfenster weiterlaufen.

Dieselbe Windows-Anwendung kann optional einen KI-Metadatenarbeitsprozess mit niedriger Priorität betreiben. Er übernimmt einen authentifizierten Analyseauftrag, lädt nur die übernommene Ressource herunter, erzeugt interne durchsuchbare Metadaten mit der ausgewählten lokalen Backend-/Modellgeneration und liefert das Ergebnis an den bestehenden Uploadautomatisierungsendpunkt zurück. Dieser Arbeitsprozess ist von menschlichen Bildunterschriften getrennt. Die galeriebezogene Aktion \ui{Neuerzeugung von KI-Metadaten erzwingen} löscht gespeicherte erzeugte Zeilen und den alten Warteschlangenzustand für diesen Galeriezweig und stellt neue Arbeit ein; die aufwendige Analyse läuft weiterhin im Windows-Arbeitsprozess.

Mobile WebDAV-Zugangsdaten sind galeriegebunden und verwenden HTTP-Basic-Authen\-tifizierung sowie ein nicht erratbares Pfadtoken. Bei der Erstellung einer Verbindung wird ihr erzeugtes Passwort einmal angezeigt; anschließend hält der Server nur einen Passworthash vor. Die Administrationsseite führt Verbindungsbezeichnung, Zielgalerie, Server-URL und letzte Nutzungszeit auf und kann die Verbindung widerrufen. Der WebDAV-Endpunkt unterstützt das minimale Verhalten von \codeword{OPTIONS}, \codeword{PROPFIND}, \codeword{MKCOL} und \codeword{PUT}, das Fototransferclients wie PhotoSync benötigen. Die aktuelle WebDAV-Aufnahme akzeptiert JPG, PNG, GIF und WebP. Konfigurieren Sie den mobilen Client nach Möglichkeit so, dass er HEIC vor der Übertragung in JPEG umwandelt. Hochgeladene Medien gelangen anschließend in dieselbe Pipeline für Scan, optionale automatische Umbenennung, Datenbank und Vorschaubilder wie andere Uploads.
\section{Vorschaubilderzeugung und öffentliche Darstellung}
\label{sec:thumbnails}
\index{v00262@Vorschaubilder}
\index{w00271@WebP}
\index{j00152@JPEG}
\index{b00048@Bildvorschauen}
\index{s00215@Seitengeschwindigkeit!v00262@Vorschaubilder}

\subsection{Was ein Vorschaubild ist}
\index{v00261@Vorschaubild!d00072@Definition}

Ein Vorschaubild ist eine kleinere Anzeigekopie eines Fotos. Das Originalfoto kann mehrere tausend Pixel breit sein und viele Megabytes belegen. Eine Galeriekarte auf einem Telefon benötigt möglicherweise nur einige hundert Pixel. Würde für jede kleine Karte das vollständige Original gesendet, müssten Besucher auf Details warten, die sie bei dieser Größe nicht sehen können.\footnote{Bei der Vorschaubilderzeugung bleibt das Originalfoto die Ressource für Archivierung und vollständige Betrachtung. Die erzeugte Kopie ist eine austauschbare Navigationshilfe.}

Stellen Sie sich eine Galerie mit 60 Fotos vor, deren Originaldateien jeweils 12~Megabytes groß sind. Das Laden sämtlicher Originale für das erste Raster könnte hunderte Megabytes erfordern. Kleine Vorschaubilder können diesen ersten Betrachtungsaufwand erheblich reduzieren. Das Original bleibt für autorisierte vollständige Betrachtung oder Downloads verfügbar, während das Raster effiziente Vorschauen erhält.\footnote{Die tatsächliche Einsparung hängt von Bildinhalt, Qualität, Format, Abmessungen und dem vom Browser gewählten Kandidaten ab. Das Beispiel veranschaulicht die Größenordnung und verspricht kein festes Kompressionsverhältnis.}

Vorschaubilder erfüllen mehrere alltägliche Zwecke:

\begin{itemize}
  \item Galerieraster erscheinen früher;
  \item Scrollen beansprucht weniger Speicher und Netzwerkkapazität;
  \item mobile Besucher erhalten ein Bild näher an der Größe ihres Bildschirms;
  \item Galerietitelbilder und Suchergebnisse bleiben kompakt;
  \item der große Betrachter kann mit einer geeigneten Vorschau beginnen, während eine größere Quelle vorbereitet wird;
  \item wiederholte Besuche können bereits im Browsercache gespeicherte Dateien wiederverwenden.\footnote{Ein Browsercache speichert kürzlich verwendete Webressourcen gemäß Server- und Browserrichtlinie auf dem Gerät des Besuchers. Die Wiederverwendung kann einen späteren Besuch deutlich schneller wirken lassen.}
\end{itemize}

\subsection{Warum mehrere Vorschaubildgrößen existieren}
\index{v00261@Vorschaubild!g00145@Größen}
\index{g00138@Gerätepixelverhältnis}

Eine kleine Größe passt nicht zu jedem Bildschirm. Eine schmale Telefonkarte kann mit einem 300-Pixel-Vorschaubild klar aussehen. Eine breite Desktopkarte benötigt möglicherweise 600 oder 800~Pixel. Ein Bildschirm mit hoher Pixeldichte verwendet für dieselbe physische Fläche mehr Bildpixel und kann von einem größeren Kandidaten profitieren.

PHP Gallery kann je nach Konfiguration und Grenzen des Originalbilds übliche Größen wie 300, 600, 800, 960, 1280 und 1600~Pixel vorbereiten. Die Zahl bezeichnet die maximale lange Seite der erzeugten Kopie. Hoch- und Querformatfotos behalten ihre ursprünglichen Proportionen.\footnote{Bei einem Querformatfoto ist die lange Seite gewöhnlich die Breite. Bei einem Hochformatfoto ist sie gewöhnlich die Höhe. Die kürzere Seite wird aus dem ursprünglichen Seitenverhältnis berechnet.}

Die kleineren Kandidaten dienen Rastern und Titelbildern. Größere Kandidaten dienen breiten Karten, hochdichten Displays, Suchdarstellungen, Kartenvorschauen und dem großen Betrachter. Der Browser kann aus der verfügbaren Gruppe wählen, statt für jede Situation eine feste Größe zu erhalten.

\subsection{Warum JPEG und WebP verfügbar sind}
\index{v00261@Vorschaubild!f00105@Formate}
\index{j00152@JPEG!v00264@Vorschaubildverwendung}
\index{w00271@WebP!v00264@Vorschaubildverwendung}

JPEG ist für Fotos weit verbreitet, kompatibel und vertraut. WebP bietet häufig ähnliche sichtbare Qualität bei weniger übertragenen Bytes. PHP Gallery kann beide erzeugen, sodass ein kompatibler Browser WebP verwendet und JPEG eine verlässliche Alternative bleibt.\footnote{Die Kompressionseffizienz variiert mit Foto und Qualitätseinstellungen. Feine Textur, Rauschen, Verläufe und scharfe grafische Details können unterschiedliche Ergebnisse liefern.}

Der Besucher wählt das Format nicht manuell. Die Seite teilt dem Browser die vorhandenen Versionen mit, und der Browser wählt ein verstandenes Format. Das Originalfoto behält sein eigenes Quellformat und bleibt von diesen Betrachtungskopien getrennt.

\subsection{Wann Vorschaubilder erstellt werden}
\index{v00261@Vorschaubild!e00095@Erzeugung}
\index{v00261@Vorschaubild!w00268@Wartung}

Vorschaubilder können beim Import oder Upload von Fotos vorbereitet, durch administrative Wartung neu erstellt oder durch einen geschützten Reparaturprozess erzeugt werden, wenn die öffentliche Seite eine fehlende erwartete Kopie entdeckt. Vorbereitung nahe dem Uploadzeitpunkt bietet späteren Besuchern das flüssigste Erlebnis, weil die benötigten Dateien bereits vor dem ersten Besuch existieren.\footnote{Große Importe können erhebliche Verarbeitungszeit und Speicher benötigen. Eine bewusst ausgeführte administrative Vorbereitung gibt dem Betreiber eine bessere Gelegenheit, den Abschluss zu überwachen.}

Ein Galeriebetreiber sollte die Vorschaubildwartung nach folgenden Ereignissen prüfen:

\begin{itemize}
  \item Import eines großen vorhandenen Ordnerbaums;
  \item Änderung aktivierter Vorschaubildgrößen oder der Qualität;
  \item Wiederherstellung einer Sicherung mit Originalen und Datenbankdaten, aber unvollständigem Vorschaubildcache;
  \item Verschieben einer Installation zwischen Servern mit unterschiedlicher Bildformatunterstützung;
  \item wiederholte Ausweichanzeigen von Originaldateien oder fehlende Vorschauen im Systemzustand.
\end{itemize}

Erzeugte Vorschaubilder sind reproduzierbare Cachedaten und können aus den Originalfotos neu erstellt werden. Die Originale selbst benötigen dauerhaften Sicherungsschutz.

\subsection{DNG- und RAW-Anzeigemaster}
\index{d00078@DNG}
\index{r00200@RAW-Fotos}

DNG-Fotos können die Archivquelle bleiben, während PHP Gallery einen Browser-Anzeigemaster für Vorschaubilderzeugung und öffentliche Betrachtung erstellt. Der Anzeigemaster ist eine erzeugte Ableitung, kein Ersatz für das Original-DNG. Ist das erforderliche lokale Konvertierungswerkzeug nicht verfügbar oder scheitert die Konvertierung, kennzeichnen administrative Upload- und Vorschaubildberichte die fehlgeschlagene DNG-Ableitung. Die Vorschaubildwartung kann den Master später neu erstellen; das Löschen erzeugter Ableitungen löscht das Original-DNG nicht.

\ui{Laufzeitdiagnose} zeigt die aktive DNG-Konvertierungsrichtlinie gemeinsam mit den tatsächlichen GD-, Imagick/ImageMagick- und Formatfähigkeiten des Servers. Die Quellrichtlinie kann automatisch zuerst vollständige RAW-Dekodierung und dann die eingebettete Kameravorschau versuchen, ausdrücklich vollständige RAW-Dekodierung bevorzugen oder die eingebettete Vorschau bevorzugen. Die Farbverarbeitung kann browsersicheres sRGB erzwingen, das dekodierte Erscheinungsbild möglichst genau erhalten oder den Kameraweißabgleich verwenden, sofern verfügbar. Diese Entscheidungen betreffen nur erzeugte Browserableitungen. Sie schreiben niemals das Archiv-DNG um; ihre Verfügbarkeit hängt weiterhin von den Bilddelegaten und Speichergrenzen des Hosters ab.

\subsection{Was Besucher erhalten}
\index{v00261@Vorschaubild!b00049@Browserauswahl}

Die Galerie sendet dem Browser eine Liste der tatsächlich für eine Karte vorhandenen Vorschaubildkopien. Der Browser berücksichtigt den verfügbaren Platz auf der Seite, den Schärfebedarf des Bildschirms, das unterstützte Format und seine eigenen Ladeentscheidungen. Anschließend fordert er eine passende Datei an.

Angenommen, ein Foto besitzt Kopien mit 300, 600, 800 und 960~Pixeln:

\begin{itemize}
  \item eine kleine Telefonkarte erhält möglicherweise die 300-Pixel-Kopie;
  \item eine mittlere Karte erhält möglicherweise die 600-Pixel-Kopie;
  \item eine große oder hochdichte Karte erhält möglicherweise die 800- oder 960-Pixel-Kopie;
  \item der große Betrachter fordert möglicherweise eine eigens dafür vorbereitete, noch größere Vorschau an.
\end{itemize}

Die genaue Auswahl kann sich zwischen Browsern und Bildschirmen unterscheiden. Diese Flexibilität ist nützlich, weil der Server keine universelle Größe erraten muss.

\subsection{Responsive Browserauswahl}
\index{r00205@responsiver Renderer}
\index{s00233@srcset}
\index{v00261@Vorschaubild!r00204@responsiver Modus}

\ui{Responsive Browserauswahl} ist der unterstützte ältere Modus unter \ui{Administration > Design > Layout}. Die Seite stellt die verfügbaren Vorschaubildkandidaten sofort bereit, und der Browser wählt den am besten zur dargestellten Karte und Displaydichte passenden aus. JavaScript wird für diese Auswahl nicht benötigt.\footnote{Der zugrunde liegende Webstandard nennt die Kandidatenliste \emph{srcset}. Galeriebetreiber können sich auf das sichtbare Verhalten verlassen, ohne diesen Markupbegriff lernen zu müssen.}

Responsive Auswahl überträgt gewöhnlich ein ausgewähltes Vorschaubild je geladener Karte und ist die normale Wahl, wenn vorhersehbares browsernatives Verhalten am wichtigsten ist.

\begin{tabularx}{\textwidth}{>{\RaggedRight\arraybackslash}p{0.30\textwidth} X}
\toprule
\textbf{Situation} & \textbf{Wahrscheinliches Ergebnis}\\
\midrule
Schmale Mobilkarte & Ein kleinerer Vorschaubildkandidat genügt gewöhnlich.\\
Breite Desktopkarte & Der Browser wählt möglicherweise einen mittleren oder größeren Kandidaten.\\
Display mit hoher Pixeldichte & Für dieselbe CSS-Größe kann ein Kandidat mit höherer Auflösung gewählt werden.\\
JavaScript nicht verfügbar & Responsive Kandidatenauswahl funktioniert weiterhin.\\
\bottomrule
\end{tabularx}

Die 300-Pixel-Kopie dient als Ersatz, sofern verfügbar. Ihre Existenz zwingt einen modernen Browser nicht, sie zuerst herunterzuladen, weil die größeren Optionen gleichzeitig angeboten werden.

\subsection{Progressive Vorschaubildschärfung}
\index{p00196@progressiver Renderer}
\index{p00194@progressive Schärfung}
\index{v00261@Vorschaubild!p00195@progressiver Modus}

\ui{Progressive Vorschaubildschärfung} ist Standard und sicherer Ersatz unter \ui{Administration > Design > Layout}. Die Galerie zeigt zuerst ein kleines Vorschaubild, gewöhnlich 300~Pixel, und ersetzt relevante Karten später durch größere Kopien. Bei langsameren Verbindungen kann die Sammlung dadurch früher erkennbar werden; dafür werden möglicherweise sowohl das kleine als auch das ersetzende Vorschaubild übertragen.

\subsubsection{Was auf dem Bildschirm geschieht}

\begin{enumerate}
  \item Die Seite reserviert die endgültige Kartengeometrie.
  \item Kleine Vorschaubilder beginnen, sichtbare Karten zu füllen.
  \item Relevante Karten werden für einen größeren Kandidaten eingereiht.
  \item Die größere Kopie lädt, während die kleine sichtbar bleibt.
  \item Die Karte wird ersetzt, sobald die größere Kopie bereit ist.
\end{enumerate}

Größere Aktualisierungen werden eingereiht statt für jede Karte zugleich gestartet; sichtbare Karten haben Vorrang vor entfernten Inhalten. Die aktuelle progressive Option gilt für Fotokarten innerhalb einer geöffneten Galerie. Galerietitelbilder, Collagezellen, Galeriekarten auf der Startseite und Administrationsbilder verwenden weiterhin responsive Browserauswahl. Die Administrationsoberfläche kennzeichnet progressive Schärfung als Standard und responsive Auswahl als unterstützte ältere Option.

\subsubsection{Abwägung beim Datentransfer}

Der progressive Modus kann für ein Foto zwei Dateien übertragen: die erste kleine Kopie und den schärferen Ersatz. Der responsive Modus überträgt häufig eine vom Browser gewählte Kopie. Der progressive Modus bevorzugt daher frühe sichtbare Befüllung; der responsive Modus reduziert gewöhnlich die Gesamtübertragung für dieselbe Karte.

\subsection{Welchen Modus sollte ein Betreiber wählen?}
\index{v00261@Vorschaubild!r00203@Renderer wählen}

Wählen Sie \ui{Responsive Browserauswahl}, wenn effiziente browsernative Auswahl und gewöhnlich eine Vorschaubildübertragung je geladener Karte zur Sammlung passen.

Behalten Sie \ui{Progressive Vorschaubildschärfung} für ein kleines erstes Bild mit anschließenden begrenzten Steigerungen relevanter Karten bei. Testen Sie beide unterstützten Renderer mit repräsentativen Galerien und Geräten, wenn Sie die passende Präsentation für Ihre Besucher wählen.

Nützliche Fragen sind:

\begin{itemize}
  \item Sehen Besucher die ersten Fotos früher?
  \item Sieht die kleine Version während der Schärfung akzeptabel aus?
  \item Wie viele zusätzliche Daten werden bei einem typischen Besuch übertragen?
  \item Werden Telefone und hochdichte Displays in angenehmem Tempo schärfer versorgt?
  \item Passt der gewählte Modus zur wahrscheinlichen Verbindungsqualität der Zielgruppe?
\end{itemize}

\subsection{Ladeprioritäten und stabiles Layout}
\index{v00259@verzögertes Laden}
\index{l00161@Layoutstabilität}
\index{v00261@Vorschaubild!v00259@verzögertes Laden}

Die ersten sichtbaren Fotos erhalten frühere Ladeaufmerksamkeit, weil sie den ersten Eindruck prägen. Fotos weiter unten verwenden verzögertes Laden: Der Browser kann warten, bis sie sich dem sichtbaren Bereich nähern. Dadurch wird keine Verbindungskapazität für das Ende einer langen Galerie verbraucht, während der Besucher noch ihren Anfang betrachtet.\footnote{Der Zeitpunkt verzögerten Ladens liegt teilweise beim Browser. Browser berücksichtigen Scrollen, Abstand zum sichtbaren Bereich, Verbindungszustand und ihre eigene Implementierungsrichtlinie.}

Die Galerie kennt Breite und Höhe indizierter Fotos. Sie verwendet diese Proportionen, um vor Abschluss des Bildladens stabile Kartenformen zu reservieren. Text, Schaltflächen und benachbarte Karten behalten daher beim Erscheinen der Vorschaubilder ihre Positionen.

Besucher mit einer Präferenz für reduzierte Bewegung erhalten eine ruhige Präsentation ohne erforderliche fortlaufende Ladeanimation. Fotolinks bleiben während des Ladens und Schärfens nutzbar.

\subsection{Vorschaubildqualität und Speicher}
\index{v00261@Vorschaubild!q00198@Qualität}
\index{v00261@Vorschaubild!s00228@Speicher}

Höhere Qualität erhält mehr feine Details und kann größere Dateien erzeugen. Niedrigere Qualität spart Speicher und Übertragung, führt aber zu sichtbarer Kompression. Die beste Einstellung hängt von Motiv, Bildschirmnutzung, Hostingkapazität und Zielgruppe ab. Detailliertes Blattwerk, Haare, Text und feine Architektur können Kompression leichter offenbaren als breite weiche Hintergründe.

Jede zusätzliche aktivierte Größe und jedes zusätzliche Format benötigen Speicher, weil sie eine weitere Kopie jedes Fotos erzeugen. Der Nutzen ist eine genauere Auslieferung an unterschiedliche Bildschirme. Systemzustand und Speicherstatistik helfen dem Betreiber, den resultierenden Cache zu verstehen. Ein großes Quellarchiv kann eine umfangreiche Vorschaubildsammlung erzeugen; diese Dateien bleiben jedoch aus den Originalen reproduzierbar.

\subsection{Vorschaubildkompatibilität und Wartungsmodi}
\index{v00263@Vorschaubildkompatibilität}
\index{v00266@Vorschaubildwartung}
\index{b00051@browserseitiger Vorschaubildneuaufbau}

Der Wartungsbereich besitzt zwei Kompatibilitätsrichtlinien. \ui{Modern} erzeugt WebP-Vor\-schau\-bil\-der, wo der Server WebP schreiben kann. \ui{Legacy} erzeugt neben WebP auch JPEG-Ersatzvorschaubilder. Nach einer dauerhaften Umstellung auf Modern löscht \ui{Alte JPG-Vorschaubilder entfernen} nur erzeugte JPEG-Vorschaubilddateien; Originale, WebP-Ableitungen und DNG-Anzeigemaster bleiben erhalten.

Ein vollständiger Durchlauf von \ui{Fehlende Vorschaubilder prüfen} untersucht importierte Fotos und vermerkt Galerien mit fehlenden, veralteten, metadatenlosen oder im Seitenverhältnis falschen erzeugten Varianten, ohne etwas zu erstellen. Das gespeicherte Ergebnis ermöglicht \ui{Fehlende Vorschaubilder erstellen}, das nur die gemeldeten Bilder bearbeitet. \ui{Alle Vorschaubilder erstellen} unter Wartung umfasst alle Galerien; dieselbe Aktion im Bildereditor einer Galerie umfasst diese Galerie und bei ausgewähltem \ui{Untergalerien einschließen} ihre Nachkommen. Diese administrativen Vorgänge bleiben der primäre bewusste Wartungsablauf. Öffentliches Durchsehen startet keinen unbegrenzten Reparaturdurchlauf; der aktuelle Renderer besitzt jedoch eine kleine geschützte Hintergrundvorbereitung für bestimmte Karten, die wegen einer fehlenden erwarteten Ableitung auf autorisierte Originalmedien ausweichen mussten.

\ui{Vorschaubilddatenbank aktualisieren} gleicht Datenbankmetadaten mit vorhandenen erzeugten Dateien ab und entfernt erzeugte Dateien mit falschem Seitenverhältnis, statt sie zu veröffentlichen. \ui{Alle Vorschaubilder löschen} ist ein separat bestätigter destruktiver Cachevorgang; Originalfotos bleiben unberührt, und Ableitungen können später neu erstellt werden.

\ui{Browserseitiger Vorschaubildneuaufbau} ist bei diesen Neuerzeugungsaktionen standardmäßig ausgewählt, sofern Browserverarbeitung verfügbar ist. Deaktivieren Sie das Kontrollkästchen für Servererzeugung. Der Server sendet Originaldateien in begrenzten Quell-ZIP-Blöcken an den authentifizierten Browser; dieser erzeugt Vorschaubildvarianten, und vorbereitete Vorschaubild-ZIP-Chargen werden zur serverseitigen Prüfung und Installation zurückgeladen. Quellblock- und Chargengrenzen werden in den Browseruploadeinstellungen konfiguriert, sodass der Ablauf keine einzelne übergroße PHP-Anfrage benötigt.

\subsubsection{Geschützte öffentliche Vorschaubildvorbereitung}
\index{v00265@Vorschaubildvorbereitung}
\index{v00261@Vorschaubild!s00217@Selbstheilung}

Hat der Server ein öffentlich autorisiertes Foto mittels Originalmedienersatz dargestellt, kann er genau dieses Foto und die fehlenden unterstützten Größen als Kandidaten für Hintergrundvorbereitung markieren. Der Browser wartet auf Leerlaufzeit, hält bei verborgener Seite an, sendet höchstens zwei Kandidaten je Anfrage, führt höchstens vier Anfragen je Seitenladung aus und lässt etwa 1,8~Sekunden zwischen den Anfragen. Der Server begrenzt die normalisierte übermittelte Liste unabhängig auf 24 Bild-IDs; das tatsächliche Erzeugungsbudget ist jedoch kleiner: Eine anonyme Anfrage verarbeitet höchstens zwei Originalbilder in etwa vier Sekunden, eine authentifizierte Administratoranfrage höchstens vier in etwa acht Sekunden.

Ein Kandidat trägt ein HMAC-Token, gebunden an Bild-ID, Galerie-ID, Dateiname, relativen Pfad und Galerieordner der ursprünglichen Serverdarstellung. Vor jeder Erzeugung lädt der Endpunkt Bild/Galerie erneut, prüft das Token, wiederholt aktuelle Besucherzugriffsprüfungen, normalisiert angeforderte Größen auf den konfigurierten Vorschaubildsatz, konsultiert verfügbare Vorschaubildmetadaten und überspringt bereits darstellbare Varianten. Eine 60-Sekunden-Sperrfrist je Bild verhindert wiederholte aufwendige Versuche. Eine einzelne globale nicht wartende Arbeitsprozesssperre verhindert, dass paralleler öffentlicher Verkehr mehrere Bilddekodierungsaufträge anhäuft; eine belegte Anfrage meldet einfach einen anderen aktiven Vorbereitungsprozess. Die Vorbereitung ist standardmäßig über die interne Anwendungseinstellung \setting{thumbnail_background_warmup_enabled} aktiviert und schreibt begrenzte Diagnoseereignisse ins Administrationsprotokoll, sofern Protokollspeicher verfügbar ist.

\securitynote{Hintergrundvorbereitung macht ein privates oder anderweitig unzugängliches Foto nicht zum öffentlichen Reparaturziel. Der signierte Kandidat belegt, dass der Server dieses Bild als Ersatzkandidaten dargestellt hat; der Endpunkt wiederholt dennoch die aktuelle Zugriffsrichtlinie vor dem Öffnen des Originals oder Schreiben von Ableitungen.}

\subsection{Wenn ein Vorschaubild fehlt}
\index{v00261@Vorschaubild!f00101@fehlend}
\index{v00261@Vorschaubild!n00177@Neuaufbau}

Ein fehlendes Vorschaubild kann je nach Bild und verfügbaren Dateien als langsamerer Originaldateiersatz, leere Vorschau oder Wartungswarnung erscheinen. Öffnen Sie die Vorschaubildwartungswerkzeuge, prüfen Sie betroffene Galerie oder Größen und erstellen Sie fehlende Varianten neu. Prüfen Sie Schreibberechtigungen und Bildformatunterstützung, wenn der Neuaufbau wiederholt scheitert.

Öffnen Sie nach der Reparatur die Galerie mit leerem Browsercache und bestätigen Sie, dass Titelbilder, Karten und der große Betrachter die erwarteten Vorschauen erhalten.
\section{Lightbox, Karten, EXIF und Abstimmung}
\label{sec:lightbox}
\index{l00164@Lightbox}
\index{e00096@EXIF}
\index{g00144@GPS}
\index{a00004@Abstimmung}

\subsection{Asynchrone Lightbox-Metadaten}

Die Lightbox ist der große Fotobetrachter, der sich über der Galerieseite öffnet. Besucher können sich auf ein Foto konzentrieren und behalten dabei verfügbare Steuerungen für Vorwärts, Rückwärts, Schließen, Vollbild, Diashow, Karte und weitere Aktionen.

Die Galerie bereitet den Betrachter bei Bedarf vor und ruft Informationen für benachbarte Fotos in überschaubaren Gruppen ab. So kann ein Besucher durch eine große Galerie wechseln, ohne dass die erste Seite sämtliche Titel, Vorschauen, Kartenpunkte und Stimmen zugleich enthalten muss.\footnote{Der aktuelle Betrachter fordert Informationen gewöhnlich in Gruppen benachbarter Fotos an. Zugriffsprüfungen werden für jede Anfrage wiederholt, sodass private Galerieregeln weiterhin gelten.}

Die Pfeiltasten nach links und rechts wechseln ein Foto zurück beziehungsweise vor. \codeword{Shift+Left} und \codeword{Shift+Right} wechseln zehn Positionen für schnelleres Durchsehen großer Ergebnismengen. Der Sprung folgt derselben stabil geordneten Abfolge wie die gewöhnliche Navigation, springt an den Enden um und funktioniert auch bei einer Smart Gallery als aktiver Ergebnismenge. Die übersetzte Tastaturhilfe führt diese Kürzel auf; gesonderte grafische Sprungschaltflächen werden nicht ergänzt.

Der Betrachter beginnt gewöhnlich mit einer passend großen Vorschau. Eine größere oder originale Quelle kann verwendet werden, wenn Aktion und Zugriffsrichtlinie dies vorsehen. Benachbarte Vorschauen können unauffällig vorbereitet werden, damit das nächste Foto flüssig erscheint.

Während ein neues Foto vorbereitet wird, verwendet der Betrachter denselben Fortschrittsbereich, der das Laden voller Qualität meldet. Bei schnellem Vor- oder Zurückwechseln erhält stets das zuletzt angeforderte Foto die Zuständigkeit für diese Anzeige, sodass eine ältere Vorschau- oder Metadatenanfrage den Betrachter nicht dauerhaft beschäftigt halten oder die aktuelle Auswahl ersetzen kann. Kann ein Foto nicht angezeigt werden, zeigt der Fortschrittsbereich eine behebbare Fehlermeldung; der Besucher kann mit Vorwärts, Rückwärts, Schließen oder erneutem Öffnen fortfahren. Diese Ladezustände verändern keine Regeln für Galeriezugriff, geschützte Medien, Smart Galleries, Karten, Diashow oder Nutzung ohne JavaScript.

\subsection{Ein Foto vergrößern und verschieben}
\index{l00164@Lightbox!z00279@Zoom}
\index{z00280@Zoomsteuerungen}
\index{z00281@Zweifingergeste}

Der Zoombereich der Lightbox beträgt 100--400~\%. Verwenden Sie Minus und Plus oder wählen Sie die Prozentanzeige, um auf 100~\% zurückzukehren. Tastenkürzel sind \codeword{+} oder \codeword{=} zum Vergrößern, \codeword{-} oder \codeword{\_} zum Verkleinern und \codeword{0} zum Zurücksetzen. Fokusindikatoren und aktueller Prozentwert bleiben für assistive Technik zugänglich.

Am~Computer verschieben normales Mausrad-Scrollen und~Scrollen mit zwei Fingern auf~dem Trackpad ein~vergrößertes Foto in~beiden Achsen; bei~100~\% ändern sie den~Maßstab nicht. Eine~Spreizgeste auf~dem Trackpad oder~\codeword{Strg} mit dem~Mausrad vergrößert das~Foto um~den Zeiger; \codeword{Command} mit dem~Mausrad bleibt eine~Browserfunktion. Auf~Touchscreens vergrößert eine~Zweifingergeste um~ihren Mittelpunkt, während ein Finger das~vergrößerte Foto verschiebt. Bei~100~\% blättert horizontales Wischen mit einem Finger zwischen Fotos; der~Übergang von~Wischen zu~Pinch oder~von~Pinch zu~Verschieben löst keine Navigation aus.

Bei 100~\% ist das Foto eingepasst und im Betrachter zentriert. Vollbild verwendet dasselbe Zoom- und Verschiebeverhalten; Schließen und andere Steuerungen bleiben bei vergrößertem Bild anklickbar.

Zoom über 100~\% fordert sofort die geschützte Quelle voller Qualität für das aktive Foto an. Die vorhandene Vorschau bleibt sichtbar, bis das größere Bild bereit ist; aktuelle Zoom- und Verschiebeposition bleiben erhalten. Im normalen, Vollbild- und mobilen Vollbildmodus zeigt diese Übertragung dieselbe Fortschrittsleiste mit Prozentwert und Byteanzahl. Fotowechsel, Schließen des Betrachters oder eine neuere Anfrage voller Qualität brechen die vorherige Übertragung ab. Eine fehlgeschlagene oder abgebrochene Anfrage voller Qualität belässt die geschützte Vorschau.

Navigation, Diashow-Start, Schließen und erneutes Öffnen setzen den Betrachter auf zentrierte 100~\% zurück. Ein- oder Austritt aus dem Vollbild erhält den aktuellen Maßstab, solange das Foto selbst unverändert ist, und passt die Verschiebung an den neuen sichtbaren Bereich an. Präferenzen für reduzierte Bewegung entfernen den kurzen einzelnen Zoomübergang.

Zoom betrifft nur die Präsentation. Er verändert das gespeicherte Foto nicht und umgeht keine Medienautorisierung. Regeln für Passwort, Freigabelink, NSFW, Smart Galleries, Karten, Abstimmung, Seitennavigation und Lightbox-Metadaten gelten weiterhin.

\subsection{EXIF- und Kartenrichtlinie}

EXIF sind Informationen, die viele Kameras und Telefone in einem Foto speichern. Sie können Aufnahmedatum, Kamera, Objektiv, Belichtung, Abmessungen, Orientierung und GPS-Koordinaten enthalten. PHP Gallery kann sie für Datumsvorschläge, Karten, technische Details, Sortierhilfen und Duplikatprüfung verwenden.

Karten laden auf Anforderung des Besuchers; das spart Arbeit für Besucher, die nur Fotos betrachten möchten. Eine Galerie kann einzelne Fotoorte und eine Übersicht verfügbarer Punkte zeigen. Einstellungen übergeordneter und untergeordneter Galerien bestimmen, welche Galerien diese Informationen offenlegen. Die globale Einstellung zur öffentlichen EXIF/GPS-Anzeige dient als Vorgabe für Galerien mit geerbter Einstellung. Die Administration zeigt auch die Anzahl von Galerien mit ausdrücklicher Überschreibung und kann alle diese Überschreibungen in einer bestätigten Einstellungsspeicherung auf Vererbung zurücksetzen. Das Zurücksetzen ist schemagesichert: Lässt sich der Überschreibungszustand nicht überprüfen, verweigert die Anwendung die Änderung und verweist den Administrator zum Systemzustand, statt zu raten.

Die Auswahl von \ui{Foto öffnen} in einer Galeriekartenmarkierung öffnet dieses Foto in der aktuellen Lightbox, wenn es zur aktiven Galerie gehört. Das funktioniert auch außerhalb der aktuell angezeigten Seite: Der Betrachter fordert nur ein begrenztes Metadatenfenster um das ausgewählte Foto an. Im geteilten Vollbild-Kartenmodus bleibt die Karte geöffnet, während sich der Fotobereich ändert. Bei einer normalen Kartenüberlagerung schließt sich die Karte und zeigt das ausgewählte Foto im Betrachter. Kann der erweiterte Betrachter die Auswahl nicht auflösen, folgt dieselbe Steuerung der kanonischen Fotoseite; sie verlinkt nicht direkt auf die rohe Bilddatei.

\securitynote{GPS-Daten können sensible Orte offenlegen. Administratoren sollten geerbte Karteneinstellungen, Bildmetadaten und öffentlichen Zugriff prüfen, bevor sie persönliche Fotos veröffentlichen.}

Eine Galerie mit gespeicherter Flugroute kann ihre Karte bereits vor dem ersten Fotoupload öffnen, auch auf Seiten ohne Lightbox-Markup. Die Karte umfasst alle gespeicherten Routenpunkte und behält die normalen Zurücksetzen-Steuerelemente; Routenpunkte werden nicht als Fotomarker behandelt. Ohne gespeicherte Koordinaten erscheint die normale Leerkartenmeldung. Galeriezugriff und wirksame Karteneinstellungen gelten weiter; die Karte lädt auf Anforderung.

\subsection{Flugkarten und Navigationsdaten}
\index{s00224@SimBrief}
\index{f00104@Flugkarten}
\index{n00176@Navigationsdaten}

Flugbezogene Sammlungen können die optionale SimBrief-Integration verwenden, um den neuesten operationellen Flugplan für eine Pilot ID oder einen Pilotennamen abzurufen. Der kompakte Editor akzeptiert eine sichtbare Kennung: Nur Ziffern wählen den Pilot-ID-Weg; anderer Text den Pilotennamenweg. Vorhandene Integrationen mit getrennten Feldern bleiben kompatibel; bei beiden Angaben wird die Pilot ID bevorzugt. Der Editor erstellt einen bearbeitbaren Markdown-Beschreibungs- und Routenentwurf; OFP und verfügbare Routenkoordinaten werden erst nach dem Speichern der Galerie angehängt. Eine fehlgeschlagene SimBrief-Anfrage lässt das normale Beschreibungsfeld und den manuellen Routeneditor verfügbar. Der Entwurf vor der Erstellung ist privat und vorübergehend, wie im Abschnitt zur Galerieerstellung beschrieben.

Die SimBrief-Vorschau füllt den vorhandenen Routeneditor mit der vollständigen eingereichten Route. Fehlende Abflug- und Zielflughäfen werden jeweils einmal ergänzt; Routenanweisungen und Flughafen-/Pistenangaben bleiben erhalten. Beschreibungsentwürfe in allen gepflegten Sprachen behalten die gesamte Route einschließlich des Ziels bei langen Routen. Der Status erklärt, dass OFP und verfügbare Kartenpunkte beim Speichern der Galerie angehängt werden; der Import allein speichert sie nicht. Ändern Sie während der Anfrage Route, Kennung, Beschreibung oder Ausgangssprache, bleiben Ihre Formularwerte erhalten und die Antwort fordert einen neuen Import an. Nach dem Speichern behält die Karte die ursprünglichen OFP-Koordinaten.

Ist für~eine Galerie ein lokales SimBrief-OFP-PDF gespeichert, erscheinen \ui{Flugdurchführungsplan (OFP) anzeigen} und~\ui{OFP-PDF herunterladen} in~ihrer Beschreibung, auch ohne Fotos. Das~Dokument öffnet sich in~einer eigenen PDF-Lightbox; die~Seitennavigation wechselt niemals zu~Galeriefotos. Standardmäßig passt \ui{Ganze Seite anpassen} die~PDF-Seite gleichzeitig an~Breite und~Höhe des~verfügbaren Bereichs zwischen Kopfzeile und~Werkzeugleiste an~und zentriert sie. \ui{An Breite anpassen} erlaubt vertikales Scrollen; \ui{Originalgröße (100\%)} verwendet den~üblichen PDF-Bildschirmmaßstab. Der~aktive Modus wird hervorgehoben. Vergrößern/Verkleinern, Zurücksetzen auf~100\% des~zuletzt gewählten Anpassungsmodus, Finger-Zoom und~Verschieben, Tastaturnavigation sowie Vollbild bleiben verfügbar.

\paragraph{Historische OFP-Anhänge mit Prüfung vor dem Dispatch.}
\index{s00224@SimBrief!h00030@historischer OFP}
Öffnen Sie für~eine bestehende Galerie ohne gültiges OFP-PDF die~Registerkarte \ui{API} im~Administrator-Editor und~wählen Sie \ui{Historisches OFP erstellen / SimBrief vorausfüllen}. Zuerst erscheint ein~modales Prüfungsfenster mit der~Galeriekennung sowie bearbeitbaren ICAO-Codes für~Start und~Ziel, Datum, UTC-Abflugzeit, Flugzeugtyp, Flugangaben, Route, Flughöhe, Ausweichplatz und~Passagierzahl. Gespeicherte OFP-Daten werden als~bestätigt gekennzeichnet. Daten aus~Routenkarten und~ein~Galeriedatum sind ausdrücklich unsicher; Unbekanntes bleibt leer. Titel und~EXIF-Zeitpunkte werden nicht als~Flugdaten interpretiert. \ui{Abbrechen / Zurück} löst keine Navigation oder~Änderung aus.

Das~Prüfungsfenster funktioniert auch über dem~geöffneten rechten Admin-Seitenbereich, ohne ihn zu~schließen. Die~Felder bleiben bearbeitbar, \codeword{Tab} wechselt zwischen ihnen, und~\codeword{Escape} oder~\ui{Abbrechen / Zurück} schließt nur das~Prüfungsfenster und~bringt den~Fokus zum~Editor zurück.

Auch wenn bereits ein~gültiges PDF angehängt ist, öffnet \ui{Überarbeitetes OFP in SimBrief vorbereiten} denselben bearbeitbaren Prüfdialog, ohne das~gespeicherte PDF zu~verändern. Erzeugen und~laden Sie das~neue PDF zuerst auf~SimBrief herunter; laden Sie es~anschließend über \ui{Vorhandenes OFP-PDF ersetzen} mit ausdrücklicher Bestätigung hoch. Die~Auswahl \ui{Herkunft des hochzuladenden PDFs} erfasst ausschließlich die~Herkunft; sie~erzeugt oder~lädt kein Dokument herunter. Das~Datumsfeld erhält eine~eigene vollständige Zeile, damit die~native Datumauswahl auch bei~schmalem Bildschirm im~Dialog bleibt.

Geben Sie~die~Reiseflughöhe in~Hunderten Fuß ein, beispielsweise \codeword{350}; die~Weiterleitung sendet das~dokumentierte Format \codeword{fl=FL350}. Ohne Anmeldung erscheint zuerst die~SimBrief-Anmeldeseite. Diese Seite bestätigt nicht die~Annahme eines historischen Datums; prüfen Sie~das~Datum nach der~Anmeldung und~vor der~Generierung.

Erst \ui{SimBrief Dispatch öffnen} übergibt die~geprüften Optionen URL-kodiert an~die~offizielle HTTPS-Dispatchseite in~einem neuen, geschützten Tab. Anmeldung und~\ui{Generieren} erfolgen ausdrücklich bei~SimBrief, ohne Entwicklerschlüssel, OAuth oder~Hintergrundabfragen. Historische Daten verwenden DDMMMYY; SimBrief kann frühere Daten jedoch ablehnen oder~ändern. Ein~neu erzeugter retrospektiver Plan ist nicht das~ursprüngliche OFP vom Flugtag: Wetter, NOTAMs, AIRAC, Strecke und~Treibstoff können abweichen.

Nach dem Download fügen Sie im~Register \ui{API} mit \ui{OFP-PDF hochladen} genau ein~geprüftes PDF bis~25\,MiB an. Die~Herkunft wird als~retrospektiv oder~manuell bereitgestellt dokumentiert. Ein~gültiges vorhandenes PDF bietet \ui{OFP anzeigen}, \ui{PDF herunterladen} sowie die~separate, eingeklappte und~bestätigungspflichtige Aktion \ui{Vorhandenes OFP-PDF ersetzen}. Das~kanonische Manifest speichert die~Herkunft; ursprüngliche JSON-Daten, Fotos, Galeriedatum, Beschreibung, Route und~Sichtbarkeit bleiben unverändert. Es~gibt keinen Stapelauftrag.

Das~gespeicherte PDF wird über die~galeriegebundene Route \codeword{gallery\_ofp\_pdf} auf~derselben Browser-Herkunft wie die~Galerie ausgeliefert. Die~\codeword{index.php}-Abfragelinks funktionieren auch in~Unterverzeichnissen ohne URL-Umschreibung. Ein~neuer anonymer Besucher kann das~OFP einer öffentlichen oder~unveröffentlichten Galerie über deren direkten Link ohne Administratoranmeldung anzeigen und~herunterladen. Unveröffentlicht bedeutet hier nicht gelistet, aber nicht privat: das~Original-PDF unterliegt denselben Passwort-, Freigabetoken- und~NSFW-Zugriffsregeln wie die~Galerieseite. Wirklich private Quellgalerien ohne gültige Besucherberechtigung und~private Untergalerien mit erzeugten OFP-Seiten bleiben geschützt. Ein~berechtigter Administrator darf private Originale außerhalb der~anonymen Vorschau öffnen. Verweigerter Zugriff, fehlende und~ungültige PDF-Dateien liefern dieselbe nicht zwischengespeicherte Textantwort \codeword{404}, ohne HTML-Galerieseite oder~Dateipfad offenzulegen. Frühere \codeword{media}-URLs mit \codeword{ofp=1} bleiben kompatibel.

Die~kompakten Schaltflächen \ui{Flugdurchführungsplan (OFP) anzeigen} und~\ui{OFP-PDF herunterladen} stehen in~einer Zeile. Administratoren können über eine~separate Schaltfläche darunter optional eine~private Untergalerie der~PDF-Seiten erstellen. Eine~lokalisierte Hilfe mit Fragezeichen erläutert die~Umwandlung bei~Mauskontakt, Tastaturfokus oder~Antippen, ohne den~Beschreibungstext zu~verlängern.

Automatische Ganzseiten- und~Breitenanpassung werden für~die~tatsächlichen Maße und~die~Ausrichtung jeder PDF-Seite sowie nach Fenstergrößenänderungen, Gerätedrehung und~Vollbildwechsel neu berechnet. Manueller Zoom bewahrt dagegen seine CSS-Pixelgröße und~Scrollposition bei gewöhnlichen Größenänderungen; die~Wahl eines Anpassungsmodus oder~das Zurücksetzen stellt die~automatische Anpassung wieder her. Die automatische Berechnung berücksichtigt keine vorherigen Ränder für manuelles Verschieben; Zurücksetzen sowie Breiten- und Ganzseitenanpassung halten das PDF auch nach einer Größenänderung des gezoomten Dokuments lesbar und proportional. Bei einem Darstellungsfehler bleibt der~Link zum geschützten Original-PDF zugänglich. Sichtbarkeit, Passwort- und~NSFW-Zugriffsregeln der~Galerie gelten weiterhin für~Anzeige und~Download. Der~obligatorische Chromium-CI-Lauf bewahrt die~ausführliche Datei \codeword{browser-map.log} als~Artefakt zur~Diagnose fehlgeschlagener Dokumentanzeigetests auf.

Die~Taste \codeword{F} und~die~Vollbild-Schaltfläche steuern nur~den OFP-Dokumentbetrachter. \codeword{Escape} beendet zuerst das~Vollbild und~schließt das~Dokument erst in~der~normalen Ansicht. Eingaben in~bearbeitbaren Feldern, Tastenkombinationen mit Modifikatoren und~automatische Tastenwiederholungen lösen keinen Vollbildwechsel aus. Verweigert der~Browser die~native Vollbilddarstellung, bleibt derselbe PDF-Betrachter mit einem CSS-Vollbildlayout geöffnet. Im~Vollbild liegen Kopfzeile und~Werkzeugleiste über dem~PDF-Bereich, werden nach kurzer Zeigerinaktivität ausgeblendet und~bei Zeigerbewegung oder~Tastatureingabe wieder eingeblendet. Bedienelemente unter dem~Zeiger oder~mit Tastaturfokus bleiben sichtbar; das~Einblenden verändert weder PDF-Skalierung noch Scrollposition.

Normales Scrollen mit dem~Mausrad oder~zwei Fingern auf~dem Trackpad verschiebt ausschließlich die~PDF-Seite, auch horizontal und~im Modus \ui{An Breite anpassen}; der~Zoom von~Dokument, Bedienelementen und~Browser bleibt unverändert. \codeword{Strg} mit dem~Mausrad oder~eine~Spreizgeste auf~dem Trackpad zoomt das~PDF um~den Zeiger, während \codeword{Command} mit dem~Mausrad dem~Browser vorbehalten bleibt. Auf~Mobilgeräten verschiebt ein Finger das~PDF, zwei Finger zoomen und~ein verbleibender Finger kann nach einem Pinch weiterschieben. Die~PDF-Werkzeugleiste nutzt feste, kontrastreiche dunkle Farben unabhängig vom~Galeriedesign und~behält im~normalen und~Vollbildmodus ihre Größe. \ui{Ganze Seite anpassen}, \ui{An Breite anpassen} und~\ui{Originalgröße (100\%)} stellen die~automatische Ansicht wieder her.

Manueller Routentext wird vorrangig offline aufgelöst. Zuerst werden importierte OurAirports-Datenbankzeilen geprüft, dann die kleine mitgelieferte CSV-Ersatzquelle. Als Nächstes kann ein gültiges zwischengespeichertes Anbieterergebnis dienen; hat ein Administrator ein konfiguriertes Navigraph-Konto verknüpft und ist entfernte Suche erlaubt, kann eine unbekannte Kennung über Navigraph aufgelöst und unter dem bekannten AIRAC-Zyklus zwischengespeichert werden. Die Navigationsdatenverwaltung zeigt lokale Datenbank- und mitgelieferte Anzahlen, bietet einen Suchtester für Kennungen wie Flughäfen/Funkfeuer/Wegpunkte und kann den lokalen OurAirports-Flughafen-/Funkfeuerdatensatz aktualisieren. Navigraph-OAuth-/Kontospeicher ist optional und kann getrennt werden, ohne lokale Daten zu entfernen.

Bei aktiviertem JavaScript prüft auch die Eingabe oder Fokussierung einer nicht leeren Route, das Anfordern eines SimBrief-Entwurfs oder das Speichern eines Formulars mit Route die Aktualität lokaler Navigationsdaten im Hintergrund. Sowohl \ui{Flugkarten} als auch \ui{Navigationsdaten} müssen aktiviert sein. OurAirports-Daten werden erneuert, wenn der letzte erfolgreiche Import mindestens eine Woche zurückliegt; zwischen fehlgeschlagenen Versuchen liegt eine Stunde. Das Speichern der Galerie wartet nicht auf diese Anfrage; Editorentwurf, geöffnetes Seitenpanel und Browseradresse bleiben erhalten. Auch bei fehlgeschlagener Aktualisierung kann die Galerie gespeichert werden.

Nach der Prüfung können gespeicherte manuelle Routen fehlende Koordinaten aus aktuellen Navigationsdaten ergänzen. Vorhandene Koordinaten und SimBrief-OFP-Geometrie bleiben erhalten; eine neuere Routenänderung oder Löschung verhindert, dass das Hintergrundresultat diesen Zustand überschreibt. Nur gespeicherte Routen werden vervollständigt; ungespeicherter Browsertext wird mit der Aktualisierung nicht übertragen. Betroffene öffentliche Karten werden über den gemeinsamen Seitenpanel-Ablauf aktualisiert. Import und Routenschreibzugriffe benötigen verifizierten Datenbankspeicher; fehlender oder unklarer Speicher verhindert den optionalen Schreibzugriff, während gewöhnliche Galeriebearbeitung verfügbar bleibt.

Die öffentliche Kartendarstellung ruft niemals SimBrief, Navigraph oder einen anderen aktiven Planungsdienst auf. Sie stellt nur bereits bei der Galerie gespeicherte Koordinaten dar. Gespeicherte SimBrief-OFP-Koordinaten bleiben daher nutzbar, selbst wenn der entfernte Anbieter später ausfällt. Manueller Routentext kann die Kennungssuche auch vollständig durch ausdrückliche \codeword{NAME@latitude,longitude}-Punkte umgehen.

\subsection{Abstimmung}

Öffentliche Karten und Lightbox-Steuerungen teilen den Bildabstimmungszustand. Die aktuelle öffentliche Stimme ist eine umkehrbare positive Stimme: Erneutes Drücken der aktiven positiven Stimme entfernt sie. Die angezeigte Bewertung ist die Summe gespeicherter Bildstimmen; die aktuelle Implementierung besitzt keinen gesonderten Endpunkt für Abstimmungen über ganze Galerien. Angemeldete Stimmen sind dem Konto zugeordnet, anonyme Stimmen verwenden den datenschutzfreundlichen Besucherhash der Anwendung. Neue positive Stimmen sind ratenbegrenzt; der Widerruf einer bestehenden Stimme ist sofort erlaubt. Serverseitig dargestellte Formulare ermöglichen direkte Bedienung; JavaScript übermittelt normale Interaktionen asynchron und synchronisiert die Bewertung zwischen passenden Ansichten. Der Server prüft Bildkennung, Galerierichtlinie, Zugriff, Schemabereitschaft, CSRF und Anfrageintegrität.

\subsection{Picture Game}
\index{p00187@Picture Game}
\index{b00047@Bildvergleich}

Picture Game ist ein optionaler öffentlicher Vergleichsmodus. Ein Administrator aktiviert ihn für einen Galeriezweig; geeignete öffentliche Fotos dieser Galerie und ihrer Nachkommen können dann gemäß gewöhnlichen Zugriffs- und Sichtbarkeitsregeln teilnehmen. Das öffentliche Spiel zeigt zwei Fotos nebeneinander, vermerkt das gezeigte Paar und lässt den Besucher das bevorzugte Bild wählen, bevor der nächste Vergleich folgt. Die Galerie kann auch führende, durch das Spiel angesammelte Ergebnisse zeigen.

Weil bereits das Zeigen eines neuen Paars dauerhaften Spielzustand erzeugt, bevor der Besucher einen Sieger wählt, erfordert Picture Game ein bekanntes und verfügbares Datenbankschema. Lässt sich der erforderliche Speicher nicht prüfen, wird der Schreibvorgang verweigert, statt still auf einen nicht aufgezeichneten Vergleich auszuweichen. Der globale Funktionsschalter kann Picture Game ausblenden, ohne vorhandene Spieldaten oder galeriebezogene Aktivierungseinstellungen zu löschen.

\section{Fotoduplikatdetektor}
\label{sec:duplicates}
\index{f00107@Fotoduplikatdetektor}
\index{d00081@Duplikatprüfliste}
\index{s00219@SHA-256}

Der administrative Fotoduplikatdetektor analysiert standardmäßig einen ausgewählten Galeriezweig und kann auf globalen Umfang erweitert werden, wenn der Administrator die entsprechende Option aktiviert. Er meldet exakte und mögliche Duplikatpaare, bietet Galerie- und Fotokontextlinks, vermerkt Prüfentscheidungen und übergibt bestätigte Löschungen an den normalen Galeriebild-Löschdienst.\footnote{Der Detektor bleibt ein Prüfwerkzeug. Anhaltspunkte unterstützen das Urteil des Administrators, weil übereinstimmende Dateigröße und Fotometadaten in manchen Abläufen unterschiedliche Fotos beschreiben können.}

\subsection{Kategorien von Anhaltspunkten}

\subsubsection{Exakte und mögliche Funde}

Ein exaktes Duplikat enthält denselben Dateiinhalt. Die Anwendung erkennt dies über eine Prüfsumme, die wie ein sehr detaillierter digitaler Fingerabdruck aus der Datei berechnet wird. Übereinstimmende Fingerabdrücke sind ein starker Hinweis darauf, dass zwei Dateien dieselben Bytes enthalten.

Ein mögliches Duplikat besitzt unterstützende Ähnlichkeiten ohne exakte Fingerabdruckübereinstimmung. Der Detektor kann Dateigröße und verfügbare Kamerainformationen wie Abmessungen, Aufnahmezeit, Kamera, Objektiv, Belichtung, Blende, Brennweite, ISO und Orientierung vergleichen. Zwei getrennt exportierte Kopien eines Fotos können als Dateien unterschiedlich sein und dennoch genügend gemeinsame Informationen für eine Prüfung behalten.

Fehlende Kamerainformationen liefern dem Detektor lediglich weniger Hinweise. Die Ergebnisse bleiben Vorschläge für menschliche Prüfung. Öffnen Sie beide Fotolinks, vergleichen Sie Galeriekontext und Qualität und löschen Sie nur die Kopie, deren Entfernung zur Sammlung passt.

\subsection{Dauerhafte Prüfliste}

Die Prüfliste gehört dem authentifizierten Administrator. Paareinträge speichern kanonische Bildkennungen, die kleinere Kennung in der unteren und die größere in der oberen Spalte. Galerieeinträge gelten für die genau gespeicherte Galerie. Entscheidungen für übergeordnete und untergeordnete Galerien bleiben unabhängig.

Verfügbare Prüfaktionen sind:

\begin{itemize}
  \item das ausgewählte Paar ignorieren;
  \item alle Funde aus der genauen Galerie des linken Bilds ignorieren;
  \item alle Funde aus der genauen Galerie des rechten Bilds ignorieren;
  \item die Prüfliste des aktuellen Administrators leeren;
  \item ein bestätigtes Duplikat durch validierte normale Bildlöschung löschen.
\end{itemize}

AJAX-Aktionen aktualisieren das Detektorfragment im bestehenden rechten Administrations-Seitenbereich. Direkte POST-Anfragen verwenden Weiterleitung als Ersatzweg. Serverseitig verwaltete Scanaufträge bewahren den Umfang zwischen Fortsetzungsanfragen; eine Löschung entfernt das betroffene Bild aus den aktuellen Auftragsergebnissen.
\section{Zugriffskontrolle und Sicherheit}
\label{sec:security}
\index{s00221@Sicherheit}
\index{c00057@CSRF}
\index{a00023@Authentifizierung}
\index{n00179@NSFW}

\subsection{Mehrstufige Zugriffsauswertung}

Die Privatsphäre wird geprüft, bevor geschützte Ausgaben zusammengestellt werden. Galeriesichtbarkeit und geerbte Zugriffsregeln werden mit Fotosichtbarkeit, Passwort-/Freigabelinkzustand und gegebenenfalls Altersbeschränkungsregeln kombiniert. Dieselbe Autorisierung gilt für Galerieseiten, Vorschaubilder, Vorschauen, Originaldateien, Lightbox-Metadaten, Karten, Suchergebnisse und Downloads.

Administrative Vorgänge erfordern ein authentifiziertes Konto. Formulare enthalten ein privates Anfragetoken; der Server prüft Zielgalerie, Foto und angeforderten Vorgang, bevor er dauerhafte Daten verändert.\footnote{Der technische Name des privaten Formulartokens ist CSRF-Token. Galeriebetreiber begegnen ihm gewöhnlich nur über normale Administrationsformulare.}

\subsection{Passwörter und Freigabelinks}
\index{p00184@Passwortschutz}
\index{f00111@Freigabelinks}

Geschützte Galerien speichern Passworthashes und gewähren nach erfolgreicher Prüfung sitzungsgebundenen Zugriff. Freigabelinks verwenden serverseitig erzeugten Tokenzustand und eine Ablaufrichtlinie. Passwortzurücksetzung verwendet einmalige gehashte Tokens, eine optionale Wieder\-her\-stel\-lungs-E-Mail, eine konfigurierte Gültigkeitsdauer und den ausgewählten E-Mail-Transport. Verwenden Sie HTTPS, starke Administratorzugangsdaten, eingeschränkte Datenbankkonten und geschützte E-Mail-/API-Geheimnisse.

Der Schutz von Administrator-Cookies, das Schema erzeugter URLs und die HTTPS-Prüfung für Besucherkonten verwenden dieselbe Transportrichtlinie. Direktes TLS wird ohne Proxy-Einstellungen erkannt. Beendet ein Reverse-Proxy TLS und leitet HTTP an PHP weiter, konfigurieren Sie vor der Aktualisierung seine genaue IP-Adresse oder sein CIDR-Netz in \codeword{security.trusted_proxies} und aktivieren Sie den verwendeten Protokoll-Header in \codeword{security.trusted_proxy_protocol_headers}. Unterstützt werden \codeword{x-forwarded-proto} und \codeword{x-forwarded-ssl}; Vertrauen in Client-IP-Header wird getrennt konfiguriert. Der Proxy muss vom Client gelieferte Werte aktivierter Header überschreiben. Nicht vertrauenswürdige, ungültige, mehrfache oder widersprüchliche Protokollangaben belegen kein HTTPS. Konfigurierte Proxys und direktes HTTPS bleiben unterstützt; weitergeleitete Header allein stellen öffentliche URLs nicht mehr auf HTTPS um und setzen nicht das Secure-Attribut des Administrator-Cookies.

\subsection{Besucherkonten und Registrierung nur auf Einladung}
\index{b00036@Besucherkonten}
\index{b00034@Besuchereinladungen}
\index{b00037@Besucherkontoverwaltung}
\index{b00033@Besucheranmeldung}

Besucherkonten sind eine optionale Identitätsebene für personalisierte Galeriefunktionen. Sie sind vom Administratorkonto getrennt und entsperren allein keine private oder passwortgeschützte Galerie. Vorhandene Galeriepasswörter, Freigabelinks, Smart-Gallery-Zugriff und Medienautorisierung entscheiden weiterhin darüber, was ein Besucher sehen darf.

Das vollständige Besucherkontosystem ist durch einen Hauptschalter unter \ui{Administration > Funktionen} geschützt. Diese Hauptfunktion ist standardmäßig deaktiviert. Solange sie aus ist, zeigt die öffentliche Galerie keine Besuchersteuerungen für Anmeldung/Konto, Besucherrouten sind nicht verfügbar und der Eintrag \ui{Besucherkonten} ist im administrativen Kontomenü verborgen. Vorhandene Besucheridentitäten, Favoriten und private Sammlungen bleiben gespeichert; der Schalter steuert die Erreichbarkeit, statt Daten zu löschen.

Nach Aktivierung der Hauptfunktion kann der Administrator \ui{Konto > Besucherkonten} öffnen. Diese Seite behält den getrennten Modusschalter \ui{Besucherkonten (nur auf Einladung)}. Der Modusschalter steuert Besucheranmeldung und Einladungen innerhalb des aktivierten Systems; der Hauptschalter unter Funktionen bestimmt, ob das System überhaupt zugänglich ist.

Ein Administrator verwaltet Besucheridentitäten unter \ui{Konto > Besucherkonten}. Der bestehende Einladungsablauf bleibt verfügbar: Der Administrator kann eine Einladung optional an eine vorgesehene E-Mail-Adresse binden und ihren einmal angezeigten geheimen Link kopieren. Das Öffnen des Links erstellt kein Konto. Der Empfänger übermittelt eine E-Mail-Adresse, erhält eine Bestätigungsnachricht über den konfigurierten E-Mail-Transport, bestätigt im Browser und wählt anschließend ein Passwort mit mindestens 15 Zeichen. Es gibt keine allgemeine Seite \ui{Registrieren} oder \ui{Konto erstellen}.

Dieselbe Administrationsseite kann auch direkt ein Besucherkonto erstellen. Geben Sie die Besucher-E-Mail ein und entweder ein temporäres Passwort gemäß der Besucherpasswortrichtlinie oder lassen Sie das Passwort leer, damit die Galerie ein starkes erzeugt. Das temporäre Passwort wird nach der Erstellung einmal angezeigt. Wird die optionale Benachrichtigung zur Kontoerstellung versendet, enthält die Nachricht die vertrauenswürdige Anmeldeadresse und Anweisungen, aber niemals das temporäre Passwort; übermitteln Sie dieses separat über einen vertrauenswürdigen Kanal. Bei aktivierter Hauptfunktion bleibt direkte Kontoerstellung auch bei deaktiviertem untergeordnetem Besucher-Frontendmodus verfügbar, damit Konten vorab vorbereitet werden können. Diese Besucher können sich jedoch erst anmelden, wenn der Besuchermodus aktiviert ist.

Ein direkt angelegter Besucher kann das temporäre Passwort nicht als dauerhafte Anmeldung verwenden. Nach der ersten erfolgreichen Zugangsdatenprüfung wird der Browser zur privaten Erstanmeldeseite weitergeleitet und muss ein anderes richtlinienkonformes Passwort wählen. Bis dieser Austausch gelingt, stellt die Galerie keine normale Besucheridentität her, gibt kein \ui{Angemeldet bleiben} aus und erlaubt keine Berechtigung für Besucherfavoriten/-sammlungen. Die normale Zurücksetzung bei vergessenem Passwort kann die temporären Zugangsdaten ebenfalls über ihren unabhängig geprüften Wiederherstellungsablauf ersetzen. Eine gleichzeitige Administratoranmeldung im selben Browser bleibt unabhängig.

Vorhandene Besucherkonten werden auf der Administrationsseite mit Kontozustand und Sicherheitssteuerungen aufgeführt. Ein aktiver Besucher kann mit \ui{Sperren} gesperrt oder überall abgemeldet werden; ein gesperrter Besucher kann mit \ui{Wiederherstellen} reaktiviert werden. Das Sperren blockiert Besucherauthentifizierung und macht über den zentralen Besucherlebenszyklusdienst bestehende Besuchersitzungen, dauerhafte Anmeldung, Zurücksetzungen, ausstehende Bestätigungen/E-Mail-Änderungen und vorbereitete Sammlungsfreigabeberechtigungen ungültig. Das Wiederherstellen aktiviert das Konto erneut, belebt aber keine vor der Sperre vorhandenen Zugangsdaten oder Sitzungen wieder. Der Besucher muss sich daher erneut normal authentifizieren. Ein erzwungener Passwortaustausch bei Erstanmeldung bleibt über Sperren/Wiederherstellen hinweg erforderlich. \ui{Überall abmelden} belässt ein aktives Konto aktiv, widerruft aber seine Besucheranmeldeberechtigung auf allen Geräten. Diese Sicherheitssteuerungen bleiben dem Administrator verfügbar, solange die Hauptfunktion aktiviert und der untergeordnete Besucher-Frontendmodus deaktiviert ist. Das Ausschalten der Hauptfunktion verbirgt auch die Besucherverwaltung. Keiner dieser Schalter oder Sicherheitssteuerungen entfernt Favoriten, private Sammlungen, Sammlungseinträge, Fotos, Galerien, Smart Galleries, Galeriefreigabelinks oder Administratorauthentifizierung.

Administratives Löschen bleibt eine separate destruktive Aktion mit ausdrücklicher Bestätigung. Es entfernt Besucheridentität sowie besuchereigenen Authentifizierungs-, Lebenszyklus-, Favoriten- und Sammlungszustand über die vorbereiteten Datenbankbeziehungen. Es löscht keine Fotos, Galerien, Smart Galleries, Galeriefreigabelinks oder Administratorkonten/-sitzungen.

Die öffentliche Kopfzeile zeigt \ui{Anmelden} nur bei aktivierten Besucherkonten. Nach der Anmeldung zeigt sie \ui{Konto}. Die Kontoseite zeigt die bestätigte E-Mail, Links zu privaten Seiten \ui{Favoriten} und \ui{Sammlungen}, \ui{Passwort ändern}, \ui{E-Mail ändern}, \ui{Konto löschen} und eine ausschließlich besucherbezogene Aktion \ui{Abmelden}. Die Abmeldung eines Besuchers oder Änderung seiner Kontoberechtigung meldet keinen Administrator ab, der zufällig in derselben Browsersitzung authentifiziert ist.

\index{f00100@Favoriten}
Ein angemeldeter Besucher kann ein autorisiertes Foto über das kleine Herz auf einer Galeriekarte oder in der Lightbox als Favoriten markieren. Dasselbe Foto bleibt zwischen Karte und Lightbox synchronisiert; die private Seite \ui{Favoriten} führt gespeicherte Fotos auf, die der Besucher weiterhin öffnen darf. Das Entfernen eines Favoriten löscht nur die gespeicherte Referenz des Besuchers; das Foto wird weder gelöscht noch verändert.

Ein Favorit ist niemals eine Zugriffsfreigabe. Wird die Quellgalerie später privat, läuft ihre Passwort-/Freigabeberechtigung ab oder ist das Foto aus anderen Gründen für diese Anfrage nicht mehr autorisiert, umgeht der gespeicherte Favorit diese Regeln nicht; die Quellmetadaten erscheinen nicht auf der Favoritenseite. Ein gleichzeitig als Administrator und Besucher angemeldeter Browser benötigt weiterhin die unabhängige Quellautorisierung des Besuchers, bevor dieser das Bild als Favoriten speichern darf.

\index{b00038@Besuchersammlungen}
Ein angemeldeter Besucher kann auch private benannte \ui{Sammlungen} erstellen. Eine Sammlung ist eine geordnete Referenzliste vorhandener Fotos; sie ist keine Galerie, kopiert keine Fotos und verändert die administratorverwaltete Galeriehierarchie nicht. Auf einer autorisierten Galerie- oder Smart-Gallery-Karte oder der privaten Seite \ui{Favoriten} kann der Besucher \ui{Zur Sammlung hinzufügen} verwenden und eine eigene Sammlung wählen. Der private Sammlungsbereich unterstützt Erstellen und Umbenennen einer Sammlung, Entfernen von Fotos, Verschieben sichtbarer Fotos nach oben oder unten und Löschen der Sammlung selbst.

Sammlungseigentum und Fotozugriff sind getrennte Entscheidungen. Das Speichern eines Fotos hält nur seine interne Bildreferenz und Reihenfolge fest. Die Sammlung speichert kein Galeriepasswort, Galeriefreigabetoken, Dateisystempfad, Vorschaubildpfad oder gemerkte Zugriffsentscheidung. Bei jedem Öffnen der Sammlung bewertet die Galerie die aktuelle Quellgalerie- und Fotoautorisierung erneut. Ist ein zuvor gespeichertes Foto nicht mehr autorisiert, wird es ausgelassen, ohne Titel, Dateiname, Galerietitel, Vorschaubild, EXIF-Daten oder Ablehnungsgrund offenzulegen. Die gespeicherte Referenz kann erhalten bleiben, damit das Foto bei später wieder vorhandenem normalem Quellzugriff erneut sichtbar wird.

Diese Regel gilt auch bei gleichzeitiger Anmeldung als Administrator und Besucher im selben Browser. Administratorsichtbarkeit wird nicht zur Berechtigung für Besuchersammlungen; eine Sammlung kann niemals als zweiter Weg um Galeriepasswörter, Freigabeberechtigungen, Altersbeschränkungen oder Medienauslieferungsprüfungen dienen. Favoriten und Sammlungen sind unabhängig: Ein Foto kann in einem, beiden oder keinem Bereich liegen; das Löschen einer Sammlung löscht weder einen Favoriten noch das zugrunde liegende Foto.

\index{s00209@Sammlungsfreigabe}
Ein Besucher kann für eine eigene Sammlung einen nicht gelisteten schreibgeschützten Freigabelink erstellen. Der Link läuft nach 30~Tagen ab und kann jederzeit ersetzt oder widerrufen werden. Der vollständige geheime Link wird nur unmittelbar nach Erstellung oder Ersetzung angezeigt; später zeigt die Sammlungsseite nur die aktive Freigabe sowie Erstellungs- und Ablaufzeiten. Das Ersetzen macht den bisherigen Link unbrauchbar. Der Widerruf entfernt nur die Freigabeberechtigung und löscht weder Sammlung noch Einträge, Favoriten oder Fotos.

Ein Empfänger benötigt kein Besucherkonto. Das Öffnen des geheimen Links prüft ihn, stellt nur eine eng begrenzte Sitzungsberechtigung für diese eine Sammlung her und leitet zu einer bereinigten Adresse ohne Geheimnis weiter. Der Link bleibt bis Ablauf/Widerruf wiederverwendbar, sodass automatische E-Mail- oder Chat-Linkscanner ihn nicht verbrauchen. Die geteilte Seite ist nicht gelistet, nicht zwischenspeicherbar und gegen Suchindexierung gekennzeichnet.

Die Freigabe gewährt nur Zugriff auf den Sammlungscontainer. Sie entsperrt niemals eine Quellgalerie oder ein Foto. Bei jeder geteilten Ansicht wird jede gespeicherte Fotoreferenz erneut gegen die aktuelle Quellgalerieautorisierung des Empfängers geprüft. Eine passwortgeschützte Quelle bleibt verborgen, bis dieser Empfänger sie unabhängig über den normalen Galeriemechanismus entsperrt; eine unabhängig vorhandene Galeriefreigabeberechtigung kann sie ebenfalls normal autorisieren. Selbst bei gleichzeitiger Administratoranmeldung im selben Browser wird der Administratorausnahmepfad nicht zum Befüllen der geteilten Sammlung verwendet. Direkte Vorschaubild-/Medienrouten erzwingen weiterhin ihre eigenen bestehenden Quellregeln. Unzugängliche Einträge werden ausgelassen, ohne verborgene Titel, Dateinamen, Pfade, Galerien oder Gründe offenzulegen.

Das Sperren des Sammlungsbetreibers widerruft aktive Sammlungsfreigaben; die Wiederherstellung des Kontos belebt sie nicht wieder. \ui{Überall abmelden} widerruft nur Besucheranmeldeberechtigung und keine ansonsten gültige Freigabe. Das Löschen des Betreibers oder der Sammlung macht die Freigabe ungültig, ohne Originalfotos zu löschen. Der globale Funktionsschalter für Besucherkonten schützt auch Sammlungsfreigaben und bleibt standardmäßig deaktiviert.

\ui{Angemeldet bleiben} für Besucher verwendet eigene rotierende Besucherzugangsdaten statt des dauerhaften Administratoranmeldetokens. Ihre Wiederherstellung meldet nur den Besucher an und zählt nicht als kürzlich erfolgte Passwort-Neuauthentifizierung. Anfragen wegen vergessener Passwörter liefern eine allgemeine Antwort und verwenden zweistufige Zurücksetzungsbestätigung, damit automatische Linkscanner kein Passwort allein durch Öffnen einer URL zurücksetzen können.

Sensible Kontoänderungen erfordern einen aktuellen Nachweis des Besucherpassworts. Ein nur durch \ui{Angemeldet bleiben} wiederhergestellter Besucher wird daher aufgefordert, vor Passwortänderung, Beginn einer E-Mail-Adressänderung oder Kontolöschung das aktuelle Besucherpasswort zu bestätigen. Der Rückkehrpfad ist auf diese Kontoaktionen beschränkt und kann nicht auf eine beliebige Seite weiterleiten.

Die Änderung des Besucherpassworts verwendet dieselbe Richtlinie wie Aktivierung und Zurücksetzung. Eine erfolgreiche Änderung macht gemäß Sicherheitslebenszyklus die alte Passwortberechtigung, Besuchersitzungen und dauerhaften Besucherzugangsdaten ungültig und fordert anschließend eine erneute Anmeldung. Gespeicherte Favoriten bleiben demselben Besucherkonto zugeordnet. Ein in derselben Browsersitzung angemeldeter Administrator bleibt angemeldet.

Die Änderung der Besucher-E-Mail erfolgt stufenweise. Die Übermittlung einer neuen Adresse ersetzt nicht die aktuell bestätigte Adresse. Die Galerie sendet zunächst innerhalb der Besucher-E-Mail-Missbrauchsgrenzen einen Bestätigungslink an die neue Adresse. Das Öffnen dieses Links prüft ihn nur und bereitet einen kurzlebigen Bestätigungszustand vor; es verändert das Konto nicht. Die endgültige Änderung erfolgt erst nach Übermittlung des gesonderten geschützten Bestätigungsformulars durch den Besucher. Ein erfolgreicher Wechsel meldet den Besucher ab, damit die neue bestätigte E-Mail bei der nächsten Anmeldung verwendet wird.

\ui{Konto löschen} ist eine endgültige, ausschließlich besucherbezogene Aktion. Nach aktueller Passwortbestätigung muss der Besucher das Löschen ausdrücklich bestätigen. Der Vorgang entfernt Besucherkonto und besuchereigene Daten wie Favoriten und private Sammlungen gemäß vorbereitetem Datenbanklebenszyklus. Er löscht keine Galeriefotos, Galerien, Smart Galleries oder deren bestehende Freigabelinks und zerstört keine gleichzeitige Administratoranmeldung.

Besuchersicherheit und private Konto-/Favoriten-/Sammlungs-/Lebenszyklusseiten werden als private, nicht zwischenspeicherbare Antworten ausgeliefert. Einladungs-, Bestätigungs-, Zurücksetzungs- und E-Mail-Änderungslinks verwenden den konfigurierten vertrauenswürdigen Anwendungsursprung. Bestätigungs-, Zurücksetzungs- und E-Mail-Änderungsnachrichten für Besucher unterliegen weiterhin adress-, besucher- und globalbezogenen Missbrauchsbudgets, bevor der konfigurierte PHP-E-Mail- oder SMTP-Transport verwendet wird.

Nicht gelistete schreibgeschützte Sammlungsfreigabe ist nur innerhalb der Funktion Besucherkonten verfügbar. Öffentliche Sammlungssuche, öffentliche Besucherprofile, Bearbeitung/Zusammenarbeit durch Empfänger, Besucheruploads/-kommentare, offene Registrierung, soziale Besucheranmeldung, TOTP, Passkeys und Magic-Link-Anmeldung bleiben nicht verfügbar.

\subsection{Administratorkonto und E-Mail-Einstellungen}
\index{a00009@Administratorkonto}
\index{s00227@SMTP}
\index{g00143@Google-Anmeldung}

Der Kontobereich unterstützt Benutzername- und Passwortänderungen, eine optionale Wiederherstellungsadresse, Gültigkeitsdauer der Passwortzurücksetzung und eine Testnachricht. Wieder\-her\-stel\-lungs-E-Mail kann PHP \codeword{mail()} oder einen authentifizierten SMTP-Transport mit Einstellungen für Host, Port, Verschlüsselung, Benutzername und Passwort verwenden. Behandeln Sie SMTP-Zugangsdaten wie Datenbankzugangsdaten und nehmen Sie sie niemals in Exporte oder Diagnosebildschirmfotos auf.

Google-Anmeldung ist optional. Der Administrator konfiguriert OAuth-Client und autorisierte Callback-/Weiterleitungs-URI und verknüpft anschließend eine Google-Identität nur bei bereits vorhandener Authentifizierung durch die normale Administrator-Passwortsitzung. Eine nicht verknüpfte Google-Identität wird auf der Anmeldeseite abgelehnt und darf kein Administratorprofil erstellen oder beanspruchen. Nach der Verknüpfung kann sich nur diese gespeicherte Google-Identität über die Google-Schaltfläche für das Profil anmelden. Das Trennen erfordert das aktuelle Administratorpasswort. Externe Anmeldung ersetzt keine Passwortanmeldung; fehlt OAuth-Konfiguration oder Identitätsverknüpfungsspeicher oder lässt sich dieser nicht prüfen, bleibt gewöhnliche Passwortauthentifizierung unabhängig.

Das Anmeldeformular akzeptiert Administratorbenutzername oder konfigurierte Wieder\-her\-stel\-lungs-E-Mail. Authentifizierungsversuche werden nach Besucher und normalisierter übermittelter Kennung gedrosselt, wenn das Ratenbegrenzungsschema verfügbar ist. Die Drosselungstabelle speichert gehashte Subjekte statt unverarbeiteter IP-Adressen oder übermittelter Kennungen. Die aktuelle Richtlinie begrenzt wiederholte fehlgeschlagene Anmeldungen in einem 15-Minuten-Fenster und Passwortzurücksetzungsanfragen in längeren begrenzten Zeitfenstern; erfolgreiche Anmeldung leert den relevanten Anmeldedrosselungszustand.

Passwortwiederherstellung liefert bewusst ein allgemeines Anfrageergebnis, damit das öffentliche Formular nicht die Existenz eines übermittelten Benutzernamens/einer E-Mail bestätigt. Zurücksetzungslinks verwenden einmalige gehashte Tokens und eine konfigurierte Gültigkeit von 15 bis 1440~Minuten. Der Abschluss ändert das Passwort, markiert das Token als verwendet und widerruft dauerhafte Anmeldetokens für dieses Konto. Die Kontoseite kann eine Wiederherstellungs-Testnachricht über den konfigurierten PHP-\codeword{mail()}- oder SMTP-Transport senden.

Die Option \ui{Angemeldet bleiben} verwendet ein gehashtes Browsertoken, damit ein Administrator gewöhnliche PHP-Sitzungsbereinigung auf Shared Hosting übersteht, ohne dauerhafte Zugangsdaten im Klartext in der Datenbank zu speichern. Dauerhafte Anmeldung bleibt optional. Sie wird durch entsprechende Konto-/Abmeldesicherheitsabläufe widerrufen und kann bei nicht überprüfbarem Schema nicht verfügbar sein, ohne die aktuelle gewöhnliche PHP-Sitzung zu deaktivieren.
\subsection{Vertrag zur Datenbankbereitschaft}
\label{sec:db-readiness-contract}
\index{s00238@Systemzustand!d00063@Datenbankbereitschaft}
\index{d00062@Datenbank!b00032@Bereitschaft}
\index{a00001@503-Antwort}

Funktionen, die vom Datenbankschema abhängen, verwenden einen gemeinsamen Bereitschaftsvertrag. Entscheidend ist der Unterschied zwischen einem erfolgreich geprüften und als älter erkannten Schema und einem Schema, das derzeit nicht geprüft werden kann.

\begin{longtable}{>{\RaggedRight\arraybackslash}>{\hspace{0pt}}p{0.12\textwidth} >{\RaggedRight\arraybackslash}>{\hspace{0pt}}p{0.20\textwidth} >{\RaggedRight\arraybackslash}>{\hspace{0pt}}p{0.35\textwidth} >{\RaggedRight\arraybackslash}>{\hspace{0pt}}p{0.17\textwidth}}
\toprule
\textbf{Zustand} & \textbf{Bedeutung} & \textbf{Anwendungsverhalten} & \textbf{Betreiberaktion}\\
\midrule
\endfirsthead
\toprule
\textbf{Zustand} & \textbf{Bedeutung} & \textbf{Anwendungsverhalten} & \textbf{Betreiberaktion}\\
\midrule
\endhead
Verfügbar & Erforderlicher Speicher wurde gefunden und geprüft. & Normale Lese- und Schreibvorgänge für die Fähigkeit. & Keine.\\
Fehlend & Prüfung erfolgreich, aber ein erforderliches Objekt fehlt. & Ein dokumentierter Weg für ältere Schemata kann verfügbar bleiben; andernfalls fordert die Funktion eine Migration. & Sichern und die ausstehende Migration anwenden.\\
Unbekannt & Das Schema konnte nicht zuverlässig geprüft werden. & Geschützte Lesevorgänge können sicher verweigert oder optionale schreibgeschützte Ausgaben ausgelassen werden. Abhängige Schreibvorgänge sind blockiert. & Datenbankverbindung, ausgewähltes Schema oder Metadatenberechtigungen reparieren.\\
Deaktiviert & Eine optionale Funktion ist ausgeschaltet. & Die Funktion wird nicht verwendet. & Keine, sofern die Funktion nicht aktiviert werden soll.\\
\bottomrule
\end{longtable}

\ui{Systemzustand} und \ui{Laufzeitdiagnose} verwenden diese Zustände einheitlich. Ein Ergebnis \ui{Fehlend} weist gewöhnlich auf Migrationsarbeit hin; \ui{Unbekannt} ist ein Betriebsproblem und darf nicht als Beweis gelten, dass eine Tabelle oder Spalte nie existiert hat. Diagnosen können Fähigkeitsname, Status, validierte Datenbankobjektnamen, vorgeschlagene Prüfungen und eine Anfragereferenz zeigen. SQL-Text, unverarbeitete Datenbankausnahmen, Zugangsdaten, Tokens, Cookies und private Dateisystempfade sind ausgeschlossen.

\subsection{Sicherheitsrelevante Ausnahmen}
\index{n00180@NSFW Guard}
\index{p00185@Passwortzurücksetzung!d00063@Datenbankbereitschaft}
\index{g00143@Google-Anmeldung!d00063@Datenbankbereitschaft}
\index{f00111@Freigabelinks!d00063@Datenbankbereitschaft}

Einige Fähigkeiten besitzen strengeres oder spezifischeres Verhalten:

\begin{description}
  \item[NSFW Guard.] Bestätigt die Prüfung das ältere Schema ohne NSFW-Felder, behält die Website bis zur Migration ihr Verhalten vor Einführung der Funktion. Sind die Felder \ui{Unbekannt}, werden öffentliche Anfragen, die altersbeschränkte Medien offenlegen könnten, sicher verweigert und NSFW-Bearbeitung pausiert.
  \item[Galeriezugriff.] Eine Datenbank gilt nur dann als älteres Format vor Einführung des Schutzes, wenn die Prüfung gelingt und alle zentralen Zugriffsfelder fehlen. Ein gemischter Aufbau bedeutet eine unvollständige Migration; geschützte Routen bleiben daher bis zur Reparatur unverfügbar. Der ältere Veröffentlichungswert \codeword{draft} wird nur geschrieben, wenn die bekannte Felddefinition ihn erfordert; aktuelle Schemata verwenden \codeword{unpublished}.
  \item[Freigabelinkwiderruf.] Widerruf darf die kleinere bekannte Feldmenge verwenden, die zum Leeren des prüfenden Hashes nötig ist, weil dieser Vorgang Zugriff reduziert. Er stoppt dennoch, wenn seine eigenen erforderlichen Felder \ui{Unbekannt} sind.
  \item[Dauerhafte Administratoranmeldung.] Eine fehlende Tabelle für dauerhafte Anmeldung deaktiviert dauerhafte Browsertokens, verhindert aber nicht, dass eine gewöhnliche Passwortanmeldung die aktuelle PHP-Sitzung erstellt.
  \item[Passwortzurücksetzung.] Wiederherstellung erfordert sowohl Wieder\-her\-stel\-lungs-E-Mail-Feld als auch Zurücksetzungstokenspeicher. Fehlender Speicher liefert Migrationshinweise; \ui{Unbekannt}-Speicher liefert einen vorübergehenden Fehler statt eines falschen Konto-/Tokenergebnisses.
  \item[Google-Identitätsverknüpfungen.] Externe Anmeldung benötigt gültige OAuth-Konfiguration und bekannten Identitätsverknüpfungsspeicher. Fehlt dieser Speicher oder ist er unbekannt, bleibt Passwortanmeldung unabhängig verfügbar.
\end{description}

\subsection{Schreibvorgänge, Wartung und Zugangsdaten}
\index{d00062@Datenbank!a00014@Änderungssicherheit}
\index{u00249@Uploads!d00063@Datenbankbereitschaft}
\index{a00010@Aktualisierungen!d00063@Datenbankbereitschaft}

Vorgänge, die Dateien, Zugangsdaten oder dauerhafte Datensätze verändern, führen ihre Bereitschaftsprüfung vor dem ersten zielverändernden Schritt durch. Dazu gehören Galerie-/Fotoänderungen, Uploads, WebDAV-Schreibvorgänge, Galeriemigration, Vorschaubildwartung, Datenbankreparatur/-bereinigung und Aktivierung von Anwendungsaktualisierungen. Ein Ergebnis \ui{Unbekannt} stoppt den Vorgang vor einer Zieländerung. Ein bestätigtes Ergebnis \ui{Fehlend} darf nur dort einen Kompatibilitätsweg für ältere Schemata oder eine Migrationsinitialisierung verwenden, wo dieses Verhalten dokumentiert ist.

Zugangsdatenwiderruf folgt demselben restriktiven Prinzip wie Freigabelinkwiderruf: Eine kleinere bekannte Feldmenge darf verwendet werden, wenn diese Felder zur Zugriffsdeaktivierung ausreichen. Erstellung oder Vertrauen in Zugangsdaten erfordert weiterhin den vollständigen Speicher des betreffenden Authentifizierungswegs.

Ist ein Schreibvorgang wegen \ui{Unbekannt} blockiert, reparieren Sie vor dem erneuten Versuch die Datenbankprüfung. Wird der Zustand \ui{Fehlend}, sichern Sie und wenden Sie die erforderliche Migration an. Wird er \ui{Verfügbar}, wiederholen Sie den Vorgang einmal und prüfen Sie Dateisystem- und Datenbankergebnisse. Setzen Sie einen bestehenden wiederherstellbaren Migrations-/Aktualisierungsauftrag fort, statt einen doppelten Auftrag zu starten.

\subsection{Optionale Funktionen}
\index{k00155@Karten!d00063@Datenbankbereitschaft}
\index{p00187@Picture Game!d00063@Datenbankbereitschaft}
\index{a00004@Abstimmung!d00063@Datenbankbereitschaft}
\index{o00182@OpenAI!d00063@Datenbankbereitschaft}
\index{s00224@SimBrief!d00063@Datenbankbereitschaft}
\index{t00240@Telemetrie!d00063@Datenbankbereitschaft}

Optionale Präsentationsfunktionen können eingeschränkt werden, ohne die Hauptgalerie zu deaktivieren, wenn dabei keine geschützten Inhalte offengelegt oder Zugriffe gewährt werden. Beispielsweise darf eine Karte entfallen, wenn ihr optionaler Speicher nicht gelesen werden kann; das Speichern einer GPS-/Lightbox-Überschreibung erfordert aber weiterhin ein bekanntes relevantes Feld.

Abstimmung und Picture Game schreiben beide dauerhaften Zustand. Picture Game vermerkt ein dargestelltes Paar vor der Siegerwahl; somit ist auch das Laden eines neuen Paars eine Schreibgrenze und benötigt bekannten Speicher. KI-Warteschlangen, Telemetriewartung, Navigationstokens, gespeicherte SimBrief-Routendaten und ähnliche optionale Schreibvorgänge folgen derselben Regel. Eine Funktion darf nur dort ein schreibgeschütztes Ergebnis ohne optionale Speicherung liefern, wo die Oberfläche diese Einschränkung deutlich macht.

\subsection{Betriebliche Sicherheit}

Empfohlene Vorgehensweisen sind:

\begin{itemize}
  \item Stellen Sie die Anwendung über HTTPS bereit.
  \item Beschränken Sie Datenbankzugangsdaten auf Anwendungsschema und erforderliche Vorgänge.
  \item Halten Sie \filepath{config.php}, Sicherungen, Caches und erzeugte Archive außerhalb öffentlicher Auffindbarkeit, soweit die Hostingstruktur dies erlaubt.
  \item Wenden Sie Migrationen und Aktualisierungen über authentifizierte Wartungsabläufe an.
  \item Prüfen Sie Auditprotokolle und Sicherheitsdiagnosen.
  \item Testen Sie geschützte Galerien in anonymen Besuchersitzungen.
  \item Pflegen Sie abgestimmte Dateisystem- und Datenbanksicherungen.
\end{itemize}
\subsection{Dateigrenzen für Medien und wiederaufnehmbare Migration}
\index{Medien!Dateisystemgrenzen}
\index{Galerie-Migration!atomare Installation}

Vor der Ausgabe öffentlicher Galeriebilder, Titelbilder und galeriespezifischer Branding-Ressourcen wird die Zugriffsberechtigung erneut geprüft. Geschützte oder altersbeschränkte Galerien verwenden auch bei einer Administratoransicht private Cache-Regeln. Gespeicherte Bild- und Branding-Pfade müssen innerhalb des jeweiligen Galerieverzeichnisses bleiben; dabei wird auch das tatsächliche Ziel symbolischer Links geprüft. Globale Designhintergründe und Branding-Dateien sind auf ihre eigenen Cache-Verzeichnisse begrenzt. Öffentliche Bildantworten erlauben unterstützte Rasterformate, nicht jedoch aktive SVG-Dateien. Das ursprüngliche Designhintergrundbild ist nur für angemeldete Administratoren abrufbar; die normale öffentliche Ableitung bleibt verfügbar.

Bei der Galerie-Migration prüfen sowohl die Erstellung des Asset-Manifests als auch die authentifizierte Dateiausgabe über den API-Schlüssel die endgültige Dateiposition innerhalb der Quellgalerie. Eine empfangene Datei wird zuerst in eine eindeutige temporäre Datei im Zielverzeichnis kopiert. Vor der Umbenennung werden die Byteanzahl und ein gegebenenfalls im Manifest vorhandener SHA-256-Wert geprüft. Ein fehlgeschlagener Kopiervorgang hinterlässt keine unvollständige Zieldatei. Ein bereits installiertes identisches Asset kann bei einem erneuten Versuch wiederverwendet werden; Konflikte und symbolische Links werden abgelehnt. Vom Browser erzeugte unkomprimierte ZIP-Pakete lehnen zusätzlich doppelte normalisierte Eintragsnamen, fehlerhafte CRC-32-Werte und unvollständige oder inkonsistente Zentralverzeichnisse ab.

\section{Design, Präsentation und Lokalisierung}
\label{sec:theme}
\index{d00073@Design}
\index{l00165@Lokalisierung}

Der Designbereich steuert Websiteidentität, Farben, Abmessungen, Seitenbreite, Hintergründe, Branding, Favicon-Ressourcen, Sprache, Galeriekartenpräsentation, Lightbox-Vorgaben, GPS-Präsentation, Favoritennavigation und öffentliche Vorschaubilddarstellung. Einstellungen verwenden normalisierte skalare Werte oder gezielte Dienstmodelle.

Die Unterkarte \ui{Animationen} legt die Öffnungs- und Schließdauer für das öffentliche Informationsfenster auf Galeriekarten und das Admin-Seitenpanel fest. Beide Werte liegen zwischen 0 und 800~ms; 0 deaktiviert die jeweilige Bewegung. Die Standardwerte sind 320~ms und 260~ms. Das öffentliche Panel enthält Galerietitel, Beschreibung, Datum und die vollständige Schlagwortliste. Die native Offenlegung funktioniert auch ohne JavaScript; mit JavaScript bleibt das Panel im Ansichtsbereich und lässt sich durch Klick außerhalb oder mit Escape schließen. Die Interaktion berücksichtigt die Einstellung für reduzierte Bewegung des Browsers.

Der Designeditor verwendet keinen getrennten Hell-/Dunkelmodus-Schalter. Stattdessen bestimmen gespeicherte Farbsteuerungen und ein optionales benutzerdefiniertes CSS-Design die öffentliche Farbpalette direkt. Er legt außerdem Galeriekarten- und Untergalerielayouts, den Standardstil der Breadcrumbs und bis zu drei Kopfzeilenverknüpfungen fest.

Wählen Sie den globalen Breadcrumb-Stil unter \ui{Design > Layout > Karten und Kennzeichnungen}. Jede Option ist eine native Radiokarte mit einem echten gerenderten Breadcrumb-Beispiel, sodass die Auswahl nach dem Aussehen möglich ist. Die Karte \ui{Standard / vom Design übernehmen} in Galerie- und Smart-Gallery-Einstellungen zeigt den aktuell wirksamen Designstil. Die Radiokarten lassen sich per Tastatur bedienen und funktionieren ohne JavaScript. Die stabilen IDs lauten \codeword{minimal}, \codeword{chevron}, \codeword{pills}, \codeword{surface}, \codeword{ribbon}, \codeword{nodes}, \codeword{tabs}, \codeword{tiles} und \codeword{gradient}. Band (\codeword{ribbon}) bildet Pfeilsegmente; Verbundene Knoten (\codeword{nodes}) verknüpfen dekorative Punkte; Registerkarten (\codeword{tabs}) sind kompakt; Kacheln (\codeword{tiles}) sind erhabene Rechtecke; Farbverlauf (\codeword{gradient}) ergänzt eine Designfarbfläche mit Akzentstreifen.

Jede Kopfzeilenverknüpfung kann zur Hauptseite oder einer Galerie führen oder leer bleiben. Anonyme Besucher sehen nur konfigurierte Galerieverknüpfungen, die weiterhin öffentlich und gelistet sind; doppelte oder gelöschte Galerieauswahlen werden beim Speichern verworfen. Drei leere Plätze verbergen diese Favoritengalerie-Schaltflächen. Galeriekarten können außerdem eine konfigurierbare Kennzeichnung enthaltener Bilder anzeigen. Ist diese für eine Galerie wirksam, bleibt beim Öffnen dieselbe Zweiganzahl im Hero-Bereich sichtbar; gezählt werden Bilder der aktuellen Galerie und ihrer zugänglichen Untergalerien nach denselben Öffentlichkeits-/Zugriffsregeln wie bei der Karte. Galeriespezifische Präsentation kann Websitevorgaben erben oder die im Galerieeditor angebotenen Entscheidungen überschreiben.

\subsection{Referenz der Designsteuerungen}
\index{d00073@Design!s00236@Steuerungen}
\index{s00216@Seitennavigation!d00074@Designeinstellungen}
\index{f00099@Favicon}
\index{b00029@benutzerdefiniertes CSS}

Die aktuelle Designseite gliedert sich in \ui{Erscheinungsbild}, \ui{Branding \& Medien}, \ui{Layout}, \ui{Sprache} und \ui{Benutzerdefiniertes CSS}. Die folgende Referenz beschreibt betreibersichtbare Steuerungen und die von der aktuellen Version erzwungene Normalisierung.

\begin{longtable}{>{\RaggedRight\arraybackslash}>{\hspace{0pt}}p{0.20\textwidth} >{\RaggedRight\arraybackslash}>{\hspace{0pt}}p{0.27\textwidth} >{\RaggedRight\arraybackslash}>{\hspace{0pt}}p{0.41\textwidth}}
\toprule
\textbf{Bereich} & \textbf{Steuerung} & \textbf{Aktuelles Verhalten}\\
\midrule
\endfirsthead
\toprule
\textbf{Bereich} & \textbf{Steuerung} & \textbf{Aktuelles Verhalten}\\
\midrule
\endhead
Erscheinungsbild & Websiteidentität und Farben & Websitename und Farben für Akzent, dunklen Akzent, Seitenhintergrund, Bereich, geöffneten Galeriebereich, Kopfzeilentitel und Galerietitel. Der \glqq{}dunkle Akzent\grqq{} ist eine Hover-/Umrandungsfarbe, kein Dunkelmoduswähler.\\
Erscheinungsbild & Ecken und Schrift & Eckenradius 0--32~Pixel. Schriftmodus \ui{Klassische Serifenschrift} oder \ui{Klare serifenlose Schrift}. Die Live-Vorschau aktualisiert sich bei der Bearbeitung.\\
Erscheinungsbild & Öffentliche Seitenbreite & \ui{Standard}, \ui{Breiter}, \ui{Volle Breite} oder \ui{Benutzerdefiniert}. Benutzerdefinierte Breite wird auf 1024--2048~Pixel begrenzt und beträgt ohne gültigen Wert 1440.\\
Erscheinungsbild & GPS-Markierung & Bei aktivierter GPS-Kartenfunktion können Fotokartenmarkierungen und deren Hintergrundunterlegung separat gezeigt oder verborgen werden. Markierungsgröße 14--48~Pixel, Vorgabe 26. Der Unterlegungsregler zeigt 0--48~Pixel; die aktuelle Servernormalisierung behandelt 0 oder andere nicht positive Werte jedoch als ungesetzt und stellt 22~Pixel wieder her. Verbergen Sie den Hintergrund über das separate Unterlegungskontrollkästchen.\\
Erscheinungsbild & Öffentliche Schlagwortseiten & Eigene Spalten und Zeilen für Schlagwortseiten verwenden dieselben sicheren Höchstwerte wie andere öffentliche Raster: 1--12 Spalten und 1--50 Zeilen. Schlagwort-Ergebniskarten können unabhängig von normalen Galerielisten vertikales oder horizontales Galeriekartendesign verwenden.\\
Erscheinungsbild & Galerie-Hero-Schlagwörter & Nach Nutzung oder alphabetisch sortieren. Entweder alle sofort zeigen oder 1--200 Schlagwörter vor der Erweiterungssteuerung; Standardschwelle 20. Optionales Scrollen verwendet 1--12 sichtbare Zeilen, Vorgabe 5, und aktiviert sich nur, wenn die dargestellten Schlagwörter tatsächlich mehr konfigurierte Zeilen belegen.\\
Erscheinungsbild & Oberflächenanimationen & Dauer des öffentlichen Informationspanels auf Galeriekarten: 0--800~ms (Standard 320~ms); Dauer des Admin-Seitenpanels: 0--800~ms (Standard 260~ms). 0 deaktiviert die jeweilige Bewegung.\\
Layout & Kopfzeilenverknüpfungen & Drei geordnete Plätze. Ein Platz kann leer sein, auf die Hauptseite oder eine per Galerieauswahl gefundene Galerie zeigen. Öffentliche Besucher sehen nur öffentliche/gelistete Galerieziele.\\
Layout & Breadcrumb-Standardstil & Legen Sie den globalen Stil unter \ui{Design > Layout > Karten und Kennzeichnungen} fest. Physische Galerien können ihn übernehmen oder überschreiben; ungültige Werte fallen sicher auf den integrierten Stil \codeword{chevron} zurück.\\
Erscheinungsbild & Galeriekartendesign & \ui{Vertikales System} zeigt das Foto über der Beschreibung; \ui{Horizontales System} daneben. Neuinstallationen beginnen vertikal; die Migration erhält das Aussehen älterer Installationen. Eine physische Galerie kann die Designvorgabe bei verfügbarem Präsentationsschema erben oder überschreiben.\\
Layout & Bildanzahlkennzeichnung & Die globale Kennzeichnung enthaltener Bilder ist standardmäßig aktiviert. Einzelne Galerien können sie erben oder überschreiben. Bei Aktivierung für die geöffnete Galerie zeigt auch der Hero-Bereich deren Zweigbildanzahl einschließlich zugänglicher Untergalerien.\\
Layout & Öffentlicher Vorschaubildrenderer & \ui{Progressive Vorschaubildschärfung} ist Standard und sicherer Ersatz. \ui{Responsive Browserauswahl} bleibt die unterstützte ältere Option mit vollständigen serverseitigen Kandidaten und Verhalten ohne JavaScript. Titelbilder und Collagen behalten ihre responsive Pipeline.\\
Layout & Seitennavigation & Globale Seitennavigation ist standardmäßig aus. Spalten: Vorgabe 3, begrenzt auf 1--12; Zeilen: Vorgabe 3, begrenzt auf 1--50. Wirksame Einträge je Seite sind Spalten mal Zeilen.\\
Layout & Hauptseitenraster & Die Startseite besitzt ein eigenes Raster mit 1--12 Spalten und 1--50 Zeilen. Raster innerhalb von Galerien bleiben unverändert.\\
Layout & Galerieeigene Raster zurücksetzen & \ui{Alle benutzerdefinierten Galerieraster zurücksetzen} leert Datenbanküberschreibungen, stellt den standardmäßigen Untergalerie-Vererbungszustand wieder her und entfernt entsprechende Rasterschlüssel aus \filepath{gallery.json}-Begleitdateien, damit spätere Erkennung keine veralteten Rastereinstellungen erneut importiert.\\
Layout & Lightbox-Vorgabe & \ui{Einzelbild}, \ui{Bildstreifen} oder \ui{3D-Karussell} bei aktivierter Lightbox-Modusfunktion. Klassisches Einzelbild ist der Ersatzmodus und bei deaktivierter Funktion der einzige angebotene Modus.\\
Branding & Öffentlicher Kopfzeilenbanner & Optionaler designweiter JPG/PNG/GIF/WebP-Banner, maximal 8~MB. Er ersetzt den sichtbaren Websitetitel, wenn kein galeriespezifischer Banner konfiguriert ist.\\
Branding & Kopfzeilentrennlinie & Optionale JPG/PNG/GIF/WebP-Trennlinie, maximal 8~MB. Breite 0 folgt der responsiven Seitenbreite; positive Breite wird auf 160--3840~Pixel begrenzt. Höhe: Vorgabe 72, begrenzt auf 8--512~Pixel. Normalmodus erhält das Seitenverhältnis und behandelt Höhe als Maximum; Streckmodus verwendet genau das konfigurierte Darstellungsrechteck und kann das Bild bewusst verzerren.\\
Branding & Favicon & JPG/PNG/GIF/WebP hochladen, einen quadratischen Ausschnitt ziehen und diesen auf das 1--3-Fache vergrößern. Speichern erzeugt PNG-Faviconvarianten mit 32, 48 und 180~Pixeln.\\
Branding & Designhintergrund & Originalupload erhalten und bei verfügbarer GD/WebP-Unterstützung optional eine WebP-Ableitung erzeugen. Optimierte längste Seite 1024--3840~Pixel, Vorgabe 1920, Regelschritte 128~Pixel. Neuerzeugung oder Entfernung löscht das Original nicht.\\
Branding & Hintergrundsichtbarkeit/\allowbreak Ersatzquelle & Hintergrundsichtbarkeit 0--100~Prozent, Vorgabe 65. Galerieersatzquelle: keine, neuer Designupload, vorhandenes Galeriebild oder Collage aus öffentlichen Galerien. Separate Aktionen setzen alle galerieeigenen Hintergrundauswahlen zurück oder entfernen Designhintergrund beziehungsweise Favicon.\\
Benutzerdefiniertes CSS & Designvorlage oder Upload & Vorlagen sind die derzeit in \filepath{custom_css/} vorhandenen CSS-Dateien; die Auswahl kopiert eine nach \filepath{public/assets/custom.css}. Eine manuell hochgeladene \filepath{.css}-Datei ersetzt diese aktive Datei. Das aktive benutzerdefinierte Stylesheet lädt nach eingebautem CSS und gespeicherten Designsteuerungen.\\
Benutzerdefiniertes CSS & Rücksetzaktionen & \ui{Benutzerdefiniertes CSS zurücksetzen} löscht das aktive Stylesheet und leert dessen Vorlagenmarker. \ui{Auf CSS zurücksetzen} leert gespeicherte visuelle Überschreibungen für Palette/\allowbreak Schrift/\allowbreak Radius/\allowbreak Hintergrundquelle/\allowbreak GPS, sodass das zugrunde liegende CSS wieder maßgeblich wird; strukturelle Seitenbreiteneinstellungen bleiben bewusst erhalten.\\
Sprache & Administrations-/öffentliche Sprache & Administrationssprache wird für Sitzung/Browser des Administrators gespeichert; öffentliche Sprache ist die websiteweite anonyme Vorgabe. Die Besucherauswahl ist eine unabhängige Besucherüberschreibung, unten erläutert.\\
Sprache & Paketwartung & Unterstützte Pakete zeigen Zeichenkettenanzahl und Abdeckung. Der JSON-Editor prüft ein Objekt mit Zeichenkettenwerten, importiert/exportiert JSON und kann authentifizierte Diagnosen fehlender Schlüssel leeren. Englisch bleibt Quell-/Ersatzkatalog.\\
\bottomrule
\end{longtable}

\subsection{Lightbox-Browsingmodi}
\index{l00164@Lightbox!b00052@Browsingmodi}
\index{b00045@Bildstreifen}
\index{a00000@3D-Karussell}

Der Designeditor definiert die standardmäßige öffentliche Lightbox-Präsentation. \ui{Einzelbild} erhält den klassischen Betrachter. \ui{Bildstreifen} ergänzt benachbarte Vorschaubilder unter dem aktiven Foto. \ui{3D-Karussell} platziert eine kleine Gruppe benachbarter Fotos hinter dem aktiven Bild und bietet beim Durchsehen mehr visuellen Kontext.

Eine physische Galerie kann die Designvorgabe erben oder eine eigene Lightbox-Modus\-über\-schrei\-bung speichern. Auch Smart-Gallery-Präsentation kann aus eigenen Einstellungen und Websitevorgaben einen wirksamen Lightbox-Modus bestimmen. Die Modi teilen dieselbe autorisierte Lightbox-Metadaten- und Navigationspipeline, sodass Vollbild, Diashow, Zoom, Karten, Abstimmung und progressive Qualitätssteigerung weiterhin ihren normalen Regeln folgen. Das Deaktivieren des globalen Lightbox-Modus-Funktionsschalters verbirgt diese Betrachtungssteuerungen, ohne gespeicherte Fotos zu verändern.

\subsection{Oberflächensprache wählen}
\label{sec:language-selection}
\index{s00231@Sprache}
\index{l00165@Lokalisierung!s00230@Sprachauswahl}
\index{u00244@Übersetzungspakete}

PHP Gallery behandelt die Sprache der Programmoberfläche getrennt von selbst verfasstem Text. Eine Änderung der Oberflächensprache verändert Schaltflächen, Menüpunkte, Hilfetexte, Statusmeldungen, Dialoge und andere Anwendungstexte. Galerie- und Fototitel/-beschreibungen können optionale gepflegte Sprachversionen besitzen; diese werden jedoch getrennt bearbeitet und gespeichert. Schlagwörter und andere Betreiberinhalte bleiben unverändert.

Öffnen Sie \ui{Design > Sprache}. Zwei unabhängige Auswahlen stehen bereit:

\begin{itemize}
  \item \ui{Sprache der Administrationsoberfläche} ändert die Sprache, die der Administrator im aktuellen Browser sieht.
  \item \ui{Sprache für öffentliche Besucher} setzt die Standardoberflächensprache öffentlicher Seiten.
\end{itemize}

Die optionale \ui{Besucher-Sprachauswahl} erlaubt einzelnen Besuchern, die öffentliche Vorgabe in ihrem eigenen Browser zu überschreiben. Der Administrator wählt die gepflegten Sprachen der öffentlichen Auswahl; das Deaktivieren der Funktion bringt Besucher zur websiteweiten öffentlichen Sprache zurück. Administrative und öffentliche Sprache müssen nicht übereinstimmen.

Öffentliche Seiten verwenden eine kompakte Sprachsteuerung mit muttersprachlichen Bezeichnungen und lokal ausgelieferten Flaggensymbolen. Die Sprachwahl hält den Besucher auf derselben öffentlichen Seite und merkt sie für spätere Besuche in diesem Browser. \ui{Websitevorgabe verwenden} entfernt die persönliche Überschreibung. Der entsprechende Parameter \codeword{?lang=} bleibt für direkte Links und Nutzung ohne JavaScript verfügbar.

\subsubsection{Besucher-Sprachauswahl anpassen}

Die Auswahl besitzt fünf Vorlagen: \ui{Klassisch}, \ui{Gefüllte Kapseln}, \ui{Umrandung}, \ui{Weiche Karten} und \ui{Minimal}. \ui{Einstellungen > Allgemein} bietet grundlegende Vorlagen-/Flaggenentscheidungen; \ui{Design > Sprache} enthält detaillierte Gestaltung und Live-Vorschau. Diese Erscheinungsbildeinstellungen verändern weder angebotene Sprachen, Websitevorgabe, Administrationssprache noch die persönliche Besucherwahl.

\subsubsection{Mehrsprachige Galerie- und Fototexte}
\index{l00165@Lokalisierung!g00128@Galerieinhalt}
\index{l00165@Lokalisierung!f00110@Fototitel}

Galerie- und Fotoeditoren behalten vorhandenen Titel und Beschreibung als Ausgangsinhalt. Mit \ui{Sprache des aktuellen Titels und der Beschreibung} ordnen Sie diesen Text Englisch, Tschechisch, Deutsch oder Schwedisch zu oder lassen ihn \ui{Nicht angegeben}. Vorhandene Installationen werden nicht automatisch als Englisch eingeordnet.

Öffnen Sie \ui{Andere Sprachen} über die Globussteuerung, um optional übersetzte Titel und Beschreibungen einzugeben. Galerietitel und -beschreibungen fallen unabhängig auf den Ausgangstext zurück. Eine gespeicherte Fotosprachzeile gilt hingegen als eine Bildunterschriftsvariante: Leere Felder bleiben verborgen, statt Ausgangssprachtext unter übersetzten Fototext zu mischen. Sind beide Übersetzungsfelder leer, wird diese Sprachzeile entfernt. Gespeichert wird über das normale Galerie-/Fotoformular im bestehenden Administrations-Seitenbereich.

Öffentliche Seiten verwenden die ausgewählte gepflegte Besuchersprache. Passender Text erscheint auf Karten, in Galeriebeschreibungen, Fotoüberlagerungen, Lightbox-Metadaten, Suchergebnissen, Alternativtext und SEO-Metadaten. Fehlende Galeriefelder fallen auf Ausgangsinhalt zurück; eine vorhandene Fotobildunterschriftsvariante zeigt nur ihre gespeicherten übersetzten Felder. Übersetzungen verändern weder URLs, Slugs, Ordner, Dateinamen, Reihenfolge, Sichtbarkeit, Passwörter, NSFW-Regeln noch Medienautorisierung.

Ist OpenAI-Texthilfe konfiguriert, sendet \ui{Übersetzung mit OpenAI vorschlagen} das Quellfeld und setzt einen prüfbaren Entwurf in das Zielfeld ein. Der Entwurf wird niemals automatisch gespeichert oder veröffentlicht. Manuelle Übersetzung bleibt ohne Anbieter vollständig verfügbar.

\subsubsection{Verfügbare Sprachen und ihr aktueller Zustand}
\index{s00231@Sprache!e00086@Englisch}
\index{s00231@Sprache!t00243@Tschechisch}
\index{s00231@Sprache!d00075@Deutsch}
\index{s00231@Sprache!s00213@Schwedisch}

Englisch ist kanonische Ausgangssprache, Vorgabe neuer Installationen und Ersatzsprache. Tschechisch, Deutsch und Schwedisch werden ebenfalls als vollständige Kataloge gepflegt. Diese vier sind die derzeit auswählbaren Oberflächensprachen.

\begin{longtable}{>{\RaggedRight\arraybackslash}>{\hspace{0pt}}p{0.20\textwidth} >{\RaggedRight\arraybackslash}>{\hspace{0pt}}p{0.10\textwidth} >{\RaggedRight\arraybackslash}>{\hspace{0pt}}p{0.58\textwidth}}
\toprule
\textbf{Sprache} & \textbf{Code} & \textbf{Aktuelle Rolle}\\
\midrule
\endfirsthead
\toprule
\textbf{Sprache} & \textbf{Code} & \textbf{Aktuelle Rolle}\\
\midrule
\endhead
Englisch & \codeword{en} & Vollständige kanonische Quelle, Vorgabe neuer Installationen und Ersatzsprache.\\
Čeština & \codeword{cs} & Vollständige gepflegte auswählbare tschechische Übersetzung.\\
Deutsch & \codeword{de} & Vollständige gepflegte auswählbare deutsche Übersetzung.\\
Svenska & \codeword{sv} & Vollständige gepflegte auswählbare schwedische Übersetzung.\\
\bottomrule
\end{longtable}

\important{Derzeit sind nur Englisch, Tschechisch, Deutsch und Schwedisch auswählbar. Englisch bleibt der Ersatz, wenn eine gepflegte Übersetzung fehlt.}

\subsubsection{Funktionsweise von Übersetzungspaketen}
\index{u00244@Übersetzungspakete!e00090@Ersatzsprache}
\index{u00244@Übersetzungspakete!j00153@JSON}
\index{u00244@Übersetzungspakete!d00076@Diagnose}

Ein Übersetzungspaket ist einfach ein Wörterbuch benannter Oberflächenmeldungen. Das Programm fordert einen stabilen Meldungsschlüssel wie \glqq{}Galerie speichern\grqq{} oder \glqq{}ausgewählte Fotos löschen\grqq{} an; das aktive Sprachpaket liefert den anzuzeigenden Wortlaut. Der interne Schlüssel bleibt in jeder Sprache gleich; nur der angezeigte Text ändert sich.

Die Tabelle \ui{Erkannte Sprachpakete} zeigt die Abdeckung unterstützter Sprachen; der Unterabschnitt \ui{Diagnose} vermerkt während einer authentifizierten Administratorsitzung beobachtete fehlende Schlüssel. Englisch liefert Ersatztext, wenn einer gepflegten Übersetzung ein Schlüssel fehlt.

Der \ui{Paketeditor} dient sorgfältiger Wartung oder dem Import einer vorbereiteten Übersetzung. Sprachdaten werden als JSON gespeichert. Im Editor ist der Text links in jedem JSON-Paar der stabile Schlüssel und darf nicht übersetzt werden. Der Text rechts ist der übersetzbare Wert. Platzhalter wie \codeword{\{count\}}, \codeword{\{gallery\}} oder \codeword{\{time\}} stehen für von PHP Gallery eingesetzte aktuelle Werte und müssen in jeder Sprache dieselben Namen behalten.

Ein Sprachpaket kann auch als JSON exportiert, außerhalb von PHP Gallery übersetzt und wieder importiert werden. Importe werden nur akzeptiert, wenn die Datei ein JSON-Objekt mit Textwerten ist; Übersetzer können somit am Paket arbeiten, ohne Zugriff auf den Galerieserver zu erhalten.

Für Entwickler und Betreuer liegen die gepflegten Kataloge unter \filepath{app/lang/}. JSON ist das kanonische bearbeitbare Format. Englische, tschechische, deutsche und schwedische PHP-Wörterbücher bleiben nur als Kompatibilitätsersatz für Installationen ohne entsprechende JSON-Datei erhalten. Serverseitig dargestellte Seiten und Browser-JavaScript verwenden dasselbe Modell aus aktiver Sprache und englischem Ersatz, sodass Übersetzungen nicht in getrennten Frontend- und Backendkatalogen dupliziert werden müssen.

Benutzerdefiniertes CSS unterstützt installationsspezifische visuelle Verfeinerung. Designerzeugtes CSS stellt validierte Variablen und berechnete Selektoren über eine eigene Route bereit. Galeriebranding kann ausgewählte websiteweite Ressourcen überschreiben und dabei wirksames Ersatzverhalten erhalten.

\subsection{Galerie-Hero-Schlagwörter}
\index{s00212@Schlagwörter!g00118@Galerie-Hero-Bereich}
\index{d00073@Design!g00132@Galerieschlagwörter}

Der Unterabschnitt \ui{Erscheinungsbild > Galerieschlagwörter} steuert die kompakte Schlagwortsammlung unter dem Titel einer geöffneten Galerie. Der Schlagwortbereich nutzt die vollständige Breite des Hero-Bereichs; Schlagwörter können somit bis zum rechten Rand umbrechen, statt auf normale lesbare Absatzbreite beschränkt zu bleiben. Direkte Galerieschlagwörter und aus enthaltenen Inhalten gesammelte Schlagwörter bleiben getrennte semantische Gruppen, folgen aber derselben globalen Anzeigerichtlinie.

Derselbe Unterabschnitt steuert auch die öffentliche Schlagwort-Zielseite. \ui{Galerien je Zeile} und \ui{Zeilen je Seite} bestimmen das Galerieergebnisraster und seine Seitenkapazität; \ui{Galeriekartendesign} wählt die vertikale oder horizontale Kartenpräsentation. Diese Werte betreffen öffentliche Seiten, die über ein Schlagwort geöffnet werden, und verändern normale Galerieseiten nicht. Sind eigene Werte noch nicht gespeichert, bleiben globales Designraster und Galeriekartendesign die Vorgabe. Der globale Seitennavigationsschalter bestimmt weiterhin, ob lange Schlagwortergebnisse in Seiten geteilt werden. Unter \ui{Schlagwörter bearbeiten} öffnet \ui{Schlagwortanzeige konfigurieren} diesen Unterabschnitt direkt; \ui{Schlagwortmetadaten verwalten} bietet den umgekehrten Kontextlink.

Standardmäßig zeigt die serverseitig dargestellte Seite sofort 20 Schlagwörter. Die Grenze lässt sich per Schieberegler oder Zahlenfeld von 1 bis 200 ändern. Bei mehr Schlagwörtern öffnet die Auslassungssteuerung einen Bereich \ui{Alle Schlagwörter} mit sämtlichen direkten und enthaltenen Schlagwörtern in ihrer ursprünglichen Gruppenreihenfolge. Er verwendet dieselben Animationen und Schließaktionen wie Galeriekarten, hält die Kopfbereichsmaße stabil und bricht lange Namen um. Mit JavaScript bleibt er im sichtbaren Fenster; dieselbe Steuerung, ein Klick außerhalb oder Escape schließen ihn. Escape setzt den Fokus zurück auf die Steuerung. Alle Schlagwörter stehen bereits im HTML, daher benötigt das Öffnen keine Serveranfrage. Ohne JavaScript öffnet die native Aufklappsteuerung den Bereich unter dem Kopfbereich; lange Listen bleiben durch Seitenscrollen erreichbar. \ui{Jedes Schlagwort sofort anzeigen} zeigt alle bereits in der ersten Antwort; bei genau erreichter Grenze entfällt die Aufklappsteuerung ebenfalls.

Die Reihenfolge ist als \ui{Am häufigsten verwendet zuerst} oder \ui{Alphabetisch} konfigurierbar. \ui{Am häufigsten verwendet zuerst} ist die Vorgabe. Nutzung ist die Gesamtzahl direkter Schlagwortzuordnungen zu Galerien plus direkter Zuordnungen zu Fotos in der Installation. Gleiche Nutzungszahlen werden durch natürliche alphabetische Reihenfolge ohne Beachtung der Groß-/Kleinschreibung und dann durch stabile Kennung aufgelöst. Direkte Galerieschlagwörter und enthaltene Schlagwörter werden jeweils innerhalb ihrer Gruppe sortiert; die Gruppen werden nicht gemischt.

Die Schlagwortvorschau kann eine interne Bildlaufleiste verwenden. Sie ist standardmäßig mit einer Höhenbegrenzung von fünf Zeilen aktiviert; die Zeilenzahl ist per Schieberegler oder Zahlenfeld von 1 bis 12 einstellbar. Der Server liefert die Begrenzung im anfänglichen HTML, sodass die Vorschau nicht auf JavaScript wartet und beim Modulstart ihre Maße behält. Der Umbruch folgt der verfügbaren Breite; Scrollen ist nur bei Überschreiten der Begrenzung nötig. Der offene Überlagerungsbereich wird dadurch nicht abgeschnitten. Das Deaktivieren lässt die Vorschau natürlich wachsen.


\section{Downloads, Freigabe und Übertragung}
\label{sec:sharing}
\index{d00080@Downloads}
\index{z00278@ZIP}
\index{w00270@WebDAV}

Galeriedownloads bereiten autorisierte Inhalte als ZIP-Archive vor und übertragen das fertige Archiv als Stream zum Browser. Ein normaler Galeriedownload durchläuft den erlaubten Galerieteilbaum; Smart-Gallery-Downloads bewerten die gespeicherte Abfrage und aktuelle Zugriffsregeln neu; Auswahldownloads des Picture Manager verpacken nur die aktuell ausgewählten zugehörigen Fotos; der Administrator kann außerdem ein Archiv aller Galerien anfordern. Archiveintragsnamen werden normalisiert, um Pfadtraversierung zu verhindern und deterministische Pfade zu erhalten.

In modernen Browsern fordert ein öffentlicher Galerie- oder Smart-Gallery-Download zunächst ein privates begrenztes Manifest an. Anschließend ruft der Browser jedes autorisierte Original über einen unabhängigen Medienendpunkt ab und setzt die ZIP-Datei lokal zusammen. Dabei zeigt er Fortschritt und ermöglicht Abbruch oder Wiederholung, ohne eine lange PHP-Archivanfrage offen zu halten. Jede Quellanfrage wiederholt die aktuellen Prüfungen von Galerie, Smart-Gallery-Zugehörigkeit, Besuchersichtbarkeit, Medienversion und Pfadbegrenzung; eine Bildkennung allein gewährt niemals Zugriff. Grenzen für Dateianzahl, kumulierte Quellbytes, doppelte Namen, Speicher und ZIP64 halten den Ablauf begrenzt. Direkte Links und Clients ohne JavaScript behalten den bisherigen serverseitigen ZIP-Ersatzweg.

Erzeugte Galerie-/Gesamt-/Auswahl-ZIP-Dateien verwenden eine aus relevanten Galerie-/Bildmetadaten abgeleitete Inhaltssignatur, sodass ein weiterhin aktuelles Cachearchiv wiederverwendet werden kann und eine geänderte Sammlung eine neue Cachekennung erhält. Diese \emph{Inhaltssignatur} sind Cacheinvalidierungsmetadaten, kein öffentliches HMAC-Downloadautorisierungstoken. Der Downloadcontroller prüft die Autorisierung vor der Archivauslieferung.

\subsection{Migration zwischen Galerien}
\index{g00130@Galeriemigration}
\index{a00021@API-Migration}

Der API-Bereich des Galerieeditors unterstützt beide Übertragungsrichtungen zwischen kompatiblen PHP-Gallery-Installationen. \ui{Diese Galerie in eine andere Instanz exportieren} behandelt die aktuelle Galerie als Quelle und verwendet einen galeriegebundenen API-Schlüssel für die empfangende übergeordnete Galerie. \ui{Aus einer anderen Instanz in diese Galerie importieren} behandelt die aktuelle Galerie als empfangende übergeordnete Galerie und verwendet einen API-Schlüssel für die entfernte Quellgalerie. Der Import erstellt darunter eine neue Untergalerie; er überschreibt oder verschiebt niemals die ausgewählte empfangende Galerie.

\ui{Untergalerien einschließen} ist standardmäßig aktiviert. Bei Auswahl beschreibt das versionierte Manifest die Quellgalerie und ihre vollständige untergeordnete Hierarchie in Eltern-zuerst-Reihenfolge. Das Ziel erstellt diese Hierarchie unter der neuen importierten Wurzel erneut und erhält unterstützte Galerie-/Bildmetadaten, Übersetzungen, Flugkartenzustand, Originale, vorhandene Vorschaubilder und Galeriebranding-Ressourcen. Ohne die Option wird nur die ausgewählte Quellgalerie übertragen. Die aktuelle Kompatibilitätsrichtlinie erfordert exakt dieselbe Anwendungsversion auf Quelle und Ziel.

Der empfangende Server erstellt einen durch seine Uploadgrenzen beschränkten deterministischen Paketplan. Jedes Original bleibt mit seinen bereits erzeugten Vorschaubildern gruppiert; Ressourcen auf Galerieebene können unabhängig verpackt werden. ZIP-Einträge verwenden stabile Ressourcenschlüssel, speichern bestehende Bytes ohne erneute Fotokompression und werden vor der Installation gegen Größen und SHA-256-Prüfsummen des Manifests geprüft. Genauer API-Schlüsselumfang, Galeriebaumstruktur, Zielpfade, Datenbankbereitschaft und Ressourcen-Vollständigkeit werden an den jeweiligen Änderungsgrenzen geprüft.

Der Administrator wählt ein Wiederverbindungsintervall von 5 bis 300~Sekunden. Nach Zeitüberschreitung, Neuladen oder Verbindungsabbruch fragt die aktive Seite den dauerhaften Zielauftrag nach bereits vorhandenen Ressourcenschlüsseln und überspringt fertige Pakete, statt sie blind erneut zu übertragen. Ordnerzuordnungen und empfangene Ressourcen werden schrittweise gespeichert, damit dieselbe Übertragung sicher fortsetzbar ist. Der Abschluss bleibt ein gesonderter Schritt und wird verweigert, bis jede deklarierte Ressource vorhanden ist. Der Abbruch des Browserablaufs stoppt weitere Anfragen; Zugangsdaten sollten dennoch widerrufen werden, sobald ein temporärer Migrationsschlüssel nicht mehr benötigt wird.

\subsection{Mobile WebDAV-Übertragung}
\index{w00270@WebDAV!m00173@mobiler Upload}

\ui{Einstellungen > Uploads} erstellt eine WebDAV-kompatible Verbindung für eine ausgewählte Galerie. Geben Sie eine erkennbare Bezeichnung an, wählen Sie die Galerie, erstellen Sie die Verbindung und kopieren Sie erzeugte Server-URL, Benutzername und Passwort in die mobile Anwendung. Das Passwort wird einmal angezeigt und als Hash gespeichert. Vorhandene Zeilen zeigen die letzte Nutzungszeit und können sofort in der Administration gelöscht werden.

Verwenden Sie den Zieltyp WebDAV in Clients wie PhotoSync. Die aktuelle Serveraufnahme akzeptiert JPG, PNG, GIF und WebP; aktivieren Sie daher nach Möglichkeit HEIC-zu-JPEG-Konvertierung im Client. WebDAV-Uploads umgehen keine Galerieaufnahmeregeln: Zugangsdaten sind galeriegebunden, jede Schreibanfrage benötigt Authentifizierung, Galerie-/Bildschema muss bekannt sein, und angenommene Medien werden anschließend gescannt, optional automatisch umbenannt und durch die gewöhnliche Vorschaubildpipeline verarbeitet.

\subsection{Öffentliche Auffindbarkeit und Werkzeuge für ausgewählte Fotos}
\index{r00206@robots.txt}
\index{s00225@sitemap.xml}
\index{b00041@Bild-Sitemap}
\index{p00188@Picture Manager}

Die öffentliche Antwort \codeword{/robots.txt} beschreibt Crawlergrenzen; \codeword{/sitemap.xml} führt geeignete öffentliche Galerie- und Bild-URLs mit Bildmetadaten und letzter Änderungszeit auf. Nicht veröffentlichte, private, passwortgeschützte, altersbeschränkte oder anderweitig nicht autorisierte Inhalte werden durch diese Ausgaben nicht öffentlich. Kanonische URLs, soziale Metadaten und strukturierte Bilddaten verwenden dasselbe zugriffsbewusste öffentliche Pfadmodell.

Picture Manager ist eine Überlagerung der öffentlichen Galerieansicht für angemeldete Administratoren, keine Besucherfunktion. Ein Administrator kann Fotos und physische Untergalerien auswählen, eine gemischte Auswahl löschen oder unter einer gewählten übergeordneten Galerie in eine neue physische Galerie kopieren. Reine Fotoauswahlen können außerdem auf eine sichtbare Untergalerie gezogen, über die durchsuchbare Zielauswahl verschoben oder kopiert, über den nativen Gerätefreigabedialog geteilt, sofern unterstützt, oder heruntergeladen werden. Natives Teilen fällt auf die ZIP-Datei ausgewählter Fotos zurück, wenn der Browser vorbereitete Dateien nicht teilen kann. Smart Galleries sind von physischen Teilbaumaktionen ausgeschlossen. Die Auswahl selbst ist nur Browserzustand und gewährt keine Berechtigung. Jede Änderungs-/ZIP-Anfrage prüft erneut normales angemeldetes Konto, CSRF-Token, Quellgalerie, ausgewählte Zuordnungen und gegebenenfalls Ziel. Gewöhnliche Besucher erhalten diese Steuerungen nicht. Picture Manager ist bewusst unabhängig von der Fähigkeit Inline-Administration: Deren Deaktivierung entfernt direkte öffentliche Bearbeitungs-, Hinzufüge-, Lösch- und Sortiersteuerungen; ein unabhängig aktivierter Picture Manager kann weiterhin Auswahl-, Lösch-, Verschiebe-, Kopier-, physische Galerieerstellungs- und Auswahlexportabläufe anbieten.
\section{Wartung, Aktualisierungen und Wiederherstellung}
\label{sec:maintenance}
\index{w00268@Wartung}
\index{a00010@Aktualisierungen}
\index{s00222@Sicherung}

\subsection{Migrationen}

Schemaänderungen liegen in PHP-Dateien mit Zeitstempel unter \filepath{database/migrations/}. Der Migrationsläufer prüft ausstehende Definitionen vorab, wendet sie in Dateinamenreihenfolge an und vermerkt erfolgreiche Versionen. Neue Schemaarbeit erhält eine neue Migrationsdatei, damit vorhandene Installationen einen deterministischen Verlauf behalten.\footnote{Manche Migrationen enthalten neben SQL auch Reparaturcallbacks. Kompatibilitätstests schützen Teilaktualisierungsszenarien, in denen ein älterer Läufer auf neuere Migrationsdefinitionen trifft.}

Datenbankbereitschaftsprüfungen werden nur für die aktuelle PHP-Anfrage gespeichert. Das verringert normalerweise wiederholte Metadatenabfragen; eine Migration kann die Antwort jedoch ändern, während derselbe Administrations-, Installations-, Aktualisierungs- oder Befehlszeilenprozess noch läuft. Der Migrationsläufer leert daher diese gespeicherten Antworten nach schemaändernden Anweisungen und erfolgreichen Reparaturcallbacks. Eine Prüfung nach der Migration liest dann den neuen Datenbankkatalog, statt eine Antwort von vorher wiederzuverwenden. Gewöhnliche reine Datenmigrationsanweisungen leeren den Bereitschaftscache nicht, weil sie keine Tabelle, kein Feld und keinen Index erstellen oder entfernen können.

Version~0.106 ergänzte die Tabelle \codeword{maintenance\_jobs} durch eine gewöhnliche Migration. Führen Sie diese wie üblich auf der Galerieverwaltungsseite aus, sofern noch nicht geschehen. Sie benötigt die normale Berechtigung der Installation zur Tabellenerstellung, aber keine Datenbanktrigger, gespeicherten Routinen, \codeword{SUPER}, globale Servervariablenänderung oder separate Datenbankadministratorhilfe. Bis die Migration verfügbar und geprüft ist, verweigert die Wartungszentrale einen Auftragsstart, statt Wartung ohne dauerhafte Kontrollpunkte zu versuchen. Die Migrationen für kooperative Galerien und Vorschaubildidentitäten verwenden denselben normalen Aktualisierungsablauf.

\subsection{Wie die Galerie Datenbankbereitschaft prüft}
\index{d00062@Datenbank!b00032@Bereitschaft}
\index{s00211@Schemaprüfung}
\index{m00172@Migrationen!f00102@fehlende Objekte}
\index{s00238@Systemzustand!d00064@Datenbankprüfung}
\index{f00103@Fehlersuche!d00066@Datenbankverbindung}

Die gemeinsamen Bereitschaftszustände und Betreiberaktionen sind in Abschnitt~\ref{sec:db-readiness-contract} definiert. Wartung verwendet denselben Vertrag bei der Prüfung von Tabellen, Feldern und Indizes für Passwortzurücksetzung, NSFW Guard, Uploadautomatisierung, Vorschaubildmetadaten und andere schemaabhängige Funktionen.

Während einer Aktualisierung kann \ui{Fehlend} erwartet sein, wenn neue Anwendungsdateien vor ihrer Migration eintreffen. \ui{Unbekannt} bedeutet dagegen, dass PHP Gallery die ausgewählte Datenbank nicht zuverlässig prüfen konnte. Bei \ui{Fehlend} sichern Sie und wenden die ausstehende Migration an. Bei \ui{Unbekannt} prüfen Sie vor Schemaänderungen Datenbankverbindung, ausgewähltes Schema und Berechtigungen zur Metadatenprüfung.

Laden Sie nach der Reparatur \ui{Systemzustand} erneut und testen Sie die betroffene Funktion. Bewahren Sie eine Anfragereferenz zusammen mit Zeitpunkt und Anwendungsversion auf, wenn Sie Hoster oder Betreuer um Hilfe bitten.

\subsection{Galeriepapierkorb und Wiederherstellung}
\index{g00117@Galerie!p00183@Papierkorb}
\index{g00117@Galerie!w00274@Wiederherstellung}
\index{p00183@Papierkorb!a00024@automatische Löschung}

Öffnen Sie \ui{Wartung > Papierkorb}, um wiederherstellbare Galerielöschungen zu prüfen. Jede aktive Zeile zeigt ursprünglichen Pfad, Löschzeitpunkt, Galerie-/Fotoanzahlen, gespeicherte Größe, Lebenszyklusstatus und bei aktiver Richtlinie die automatische Löschfrist. Problemeinträge bleiben zur Prüfung sichtbar; PHP Gallery nimmt bei unklarem Dateisystem-/Datenbankzustand nicht an, dass eine Vernichtung sicher ist.

Der Papierkorbhauptschalter steuert zukünftige Galerielöschungen und ist standardmäßig aktiviert. Sein Ausschalten verbirgt oder löscht keine vorhandenen Einträge: Wiederherstellung und manuelle endgültige Löschung bleiben verfügbar. Automatische Löschung ist eine separate ausdrückliche Option und standardmäßig aus. Ihre Aufbewahrung beträgt standardmäßig 30~Tage, ihre Wartungscharge ist begrenzt, und automatische Löschung pausiert bei deaktiviertem Papierkorb. Die Aktivierung nach einer Pause gibt aktuellen wiederherstellbaren Einträgen eine neue vollständige Aufbewahrungsfrist, bevor geplante Löschung beginnen darf.

Migrationen \filepath{202609070001_gallery_trash_bin.php} und \filepath{202609070002_gallery_trash_state_machine.php} ergänzen dauerhafte Vorgangsdatensätze. Bei aktiviertem Papierkorb prüft die Löschung dieses Schema vor der ersten Dateisystemverschiebung oder Datenbankzeilenlöschung. Bestätigt fehlender oder unbekannter Speicher stoppt wiederherstellbares Löschen und verweist den Administrator zur Anwendung/Prüfung von Migrationen; es erfolgt kein stiller Wechsel zu endgültiger Löschung. Geplante Websitewartung gleicht veraltete laufende Vorgänge vorsichtig ab und führt abgelaufene automatische Löschungen nur bei aktiviertem Papierkorb und aktivierter automatischer Löschung aus.

Papierkorbdaten liegen im geschützten dauerhaften Speicher \filepath{data/gallery-trash/}, nicht im Cache oder aktiven Galeriebaum. Bereitstellungspakete enthalten die Zugriffsrichtliniendatei des Verzeichnisses, schließen aber gespeicherte Papierkorbinhalte aus. Beziehen Sie dieses Laufzeitverzeichnis und die Datenbank in Sicherungs- und Wiederherstellungsplanung ein; keine Hälfte allein bildet einen vollständigen wiederherstellbaren Eintrag.

\subsection{Wartungszentrale}
\index{w00269@Wartungszentrale}
\index{w00268@Wartung!z00276@zentraler Ablauf}

Öffnen Sie \ui{Wartung > Wartungszentrale} für einen geführten Ablauf \ui{Analysieren}, \ui{Prüfen}, \ui{Ausführen} und \ui{Verifizieren} über die Verantwortungsbereiche regelmäßiger Wartung. Analysieren ist für Galerie-, Telemetrie-, Protokoll-, Papierkorb-, Cache-, Vorschaubild- und Datenbankinhalte schreibgeschützt; geschrieben wird nur der dauerhafte Kontrollpunkt der Wartungszentrale zur Fortsetzung der Prüfung. Der Prüfschritt gruppiert Aufgaben nach Bereich, kennzeichnet erforderliche automatische Schritte und fasst verfügbare optionale Aufgaben ohne erkannten Arbeitsbedarf in einem eingeklappten Abschnitt zusammen. Erforderliche, standardmäßige und optionale Aufgaben bleiben vor Beginn der Bereinigung klar erkennbar; tiefe Medienprüfung und vollständige Datenbankoptimierung bleiben bei Verfügbarkeit auffindbar, auch wenn sie nicht automatisch ausgewählt werden.

Der Server lädt vor jeder Aufgabe die geprüften Abhängigkeiten des jeweiligen Bereichs erneut für jede separate Analyse- oder Ausführungsanfrage. Die Prüfung ausstehender Migrationen liest nur die Migrationstabelle; eine fehlende Tabelle oder eine fehlgeschlagene Prüfung gilt als ausstehende Arbeit. Wenden Sie Migrationen über den ausdrücklichen Migrationsablauf an; die Analyse erstellt oder repariert keinen Migrationsspeicher.

Der geprüfte Plan ist an aktuelle Anwendungsversion, angewandte Migrationsrevision, Registrierungsrevision der Wartungszentrale und Planinhalt gebunden. Ändert sich einer dieser Eingänge, lehnt die Ausführung den veralteten Plan ab und fordert neue Analyse. Der Server prüft bei der Ausführung erneut Aufgabenabhängigkeiten, Fähigkeitszustand und Datenbanktabelleneignung; Browserkontrollkästchen sind Präsentationseingaben, keine Autorisierung.

Die Ausführung schreitet durch begrenzte Kontrollpunkte fort, sodass Sie die Seite neu laden oder schließen und denselben Auftrag später fortsetzen können. Nur ein laufender oder pausierter zentraler Auftrag hält den dauerhaften Änderungsanspruch; zwei Browserregister können einen Auftrag nicht gleichzeitig vorantreiben. Automatische Websitewartung gibt vor ihrer ersten Änderung nach, solange dieser Anspruch gehalten wird. Öffentliche schreibgeschützte Galerieanfragen bleiben verfügbar.

Pausieren und Abbrechen erfolgen kooperativ: Sie wirken zwischen begrenzten Vorgängen und machen bereits festgeschriebene Arbeit nicht rückgängig. Ein bereits laufendes Datenbank-\codeword{ANALYZE TABLE} oder -\codeword{OPTIMIZE TABLE} muss zurückkehren, bevor der Abbruch fortgesetzt werden kann. Die Zentrale bearbeitet höchstens eine aktuell geeignete Tabelle je Browserschritt, überspringt geschützte, unbekannte, übergroße oder neu ungeeignete Tabellen und vermerkt ein unsicheres atomares Ergebnis, statt eine Tabelle nach verlorener Antwort blind erneut zu bearbeiten.

Der Abschlussbericht trennt entfernte Zeilen und Dateien, analysierte oder optimierte Tabellen, Übersprünge, Warnungen und Fehler. Er meldet physischen Speicher vor/nach der Arbeit nur bei ausreichenden Bestandsnachweisen. Der allgemeine Ablauf wendet niemals Migrationen an, führt Schemareparaturen aus, installiert Aktualisierungen, löscht Administrationsprotokollarchive, leert alle Papierkorbeinträge oder erzeugt Medien neu; verwenden Sie Spezialseiten für diese bewussten Vorgänge.

Download-/Cachebereinigung durchläuft Manifestcache, bisherigen Download-Artefaktcache und Cache erzeugter ZIP-Dateien als getrennte begrenzte Schritte. Diese Caches sind unabhängige Optimierungen: Ist einer auf dem aktuellen Host nicht verfügbar, bewahrt die Zentrale erledigte Bereinigung, vermerkt eine sichere Warnung und fährt mit anderen Verantwortungsbereichen und späteren Aufgaben fort. Die Warnung darf Vorgang und stabile Fehlerkategorie benennen, aber keine unverarbeitete Ausnahmemeldung, SQL, Zugangsdaten, privaten Pfad oder Stacktrace offenlegen. Ändert eine Aktualisierung diesen Ausführungsvertrag, ist ein älterer fertiger Plan bewusst veraltet; führen Sie vor dem Fortsetzen Analysieren erneut aus.

\subsection{Vorschaubildwartung}

Wartungswerkzeuge prüfen erzeugte Varianten, erkennen fehlende oder veraltete Ableitungen und erstellen ausgewählte Größen neu. Browsergestützter Neuaufbau kann bei entsprechender Konfiguration Quellblöcke vorbereiten. Hintergrundvorbereitung behandelt kleine Lücken öffentlicher Darstellung mit begrenzten Anfragen; ausdrückliche administrative Wartung bleibt der primäre Reparaturweg in großem Umfang.

\subsection{SEO-Anfrageschutz und öffentliche Abfragehygiene}
\index{s00218@SEO-Anfrageschutz}
\index{a00002@Abfrageparameter-Spam}
\index{k00154@kanonische URL}

Der Wartungsbereich enthält einen optionalen, standardmäßig aktivierten öffentlichen SEO-Anfrageschutz. Für anonyme öffentliche GET-Anfragen hält er je Route eine ausdrückliche Positivliste von Abfrageparametern. Bekannte Analyseparameter wie UTM-Felder und übliche Anzeigenklickkennungen werden ignoriert, weil sie den dargestellten Inhalt nicht verändern. Unerwartete Parameter werden vor dem normalen öffentlichen Renderer mit 404, \codeword{noindex,nofollow} und einer nicht zwischenspeicherbaren Antwort abgelehnt, damit crawlererzeugter Parameterspam keine indexierbaren Galerie-URLs vervielfacht.

Administrationsrouten und maschinelle Endpunkte wie Einrichtung, WebDAV, Uploadautomatisierung, Galeriemigrationsempfänger und Wartungs-Cron-Endpunkt liegen außerhalb dieses öffentlichen Abfrageschutzes. Angemeldete Administratoren sind ebenfalls ausgenommen, damit Diagnose-/Administrationsabläufe nicht von öffentlichen Crawlerregeln blockiert werden. Die Protokollierung abgelehnter Abfragen ist unabhängig aktivierbar und je UTC-Tag stichprobenbegrenzt, damit eine Crawlerflut das Administrationsprotokoll nicht füllt. Derselbe Dienst liefert kanonische Ersatz-URLs, wenn eine öffentliche Seite noch kein eigenes kanonisches Tag ausgegeben hat.

\subsection{Datenbankwartung}

Die Datenbankprüfung erfasst Schemaobjekte, Referenzen, Zuordnungskategorien, Aufbewahrungsrichtlinie und Bereinigungssicherheit. Logische Bereinigung verwendet Regeln einer Positivliste mit deterministischer Auswahl, begrenzten Chargen und Auditdatensätzen. Schemareparatur bietet Probelauf und idempotenten Migrationsweg. \codeword{ANALYZE TABLE} kann Optimiererstatistiken ausgewählter Tabellen aktualisieren. Physisches \codeword{OPTIMIZE TABLE} ist nur als Vorschau mit anschließend separat bestätigtem Vorgang für ausgewählte Tabellen verfügbar; es ist niemals ein automatischer Bereinigungsschritt.\footnote{Zuweisungswerte der Speicherengine wie Schätzungen freien Datenspeichers beschreiben den Enginezustand und garantieren keine genaue Anzahl nach Optimierung freigegebener Dateisystembytes.}

\subsection{Speicherstatistik und Gesamtbericht}
\index{s00229@Speicherstatistik}
\index{g00139@Gesamtbericht}
\index{g00120@Galeriebericht}

Detaillierte \ui{Speicherstatistik} wird nur auf Anforderung eines Administrators berechnet, damit die normale Übersicht nicht bei jedem Besuch das gesamte Ableitungsdateisystem scannt. Die Dateiansicht trennt indizierte Originalfotos von erzeugten JPEG/WebP-Vor\-schau\-bil\-dern und DNG-Browser-Anzeigemasters, zeigt Quell-Dateityp- und Größenverteilungen, größte Galerien und Originale, unbekannte Größen und Scanwarnungen. Eine zwischengespeicherte Momentaufnahme wird bei Galeriedatenänderungen als veraltet markiert. Die Aktualisierung scannt erwartete erzeugte Dateien in kleinen browsergesteuerten Chargen. Derselbe Bereich verlinkt Datenbankspeicherdetails und Datenbankwartung.

Das Register Datenbank liest Information-Schema-Speicherstatistiken getrennt vom Dateisystemscan. Es unterscheidet nach Möglichkeit galerieeigene Tabellen von der gesamten ausgewählten Datenbank und zeigt die größten Tabellen, statt der normalen Übersichtsanfrage eine vollständige Tabellenprüfung aufzuzwingen. \ui{Alles aktualisieren} erneuert den Dateiscan, aktualisiert Tabellenschätzungen mit begrenzten \codeword{ANALYZE TABLE}-Anfragen und führt anschließend die schreibgeschützte Datenbankprüfung aus. Der Fortschritt bleibt in einem serverseitigen Ablauf gespeichert und die angezeigte Registerkarte wird nach Abschluss aktualisiert, ohne die Seite zu verlassen; ein erneuter Start während des aktiven Ablaufs desselben Administrators setzt bei der gespeicherten Phase fort. Teilfehler werden je Phase gemeldet und im Administrationsprotokoll vermerkt. Die Aktion löscht keine Daten, repariert kein Schema, optimiert keine Tabellen und führt keinen ausgewählten Plan der Wartungszentrale aus.

\ui{Wartung > Gesamtbericht} erstellt einen eigenständigen HTML-Bericht über Speicher, Datenbankstruktur und -größe, Galeriestruktur, Bild- und EXIF-Statistik, GPS-Übersicht, Telemetrie, Betriebsprotokolle, Funktions-/Einstellungszustand und Laufzeitdiagnose. Bild- und erzeugte Medienarbeit erfolgt in begrenzten browsergesteuerten Chargen, damit der Bericht auf Shared Hosting praktikabel bleibt. Vorübergehende Zeitüberschreitungen, Ratenbegrenzungs- und Serverfehler werden mit kurzen zunehmenden Verzögerungen wiederholt; EXIF- und GPS-Zusammenfassungen mit vielen unterschiedlichen Werten bleiben begrenzt, um Sitzungsspeicher und Arbeitsspeicher großer Bibliotheken zu kontrollieren. Der gewählte Telemetriezeitraum betrifft nur den Telemetrieabschnitt; Bildstatistik umfasst stets die vollständige Bildbibliothek. GPS-Ortszusammenfassungen verwenden lokale ungefähre Gruppierung und Offline-Referenzdaten; wahrscheinliche Simulator- oder Spielbildschirmfotos werden bei ausreichenden Anhaltspunkten aus realen Ortshäufigkeiten ausgeschlossen.

Der fertige Bericht wird als Download zum Browser zurückgegeben und nicht als Berichtsdatei auf dem Server gespeichert. Für eine PDF-Kopie öffnen Sie das erzeugte HTML und verwenden Drucken oder Als PDF speichern im Browser; Druckformatierung ist enthalten. Behandeln Sie jeden exportierten Bericht als Verwaltungsdaten, weil er detaillierte Betriebs- und Sammlungsstatistiken enthalten kann.

\subsection{Telemetrie und Datenschutzsteuerungen}
\index{t00240@Telemetrie}
\index{d00068@Datenschutz!t00240@Telemetrie}

Der optionale Bereich \ui{Telemetrie} erfasst datenschutzkontrollierte Betriebsmessungen und zeigt sie als Administrationsberichte. Steuerungen umfassen Telemetriehauptschalter, öffentliche Nutzungsmessung, Leistung, Cache- und Datenbankereignisse, Berücksichtigung von Do Not Track, Ausschluss von Administratoren, maximale gezählte Fotobetrachtungsdauer und Aufbewahrung roher, stündlicher und täglicher Daten. Berichte enthalten bei aktivierten Kategorien beispielsweise anonyme Sitzungen, Seitenaufrufe, Fotoöffnungen und begrenzte Betrachtungszeit, Clientfehler, gemessene Medienbytes, Browserverteilung und Cache-/Datenbankaktivität.

Das Telemetriesystem ist so ausgelegt, dass es keine unverarbeiteten IP-Adressen, Browser-User-Agent-Zeichenfolgen, Referrer-URLs, Namen, E-Mail-Adressen, Kontokennungen, Anfrageinhalte oder genauen Besucherorte speichert. Exporte enthalten aggregierte anonyme Telemetrie und Sitzungshashes statt ursprünglicher Besucherkennungen. Wartung aggregiert und bereinigt Telemetrie gemäß konfigurierter Aufbewahrungsrichtlinie; die Administrationsansicht kann einen eigenständigen HTML-Bericht zur Offline-Prüfung exportieren. Prüfen Sie Datenschutzsteuerungen vor Aktivierung der Erfassung auf einer öffentlichen Website.

\subsection{Laufzeitdiagnose, Integrität und Galeriebenchmark}
\index{l00160@Laufzeitdiagnose}
\index{i00151@Integritätsprüfung}
\index{g00119@Galeriebenchmark}

\ui{Laufzeitdiagnose} bietet eine kopierbare Umgebungs- und Supportansicht für PHP, Datenbankbereitschaftsklassen, Hostinggrenzen, GD, Imagick/ImageMagick und relevante Medienformatunterstützung. Sie zeigt auch die zuvor erläuterte DNG-Konvertierungsrichtlinie. Bereitschaftsbereiche trennen Präsentations-/Berichtsabhängigkeiten, Authentifizierungs-/Sicherheitsabhängigkeiten und änderungssensible Abhängigkeiten, damit Administratoren eine optional fehlende Fähigkeit von einem zwingend schreibblockierenden Zustand unterscheiden können.

Dasselbe begrenzte Modell meldet Bereitschaft der Galeriebearbeitungsrevision und ungelöste Bildverschiebungen. Ein fehlendes oder unbekanntes erforderliches Schema ist ein Aktionszustand und stoppt die betroffene Änderung vor Veränderungen an Zieldateien oder -zeilen. Zusammenfassungen ausstehender Verschiebungen enthalten nur begrenzte Vorgangs-/Galeriekennungen und sichere Zustandskategorien; verwenden Sie den dokumentierten Befehlszeilen-Wiederherstellungsablauf für einen ausdrücklichen Vorgang, statt Dateiverschiebungen von der Diagnoseseite zu erwarten.

Der getrennte Schalter \ui{Entwicklungsdiagnose} ist nur für Administratoren. Bei Aktivierung für einen angemeldeten Administrator ergänzt er Betrachterinstrumentierung für Vorladen, Cache, Speicher, Netzwerk und Bildzeitmessung im öffentlichen/Lightbox-/Vollbilderlebnis. Er aktiviert diese Diagnosen nicht für gewöhnliche anonyme Besucher. Lassen Sie ihn im normalen Betrieb aus, außer zur Untersuchung von Darstellungs- oder Leistungsverhalten; der Galeriebenchmark hängt ebenfalls von diesem Entwicklungsdiagnosekontext ab.

Das DEV-Dashboard im Lightbox-Betrachter öffnet eine kompakte Übersicht. Einfrieren und Fortsetzen steuern die Beobachtung; Einklappen behält die Übersicht und Erweitern öffnet einzelne Seiten für Medien, Vorladen, Diagramme, Ansichtsgröße und Lebenszyklus. Die Seiten passen ohne Scrollleisten in den verfügbaren Platz, behalten die Auswahl bei Aktualisierungen und bleiben auch im nativen Vollbild innerhalb des Betrachters. Die Medienseite unterscheidet das angezeigte Foto vom noch ausstehenden Navigationsziel und zeigt geladene Abmessungen sowie angezeigte und angeforderte Qualität. Quellbeschriftungen verbergen Zugangsdaten, Abfrageparameter und URL-Fragmente.

Drei Diagramme besitzen eigene beschriftete Skalen: geschätzter Pixelspeicher gehaltener dekodierter Bilder, Cacheeinträge und aktive separate Ladevorgänge sowie Animationsbildintervalle mit einer 16,7~ms-Referenz. Höchstens 90 Messpunkte im Abstand von 350~ms umfassen ungefähr 31 Sekunden. Die Pixelschätzung schließt das aktive Bild, GPU-Speicher und den Browserheap aus; Bildintervalle enthalten auch Ablaufplanung und Hintergrunddrosselung. Früheres erfolgreiches Dekodieren beweist keinen aktuellen Cachebesitz. Nicht unterstützte Messungen erscheinen als nicht verfügbar. Die Diagnose beobachtet vorhandene Arbeit ohne zusätzliche Medien- oder Metadatenanfragen; Schließen stoppt die Beobachtung und Aufräumen entfernt das Dashboard. Bei ausgeschaltetem DEV wird dessen Modul nicht geladen.

Der administrative Ablauf \ui{Integrität} prüft integritätsverwaltete Kerndateien gegen das installierte Manifest und hilft, unvollständige oder unerwartete Bereitstellungsänderungen zu erkennen. Er ist ein Laufzeitprüfwerkzeug; abgestimmte Sicherungen bleiben erforderlich, weil ein Integritätsmanifest keine Benutzerfotos enthält und keine Wiederherstellungsmedien ersetzt.

Bei verfügbarer Entwicklungsdiagnose kann ein angemeldeter Administrator für die aktuelle Galerie einen \ui{Benchmark zum Laden der öffentlichen Galerie} starten. Der Browser lädt die Galerie mehrfach in einem verborgenen Frame desselben Ursprungs mit anonymer Vorschau, kombiniert PHP-Darstellungszähler mit Browsernavigationstiming und bietet anschließend das Benchmarkprotokoll als JSON an. Verwenden Sie zum Vergleich zweier Änderungen dieselbe Galerie, denselben Präsentationsmodus und dieselben Netzwerk-/Cachebedingungen; der Benchmark ist als Vorher-/Nachhermessung nützlicher denn als universelle Punktzahl.

\subsection{Geplante Websitewartung}
\index{g00137@geplante Wartung}
\index{c00056@Cron!w00273@Websitewartung}

Regelmäßige Wartung kann in begrenzten Teilabschnitten statt einer langen Anfrage laufen. Der Wartungsdienst verkettet sichere Webanfragen bei vorhandenem Verkehr und kann auch über Hosting-Cron mit dem in der Administration gezeigten CLI-Befehl aufgerufen werden. Er kann Vorschaubildzustand prüfen, Metadaten aktualisieren, Bereinigungsaufgaben bearbeiten, Telemetrie aggregieren und geeignete Tage des Administrationsprotokolls gemäß konfigurierten Richtlinien archivieren. Das verborgene Web-Cron-Token ist austauschbar; ein Austausch macht zuvor veröffentlichte Web-Cron-URLs ungültig. Wartung ist standardmäßig nicht destruktiv und ersetzt keine ausdrückliche Bestätigung von Datenbankoptimierung oder anderen folgenreichen Aktionen.
\subsection{Administrationsprotokollverlauf und Archive}
\index{a00006@Administrationsprotokolle}
\index{p00197@Protokollarchiv}

Alltägliche Verwaltungsarbeit erzeugt einen nützlichen Verlauf: Uploads, Galerieänderungen, Sicherheitsentscheidungen, Wartungsergebnisse, Warnungen und andere Betriebsereignisse können im Administrationsprotokoll vermerkt werden. Dieser Verlauf hilft zu verstehen, was wann geschah und welches Konto oder welche Anfrage beteiligt war. Auf einer aktiven Installation wächst die Tabelle normalerweise stark. Version~\version{} behandelt Protokollpflege daher als eigene ausdrückliche Wartungsaufgabe, statt allgemeine Datenbankbereinigung unerwartet Auditverlauf entfernen zu lassen.

Die Administrationsprotokollansicht kann einzelne Einträge zeigen oder wiederholte Ereignisse zusammenfassen. Eine Gruppe kann wiederholte Vorschaubildwarnungen, Uploadresultate oder Wartungsmeldungen darstellen; beim Öffnen erscheinen die einzelnen Datensätze. Filter und Ablaufstatussteuerungen grenzen die Betriebsprüfung ein; begrenzte Exporte erhalten ausgewählten Verlauf, ohne die gesamte Tabelle in einer Anfrage zu laden. Große Listen, das Aufklappen einzelner Rohereignisse und Exporte werden in begrenzten Teilen verarbeitet, sodass die Anforderung eines vollständigen Verlaufs weder Browser noch PHP zwingt, die ganze Tabelle gleichzeitig im Speicher zu halten.

Der Betreiber kann wählen, wie lange aktuelle Protokolle in der aktiven Datenbank bleiben. \emph{Für immer} deaktiviert automatische Archivierung. Eine endliche Aufbewahrungsdauer verwirft ältere Datensätze nicht sofort. Der Wartungsablauf wählt zunächst vollständige Kalendertage außerhalb des aktiven Aufbewahrungsfensters, schreibt ein geschütztes Archiv mit maschinenlesbarem JSON und lesbarem statischem HTML-Bericht, prüft die erwartete Datensatzanzahl und entfernt erst dann genau diese Zeilen aus der aktiven Tabelle. Das Archiv liegt unter \filepath{data/admin-log-archives/}, einem Anwendungsdatenbereich statt eines öffentlichen Galeriebereichs, und besitzt eine zusätzliche Webserver-Schutzregel.

Archivwartung ist fortsetzbar. Sie schreibt vor der Bereitstellung eines fertigen Archivs temporäre Dateien, vermerkt eine kleine Zustandsdatei und verwendet eine Sperre, damit zwei Wartungsanfragen denselben Verlauf nicht gleichzeitig bearbeiten. Stoppen Hostinggrenzen, Verbindungsabbruch oder ein anderes Problem die Arbeit, kann die Administrationsansicht den unvollständigen Zustand zeigen; der nächste Wartungsversuch kann ihn prüfen, fortsetzen oder wiederherstellen. Ein Archiv gilt niemals allein wegen einer Datei mit ähnlichem Namen als Grund zum Löschen aktiver Datensätze; Zeilenanzahl- und Archivmetadatenprüfungen gehören zur Sicherheitsgrenze.

Administratoren sollten weiterhin Datenbank und Verzeichnis \filepath{data/admin-log-archives/} in abgestimmte Sicherungen einbeziehen. Die aktive Datenbank enthält neueren Betriebsverlauf; das Archivverzeichnis enthält durch die ausdrückliche Aufbewahrungsrichtlinie ausgewählten älteren Verlauf. Beide zusammen bieten dem Betreiber eine durchgehende Aufzeichnung und erleichtern die Untersuchung einer wiederhergestellten Installation nach Aktualisierung, Migration oder unerwartetem Ereignis.
\subsection{Vollständige aktuelle Funktionsreferenz}
\label{sec:complete-feature-reference}
\index{f00113@Funktionsreferenz}
\index{f00115@Funktionsverzeichnis}

Diese Referenz ist die Vollständigkeitscheckliste des Handbuchs für Version~\version{}. Sie führt die aktuellen Produktfähigkeiten auf, die öffentliche Routen, Administrationsnavigation, Galerie-/Fotoeditoren, Wartungswerkzeuge und optionale Integrationen bereitstellen. Interne Hilfsfunktionen erscheinen nicht als eigene Funktionen, wenn sie lediglich eine der folgenden Zeilen umsetzen.

\begin{longtable}{>{\RaggedRight\arraybackslash}>{\hspace{0pt}}p{0.23\textwidth} >{\RaggedRight\arraybackslash}>{\hspace{0pt}}p{0.66\textwidth}}
\toprule
\textbf{Funktion} & \textbf{Implementiertes Verhalten}\\
\midrule
\endfirsthead
\toprule
\textbf{Funktion} & \textbf{Implementiertes Verhalten}\\
\midrule
\endhead
Dateisystemorientierte Galerien & Verschachtelte physische Ordner mit gespiegelten Datenbankdatensätzen, praktisch unbegrenzte Hierarchie, Erkennung/Import, manuelle Erstellung, physisches Verschieben, sprechende öffentliche Pfade, Begleitmetadaten und stabile Geschwister-/Fotoreihenfolge.\\
Galeriemetadaten & Titel, Beschreibung, gepflegte Übersetzungen, manuelles Datum/Datumsbereich, EXIF-Datumsvorschläge, Slug, Ordnername, übergeordnete Galerie, Reihenfolge, Schlagwörter, Titelbild, hochgeladenes Galerievorschaubild, Branding, Hintergrund und Präsentationsüberschreibungen.\\
Galeriezugriff & Öffentlich, nicht veröffentlicht, privat, Passwortschutz, ablaufende/unbefristete Freigabelinks, geerbte Zugriffsauswertung, Administratorzugriff und NSFW-/18+-Bestätigungsgrenzen.\\
Sammelgalerieverwaltung & Zweigbewusstes Löschen, Sichtbarkeitsänderungen, EXIF/GPS-Kartenüberschreibungen, Abstimmungsänderungen, Dateinamensanzeigeänderungen, Picture-Game-Änderungen, Pfadneuerzeugung, Vorschaubild-/Scanaktionen und Abläufe für Geschwisterreihenfolge/Datumssortierung.\\
Smart Galleries & Gespeicherte verschachtelte boolesche Abfragen über physische Galerien, Schlagwörter, Daten, EXIF/GPS, Text, KI-Metadaten, Duplikatstatus, Dateieigenschaften und private redaktionelle Bewertung; öffentliche/nicht gelistete/Stamm-/eingebundene Platzierung; ober-/unterhalb-Reihenfolge je übergeordneter Galerie; Zyklenschutz; unabhängige Präsentationsüberschreibungen; vollständiger seitenübergreifender Lightbox-Durchlauf; begrenzter ZIP-Download.\\
Fotoeditor & Titel/Beschreibung/Übersetzungen, Sichtbarkeit, NSFW, Sortierreihenfolge, private redaktionelle Bewertung, Schlagwörter, fotoeigene Vorschaubildgrenzen, EXIF/GPS-Prüfung, interne KI-Metadatenprüfung und OpenAI-Entwurfswerkzeuge.\\
Sammelfotoverwaltung & Verschieben in vorhandene/neue Galerie, Löschen, Titelbildzuweisung, Sichtbarkeit, NSFW, Vorschaubilderstellung und Sortierung durch Ziehen/Dateiname.\\
Schlagwörter & Gemeinsame Galerie-/Fotoschlagwörter, normalisierte wiederverwendbare Namen/Slugs, öffentliche Beschreibungen, Nutzungszahlen, gewichtete Editorvorschläge, administratives Umbenennen/Bearbeiten/Löschen, öffentliche Schlagwort-Zielseiten, eigene Raster-/Karteneinstellungen sowie Reihenfolge/Einblenden/Scrollen der Hero-Schlagwörter.\\
Öffentliche Start- und Galerieseiten & Hierarchische Karten, Brotkrümelnavigation, Titelbilder/Collagen, Seitennavigation, direkte Fotoraster, Galerieanzahlen, Datumsangaben, Beschreibungen, Branding/Hintergründe, Favoriten-Kopfzeilenverknüpfungen, kontextbezogene Bearbeitungs-/Entfernungs-/Sortiersteuerungen für angemeldete Administratoren, schwebende Nach-oben-Hilfe bei langen Listen, responsive öffentliche Pfade und Abfrageparameter-Ersatz bei deaktivierter Umschreibung.\\
Öffentliche Suche & Optionale Website-/Zweigsuche über Galerie-/Fototitel, Beschreibungen, Dateinamen, Schlagwortmetadaten, gepflegten Sprachtext und verfügbaren internen durchsuchbaren KI-Text; Mindestanfragelänge und begrenzte Ergebniszahl; Zugriffs-/Listungsrichtlinie vor Ergebnisausgabe angewandt.\\
Robots, Sitemap und SEO & Zugriffsbewusstes \codeword{robots.txt}, bildbewusstes \codeword{sitemap.xml}, kanonische/soziale/strukturierte Metadaten und der oben beschriebene routenspezifische Abfrageparameterschutz.\\
Öffentliche Medienautorisierung & Eigene Routen für Originale/Vorschaubilder/Titelbilder/Branding/Hintergrund/Favicon bewerten Galerie-/Bildsichtbarkeit und geschützten Zugriff neu, statt beliebige Dateisystempfade offenzulegen.\\
Vorschaubilderzeugung & Responsive Größenvarianten, JPEG/WebP-Kompatibilitätsrichtlinie, responsive und progressive Renderer der ausgewählten Galerie, verzögertes/viewportnahes Verhalten, galerie-/fotoeigene Größengrenzen, DNG-Anzeigemaster, Metadatenabgleich, Administrations-/Browserneuaufbau und signierte zugriffsgeprüfte öffentliche Ersatzvorbereitung mit kleinem Budget.\\
Lightbox & Autorisierte asynchrone Metadaten, Tastatur-/Berührungsnavigation, Diashow/Vollbild, Einzelbild-/Bildstreifen-/3D-Karussellmodi, EXIF-/Kartenüberlagerung, Abstimmungszustand, Zoom von 100--400~\% mit Verschieben/Zweifingergeste/Cursorverankerung, sofortige Originalqualitätssteigerung beim Zoom und Umschalt+Pfeil-Sprünge um 10 Fotos.\\
EXIF und GPS & Kamera-/Objektiv-/Belichtungs-/Datums-/GPS-Extraktion, globale Vorgabe mit galerieeigenem Erben/erzwungenem Ein-/Ausschalten, Fotomarkierungen, Galeriekartendaten, Datumsvorschläge, Duplikathinweise und datenschutzrelevante öffentliche Unterdrückung.\\
Abstimmung & Umkehrbare positive Bildstimme mit synchronisierter Bewertung, anonymem Besucherhash oder angemeldeter Zuordnung, Ratenbegrenzung neuer Likes und galeriebezogenem Ein-/Ausschalten.\\
Picture Game & Optionaler Paarvergleich auf Zweigebene mit geeigneten öffentlichen Bildern, dauerhaftem Paar-/Siegerzustand, Ranglisten-/Ergebnisaggregation und schemabewusster Verweigerung bei unbekanntem Speicher.\\
Picture Manager & Verwaltungsüberlagerung der öffentlichen Ansicht für Angemeldete mit gemischter Foto-/physischer Untergalerieauswahl, wiederherstellbarem Löschen, physischer Galerieerstellung unter gewählter übergeordneter Galerie mit Teilbaumkopie und Rücknahme, reinen Fotoaktionen für Verschieben/Kopieren/Ziehen/Teilen, nativem Mehrdatei-Web-Share, sofern unterstützt, und ZIP-Ersatz/Download ausgewählter Fotos.\\
Downloads & Autorisierte Galerie-ZIP-Dateien, Smart-Gallery-ZIP-Dateien, ZIP-Dateien ausgewählter Fotos, administratives Gesamtgalerie-ZIP, normalisierte Archivpfade, Cacheinhaltssignaturen, begrenzte Smart-Gallery-Quellanzahl/-Bytes und sichere Cachebereinigung.\\
Browseruploads & Uploadabläufe für vorhandene/neue Galerien, ausdrückliche Browser-/Serververarbeitungswahl, vollständige Prüfung vorbereiteter Vorschaubilder, begrenzte Arbeitsprozesse/Chargen mit einzelnen übergroßen Paketen unter der harten PHP-Grenze, Formatrichtlinienhinweise, lokaler ZIP-Import und automatische Umbenennung nach dem Scan.\\
Dateisystemerkennung & Sitzungsgestützter begrenzter Ordnerscan; überspringt symbolisch verknüpfte/verborgene/erzeugte Verzeichnisse; importiert bildhaltige Hierarchie an Ort und Stelle; kann nicht verwaltete Fotos in vorhandene Galerien verschieben oder sichere nicht verwaltete Kandidatenbäume löschen; meldet reine Metadaten\-begleit\-dateien separat.\\
Dateisystem-Metadaten\-begleit\-dateien & Optionale galerieeigene \filepath{gallery.json} aus bearbeitbarem Galeriezustand geschrieben und bei Erkennung/Elternreparatur für unterstützte Metadaten gelesen; nützlich für Dateisystemportabilität, aber kein Ersatz für Datenbanksicherung.\\
Upload-API & Galeriegebundene einmal sichtbare API-Schlüssel, zentraler API-Manager, Schlüsselwiderruf, gemeinsame Aufnahmepipeline und authentifizierter Automatisierungsendpunkt.\\
Windows-Begleitanwendung & Uploads aus überwachten Ordnern, manueller Sammelupload, optionale lokale responsive Vorschaubilderzeugung, parallele/Pipelinearbeit, Wiederholungs-/Zustandsdateien, Infobereich, optionale SimConnect-Kamerakoordinaten, Quelldateilöschung nach bestätigtem Erfolg und optionaler lokaler KI-Metadatenarbeitsprozess.\\
Mobiles WebDAV & Galeriegebundene Zugangsdaten mit einmaligem Passwort, Basic-Authentifizierung, von mobilen Transferclients benötigte WebDAV-Erkennungs-/Schreibmethoden, JPG/PNG/GIF/WebP-Aufnahme, Nachverfolgung letzter Nutzung und Widerruf.\\
Medienumbenenner & Probelaufplan physischer Dateinamen, Vorgabe \codeword{\{gallery_path\}_\{seq4\}} und dokumentierte Platzhalter, Kollisions-/Sicherheitsbehandlung, Galeriereihenfolgekontext, abgestimmte Datenbank-/öffentliche Pfad-/Ableitungs-/Cacheaktualisierungen, begrenzte websiteweite Anwendung und automatische Umbenennung nach Upload.\\
Metadatenorganizer & Reine Vorschau von EXIF-Datumsgruppierung mit frühestem/spätestem Datumsfilter, ISO-\codeword{YYYY-MM-DD}-Unterordnern, Wiederverwendung passender direkter Datumsuntergalerien, bestätigtem physischem Verschieben samt Ableitungen und begrenzten Übernahmechargen; Ortsgruppierung bleibt geplant statt aktiv.\\
Fotoduplikatdetektor & Exakte SHA-256- und metadatenunterstützte mögliche Paare, Kontextlinks, Löschablauf, administratorbezogene Prüfliste für geprüfte Paare/genaue Galerien und dauerhafter Prüfzustand.\\
OpenAI-Texthilfe & Kontoindividuelle verschlüsselte API-Schlüssel-/Modelleinstellungen, optionaler ausdrücklicher Vorschaubild-Zustimmungsschalter, Galerie-/Fotobeschreibungsentwürfe, Textbereinigung, Kontextzusammenfassung, Beschreibung sichtbarer Inhalte aus kleinen erzeugten Vorschaubildern, Übersetzungsvorschläge und ausschließlich manuell zu prüfende Einfügung.\\
Interne KI-Metadaten & Getrennte durchsuchbare Maschinenmetadaten und Warteschlangentabellen, authentifizierte Windows-Arbeitsprozessübernahmen/-ergebnisse, schreibgeschützte Fotoeditorprüfung, Smart-Gallery-/Suchhinweise, Modellgenerationsnachverfolgung und erzwungene Zweigneuverarbeitung.\\
Flugkarten und SimBrief & Manueller Routentext, ausdrückliche Koordinatensyntax, SimBrief-Abruf des neuesten OFP, bearbeitbares erzeugtes Markdown, gespeicherte OFP-Daten und Koordinaten sowie öffentliche Darstellung ausschließlich aus gespeicherten Punkten.\\
Navigationsdaten & Lokaler OurAirports-Datenbankimport, mitgeliefertes Offline-CSV, zwischengespeicherte Anbieterpunkte, optionale Navigraph-OAuth-/Paket-/AIRAC-Integration, vorrangig offline arbeitender Auflöser und administrativer Suchtester.\\
Galeriemigration & API-Übertragung zwischen Installationen derselben Version per Push von der Quelle oder Pull durch das Ziel, versionierte Manifeste, Metadaten-/Ressourcenprüfsummen, auf eine Ressource begrenzte Anfragen, dauerhafter Zielstatus, erneute Verbindung/Fortsetzung und separater Abschluss.\\
Design und Layout & Direkte Palettenfarben, Schriftmodus, abgerundete Ecken, Seitenbreite, Galerie-/Kartenraster, Seitennavigation, Kartenlayout, Anzahlkennzeichnung, öffentlicher Vorschaubildrenderer, Lightbox-Vorgabe, GPS-Markierungspräsentation, drei Favoriten-/Kopfzeilenverknüpfungen und Live-Vorschau. Die aktuelle Designoberfläche besitzt keinen getrennten Hell-/Dunkelmodus-Schalter.\\
Branding und Browseridentität & Website-Kopfzeilenbanner, Trennlinienabmessungen/-streckung, optimierter Hintergrund und Deckkraft, galerieeigene Überschreibungen, öffentliche Titelbildressourcen, quadratisch zugeschnittenes Favicon mit erzeugten 32-/48-/180-Pixel-Varianten und benutzerdefinierte CSS-Vorlagen/-Editor.\\
Lokalisierung & Kanonischer englischer Ersatz und gepflegte auswählbare tschechische, deutsche und schwedische Oberflächenpakete, besucherindividuelle Browserpräferenz, Auswahlgestaltungen mit Flaggen/muttersprachlichen Namen, gemeinsame Browser-/Serverkataloge, Paketdiagnose/-editor/-import/-export und gepflegte galerie-/fotoeigene Sprachinhalte.\\
Funktionsschalter & Globale routenbewusste Schalter für öffentliche Suche, Lightbox-Modi, Picture Manager, Bildabstimmung, Picture Game, Downloads, Besucherkonten und private Sammlungen, Galeriekarten, Flugkarten, Navigationsdaten, SimBrief, OpenAI, lokale KI-Metadaten, Upload-API, mobiles WebDAV, Galeriemigration, Medienumbenenner und Telemetrie.\\
Zentrale Einstellungen & Geführter stufenweiser Einrichtungsassistent mit abschließender Zustimmung, Live-Vorschau und sicheren Verantwortungsadaptern; durchsuchbare kategorisierte Registrierung, tastaturgesteuerte Spotlight-Navigation, kanonische Speicherdelegation genehmigter globaler Werte, Spezialtiefenlinks, vererbungsbewusste Zuständigkeit, Vorgaben-/Quellkennzeichnung und Geheimnisausblendung.\\
Administrativer Galeriearbeitsbereich & Sichtbarkeitsfilter mit Auswahlaufhebung verborgener Zeilen, dauerhafter Klappzustand je Zweig, Alle zuklappen/aufklappen, Hierarchiesortierung, öffentliche Sortierwerkzeuge der sichtbaren Seite, Seitenbereichseditoren und kontextbezogene öffentliche Bearbeitungs-/Entfernungs-/Sichtbarkeitssteuerungen für angemeldete Administratoren.\\
Administratorauthentifizierung & Anmeldung per Benutzername oder Wieder\-her\-stel\-lungs-E-Mail, CSRF-/Sitzungsschutz, gehashte Ratenbegrenzungssubjekte, optionales dauerhaftes Anmeldetoken, Passwortzurücksetzung per PHP-E-Mail/SMTP, Konto-/Profiländerungen mit aktuellem Passwortschutz und optionale Google-Anmeldung, die zuerst aus einer authentifizierten Passwortsitzung verknüpft werden muss und zum Trennen das aktuelle Passwort benötigt.\\
Besucherkonten, Lebenszyklus, Favoriten, private Sammlungen und nicht gelistete Freigabe & Direkte administrative Besuchererstellung/-löschung und Registrierung nur auf Einladung, administrative Steuerungen für Sperren/Wiederherstellen/Überall abmelden mit zentralem Besucherberechtigungswiderruf, erzwungener Austausch administrativ bereitgestellter temporärer Passwörter vor normaler Besucherberechtigung, scannersichere E-Mail-Bestätigung und Aktivierung, getrennte Besucheranmeldung/-abmeldung, eigene rotierende dauerhafte Zugangsdaten, allgemeine scannersichere Passwortzurücksetzung, aktuelle Passwort-Neuauthentifizierung für sensible Kontoaktionen, Passwort-/E-Mail-Änderung, Besucherselbstlöschung, private nicht zwischenspeicherbare Konto-/Favoriten-/Sammlungs-/Lebenszyklusantworten, Favoritensteuerungen auf autorisierten Karten/Lightbox-Bildern, besuchereigene geordnete Sammlungen mit aktueller Quellautorisierung, ein widerrufbarer nicht gelisteter schreibgeschützter 30-Tage-Sammlungslink mit einmal gezeigtem Geheimnis und bereinigter Weiterleitung, Schnittmenge mit Quellzugriffsregeln im Empfängerkontext ohne Administratorausnahme und strikte Trennung von Galerieautorisierung; keine offene Registrierung, öffentliche Sammlungssuche/Profile, Empfängerzusammenarbeit, Kommentare oder Uploads.\\
Datenbankbereitschaft/\allowbreak Systemzustand & Dreistufige Schemaprüfung mit Trennung von Verfügbar, Fehlend, Unbekannt und gegebenenfalls Deaktiviert; getrennte Bereiche für Sicherheits-/Authentifizierungs-, änderungssensible und optionale Präsentations-/Berichtsfähigkeiten; Schreibvorgänge werden bei unbekanntem erforderlichem Schema sicher verweigert.\\
Datenbankwartung & Ausdrücklicher Schema-/Migrations-/Codereferenzbestand, Speicherinformationen, Probelaufbereinigungskandidaten, begrenzte logische Bereinigung, eigene Bereitschaft für ältere Reparaturen, ausgewähltes ANALYZE, Optimierungsvorschau ausgewählter Tabellen und separat bestätigtes OPTIMIZE TABLE.\\
Speicherstatistik & Quell-/erzeugte-Medien-Scan auf Anforderung, Dateityp-/Größenverteilung, größte Galerien/Dateien, DNG-Master- und Vorschaubilderfassung, Behandlung veralteter Momentaufnahmen, getrennte Datenbank-/Tabellennutzungsansicht, kombinierte Aktion \ui{Alles aktualisieren} für Dateien, Datenbankschätzungen und schreibgeschützte Prüfung sowie begrenzter Fortschritt.\\
Gesamtbericht & Eigenständiger HTML-Export über Speicher, Datenbank, Galerie-/Bild-/EXIF-/GPS-Statistik, Funktions-/Einstellungszustand, Telemetrie, Protokolle und Laufzeitdiagnose mit begrenzter Bild-/erzeugte-Medien-Verarbeitung, Wiederholungen vorübergehender Anfragefehler und begrenzten Zusammenfassungen vieler verschiedener Werte.\\
Telemetrie & Optionale datenschutzkontrollierte anonyme Nutzungs-/Leistungs-/Cache-/Datenbankmessungen, DNT-Berücksichtigung, Administratorausschluss, begrenzte Betrachtungszeit, rohe/stündliche/tägliche Aufbewahrung, Aggregation/Bereinigung, Administrationsübersicht und HTML-Export.\\
Administrationsprotokolle & Strukturierte Ereignisse, Filter/Status, Gruppierung wiederholter Ereignisse mit Einzelansicht, begrenzte Exporte/ZIP-Export, konfigurierbare aktive Aufbewahrung, geschützte tägliche JSON+HTML-Archive, Fortsetzungs-/Wiederherstellungszustand und Archivdownload-/Anzeigeinstrumente.\\
Wartungszentrale & Dauerhafter Ablauf Analysieren--Prüfen--Ausführen--Verifizieren, gruppierte Aufgabenprüfung mit eingeklapptem Abschnitt für optionale Aufgaben ohne erkannten Arbeitsbedarf, auffindbare tiefe Medienprüfung und vollständige Optimierung, serverautorisierter Aufgabenplan, begrenzte Kontrollpunkte, Pausieren/Fortsetzen/Abbrechen, zentraler Änderungsanspruch, physische Datenbankschritte je Einzeltabelle und geprüfter Vorher-/Nachherbericht.\\
Geplante Websitewartung & Tägliches UTC-Fenster, optionale anfrageausgelöste Fortsetzung, begrenztes Bildchargen-/Zeitbudget, CLI-/Web-Cron-Unterstützung, Aktion für einen Teilabschnitt, Rücksetzung unterbrochenen Zustands, austauschbares verborgenes Web-Cron-Token, Vorschaubild-/Metadaten-/Bereinigungs-/Telemetrie-/Protokollarchivaufgaben und Nachgeben bei zentralem Änderungsanspruch der Wartungszentrale.\\
Laufzeitdiagnose & PHP-/Hosting-/Datenbankbereitschaft, GD-/Imagick-/ImageMagick-/Formatunterstützung, DNG-Richtlinie, kopierbare Supportzusammenfassung, administratorbeschränkte Entwicklungsdiagnose für Betrachtervorladen/Cache/Speicher/Netzwerk/Bildzeitmessung und entwicklungsbeschränkte Benchmarkeinstiege.\\
Integrität & Hashprüfung des installierten Kernmanifests für verwaltete Anwendungsdateien und Manifestgenerierung/-prüfung bei Veröffentlichung; getrennt von Benutzerdatensicherung.\\
Galeriebenchmark & Entwicklungsbeschränkter authentifizierter Browserbenchmark mit anonymen Vorschauladungen, verborgenen Frames desselben Ursprungs, Serverdarstellungszählern und Navigationstiming, anschließendem JSON-Export für Vorher-/Nachhervergleiche.\\
Aktualisierungssystem & Stable-/Beta-Erkennung, zwischengespeicherte Versionshinweise, fortsetzbare dauerhafte Aufträge, Downloadfortsetzung per Range, Archiv-/Manifestprüfung, Vorbereitung, Rücknahmemomentaufnahme, Aktivierung, Migrationen, Bereinigung, Sperren/Wiederherstellung veralteter Sperren, Wiederholung/Abbruch vor Aktivierung, Dateirücknahme nach Aktivierung, Stable-Wiederherstellung/saubere Neuinstallation/Beta-Installation, CLI-Fortsetzung und Notfall-\filepath{reset.php}.\\
Bereitstellungshelfer & Windows-/Linux-Bereitstellungsskripte, Veröffentlichungsmetadaten, Kernmanifestgenerierung und Paketprüfung; die Bereitstellung ist darauf ausgelegt, lokal installiertes PHP nicht zur unbedingten Laufzeitvoraussetzung der Galerie selbst zu machen.\\
\bottomrule
\end{longtable}
\subsection{Anwendungsaktualisierungen}

Die kompakte Versionsübersicht fasst installierte und verfügbare Version, Kanal, Status und Hauptaktionen zusammen; Repository- und GitHub-API-Details bleiben in einem eingeklappten Diagnosebereich. Nach Abschluss eines Auftrags aktualisiert der Browser die Übersicht direkt aus lokalen Metadaten, hält das Panel offen und die URL unverändert und blendet eine überholte Aktion für stabile Updates aus. Dabei erfolgt keine weitere GitHub-Prüfung. Schlägt das Laden fehl, wiederholt die Schaltfläche zur Statusaktualisierung nur den lokalen Lesevorgang, ohne die Installation zu wiederholen.

Unter Windows sind für das atomare Ersetzen eines Aktualisierungs-Kontrollpunkts bis zu zehn Versuche mit Pausen von 50~ms und einer gesamten Pausengrenze von 450~ms zulässig, um kurze Dateifreigabesperren zu überbrücken. Scheitert der Schreibvorgang weiterhin, bleibt der vorherige Kontrollpunkt erhalten. Bei dauerhaften Berechtigungsproblemen muss vor einem erneuten Versuch der Speicherzugriff korrigiert werden. Windows verwendet dieselbe begrenzte Wiederholungsregel bei der Migration von Ausrichtungsmetadaten in \filepath{gallery.json}; eine gleichzeitige Änderung oder dauerhafte Verweigerung bewahrt das vorhandene Dokument und stoppt die Migration.

Das Aktualisierungssystem prüft konfigurierte Kanäle, validiert Pakete, bereitet Inhalte vor, prüft erforderliche Laufzeitdateien und Größen, bewahrt überschriebene Dateien im Wiederherstellungsspeicher auf und aktiviert die Aktualisierung. Ratenbegrenzungsbehandlung und Wiederholungszustand verhindern wiederholte entfernte Anfragen während Anbieterwartezeiten. Eine fortsetzbare Auftragskarte erscheint sowohl unter \ui{Status} als auch unter \ui{Erweiterte Werkzeuge}, sodass Beta-Installation, Wiederherstellung oder Neuinstallation ihren Fortschritt weiterhin am Ausgangsort zeigt. Die Seite darf geschlossen und wieder geöffnet werden, ohne abgeschlossene Kontrollpunkte zu verwerfen.

Der administrative Aktualisierungsbereich zeigt außerdem zwischengespeicherte Veröffentlichungsinformationen und versionierte Versionshinweise für den verfügbaren Kanal, damit der Betreiber eine Veröffentlichung vor Auftragsbeginn prüfen kann. Stable- und Beta-Kanäle bleiben getrennte Entscheidungen; die Beta-Bezeichnung gehört zur Veröffentlichungsauswahl der Administrationsoberfläche und verändert nicht die Terminologie dauerhafter Aktualisierungsaufträge oder Manifeste.

Versionshinweise werden als rohes Markdown in~der~installierten Datei \filepath{PATCH_NOTES.md} und~im~versionsbezogenen Cache gespeichert. Die~Administrationsansicht rendert die~ausgewählte Version erst beim Anzeigen, auch wenn Einträge aus~älteren Caches stammen; früher zwischengespeichertes HTML wird verworfen statt als vertrauenswürdig übernommen. Das~bloße Öffnen verwendet lokale oder bereits zwischengespeicherte Inhalte ohne weitere GitHub-Anfrage. Ein~separater zeitlich begrenzter Schritt der~Versionssuche kann entfernte Hinweise aktualisieren; bei einem Fehler bleiben die~mitgelieferten Hinweise verfügbar.

Markdown-Links wie~\codeword{[Aufgabe](https://example.org/issue/123)}, Autolinks in~spitzen Klammern und~reine HTTP(S)-Adressen werden anklickbar. Sie öffnen einen neuen Tab mit~den~Schutzattributen \codeword{noopener} und~\codeword{noreferrer}. Andere URL-Schemata bleiben maskierter, inaktiver Text. Die~Darstellung bewahrt Überschriften, Listen, Hervorhebungen, Codeblöcke und~wörtlichen Inline-Code; Adressen innerhalb eines Codespans werden nicht verlinkt. Interpunktion, Unterstriche und~Abfrageparameter von~Links bleiben erhalten, ohne beliebiges HTML zuzulassen.


Vorübergehende Download- oder Hostingfehler können erneut versucht werden. Ein Paket, dessen Dateien nicht zum eigenen Integritätsmanifest passen, ist anders: Erneutes Herunterladen derselben veröffentlichten Version verändert deren Bytes nicht. PHP Gallery fordert daher den Administrator auf, diesen vorbereiteten Auftrag abzubrechen und eine neuere Veröffentlichung oder Beta-Code zu wählen, statt eine endlose Wiederholungsschleife anzubieten.

Stable-Erkennung und Paketverarbeitung sind getrennte begrenzte Vorgänge. Eine Erkennungsanfrage kann einen dauerhaften Hintergrundauftrag erstellen; eine spätere sichere Anfrage oder der CLI-Läufer führt einen kurzen Teilabschnitt über Download, Archivprüfung, Entpacken, Manifestprüfung, Vorbereitung, Rücknahmesicherung, Aktivierung, Migrationen und Bereinigung aus. Range-Anfragen können unterbrochene Downloads fortsetzen, Arbeitsprozesssperren verhindern paralleles Fortschreiten, und eine veraltete Sperre kann nach Ablauf ihres Besitzers wiederhergestellt werden. Planen Sie auf wenig frequentiertem Shared Hosting \codeword{php scripts/application_update.php} per Cron, beispielsweise alle fünf Minuten. Jeder Aufruf führt begrenzte Erkennung aus oder setzt vorhandene Hintergrundarbeit fort; die Korrektheit hängt nicht von \codeword{set_time_limit()} oder \codeword{ignore_user_abort()} ab.

Aktuelle Veröffentlichungen gleichen außerdem die kleine Menge anwendungseigener Serverrichtliniendateien während Vorbereitung und Aktivierung ab. Das schützt Installationen, deren älteres Aktualisierungssystem eine Richtliniendatei übersprungen hat: Die validierte Veröffentlichungskopie wird verglichen und atomar bereitgestellt, Rücknahmemetadaten werden vor dem Austausch erweitert, und ein migrationsgestützter Initialisierungsweg kann einen bereits aktivierten älteren Baum reparieren. Der Abgleich ist auf bekannte Richtlinienpfade begrenzt und verändert weder Galeriemedien noch Benutzerkonfiguration.

Vor Aktivierung kann ein Administrator einen wiederherstellbaren Auftrag abbrechen oder erneut versuchen. Danach kann die administrative Wiederherstellungsaktion verwaltete Anwendungsdateien aus der gespeicherten Rücknahmemomentaufnahme wiederherstellen; Dateirücknahme kehrt Datenbankmigrationen jedoch nicht um. Stable-Wiederherstellung, saubere Neuinstallation und Beta-Installation bleiben fortsetzbare Aufträge mit eigener Paketprüfung. Der eigenständige Notfallweg \filepath{reset.php} ist ein getrenntes Wiederherstellungswerkzeug zur Rückkehr auf den Stable-Zweig, wenn die normale Administrationsroute nicht nutzbar ist; er erfordert weiterhin dieselbe sorgfältige Sicherungs- und Integritätsprüfung.

Während des kurzen kritischen Aktivierungsabschnitts verhindert eine abhängigkeitsfreie frühe Laufzeitsperre, dass neue gewöhnliche Anfragen eine teilweise ersetzte Anwendung laden. Solche Anfragen erhalten eine private nicht zwischenspeicherbare \codeword{503 Service Unavailable}-Antwort, bis der Aktivierungsabschluss dauerhaft vermerkt ist. Der authentifizierte Aktualisierungs-Wiederherstellungs-/Statusweg bleibt verfügbar, sodass eine unterbrochene Aktivierung fortgesetzt werden kann. Beschädigter oder unvollständiger Sperrzustand verweigert sicher den Zugriff, statt die Installation als fehlerfrei darzustellen.

Dieselbe frühe Laufzeitgrenze wandelt nicht abgefangene und abfangbare fatale PHP-Fehler in begrenzte \codeword{500}-Antworten um, wenn PHP die Antwort noch steuern kann. Produktivhosting muss \codeword{display\_errors} deaktiviert halten; andernfalls kann die PHP-Laufzeit selbst Warnungs- oder fatale Fehlerdetails ausgeben, bevor die Anwendung eine sichere Antwort formatieren kann.

\important{Erstellen Sie vor einer Anwendungsaktualisierung eine abgestimmte Sicherung. Halten Sie sie verfügbar, bis Laufzeitprüfungen, Migrationen, öffentliche Galerien, geschützter Zugriff, Uploads und administrative Vorgänge geprüft wurden.}

Vor der Aktivierung verlangt die Archivprüfung des Aktualisierers auch den vollständigen neuen Anfragekernel und seinen geprüften Modulplan. Ein unvollständiges Archiv wird abgelehnt, während vorbereitete Dateien außerhalb der aktiven Installation bleiben. Halten Sie bei manueller Bereitstellung alle ausgelieferten Laufzeitdateien zusammen; ein Austausch nur des Bootstrap würde eine unvollständige Anwendung hinterlassen.

Die Update-Aktivierung verwendet die geprüfte Updater-Teilmenge des Produktionsinventars. Eingehende Inventare führen WinApp-Quellen getrennt auf, damit ältere installierte Updater das Release vor der Installation seines Lesecodes prüfen können. Der CMS-Updater installiert diese Quellen nicht. Ältere Quellarchive ohne Inventar behalten einen begrenzten Kompatibilitätspfad anhand ihres eigenen Manifests; ein vorhandenes ungültiges Inventar wird abgelehnt. Das Entfernen veralteter Dateien erfordert frühere Eigentumsnachweise und bewahrt unbekannte lokale Dateien. Die Dateiwiederherstellung nutzt den serverseitig geschriebenen Sicherungsindex und lehnt nicht indexierte Dateien sowie durch symbolische Links oder Windows-Junctions umgeleitete Pfade ab. Diese Dateisicherung macht Datenbankmigrationen nicht rückgängig.

Wenn ein Update vor der Aktivierung gestoppt wurde, brechen Sie den alten Auftrag ab, erzwingen Sie eine neue Release-Prüfung und starten Sie ein neues stabiles Update, sobald das korrigierte Release verfügbar ist. Ein erneuter Versuch kann die ursprüngliche Zielversion des Auftrags beibehalten.

\subsection{Integritätsmanifest}
\index{k00156@Kernmanifest}

\filepath{app/core-manifest.json} vermerkt Anwendungsversion und Hashes integritätsverwalteter Kerndateien. Der lokale Generator unter \filepath{scripts/generate_manifest.php} berechnet nach Laufzeitänderungen endgültige Hashes. Die Einbeziehung von Dokumentation und Benutzerdaten folgt der etablierten Generatorrichtlinie.
Für~die~deterministische Release-Vorbereitung akzeptiert der~Manifestgenerator den~optionalen Build-Zeitstempel \codeword{SOURCE_DATE_EPOCH} für~das~Feld \codeword{generated_at}; ohne ihn gilt die~aktuelle Uhrzeit. Der~gehostete Workflow leitet diesen Zeitpunkt aus~den~Release-Metadaten ab und~verwendet ihn für~die~Handbuch-PDFs und~das~abschließende Manifest. Bei identischen vorbereiteten Eingaben bleiben die~Zeitstempel dadurch reproduzierbar. Nach Änderungen an~integritätsverwalteten Anwendungsdateien muss das~Manifest dennoch neu erzeugt werden.

\section{Optionaler Hintergrund: Wie PHP Gallery funktioniert}
\label{sec:architecture}
\index{a00022@Architektur}
\index{c00054@Controller}
\index{d00077@Dienste}
\index{a00018@Ansichten}

Die folgenden Seiten beschreiben den Anfrageweg der Anwendung. Sie betreffen hauptsächlich Wartung und Entwicklung; normale Galerieverwaltung erfordert sie nicht.

Der generierte Modulplan erklärt gemeinsame Abhängigkeiten und eine kanonische Dateireihenfolge, die die vollständigen Routenabhängigkeiten und Ladereihenfolge erhält. Halten Sie bei Routenänderungen Abhängigkeitseingaben und generierten Plan zusammen. Request-/Session-Adapter besitzen den Transportzugriff; Sprache, Navigationstoken, Galeriezugriffe, NSFW-Bestätigung und Flash-Nachrichten bleiben bei ihren bisherigen Verantwortlichen. Die Lightbox trennt Navigationsgenerationen und dekodierte Ressourcen; der Viewer behält Darstellung und autorisierte Qualitätsauswahl. Diese Trennung schafft Regressionsgrenzen, belegt jedoch keinen schnelleren Browserbetrieb. Siehe \filepath{docs/BROWSER_LIFECYCLE.md} und \filepath{docs/COMPATIBILITY_LIFECYCLE.md}.

\subsection{Initialisierung, Routing und Dispatch}
\index{i00148@Initialisierung}
\index{r00207@Routing}
\index{c00055@Controller-Dispatch}
\index{a00016@Anfrage!l00162@Lebenszyklus}

Gewöhnliche Anfragen gehen über \filepath{public/index.php} ein; \filepath{index.php} im Repository-Stamm ist ein unterstützter Kompatibilitätseinstieg für andere Hostingstrukturen. Der öffentliche Einstieg lädt \filepath{app/early\_runtime.php} vor \filepath{app/bootstrap.php}; danach erreichen beide Strukturen dieselbe normale Anfragepipeline.

Der kleine anfragelokale Kernel unter \filepath{app/runtime/} steuert den vorhandenen Lebenszyklus über Request, Router, RouteRegistry und ModuleLoader und ruft anschließend den prozeduralen Controller auf. Die kanonische Routen- und Sicherheitstabelle bleibt in \filepath{app/bootstrap/dispatch.php}. Ein geprüfter generierter Plan lädt logische Module für die gewählte Route; Produktionsanfragen scannen keine Quelldateien, kompilieren keine Abhängigkeiten und besitzen keinen Ersatzweg zum Laden der gesamten Anwendung. PHP~8.1 und bisheriges Shared Hosting reichen weiterhin aus, ohne Composer oder Node-Build. Ältere CLI-/Testverbraucher können \filepath{app/bootstrap\_full.php} für die vollständige prozedurale API verwenden.

Die Initialisierung liest lokale Konfiguration und stellt den Anfragekontext her: Datenbankzugriff bei Bedarf, Sprache, Sitzungszustand, Sicherheitsvoraussetzungen und gemeinsame Laufzeitdienste. Routing übersetzt anschließend die eingehende URL oder den Ersatzweg mit Abfrageparametern in eine interne Route mit Parametern wie Galeriepfad, Seitennummer, Mediengröße oder Freigabetoken. Dispatch ruft den für diese Route registrierten Controller auf.

Die Reihenfolge ist bewusst gewählt. Eine Anfrage muss klassifiziert werden, bevor Code über ihre Bedienung entscheidet; Zugriffsprüfungen müssen vor Zusammenstellung geschützter Ausgaben laufen. Eine direkt wirkende Bild- oder Vorschaubild-URL ist daher weiterhin eine Anwendungsanfrage; der Mediencontroller autorisiert sie vor Ausgabe von Bytes.

\begin{lstlisting}
browser request
  -> index.php or public/index.php
  -> app/bootstrap.php
  -> cms_run() -> Gallery\Core\Kernel
  -> Request -> Router -> RouteRegistry
  -> ModuleLoader -> app/bootstrap/dispatch.php
  -> controller function
  -> service functions
  -> HTML, JSON, media, archive, or redirect response
\end{lstlisting}

\subsection{Controller}
\index{c00054@Controller}

Controller verantworten die Koordination auf HTTP-Ebene. Sie erhalten aufgelöste Route und Anfragekontext, prüfen den Vorgang, rufen erforderliche Dienste auf und wählen den Antworttyp. Eine öffentliche Galerieroute liefert gewöhnlich HTML; eine administrative Seitenbereichsaktion kann JSON oder ein HTML-Fragment liefern; eine Medienroute liefert eine autorisierte Datei; eine Downloadroute kann ein Archiv streamen.

Öffentliche Galeriedarstellung beginnt in \filepath{app/controllers/public_gallery.php}; verwandte Verantwortlichkeiten sind auf gezielte Controllerdateien verteilt. Uploads, Wartung, Aktualisierungen, Downloads, Medienzugriff und Administrationsfunktionen besitzen eigene Einstiege, während gemeinsame Domänenregeln in Diensten bleiben.

Controller sind keine Abkürzung bei der Autorisierung. Vorschaubild-, Karten-, Lightbox-, Such-, Download- und direkte Medienrouten wiederholen die relevante Zugriffsentscheidung, statt anzunehmen, dass ein Besucher zuvor die Galerieseite erfolgreich geöffnet hat.

\subsection{Dienste}
\index{d00077@Dienste}

Wiederverwendbare Anwendungsregeln liegen in \filepath{app/services/}. Dienste lösen Galerien und Bilder auf, bewerten Zugriff, normalisieren Einstellungen, verwalten Uploads und Änderungen, bereiten Vorschaubild-/Medienmanifeste vor, führen Suchen aus, pflegen Metadaten und koordinieren weitere Domänenarbeit. Ein Controller sollte diese Vorgänge anfordern können, ohne ihre Regeln neu zu implementieren.

Die Trennung ist wichtig, weil dasselbe Foto über mehrere Routen erreichbar ist. Zugriffsrichtlinie, Vorschaubildauswahl oder Löschsemantik sollten nicht davon abhängen, ob die Anfrage von Galeriekarte, Suchergebnis, Lightbox, Administrationswerkzeug oder direktem Link stammt. Originalfotos bleiben im Galeriebaum; Dienste koordinieren die Datenbankdatensätze und erzeugten Ableitungen, die sie beschreiben oder unterstützen.

\subsection{Ansichten und Browsermodule}
\index{a00018@Ansichten}
\index{b00050@Browsermodule}

Ansichten stellen vorbereitete Daten als HTML dar. Sie verantworten Struktur und Präsentation, nicht die Entscheidung über Besucherautorisierung oder Vorschaubilderzeugung. Serverausgabe enthält semantische Links, Formulare, Titel, Alternativtext und anfängliche Medienkandidaten vor Beginn der JavaScript-Erweiterung.

Browsermodule ergänzen Lightbox-Verhalten, Suchvorschläge, Abstimmung, Uploadfortschritt, Drag-and-drop-Aktionen, Vorschaubildsteigerungen, Diagnose und Administrations-Seitenbereich. Der Code bleibt frameworkfrei. Öffentliche Seiten behalten somit eine nützliche serverseitige Grundlage; authentifizierte Abläufe können reichhaltigere direkte Interaktionen ergänzen.

Der Administrations-Seitenbereich ist dynamisch: Sein Fragment kann ersetzt werden, ohne die Seite zu verlassen. Module für seine Steuerungen verwenden daher delegierte oder lebenszyklusbewusste Ereignisbindung, damit eine frisch dargestellte Kopie weiter funktioniert und ein altes Fragment keine Ereignislistener mehr besitzt.

\subsection{Öffentliche Darstellungspipeline der ausgewählten Galerie}
\index{d00058@Darstellungspipeline}

Eine Anfrage für eine ausgewählte Galerie löst zunächst Galerie und wirksame Zugriffsrichtlinie auf, einschließlich geerbten Schutzes, Passwort-/Freigabezustands, Veröffentlichungszustands und Altersbeschränkungen. Nur erlaubte Inhalte gelangen zur Darstellung.

Der Controller lädt dann direkte Untergalerien und die aktuelle Fotoseite samt Metadaten für Titel, Beschreibungen, Schlagwörter, Daten, Stimmen, Reihenfolge und Navigation. Vorschaubildmetadaten werden in Chargen vorbereitet; anfragelokale Medienmanifeste beschreiben autorisierte Kandidaten jedes sichtbaren Bilds.

Die Ansicht erhält dieses vorbereitete Modell und stellt Seitentitel, Branding, Navigation, Brotkrümelnavigation, Karten, Seitennavigation, Lightbox-Konfiguration und Fußzeilenressourcen dar. Browsermodule erweitern das Ergebnis nach der Auslieferung. Zugriff bleibt Serververantwortung; JavaScript kann eine Interaktion schneller oder bequemer machen, verwandelt aber keinen nicht autorisierten Kandidaten in einen autorisierten.

Diese serverorientierte Trennung erhält auch nützliches Verhalten bei blockierten oder verspäteten Skripten. Direkte Links, semantischer Inhalt und anfängliche Bildauswahl stehen bereits im HTML; JavaScript übernimmt Funktionen, die tatsächlich einen interaktiven Client benötigen.

\section{Optionaler Hintergrund: Was die Datenbank speichert}
\label{sec:database}
\index{d00065@Datenbankschema}
\index{m00172@Migrationen}

Die Datenbank speichert Anwendungszustand, der nicht in Bilddateinamen gehört: Beschreibungen, Privatsphäre, Reihenfolge, Schlagwörter, Stimmen, Daten, öffentliche Adressen, Konten und Wartungsverlauf. Die folgende Übersicht genügt für Sicherungsplanung und Hostinggespräche; Spaltendefinitionen verbleiben in der Datenbankreferenz \cite{database}.

\subsection{Wesentliche Domänen}

\begin{longtable}{>{\bfseries}>{\hspace{0pt}}p{0.25\linewidth}>{\hspace{0pt}}p{0.67\linewidth}}
\toprule
Domäne & Dauerhafte Verantwortung\\
\midrule
\endhead
Benutzer und Authentifizierung & Administratoridentität, Passworthashes, Wieder\-her\-stel\-lungs-E-Mail, externe Identitätsverknüpfungen, Zurücksetzungstokens und bereichsgebundene Zugangsdaten.\\
Galerien & Ordnerpfad, Hierarchie, öffentlicher Pfad, Metadaten, Sichtbarkeit, Zugriff, Branding, Präsentationsüberschreibungen und betriebliche Reihenfolge.\\
Bilder & Galeriezuordnung, relativer Pfad, Dateiname, Metadaten, Abmessungen, Sichtbarkeit, Prüfsumme, EXIF-Daten, Reihenfolge und abgeleiteter Zustand.\\
Schlagwörter & Wiederverwendbare Schlagwortidentität und Beziehungen zu Galerien und Bildern.\\
Stimmen & Galerie- und Bildbewertungsdatensätze mit Besucherzuordnung.\\
Vorschaubilder & Bestand gültiger erzeugter Varianten, Abmessungen, Format, Status und Ableitungsversion.\\
Einstellungen & Skalare Werte für Anwendung, Design, Funktionen, Aktualisierung, Wartung und Darstellung.\\
Audit und Telemetrie & Administratorereignisse, Betriebsdiagnose, datenschutzkontrollierte Messungen und Aggregationen.\\
Duplikatprüfliste & Administratorbezogene kanonische Paarentscheidungen und Unterdrückungsregeln für genaue Galerien.\\
Aufträge und Tokens & Begrenzter Browser-, Detektor-, Migrations-, WebDAV-, Zurücksetzungs-, Freigabe- und Wartungszustand bei erforderlicher Speicherung.\\
\bottomrule
\end{longtable}

\subsection{Beziehungen und Kaskaden}

Fremdschlüssel drücken Zuordnung aus, soweit Schemaentwicklung und Kompatibilität dies erlauben. Kaskaden entfernen abhängige Ablaufdatensätze beim Wegfall ihrer übergeordneten Identität. Inhaltslöschung läuft weiterhin über Anwendungsdienste, um Dateisystem- und Datenbankwirkungen abzustimmen. Eindeutigkeitsbedingungen kodieren Identitäten wie kanonische Duplikatpaare, genaue Administrator-Galerie-Regeln, Slugs, Tokens und Metadatenvarianten.

\subsection{Anwendungseinstellungen}

Die Tabelle \codeword{app_settings} speichert Schlüssel-Wert-Konfiguration. \codeword{app_setting()} liest Werte, \codeword{set_app_setting()} speichert normalisierte Werte, und gezielte Dienste interpretieren domänenspezifische Einstellungen. Öffentliche Sprachkonfiguration verwendet \setting{public_language} für die Websitevorgabe, \setting{public_language_selector_enabled} für die standardmäßig aktivierte Besucherüberschreibung und \setting{public_language_selector_languages} für deren geordnete nicht leere Sprachliste. Der öffentliche Vorschaubildrenderer verwendet \setting{public_thumbnail_rendering_mode} mit den Werten \setting{responsive} und \setting{progressive}.

\section{Optionale Referenz für Entwickler}
\label{sec:development}
\index{e00087@Entwicklung}
\index{c00053@Codestil}

Diese Referenz richtet sich an Personen, die PHP Gallery selbst verändern. Betreiber, die nur Veröffentlichungs- und Wartungsprüfungen benötigen, können mit Abschnitt~\ref{sec:testing} fortfahren.

\subsection{Repositoryorganisation}

\begin{longtable}{>{\ttfamily}>{\hspace{0pt}}p{0.27\linewidth}>{\hspace{0pt}}p{0.65\linewidth}}
\toprule
\textrm{Ort} & \textrm{Verantwortung}\\
\midrule
\endhead
app/controllers/ & HTTP-Controller und Antwortkoordination.\\
app/services/ & Wiederverwendbare Domänen- und Infrastrukturdienste.\\
app/views/ & Layout und wiederverwendbare serverseitige Ansichtshilfen.\\
app/lang/ & Server- und Browserübersetzungskataloge.\\
database/migrations/ & Schemaentwicklung mit Zeitstempeln.\\
public/assets/ & Frameworkfreies JavaScript, CSS und öffentliche Ressourcen.\\
scripts/ & Migrations-, Integritäts-, Bereitstellungs- und Wartungswerkzeuge.\\
tests/ & Eigenständige PHP- und gezielte JavaScript-Tests.\\
winapp/ & Windows-spezifische Werkzeuge und Integrationen.\\
\bottomrule
\end{longtable}

\subsection{Codekonventionen}

Neue PHP-Dateien verwenden strikte Typen und Einrückung mit vier Leerzeichen. Controller koordinieren Anfrageverhalten; Dienste besitzen wiederverwendbare Logik. Neues JavaScript und CSS bleiben frameworkfrei und gehören zu den öffentlichen Ressourcen. Migrationen verwenden Dateinamen mit Zeitstempelpräfix. Bestehende erklärende Kommentare, Docstrings und PHPDoc werden bei Verhaltensentwicklung erhalten und korrigiert.

Docstrings sind für benannte Funktionen, Methoden und Klassen oder vergleichbare Deklarationen einschließlich Interfaces, Traits und Enums erforderlich. PHPDoc/JSDoc steht unmittelbar über der Deklaration; Python-Docstrings stehen zuerst im~Deklarationskörper. Benannte Funktionen und Methoden benötigen eine sachliche Zweckbeschreibung, einen typisierten und beschriebenen Eintrag für jeden expliziten Parameter sowie einen typisierten Rückgabeeintrag einschließlich void. Klassen benötigen eine sachliche Zweckbeschreibung. Anonyme Funktionen, Closures, Pfeilfunktionen, Callbacks, Variablen, Konstanten und Klasseneigenschaften benötigen keine Deklarations-Docstrings; benannte Deklarationen innerhalb von Callbacks werden weiterhin geprüft. Dateiköpfe mit Autorenangaben, native Signaturtypen, Kommentare zu Betriebsrichtlinien und Parserabdeckung bleiben eigenständige Anforderungen.

Atomare Module sollten eine zusammenhängende Verantwortung besitzen. Gemeinsame Verträge verdienen normalisierte Eingaben, ausdrückliche Ausgaben und gezielte Tests. Dynamische administrative und öffentliche Fragmente benötigen lebenszyklussichere Einrichtung und Abbau.

\subsection{Code vor der Quellprüfung schreiben}

Die~kanonische Datei \filepath{AGENTS.md} verlangt benannte Deklarationen und~ihre vollständige Dokumentation in~derselben Änderung. Vor den~Tags steht ein sachlicher Zweck; beschreiben Sie jeden expliziten Parameter einschließlich Standardwerten und~neuen Parametern, tatsächliche Rückgabevarianten sowie konkrete Sammlungselemente und~Objektfelder. Native PHP/Python-Typen bleiben eigenständig. Anonyme Callbacks benötigen keinen Deklarations-Docstring; benannte Deklarationen darin schon. Laufzeitrichtlinien benötigen unmittelbar am~betroffenen Ort zusätzlich Zweck, Typ, Einheiten, Geltungsbereich, Verbraucher und~Begründung. Die~Beispiele in~\filepath{docs/AGENT_AUTHORING.md} werden durch bestehende Scanner in~einem registrierten Regressionstest geprüft; negative Fixtures erhalten Fehler für unvollständige Verträge.

GitHub Actions führt die Quellvorprüfung vor der gehosteten Vorbereitung aus. Agenten folgen CI-first: aktuellen autorisierten Arbeitszweig, HEAD und ausdrückliche Push-Berechtigung prüfen, Quellen/Tests/Dokumentation fertigstellen, committen und genau diesen Zweig pushen. \codeword{develop} und \codeword{main} niemals direkt ändern. Fehlender Zugriff bedeutet BLOCKED. Lokale Audits, Generatoren und Testschleifen sind Ausnahmediagnosen. CI behält die unveränderliche Merge-Base mit \codeword{origin/main}; Releases verwenden den vorherigen stabilen Tag. Rotes oder langsames CI rechtfertigt keine rein lokale Qualifikation.

Berichte erfassen HEAD, aufgelöstes Vergleichs-SHA und~Arbeitsbaumzustand. Ein geänderter Checkout liefert Rückmeldung, keine Qualifikation eines sauberen Commits. Die~verknüpfte Datei \filepath{source-failures.md} enthält alle Befunde fehlgeschlagener Änderungsprüfungen, Blockaden und~Befunde überzogener Kategorien; vollständige wertfreie JSON-Berichte erhalten Belege und~kennzeichnen historische Schulden. Kandidaten-Workflows laden diese gemeinsam hoch. Reparieren Sie die~gesamte Änderungsgruppe vor Wiederholung; unveränderte Eingaben sind keine Reparatur. Vervollständigen Sie alle vier Handbuchquellen, bereiten Sie generierte Artefakte nach der~letzten Änderung vor und~führen Sie den~vollständigen Audit einmal für den~exakten Kandidaten aus. Erfassen Sie SHA, Basis, Profil/Ergebnis und~wesentliche Abdeckungslücken. Messen Sie Erstversuchsergebnisse und~Reparaturzyklen aus tatsächlichen Berichten; erwartete Einsparungen bleiben bis zur Beobachtung unbekannt.

\subsection{Eine Funktion hinzufügen}

Eine typische Funktionsänderung folgt dieser Abfolge:

\begin{enumerate}
  \item Verfolgen Sie die aktuelle Implementierung und bestimmen Sie zuständige Controller, Dienste, Ansichten, Browsermodule, Speicherung und Tests.
  \item Definieren Sie Verträge für Zugriff, CSRF, Umfang, Ersatzwege, Migration und Nutzung ohne JavaScript.
  \item Ergänzen Sie gezielte Dienste und Routenintegration.
  \item Ergänzen Sie nur neue Migrationen, wenn dauerhafte Schemaänderungen erforderlich sind.
  \item Ergänzen Sie übersetzte Oberflächentexte.
  \item Ergänzen Sie gezielte automatisierte Tests und dokumentierte manuelle Prüfung.
  \item Aktualisieren Sie dauerhafte Markdown-Referenzen und~alle vier TeX-Handbuchquellen, wenn sich Benutzer-, Admin-, Architektur- oder Entwickleranweisungen ändern.
  \item Gehostete Candidate preparation führt \filepath{scripts/prepare_candidate.php} und die schreibgeschützte Prüfung aus: Laufzeitmodule, indexbasiertes Produktionsinventar und Hashes vollständiger Dateien entstehen in Abhängigkeitsreihenfolge. Nur geänderte Artefakte werden auf denselben autorisierten Zweig committet; konkurrierende Änderungen werden abgewiesen und genau das vorbereitete SHA qualifiziert. Agenten pushen Quellcommits und prüfen CI statt lokal erneut zu generieren.
  \item Die Übergabe nennt GitHub-Actions-URL, Arbeitszweig, genaues vorbereitetes SHA, unveränderliche Vergleichsbasis, Pflichtjobs und manuelle/fehlende Abdeckung. Veraltete, fehlende, übersprungene, abgebrochene, rote oder blockierte Pflichtprüfungen sind keine grüne Qualifikation. Quellen korrigieren und denselben Zweig für neues CI pushen.
\end{enumerate}

\subsection{Quellkandidaten für CI vorbereiten}

Während der Entwicklung sind kleine, getrennte Commits erlaubt. Die Qualifikation muss sich jedoch auf~eine konkrete Quellrevision beziehen, zu~der alle generierten Artefakte passen. Dauerhafte technische Dokumentation gehört zur~gleichen Änderungsgruppe. Änderungen an~Administration, Architektur, Laufzeitabhängigkeiten, Sicherheit, Paketierung oder CI erfordern die~Prüfung aller vier TeX-Handbuchquellen, nicht nur der~sichtbaren Benutzeroberfläche. Bei normaler Entwicklung bleiben \filepath{PATCH_NOTES.md} und~die gespeicherten Handbuch-PDFs unverändert.

\begin{enumerate}
  \item Gehostete Candidate preparation führt \filepath{scripts/prepare_candidate.php} und die schreibgeschützte Prüfung aus: Laufzeitmodule, indexbasiertes Produktionsinventar und Hashes vollständiger Dateien entstehen in Abhängigkeitsreihenfolge. Nur geänderte Artefakte werden auf denselben autorisierten Zweig committet; konkurrierende Änderungen werden abgewiesen und genau das vorbereitete SHA qualifiziert. Agenten pushen Quellcommits und prüfen CI statt lokal erneut zu generieren.
  \item Die Übergabe nennt GitHub-Actions-URL, Arbeitszweig, genaues vorbereitetes SHA, unveränderliche Vergleichsbasis, Pflichtjobs und manuelle/fehlende Abdeckung. Veraltete, fehlende, übersprungene, abgebrochene, rote oder blockierte Pflichtprüfungen sind keine grüne Qualifikation. Quellen korrigieren und denselben Zweig für neues CI pushen.
\end{enumerate}

Der~stabile GitHub-Workflow \filepath{.github/workflows/candidate-preparation.yml} kann Feature- und~Fix-Branches vorbereiten und~die exakte SHA des~veröffentlichten Kandidaten qualifizieren. Er~prüft die~Branchidentität, verweigert veraltete Schreibvorgänge und~führt weder Merge noch Squash oder erzwungene Pushes aus. Das~zentrale schreibgeschützte Gate \codeword{--profile=candidate-preflight} erkennt veraltete Manifeste, Laufzeitpläne, Inventare und~Quellrichtlinienverstöße vor der~aufwendigen CI-Matrix. Prüfungsrunner dürfen fehlende Artefakte nicht stillschweigend reparieren. Die~spätere Release-Vorbereitung aktualisiert nur Handbuchversionen und~Datumsangaben und~kompiliert die~vorhandenen TeX-Quellen zu~PDF; fehlende Fachtexte erstellt sie nicht.

\section{Ihre Galerie vor der Veröffentlichung prüfen}
\label{sec:testing}
\index{t00241@Tests}
\index{q00199@Qualitätssicherung}

Für Galeriebetreiber bedeutet Testen, die Website als vorgesehene Zielgruppe zu durchlaufen. Eine technisch fehlerfreie Seite kann dennoch einen privaten Ort, unklaren Titel, falsches Datum, versehentliche Downloaderlaubnis oder unhandliches mobiles Layout enthalten. Eine kurze menschliche Prüfung erkennt solche Präsentations- und Datenschutzprobleme.

\subsection{Eine besucherorientierte Veröffentlichungsprüfung}
\index{v00253@Veröffentlichungscheckliste}

Öffnen Sie ein privates Browserfenster und folgen Sie dem Link, den Sie teilen möchten. Prüfen Sie:

\begin{enumerate}
  \item Titel und erster Absatz erklären die Sammlung.
  \item Das gewählte Titelbild bleibt bei kleiner Größe klar.
  \item Untergalerien erscheinen in sinnvoller Reihenfolge.
  \item Fototitel und Beschreibungen passen zu sichtbaren Motiven.
  \item Private Fotos und Galerien bleiben verborgen.
  \item Passwort- oder Freigabezugriff funktioniert in einer neuen Sitzung.
  \item GPS-Karten zeigen nur beabsichtigte Orte.
  \item Die erste Seite lädt auf einem Telefon angenehm.
  \item Großer Betrachter, Vorwärts-/Rückwärtssteuerungen und Schließen wirken natürlich.
  \item Suche, Schlagwörter, Abstimmung und Downloads passen zum Galeriezweck.
\end{enumerate}

Bitten Sie eine andere Person, die Galerie ohne Anleitung auszuprobieren. Ihr erster Klick und ihre erste Frage zeigen häufig Navigation oder Formulierungen, die nur dem Betreiber offensichtlich erschienen.

\subsection{Automatisierte Prüfungen für Softwarebetreuer}

Der zentrale Audit prüft MIME-Typen und Namen aus der Zwischenablage, synthetische Einfügeereignisse im Browser, normales Multipart-Hochladen und den Erhalt des Seitenbereichs. Diese Tests lesen nicht die Betriebssystem-Zwischenablage. Erfassen Sie tatsächliche Ergebnisse für Windows Chrome/Edge/Firefox, Linux Chromium/Firefox und macOS Safari/Chrome/Firefox anhand der manuellen Matrix in \filepath{TESTING.md}, mit Browser- und Servervorbereitung auf direkter Seite und im Seitenbereich. Nicht ausgeführte Prüfungen bleiben offen.

Die registrierten Miniaturtests prüfen begrenzte Rasterheader, Abbruch und Ressourcenfreigabe. Chromium-Fixtures prüfen außerdem FileList-Gleichheit bei 100 Einträgen, einen 320 Pixel breiten Seitenbereich, Bereinigung bei Zurücksetzen und Navigation, Fokus sowie klassische und vorbereitete Wiederholungen ohne doppelte Übernahme. URL-/Decoderzahlen sind keine Messung des gesamten Prozessspeichers; echte Betriebssystem-Zwischenablage sowie Firefox/Safari bleiben manuell zu prüfen.

GitHub Actions ist für Agenten die normale und maßgebliche Prüfumgebung (CI-first). Die folgenden zentralen Audit-Befehle beschreiben den gehosteten Runner; Agenten führen sie standardmäßig nicht lokal aus. Lokale Reproduktion bleibt eine Ausnahme bei einem beobachteten umgebungsabhängigen CI-Fehler, tatsächlich unerreichbarem CI oder ausdrücklicher Anforderung. Lokale Nachweise ersetzen keine gehostete Pflichtabdeckung.

\begin{lstlisting}[language=bash]
# Hosted source feedback (exceptional local diagnosis only):
php scripts/audit.php --profile=release-preflight
php scripts/audit.php --profile=quick
# Hosted full source qualification:
php scripts/audit.php --profile=full
# Hosted authoritative release job:
php scripts/audit.php --profile=release
\end{lstlisting}

Der~gehostete Veröffentlichungsablauf führt vor aufwendigen AI- oder~TeX-Schritten zunächst das~besondere Richtlinien-Gate \codeword{--profile=release-preflight} aus. Diese Vorprüfung ersetzt nicht die~vollständige Release-Matrix für~den~exakten Kandidaten.

Der registrierte Breadcrumb-Designtest vergleicht die Stillisten für Design und physische Galerien in \filepath{docs/ADMIN_SETTINGS_INVENTORY.md} mit \codeword{breadcrumb_style_registry()}. Er schützt außerdem den Fallback \codeword{chevron} und den nur für Galerievererbung vorgesehenen Wert \codeword{inherit}. Ergänzen Sie bei einem neuen Stil beide dokumentierten Listen; \codeword{inherit} ist kein visueller Designstil.

\subsection{Manuelle Browserprüfung}

\subsubsection{Netzwerk- und Interaktionsprüfungen}

Browserabhängige Softwareprüfungen verwenden repräsentative Daten, anonyme und authentifizierte Sitzungen, leeren Cache, mobile und Desktopansichten, gegebenenfalls Betrieb ohne JavaScript, reduzierte Bewegungspräferenzen und Netzwerkdrosselung.

Prüfen Sie für öffentliche Vorschaubilddarstellung die Bereiche Elemente und Netzwerk. Responsiver Modus sollte seinen vollständigen aktiven Kandidatensatz bereitstellen. Progressiver Modus sollte vor Aktivierung einen kleinen aktiven Kandidaten und inaktive größere Kandidaten bereitstellen und danach höchstens die dokumentierte Zahl gleichzeitiger Steigerungen für relevante Karten ausführen. Prüfen Sie Layoutstabilität, Cachewiederverwendung, Fehlerersatz, Lightbox-Interaktion, Karten, Abstimmung, geschützte Medien und direkte Links.

\subsection{Sorgfalt bei Veröffentlichungen}

Der vollständige Release-Ablauf steht in \filepath{RELEASE.md}. Vergleichen Sie den genauen vorherigen stabilen Tag, vervollständigen Sie Versionshinweise und dauerhafte Dokumentation und gleichen Sie alle vier Handbuchausgaben ab. Erstellen Sie jedes PDF mit der indizierten Mehrpassfolge aus \filepath{docs/LATEX_BUILD.md}; prüfen Sie Compiler-Erfolg, Verweise, Indexeinträge und Layoutwarnungen. Visuelle PDF-Prüfung ist ausdrücklichen Anforderungen oder der Diagnose eines konkreten Layoutproblems vorbehalten; ein erfolgreicher Routinebau benötigt keine zusätzliche visuelle Freigabe.

Die Befehle dieses Abschnitts beschreiben gehostete Releases und den erhaltenen lokalen Wiederherstellungspfad für Maintainer. Normale Vorbereitung und Qualifikation laufen auf GitHub; Agenten führen diese Befehle standardmäßig nicht lokal aus. Lokale Wiederherstellung setzt eine ausdrückliche Anforderung oder tatsächlich unerreichbares CI voraus und bleibt ohne gehostete Nachweise BLOCKED. Geschützte Promotion und Veröffentlichung folgen dem später beschriebenen release-promotion-Workflow.

Erzeugen Sie nach den letzten Änderungen das Manifest, bestehen Sie Release-, Manifest- und Whitespace-Vorprüfungen, initialisieren Sie den Qualifizierungsfingerabdruck und frieren Sie die Eingaben ein. Führen Sie genau ein zentrales Audit \codeword{--profile=release} aus, direkt oder über die konfigurierte isolierte Testinstallation. Das Audit besitzt Syntax-, Regressions-, Vertrags-, Browser- und Release-Prüfungen; ergänzen Sie keinen zweiten manuellen Testbaumdurchlauf. Binden Sie seinen Bericht an den unveränderten Fingerabdruck und melden Sie Abdeckungslücken getrennt von ausstehenden Betriebsprüfungen.

Erstellen Sie Bereitstellungsordner oder ZIP nur auf ausdrückliche Anforderung. Die Werkzeuge verlangen aktuelle Integritätsdaten und prüfen das kanonische Produktionsinventar; kontrollieren Sie das angeforderte Artefakt vor Gebrauch. Ein reines CMS-Release benötigt keinen neuen WinApp-Installer. Commit, Tag, Veröffentlichung, Upload und Bereitstellung erfordern ausdrückliche Genehmigung. Prüfen Sie nach Veröffentlichung Update-Erkennung, Migrationen, Anmeldung, öffentliche Galerien, geschützte Medien und Integrität/Status auf der veröffentlichten Installation.

Während der~gewöhnlichen Entwicklung bearbeiten und~prüfen Sie die~Inhalte aller vier \codeword{.tex}-Quelldateien gemeinsam; erstellen Sie die~verfolgten PDFs nicht nach jeder Änderung neu. Die~PDFs dürfen bis~zur endgültigen Release-Vorbereitung hinter ihren Quellen zurückliegen. Automatische Versions- und~Datumsänderungen erzeugen keine sachliche Dokumentation: Neue Funktionen, Sicherheitsaspekte, Kompatibilität und~Betriebsverfahren müssen in~den~englischen, tschechischen, deutschen und~schwedischen Handbüchern redaktionell ergänzt werden.

Für~die~Vorbereitung auf~GitHub erstellen und~pushen Sie einen Zweig \codeword{release/v_X.Y} oder~\codeword{release/v_X.Y.Z} oder~starten die~Release-Qualifikation manuell von~diesem Zweig. Der~Ablauf erstellt und~committet einen geprüften deterministischen Kandidaten zurück in~denselben Zweig, baut alle vier PDFs gemeinsam und~qualifiziert genau diesen Commit. Er~führt weder einen Merge nach~\codeword{main} aus, noch erstellt er~einen Tag oder~veröffentlicht ein~Release.

\section{Die Galerie schnell wirken lassen}
\label{sec:performance}
\index{l00163@Leistung}
\index{g00140@Geschwindigkeit}

Wahrgenommene Geschwindigkeit ist keine einzelne Zahl. Besucher bemerken, wann die Seite erscheint, wann erste Fotos erkennbar werden, ob Scrollen reaktionsfähig bleibt, wie schnell Steuerungen reagieren und wie lange ein angefordertes größeres Bild benötigt. Verbindungsqualität, Serverentfernung, Seitengröße, Vorschaubildrichtlinie, Branding-Ressourcen und gewählter Renderer tragen dazu bei.

\subsection{Praktische Wege zu höherer Geschwindigkeit}

\begin{itemize}
  \item Halten Sie die Kartenanzahl je Seite für Zielgruppe und Hostingpaket angemessen.
  \item Erzeugen Sie konfigurierte Vorschaubildgrößen großer öffentlicher Sammlungen vorab.
  \item Vermeiden Sie unnötig große Hero-, Hintergrund-, Logo- und andere Designressourcen.
  \item Bevorzugen Sie den responsiven Vorschaubildrenderer, wenn Vorhersehbarkeit und Verhalten ohne JavaScript besonders zählen; verwenden Sie den progressiven Renderer, wenn dessen schrittweise Schärfung zu Sammlung und Browserverteilung passt.
  \item Hosten Sie die Galerie nach Möglichkeit geografisch nahe ihrer Hauptzielgruppe.
  \item Testen Sie mit leerem Cache. Ein aufgewärmter Entwicklerbrowser kann aufwendige Erstbesuchsübertragungen verbergen.
  \item Prüfen Sie PHP-/Datenbanklatenz getrennt von Bildübertragungszeit, bevor Sie jedes langsame Galerieverhalten als Vorschaubildproblem einstufen.
\end{itemize}

\subsection{Test bei langsamer Verbindung}

Öffnen Sie die Galerie in einem privaten Browserfenster, wählen Sie eine repräsentative mobile Ansicht, aktivieren Sie Netzwerkdrosselung in den Entwicklerwerkzeugen und laden Sie mit leerem Cache neu. Beobachten Sie die Reihenfolge von Seite, Titelbildern, sichtbaren Karten und späteren Bildsteigerungen. Wiederholen Sie mit der zu vergleichenden Renderer- oder Seitengrößeneinstellung; andernfalls machen zwischengespeicherte Dateien den Vergleich unzuverlässig.

\subsection{Technische Messungen für Betreuer}

Profilierung öffentlicher Darstellung und Browserdiagnose helfen, wenn ein subjektiver Bericht Zahlen benötigt. Vermerken Sie mindestens:

\begin{itemize}
  \item Zeit bis zum ersten Byte und serverseitige Darstellungsdauer;
  \item HTML-Übertragungsgröße;
  \item Anfrageanzahl und übertragene Bytes anfänglicher Vorschaubilder;
  \item Zeit bis zum stabilen größten sichtbaren Inhalt;
  \item Layoutverschiebungen beim Laden von Karten und Medien;
  \item Browsermodulen zurechenbare Arbeit des Hauptthreads;
  \item Anzahl und Größe progressiver Kandidatensteigerungen;
  \item Datenbankabfrageanzahl des Entwicklungsprofilers, soweit verfügbar;
  \item Bilddekodierungs-/Speicherdruck bei sehr großen Originalen oder starkem Lightbox-Zoom.
\end{itemize}

Entwicklungsprofiler und erfasste Benchmarkdatensätze dienen Vorher-/Nachhervergleichen zwischen Renderermodi und Galerielayouts. Verwenden Sie repräsentative hochdichte Bildschirme und eingeschränkte Verbindungen, statt sich nur auf einen schnellen Desktop mit warmem Cache zu verlassen.
\section{Fehlersuche}
\label{sec:troubleshooting}
\index{f00103@Fehlersuche}

\subsection{Anwendung startet nicht}

Prüfen Sie PHP-Version, erforderliche Erweiterungen, Konfigurationsort, Datenbankverbindung, beschreibbare Pfade und Webwurzel-Routing. Prüfen Sie PHP- und Webserverprotokolle. Führen Sie Syntaxprüfungen kürzlich geänderter PHP-Dateien aus und prüfen Sie ausstehende Migrationen.

\subsection{Galerie oder Bilder fehlen}

Prüfen Sie Dateisystempfade, Galerieerkennungszustand, Datenbankdatensätze, Sichtbarkeit, Zugriffsregeln übergeordneter Galerien, Passwortentsperrung, NSFW-Richtlinie, unterstützten Dateityp und relativen Bildpfad. Testen Sie als Administrator und anonymer Besucher.

\subsection{Vorschaubilder fehlen oder sind unscharf}

Prüfen Sie Vorschaubildmetadaten und Wartungsberichte. Bestätigen Sie erzeugte Größen, Formatunterstützung, konfigurierte Grenzen, Quellgeometrie, Ableitungsstatus und öffentliche Endpunktautorisierung. Prüfen Sie im responsiven Modus den browsergewählten Kandidaten und \codeword{sizes}. Prüfen Sie im progressiven Modus kleinen Kandidaten, Warteschlangenzustand, ausgewählten Ersatz und Dekodierungsergebnis.

\subsection{Administrationsaktionen verlassen den Seitenbereich}

Prüfen Sie, ob das passende JavaScript-Modul geladen wurde, delegierte Handler das Formular erkannten, AJAX-Header und Antwortverträge stimmen, das zurückgegebene Fragment erwartete Lebenszyklusmarker enthält und allgemeine Administrationsformularhandler funktionsspezifische Formulare ausnehmen. Direktes POST-Verhalten kann weiterhin den dokumentierten Ersatzweg bereitstellen.

\subsection{Migration scheitert}

Halten Sie eine geprüfte Datenbanksicherung vor und vermerken Sie Migrationsdateiname, Gallery-Version, Datenbankserverversion, begrenzte Bildschirmmeldung, Anfragereferenz und vorhandenen Schema-/Migrationslisten-Zustand. Wiederholen Sie über die Galerieverwaltungsseite, damit der normale idempotente Kompatibilitätsweg eine unterbrochene Migration abschließen kann. Version~\version{} benötigt keine Trigger, gespeicherten Routinen, \codeword{SUPER}, \codeword{log\_bin\_trust\_function\_creators} oder globale Datenbankänderung. Bleibt der begrenzte Fehler bestehen, prüfen Sie die Laufzeitdiagnose und geben Sie die Anfragereferenz an Betreuer oder Hoster weiter; fügen Sie keine Datenbankzugangsdaten oder unverarbeiteten Produktionsdaten in Supportberichte ein.

\subsection{PDF-Handbuchkompilierung scheitert}

Führen Sie die Befehle in \filepath{docs/LATEX_BUILD.md} aus. Prüfen Sie die Installation von \codeword{pdflatex} und \codeword{makeindex}. Ein erster MiKTeX-Lauf kann übliche Pakete anfordern. Löschen Sie nach einem fehlgeschlagenen Paketwechsel ignorierte Zwischen-Builddateien und wiederholen Sie die dokumentierte Abfolge.

\section{Checkliste für Sicherung und Wiederherstellung}
\label{sec:backup}
\index{w00274@Wiederherstellung}

\begin{enumerate}
  \item Exportieren Sie die vollständige Anwendungsdatenbank.
  \item Kopieren Sie \filepath{galleries/} samt Originaldateien und galerieeigenen Ressourcen.
  \item Kopieren Sie lokale Konfiguration und relevante Installationsgeheimnisse sicher.
  \item Bewahren Sie benutzerdefinierte Designressourcen und CSS auf.
  \item Vermerken Sie Anwendungsversion und ausstehenden Migrationszustand.
  \item Lagern Sie Sicherungen außerhalb des aktiven Webbaums.
  \item Testen Sie regelmäßig die Wiederherstellung in einer isolierten Umgebung.
  \item Prüfen Sie nach Wiederherstellung öffentliche Routen, geschützten Zugriff, Administratoranmeldung, Vorschaubilder, Uploads und Migrationen.
\end{enumerate}

Vom Aktualisierungssystem erstellte Sicherungen bewahren überschriebene Anwendungsdateien zur Aktualisierungswiederherstellung. Eine vollständige Betriebssicherung umfasst außerdem Datenbank, Galeriebaum, lokale Konfiguration und installationseigene Ressourcen.

\subsection{Einen Wiederherstellungssatz prüfen}
\index{w00274@Wiederherstellung!a00003@Absicherung}

Halten Sie Datenbankexport, Originalfotos, lokale Konfiguration, relevanten Speicher \filepath{data/}, Papierkorbinhalte, benutzerdefinierte Ressourcen und genaue Anwendungsversion als einen abgestimmten Wiederherstellungssatz zusammen. Eine Anbietermomentaufnahme oder ein begrenztes Wartungsfenster kann verhindern, dass Schreibvorgänge Datenbank- und Dateizustand auseinanderbringen. Lagern Sie eine Kopie außerhalb des Hosts und vereinbaren Sie, wie viel aktuelle Arbeit verloren gehen darf und wie schnell der Dienst zurückkehren muss.

Der CLI-Einstieg \filepath{scripts/recovery.php} kann diesen Satz inventarisieren und wiederhergestellte Dateistichproben sowie Manifesthashes an einem getrennten Ort prüfen. Folgen Sie \filepath{docs/RECOVERY_ASSURANCE.md} für Nachweisformat und Befehle. Das Werkzeug führt keinen Datenbankdump aus, stellt keine Produktionsumgebung wieder her und prüft keinen Sicherungsanbieter. Seine synthetische Beschädigungsübung belegt, dass die Prüfung Schäden an Testdaten erkennt, nicht dass Ihre tatsächliche Sicherung die Website wiederherstellen kann.

Isolieren Sie vor Öffnen einer wiederhergestellten Installation E-Mail, geplante Aufträge, Integrationen und Schreibzugriff von Produktionszielen. Prüfen Sie wiederhergestellte Datenbank, Beziehungen, Administratoranmeldung, Verweigerung privater Medien und Papierkorbwiederherstellung. Vermerken Sie Sicherungsalter, vergangene Wiederherstellungszeit und Beobachtungen des Bedieners, ohne Geheimnisse in Berichte zu kopieren. Verwenden Sie \filepath{docs/RECOVERY_OFF_HOST.md} für das vollständige Verfahren. Dateirücknahme des Aktualisierungssystems kehrt Datenbankmigrationen weiterhin nicht um.

\section{Betreuerprüfung und Veröffentlichungsnachweise}
\label{sec:release-evidence}
\index{v00254@Veröffentlichungsqualifikation}
\index{t00241@Tests!v00258@verwerfbare Abläufe}

Der zentrale Einstieg \filepath{scripts/audit.php} bietet drei Regressions-/Qualifikationsprofile sowie Quellvorprüfungen. \codeword{--profile=quick} verwendet eine ausdrückliche PHP-Testteilmenge, strenge Quell- und MVC-Verträge, die Python-Importrichtlinie, schnelle Node-Fixtures und Syntaxprüfung geänderter Dateien. Es lässt die vollständige PHP-Suite, Gesamtinventar und~Kategorienbudgets, WinApp-Tests, langsame Node-Fixtures und Chromium aus. Verwenden Sie \codeword{--profile=full} für die vollständige Übergabeprüfung und \codeword{--profile=release} für die Veröffentlichungsvorbereitung; ein schnelles PASS ersetzt keines dieser Profile.

PHP-Suiten verwenden standardmäßig vier Arbeitsprozesse. Setzen Sie \codeword{PHP\_GALLERY\_AUDIT\_WORKERS} auf eine ganze Zahl von 1 bis 8; ungültige Werte blockieren die Ausführung. Registrierte exklusive Tests laufen allein, nachdem aktive Prozesse beendet sind. Jeder Kindprozess behält getrennte Diagnosen und eine ab seinem tatsächlichen Start gemessene Zeitgrenze; Ergebnisse bleiben in stabiler Eingabereihenfolge. Jedes Profil misst außerdem Include-Phasen in frischen Prozessen vor \codeword{cms\_run()}: die frühe Laufzeit erlaubt eine PHP-Datei und 16~MiB Spitzenverbrauch, der Anwendungsbootstrap 40 Dateien und 16~MiB. Die Laufzeit ist nur eine Beobachtung; diese Sonden führen keine Datenbank-, Sitzungs- oder Routingarbeit aus.

PHP- und JavaScript-Syntaxprüfungen sowie unabhängige Node/Chromium-Fixtures verwenden dieselbe Prozessgrenze. Parallele PHP-Tests laufen vor exklusiven Fixtures, sodass der Pool nur einmal geleert wird; Berichte behalten ihre ursprüngliche Reihenfolge. Syntaxprüfungen führen den Quellcode nicht aus. Die vollständige Datenbank- und Konkurrenzqualifizierung bleibt vom Aufwand der Fixture und der Rechnerleistung abhängig. Prüfungen der Produktionsliste akzeptieren LF- und CRLF-Zeilenenden aus Git-Checkouts, lehnen aber Änderungen am Inhalt oder an der Dateimitgliedschaft weiterhin ab. Die MVC-Prüfung ermittelt Namespace- und Importkontexte einmal je Datei und erhält Alias- und PDO-Prüfungen, ohne vorherige Tokens bei jedem Aufruf erneut zu durchsuchen. Die vollständige Compiler- und Ladematrix für jedes Modul mit zusammengesetztem und ursprünglichem Loader läuft in den Profilen full und release. Schnelles Feedback behält Abhängigkeits-, Graph-, Kernel- und Bootstrap-Verträge bei, ohne jedes Modul zweimal zu starten.

Dieselbe Laufzeit-Leistungssuite misst neun echte Routen nur mit einer eigenen isolierten Datenbank-/HTTP-Testinstallation: robots, Startseite, Galerie, Vorschaubild, Medium, angemeldete Administration und Telemetrie sowie beide anonymen Administrationsverweigerungen. Routenlimits liegen bei 160--300 eingebundenen Dateien und 32--48~MiB Spitzenspeicher; echte Produktantworten und Verweigerungen werden ebenfalls geprüft. Ohne diese Installation werden Routenmessungen ausdrücklich übersprungen. Laufzeit bleibt eine Beobachtung.

CI führt erforderliche Chromium-Fixtures in einem eigenen Job mit \codeword{PHP\_GALLERY\_BROWSER\_REQUIRED=1} aus; fehlende, nicht startbare oder übersprungene erforderliche Browserabdeckung lässt den Job fehlschlagen. Datenbank- und Laufzeitjobs deaktivieren Browser nur in ihrer eigenen Umgebung. Die lokale Browsererkennung bleibt optional. Die Option \codeword{--quick} der Workflow-Starter delegiert direkt an das schnelle Profil, ohne Datenbank oder Daemon anzulegen. Grenzen stehen in \filepath{TESTING.md}, isolierte Qualifizierung in \filepath{docs/GALLERY_WORKFLOWS.md}.

Bei Browserfehlern zeigen Audit-Konsole und~Bericht den~Fixture-Namen und~die~konkrete fehlgeschlagene Prüfung oder~den~Laufzeitfehler direkt an, höchstens fünf Probleme je~Suite. Der~vollständige Auszug bleibt sieben Tage in~\codeword{browser-map.log} im~Actions-Artefakt \codeword{required-chromium}. Der~angezeigte Berichtspfad gehört zum~Runner; laden Sie das~Artefakt im~Abschnitt Artifacts des~Laufs herunter. Der~gemeinsame Admin-Browserstarter bewahrt bei Fehlern einen begrenzten stderr-Auszug und~Seitenausnahmen, bei Zeitüberschreitung den~letzten Prüfungsfortschritt. Asynchrone Oberflächentests warten innerhalb fester Fristen auf~beobachtbaren Abschluss statt auf~eine kurze Pause. Der~SimBrief-Dialogvertrag wartet höchstens zwei Sekunden auf~die~Bereinigung nach dem~nativen Schließereignis und~prüft Dialogschluss, Editorsichtbarkeit, unveränderte URL, Fokus und~Hintergrundisolation getrennt. Wiederholte native Cancel-Klicks decken 320/390/1280px und~einen neu gerenderten Editor ab; Abbrechen darf nichts absenden. Fehlgeschlagene Tests werden nicht automatisch bis zum~Erfolg wiederholt.

Automatisierte Regression und betriebliche Abnahme beantworten unterschiedliche Fragen. Das zentrale Audit prüft Code, Verträge, Syntax, verfügbare Browserfixtures und Veröffentlichungskonsistenz. Die isolierte Ablauffixture ergänzt echte Datenbank-, HTTP- und Chromium-Abläufe mit erzeugten Fotos und Konten. Sie benötigt eine ausdrücklich ausschließlich verwerfbare Konfiguration und darf niemals Datenbank oder Medien der Installation verwenden. Folgen Sie \filepath{docs/GALLERY_WORKFLOWS.md} für Voraussetzungen; eine fehlende Umgebung ist eine Abdeckungslücke, kein Beweis einer funktionierenden Funktion.

Folgen Sie für Veröffentlichungsvorbereitung \filepath{RELEASE.md}: Schließen Sie Quellcode und Dokumentation ab, erstellen Sie alle vier PDF-Handbücher neu und prüfen Sie die Compilerdiagnosen, erzeugen Sie das Manifest und initialisieren Sie den Qualifikationsdatensatz. Führen Sie ein zentrales Veröffentlichungsaudit aus, keine Quick-/Full-/Release-Abfolge. Bei konfigurierten isolierten Voraussetzungen akzeptiert \filepath{scripts/gallery_workflow_mysql.php} die Option \codeword{--release}, um seine Fixture um dasselbe Veröffentlichungsaudit herum bereitzustellen.

Das Qualifikationswerkzeug \filepath{scripts/release_qualification.php} vermerkt den genauen Quellcode- und PDF-Fingerabdruck. Sein Befehl \codeword{render} erzeugt begrenzte Seitenvorschauen; Darstellung allein genehmigt deren Erscheinungsbild nicht. Der Befehl \codeword{record-audit} akzeptiert nur den passenden zentralen Veröffentlichungsbericht, kein älteres vollständiges Audit. Eine geänderte Quelldatei oder ein neu erstelltes PDF wählt neue ausstehende Nachweise, statt frühere Zustimmung still wiederzuverwenden.

Prüfen Sie das PDF visuell nur auf ausdrückliche Anfrage oder zur Diagnose eines konkreten Layoutproblems; optionale PDF-Prüffelder dürfen ausstehend bleiben. Prüfen Sie repräsentative Browserabläufe. Vermerken Sie, was tatsächlich von wem und mit welchem Ergebnis geprüft wurde. Halten Sie ungelöste Prüfungen ausstehend und den kurzen Aktualisierungsfunktionstest nach Veröffentlichung getrennt. Ein grüner automatisierter Bericht erlaubt keine Veröffentlichung; das Werkzeug erstellt niemals Commits oder Tags, lädt nicht hoch und veröffentlicht keine Version. Datensatzbefehle finden Sie in \filepath{docs/RELEASE_QUALIFICATION.md}.

\subsection{GitHub-gehostete Release-Vorbereitung und~Qualifikation}

Der Maintainer startet \filepath{.github/workflows/release-promotion.yml} von \codeword{main} mit abgeschlossenem Qualifikationslauf, genauem SHA und Release-Zweig. Plan liest nur Nachweise. Promote benötigt die unabhängig genehmigte Umgebung \codeword{release-promotion}, aktive PR/Statusregeln und aktiviert geschütztes Zusammenführen des Release-PR. Publish prüft nach dem Merge identische Kandidat/main-Bäume, erstellt einen unveränderlichen Tag und veröffentlicht geprüfte Pakete samt dauerhaften Nachweisen. Override benötigt Begründung und genaue Namen aller roten/übersprungenen Jobs; ursprüngliche Fehler bleiben rot. Menschliche Abnahme und Produktionshost-Smoke bleiben getrennt. Fehlende Serverregeln blockieren Schreibaktionen; Agenten pushen nie direkt nach develop/main.

Der~Workflow \filepath{.github/workflows/release-qualification.yml} bereitet den~Kandidaten auf~GitHub-Runnern vor. Er~veröffentlicht kein Release. Seine Schritte folgen diesem Vertrag:

\begin{enumerate}
  \item \textbf{Zweig und~Vorprüfung.} Ein~Push nach~\codeword{release/v_X.Y} oder~\codeword{release/v_X.Y.Z} oder~ein~manueller Start auf~diesem Zweig löst die~Vorbereitung aus. Eine~optional angegebene Version muss zum Zweignamen passen. Ein~bereits vorhandener Tag \codeword{v_X.Y[.Z]} führt zur Ablehnung. Der~vorige Release-Tag liefert die~Vergleichsbasis. Vor AI oder~TeX prüft \codeword{release-preflight} Quellsyntax, Deklarationsdokumentation, Dateiinventar, Python-Importrichtlinie und~Workflow-Verträge.
  \item \textbf{Deterministische Metadaten.} Der~Helfer gleicht Anwendungsversion, Release-Metadaten, Versionshinweis-Gerüst sowie Versions- und~Datumsmarkierungen aller vier Handbücher ab und~aktualisiert das~Produktionsinventar. Er~verfasst keine inhaltliche TeX-Dokumentation. Fertige redaktionelle Versionshinweise bleiben erhalten. Gleichen Sie vor der~Freigabe die~Inhalte aller vier Handbücher ab.
  \item \textbf{Redaktionelle Hinweise.} Eine~vollständige von~Betreuern verfasste Versionssektion bleibt unverändert. Fehlt sie oder~ist sie unvollständig, nutzt der~Workflow GitHub Copilot CLI mit~\codeword{gpt-6-luna}; nur bei fehlender Verfügbarkeit dieses Modells wechselt er~zur automatischen Modellauswahl. Für~die~Erzeugung ist \codeword{COPILOT_GITHUB_TOKEN} erforderlich. Der~vollständige Git-Textdiff wird ohne Kürzung nach Bytezahl übergeben; generierte und~binäre Release-Artefakte bleiben ausgeschlossen. Getrennte temporäre Ausgabedateien verhindern blockierte Pipes. Jeder Git-Befehl hat~eine Frist von~30 Sekunden; der~gesamte Belegschritt ist auf~zwei Minuten begrenzt. Überschrittene Metadaten- oder~Diagnoselimits stoppen die~Vorbereitung ausdrücklich. Strenge Markdown-Prüfung verhindert Platzhalter, falsche Versionen, zusätzliche Abschnitte und~Assistentenkommentare. Fehler bei Erzeugung oder~Validierung blockieren die~Qualifikation, statt unvollständige Hinweise zu~veröffentlichen. Der~Helfer \filepath{.github/scripts/patch_notes_ai_chunks.php} übergibt jedes Originalbyte des~Git-Diffs genau einmal in~aufeinanderfolgenden, UTF-8-sicheren Teilen von~höchstens 100.000~Bytes. Ein~SHA-256-Manifest prüft alle Teile und~den~vollständigen ursprünglichen Beleg; jeder Teil benötigt eine~gültige sachliche Zusammenfassung von~höchstens 6.500~Bytes. Überschreiten die~Zusammenfassungen das~Limit von~180.000~Bytes für~die~Endanfrage, fassen weitere begrenzte Durchläufe alle Nachweise zusammen, ohne Teile wegzulassen. Copilot-Zeitüberschreitungen, fehlende Antworten, veränderte Prüfsummen oder~ungültiges endgültiges Markdown stoppen die~Qualifikation; der~Originaldiff wird nie stillschweigend gekürzt. Prompt und~Validator teilen die~verpflichtenden Hauptüberschriften \codeword{Highlights}, \codeword{Technical Details} und~\codeword{User Impact}, auch bei reinen Test-, Dokumentations- oder~Werkzeugveröffentlichungen. Nur irrelevante optionale Unterabschnitte dürfen entfallen. Beschreiben Sie belegte Auswirkungen für~Betreuer oder~unverändertes öffentliches/Admin-Verhalten ohne erfundene Produktvorteile. Erwähnungen im~Fließtext oder~falsche Überschriftenebenen werden abgelehnt. Bei Generierungs- oder~Validierungsfehlern bleibt verfügbarer ursprünglicher Modelltext sieben Tage im~Actions-Artefakt \codeword{release-notes-rejected-response}; Prompts und~Authentifizierungsdiagnostik sind ausgeschlossen. Ungültiger Text blockiert weiterhin die~Vorbereitung und~wird niemals angewendet.
  \item \textbf{Reproduzierbare Handbücher.} Der~Workflow verwendet durch Prüfsummen gesperrtes TinyTeX und~eine festgelegte TeX-Live-Paketliste. Ein~gültiger Cache wird anhand der~Lockdatei und~des~Paketinventars wiederverwendet. Englische, tschechische, deutsche und~schwedische Ausgaben werden jeweils mit~\codeword{pdflatex}, \codeword{makeindex} und~zwei weiteren \codeword{pdflatex}-Durchläufen gebaut. Der~Zeitpunkt aus~den~Release-Metadaten steuert reproduzierbare PDF- und~Manifest-Zeitstempel. Danach werden Produktionsinventar, Integritätsmanifest und~Release-Konsistenz geprüft.
  \item \textbf{Sicheres Zurückschreiben.} Nur freigegebene Vorbereitungsdateien werden gestaged. Eine~zwischenzeitliche Änderung am~Remote-Zweig verhindert einen veralteten Push; ein~identischer bereits vorbereiteter Tree darf den~aktuellen Commit wiederverwenden. Geschrieben wird nur in~denselben Release-Zweig, nicht nach~\codeword{main}.
  \item \textbf{Gate für~den~exakten Kandidaten.} Der~wiederverwendbare zentrale CI-Workflow checkt die~genaue SHA des~vorbereiteten Commits aus und~führt die~vollständige Release-Matrix einschließlich verpflichtender Chromium-Prüfung aus. Das~Gesamt-Gate besteht nur bei erfolgreicher Vorbereitung, vollständigen Hinweisen, erfolgreicher Matrix und~unveränderter Zweig-SHA. Merge, Tag und~Veröffentlichung bleiben Entscheidungen des~Betreuers.
\end{enumerate}

Weitere Einzelheiten stehen in~\filepath{RELEASE.md}, \filepath{docs/GALLERY_WORKFLOWS.md} und~\filepath{docs/LATEX_BUILD.md}. Eine~gewöhnliche TeX-Quelländerung erfordert weder sofortige PDF-Kompilierung noch ein~separates Release-Audit.

\section{Glossar}
\label{sec:glossary}
\index{g00142@Glossar}

\begin{description}
  \item[AJAX] Browseranfrage zur Aktualisierung eines Seitenbereichs unter Erhalt der umgebenden Oberfläche.
  \item[Zweigumfang] Eine Galerie und ihre Nachkommen, wenn die ausgewählte Funktion rekursiven Umfang definiert.
  \item[Kanonisches Paar] Zwei Bildkennungen, gespeichert in stabiler Reihenfolge von klein nach groß.
  \item[CSRF-Token] Sitzungsgebundener Wert zur Prüfung absichtlicher zustandsändernder Formularübermittlungen.
  \item[Ableitung] Erzeugte Mediendarstellung, beispielsweise ein JPEG- oder WebP-Vorschaubild.
  \item[EXIF] In unterstützte Bilddateien eingebettete Kamera- und Aufnahmemetadaten.
  \item[Integritätsmanifest] Versionierte Liste von Kernpfaden und erwarteten Hashes.
  \item[Lightbox] Vollbild- oder überlagernder Fotobetrachter mit Navigation und zugehörigen Steuerungen.
  \item[Medienmanifest] Anfragelokale Darstellungsdaten, aus Vorschaubildmetadaten für sichtbare Bilder vorbereitet.
  \item[NSFW Guard] Alters- und Galerie-/Bildrichtlinie zur Zugriffskontrolle beschränkter Medien.
  \item[Progressive Schärfung] Vorschaubildpipeline mit kleinem Anfangsbild, die später einen dekodierten größeren Kandidaten aktiviert.
  \item[Responsive Auswahl] Browsernative Kandidatenauswahl aus einem sofort aktiven Quellensatz.
  \item[Ausgewählte Galerie] Durch die aktuelle öffentliche Route dargestellte Galerie.
  \item[Seitenbereich] Rechte kontextbezogene Administrationsoberfläche, durch Fragmentanfragen geladen und aktualisiert.
  \item[Vorschaubildgrenzen] Galeriebewusste minimale und maximale erzeugte Größen, die zur Auswahl infrage kommen.
\end{description}

\clearpage
\section{Quellen}
\label{sec:references}

\begin{thebibliography}{99}

\bibitem{readme}
PHP-Gallery-Projekt, \emph{README.md}, lokale Produktübersicht, Funktionsleitfaden, Anforderungen und Installationsdokumentation, Version~\version.

\bibitem{architecture}
PHP-Gallery-Projekt, \emph{ARCHITECTURE.md}, lokale Anwendungsarchitektur und Referenz des Teilsystemlebenszyklus, Version~\version.

\bibitem{database}
PHP-Gallery-Projekt, \emph{DATABASE.md}, lokale Migrations- und dauerhafte Schemareferenz, Version~\version.

\bibitem{codemap}
PHP-Gallery-Projekt, \emph{CODEMAP.md}, lokale Quellcodezuständigkeit und Wartungsübersicht, Version~\version.

\bibitem{settingsinventory}
PHP-Gallery-Projekt, \emph{docs/ADMIN\_SETTINGS\_INVENTORY.md}, kanonisches Inventar globaler Einstellungszuständigkeiten, Ersatzwerte, Sensibilität und Migrationen, Version~\version.

\bibitem{testing}
PHP-Gallery-Projekt, \emph{TESTING.md}, lokaler Leitfaden automatisierter und manueller Prüfung, Version~\version.

\bibitem{agents}
PHP-Gallery-Projekt, \emph{AGENTS.md}, lokale Repository-Beitrags- und Sicherheitsrichtlinien.

\bibitem{phpmanual}
The PHP Documentation Group, \emph{PHP-Handbuch}, \url{https://www.php.net/manual/en/}.

\bibitem{mysqlmanual}
Oracle Corporation, \emph{MySQL-Referenzhandbuch}, \url{https://dev.mysql.com/doc/}.

\bibitem{mariadbmanual}
MariaDB Foundation, \emph{MariaDB-Serverdokumentation}, \url{https://mariadb.com/docs/server/}.

\bibitem{htmlimages}
WHATWG, \emph{HTML Living Standard: Bilder}, \url{https://html.spec.whatwg.org/multipage/images.html}.

\bibitem{owasp}
OWASP Foundation, \emph{Leitlinien zur Webanwendungssicherheit}, \url{https://owasp.org/}.

\end{thebibliography}

\clearpage
\phantomsection
\addcontentsline{toc}{section}{Index}
\printindex

\end{document}
